\documentclass[12pt,a4paper]{report}

% --- CÁC GÓI LỆNH CƠ BẢN VÀ FONT GỐC (COMPUTER MODERN) ---
\usepackage[utf8]{inputenc}
\usepackage[T5]{fontenc}
\usepackage[vietnamese]{babel}
\usepackage[left=2.2cm, right=2.2cm, top=2.2cm, bottom=2.2cm]{geometry}
\usepackage{setspace}
\onehalfspacing

% --- CÁC GÓI TOÁN HỌC VÀ BẢNG BIỂU (CHỐNG TRÀN LỀ) ---
\usepackage{amsmath, amssymb, amsfonts}
\usepackage{array, booktabs, multirow, tabularx}
\newcolumntype{L}[1]{>{\raggedright\arraybackslash}p{#1}}
\newcolumntype{C}[1]{>{\centering\arraybackslash}p{#1}}
\newcolumntype{R}[1]{>{\raggedleft\arraybackslash}p{#1}}
\newcolumntype{Y}{>{\raggedright\arraybackslash}X}
\newcolumntype{Z}{>{\centering\arraybackslash}X}

% --- CÁC GÓI VẼ HÌNH VÀ MẠCH LOGIC (DÙNG CHO CHƯƠNG 2, 3, 4) ---
\usepackage{tikz}
\usepackage{circuitikz}
\usetikzlibrary{shapes.gates.logic.US, positioning, arrows.meta, calc}

% --- GÓI ĐỊNH DẠNG DANH SÁCH VÀ CODE HỢP NGỮ ---
\usepackage{enumitem}
\setlist[itemize]{leftmargin=1.5em, itemsep=2pt, topsep=4pt}
\setlist[enumerate]{leftmargin=1.8em, itemsep=2pt, topsep=4pt}

% --- GÓI ĐÓNG KHUNG: NỀN TRẮNG HOÀN TOÀN, KHÔNG BỊ CẮT ĐÔI QUA TRANG ---
\usepackage[most]{tcolorbox}
\tcbset{
    bwbox/.style={
        colback=white,
        colframe=black,
        colbacktitle=white,
        coltitle=black,
        boxrule=0.6pt,
        titlerule=0.4pt,
        arc=0pt,
        outer arc=0pt,
        left=8pt,
        right=8pt,
        top=6pt,
        bottom=6pt,
        before skip=10pt,
        after skip=10pt,
        fonttitle=\bfseries
    }
}

% --- CẤU HÌNH KHUNG CODE HỢP NGỮ: KHÔNG ĐÈ CHỮ, KHÔNG BỊ CẮT ĐÔI ---
\usepackage{listings}
\usepackage{etoolbox}
\lstset{
    basicstyle=\ttfamily\small,
    columns=flexible,
    keepspaces=true,
    breaklines=true,
    breakatwhitespace=true,
    frame=none,
    numbers=none,
    xleftmargin=0pt,
    xrightmargin=0pt,
    aboveskip=0pt,
    belowskip=0pt,
    showstringspaces=false,
    tabsize=4
}

% Tự động xuống dòng tạo khoảng cách 6pt trước mọi khung code và bọc khung tcolorbox không cắt trang
\BeforeBeginEnvironment{lstlisting}{%
    \par\vspace{6pt}\noindent%
    \begin{tcolorbox}[
        colback=white,
        colframe=black,
        boxrule=0.5pt,
        arc=0pt,
        outer arc=0pt,
        left=6pt,
        right=6pt,
        top=4pt,
        bottom=4pt,
        before skip=0pt,
        after skip=6pt
    ]%
}
\AfterEndEnvironment{lstlisting}{\end{tcolorbox}\par\vspace{2pt}}

% --- TIÊU ĐỀ ĐẦU TRANG VÀ CHÂN TRANG ---
\usepackage{fancyhdr}
\pagestyle{fancy}
\fancyhf{}
\fancyhead[L]{\small Tài liệu ôn tập Kiến trúc máy tính (RISC-V)}
\fancyhead[R]{\small IT3030 / IT3034 / IT3283}
\fancyfoot[C]{\small Trang \thepage}
\renewcommand{\headrulewidth}{0.4pt}
\renewcommand{\footrulewidth}{0pt}

% --- ĐỊNH DẠNG TIÊU ĐỀ CHƯƠNG MỤC (GIỐNG BẢN CŨ) ---
\usepackage{titlesec}
\titleformat{\chapter}[display]
  {\normalfont\huge\bfseries}{CHƯƠNG \thechapter}{15pt}{\Huge}
\titlespacing*{\chapter}{0pt}{-10pt}{20pt}

\usepackage[hidelinks]{hyperref}

\setlength{\parindent}{0pt}
\setlength{\parskip}{5pt}

\begin{document}

% --- TRANG BÌA ---
\begin{titlepage}
    \begin{center}
        \textbf{\large ĐẠI HỌC BÁCH KHOA HÀ NỘI}\\
        \textbf{\large TRƯỜNG CÔNG NGHỆ THÔNG TIN VÀ TRUYỀN THÔNG}\\[2.5cm]
        
        {\Huge \textbf{TÀI LIỆU ÔN TẬP TRỌNG TÂM}}\\[0.5cm]
        {\LARGE \textbf{KIẾN TRÚC MÁY TÍNH (PHIÊN BẢN RISC-V)}}\\[0.5cm]
        {\large \textbf{Mã học phần: IT3030, IT3034, IT3283}}\\[2cm]
        
        \begin{tcolorbox}[bwbox, title={Phạm vi và Cấu trúc tài liệu (Chương 1 -- Chương 7)}]
            \begin{itemize}
                \item \textbf{Chương 1:} Giới thiệu chung và Đánh giá hiệu năng máy tính.
                \item \textbf{Chương 2:} Hệ thống máy tính (Cấu trúc CPU, Bộ nhớ, Vào-Ra, Ngắt và Liên kết Bus).
                \item \textbf{Chương 3:} Kiến trúc tập lệnh RISC-V (Thanh ghi, Hợp ngữ, Mã máy 32-bit, Bản đồ bộ nhớ).
                \item \textbf{Chương 4:} Số học máy tính và Logic số (Hệ đếm, Cổng logic, Mã bù 2, Tràn số, Chuẩn IEEE 754).
                \item \textbf{Chương 5:} Bộ xử lý (Đường dữ liệu đơn/đa chu kỳ RV32I, Control Unit, Đường ống lệnh Pipeline, Hazards, Forwarding và Dự đoán rẽ nhánh).
                \item \textbf{Chương 6:} Bộ nhớ máy tính (Phân cấp bộ nhớ, SRAM/DRAM/ROM, Bộ nhớ đệm Cache, AMAT, Bộ nhớ ảo Paging/TLB và Bộ nhớ ngoài HDD/SSD/RAID).
                \item \textbf{Chương 7:} Hệ thống Vào-Ra (Mô-đun Vào-Ra, Định địa chỉ MMIO/Isolated I/O, Polling, Ngắt/PIC/CSRs, DMA và Thực hành MMIO RARS).
            \end{itemize}
            Tài liệu phân tách nghiêm ngặt kiến thức theo đúng giới hạn từng chương của bài giảng chính thức, kèm quy trình giải từng bước và lời giải chi tiết cho toàn bộ ngân hàng câu hỏi Giữa kỳ và Cuối kỳ.
        \end{tcolorbox}
        
        \vfill
        {\large \textbf{Hà Nội -- 2026}}
    \end{center}
\end{titlepage}

\tableofcontents
\newpage

% --- NẠP 7 FILE CHƯƠNG ---
\chapter{GIỚI THIỆU CHUNG VÀ HIỆU NĂNG MÁY TÍNH}

\section{Phần 1: Lý thuyết trọng tâm}

\subsection{Máy tính, chương trình và phân loại máy tính}

\textbf{1. Khái niệm cơ bản}
\begin{itemize}
    \item \textbf{Máy tính (Computer):} Là thiết bị điện tử hoạt động theo chương trình, có khả năng nhập dữ liệu (input data), lưu trữ dữ liệu, xử lý dữ liệu và xuất ra thông tin (output information).
    \item \textbf{Chương trình (Program):} Là dãy các lệnh được lập trình sẵn và được lưu trong bộ nhớ để điều khiển máy tính thực hiện theo.
\end{itemize}

\textbf{2. Bốn thành phần cơ bản của mọi máy tính}\\
Mọi máy tính đều được cấu thành từ 4 bộ phận cơ bản:
\begin{enumerate}
    \item \textbf{Bộ xử lý trung tâm (CPU -- Central Processing Unit):} Điều khiển hoạt động của máy tính và xử lý dữ liệu.
    \item \textbf{Bộ nhớ chính (Main Memory):} Chứa các lệnh và dữ liệu của chương trình hiện đang được thực hiện. (Các thiết bị lưu trữ như ổ cứng HDD, SSD thuộc bộ nhớ ngoài, dùng để lưu trữ phần mềm và dữ liệu lâu dài khi chương trình chưa chạy).
    \item \textbf{Hệ thống vào-ra (Input/Output -- I/O):} Trao đổi thông tin giữa máy tính với môi trường bên ngoài.
    \item \textbf{Liên kết hệ thống (System Interconnection):} Kết nối và vận chuyển thông tin giữa CPU, bộ nhớ chính và hệ thống vào-ra.
\end{enumerate}

\textbf{3. Phân loại máy tính}\\
Dựa theo quy mô, công dụng và giá thành, máy tính được chia thành 5 nhóm:
\begin{itemize}
    \item \textbf{Máy tính cá nhân (PC -- Personal Computers):} Gồm máy tính để bàn (Desktop) và máy tính xách tay (Laptop), phục vụ mục đích đa dụng cho cá nhân.
    \item \textbf{Thiết bị di động cá nhân (PMD -- Personal Mobile Devices):} Điện thoại thông minh (Smartphones), máy tính bảng (Tablets).
    \item \textbf{Máy tính nhúng / IoT (Embedded Computers / Internet of Things):} Máy tính được đặt ẩn bên trong một thiết bị khác và được thiết kế chuyên dụng cho một chức năng cụ thể. Máy tính nhúng có kết nối Internet được gọi là thiết bị IoT.
    \item \textbf{Máy chủ (Servers):} Dùng trong mạng máy tính để quản lý tài nguyên và cung cấp dịch vụ cho nhiều người dùng; có hiệu năng và độ tin cậy cao, giá thành từ hàng nghìn đến hàng triệu USD.
    \item \textbf{Máy tính quy mô kho / Siêu máy tính (WSC -- Warehouse-Scale Computers / Supercomputers):} Kết nối hàng nghìn máy chủ với nhau (Clusters). Siêu máy tính là một dạng WSC dùng cho các bài toán tính toán lớn trong khoa học và kỹ thuật, giá thành từ hàng triệu đến hàng trăm triệu USD.
\end{itemize}

\subsection{Kiến trúc máy tính, phân lớp và các mức ngôn ngữ}

\textbf{1. Ba khía cạnh của Kiến trúc máy tính}
\begin{itemize}
    \item \textbf{Kiến trúc tập lệnh (ISA -- Instruction Set Architecture):} Nghiên cứu máy tính theo cách nhìn của người lập trình hợp ngữ; là giao diện giữa phần mềm và phần cứng. Các kiến trúc tập lệnh thông dụng hiện nay gồm:
    \begin{itemize}
        \item \textit{Intel x86:} Ra đời từ năm 1978, dùng phổ biến trong Laptop, Desktop, Servers.
        \item \textit{ARM:} Ra đời từ năm 1985, dùng phổ biến trong Smartphones, Tablets.
        \item \textit{RISC-V:} Kiến trúc tập lệnh mở (miễn phí bản quyền), phát triển từ năm 2010 tại Đại học UC Berkeley.
    \end{itemize}
    \item \textbf{Tổ chức máy tính (Computer Organization) hay Vi kiến trúc (Microarchitecture):} Nghiên cứu việc thiết kế máy tính ở mức cao để thực thi kiến trúc tập lệnh (thiết kế CPU, hệ thống nhớ, cấu trúc bus).
    \item \textbf{Phần cứng (Hardware):} Nghiên cứu thiết kế mạch chi tiết và công nghệ đóng gói vật lý của máy tính. Cùng một kiến trúc tập lệnh có thể có nhiều sản phẩm phần cứng và tổ chức khác nhau.
\end{itemize}

\textbf{2. Bảy ý tưởng lớn trong thiết kế Kiến trúc máy tính}
\begin{enumerate}
    \item Sử dụng trừu tượng hóa để đơn giản hóa thiết kế (Use abstraction to simplify design).
    \item Tối ưu hóa cho trường hợp thường gặp (Make the common case fast).
    \item Tăng hiệu năng nhờ xử lý song song (Performance via parallelism).
    \item Tăng hiệu năng nhờ kỹ thuật đường ống lệnh (Performance via pipelining).
    \item Tăng hiệu năng nhờ dự đoán (Performance via prediction).
    \item Phân cấp bộ nhớ (Hierarchy of memories).
    \item Đảm bảo độ tin cậy nhờ dự phòng (Dependability via redundancy).
\end{enumerate}

\textbf{3. Các yếu tố quyết định hiệu năng của chương trình}
\begin{itemize}
    \item \textbf{Thuật toán:} Quyết định số phép toán cần thực hiện.
    \item \textbf{Ngôn ngữ lập trình, trình biên dịch và kiến trúc tập lệnh:} Quyết định số lệnh máy (Instruction Count -- $IC$) được tạo ra để thực hiện chương trình.
    \item \textbf{Bộ xử lý và hệ thống bộ nhớ:} Quyết định tốc độ thực hiện các lệnh máy.
    \item \textbf{Hệ thống vào-ra và hệ điều hành:} Quyết định tốc độ thực hiện các thao tác vào-ra dữ liệu.
\end{itemize}

\textbf{4. Phân lớp máy tính và các mức chuyển đổi mã chương trình}
\begin{itemize}
    \item \textbf{Ba lớp của hệ thống máy tính:} Từ trong ra ngoài gồm Phần cứng (Bộ xử lý, bộ nhớ, mô-đun vào-ra) $\rightarrow$ Phần mềm hệ thống (Hệ điều hành quản lý tài nguyên và Chương trình dịch) $\rightarrow$ Phần mềm ứng dụng.
    \item \textbf{Quá trình dịch chương trình qua 3 mức ngôn ngữ:}
    \begin{enumerate}
        \item \textbf{Ngôn ngữ bậc cao (High-Level Language -- HLL):} Như ngôn ngữ C, gần với bài toán cần giải quyết, độc lập với phần cứng.
        \item \textbf{Trình biên dịch (Compiler):} Dịch chương trình từ Ngôn ngữ bậc cao sang Hợp ngữ (Assembly Language -- mô tả lệnh dưới dạng văn bản gợi nhớ, phụ thuộc vào từng kiến trúc tập lệnh cụ thể).
        \item \textbf{Trình hợp dịch (Assembler):} Dịch từ Hợp ngữ sang Ngôn ngữ máy (Machine Language / Mã máy -- được mã hóa hoàn toàn bằng các bit nhị phân 0 và 1 để phần cứng trực tiếp thực thi).
    \end{enumerate}
\end{itemize}

\textbf{5. Sự tiến hóa của công nghệ máy tính}
\begin{itemize}
    \item \textbf{Thế hệ đèn điện tử chân không (1950s):} Máy tính ENIAC (1946 -- xử lý theo số thập phân) và máy tính IAS (1952 -- do John von Neumann thiết kế theo nguyên lý chương trình lưu trữ trong bộ nhớ ``stored program'', xử lý theo số nhị phân).
    \item \textbf{Thế hệ Transistors (1960s)} $\rightarrow$ \textbf{Vi mạch SSI, MSI, LSI (1970s)} $\rightarrow$ \textbf{Vi mạch VLSI (1980s)} $\rightarrow$ \textbf{Vi mạch ULSI (1990s -- nay)}.
    \item \textbf{Sự phát triển độ rộng bit của vi xử lý:} 4-bit (Intel 4004 năm 1971) $\rightarrow$ 8-bit (1972) $\rightarrow$ 16-bit (1978, máy tính IBM-PC ra đời năm 1981) $\rightarrow$ 32-bit (1985) $\rightarrow$ 64-bit (2001) $\rightarrow$ Đa lõi (Multicores -- từ năm 2006, tích hợp nhiều nhân CPU trên cùng một chip).
\end{itemize}

\subsection{Đánh giá hiệu năng máy tính (Trọng tâm tính toán)}

\begin{tcolorbox}[bwbox, title=Bảng quy đổi đơn vị chuẩn khi giải bài tập Hiệu năng]
Khi thay số vào công thức, luôn đổi tất cả đại lượng về đơn vị chuẩn: \textbf{giây (s)} đối với thời gian, \textbf{Hertz (Hz)} đối với tần số, và \textbf{số lệnh} đối với khối lượng lệnh.
\begin{center}
\renewcommand{\arraystretch}{1.25}
\begin{tabularx}{\linewidth}{|L{4.2cm}|L{4.3cm}|Y|}
\hline
\textbf{Đại lượng} & \textbf{Đơn vị trong đề bài} & \textbf{Quy đổi về đơn vị chuẩn} \\ \hline
\multirow{2}{4.2cm}{Tần số xung nhịp ($f_0$)} & $1\text{ MHz}$ (MegaHertz) & $= 10^6\text{ Hz}$ ($1$ triệu chu kỳ/giây) \\ \cline{2-3}
 & $1\text{ GHz}$ (GigaHertz) & $= 10^9\text{ Hz}$ ($1$ tỷ chu kỳ/giây) \\ \hline
\multirow{4}{4.2cm}{Thời gian ($T_0, t_{CPU}$)} & $1\text{ ms}$ (mili-giây) & $= 10^{-3}\text{ s}$ \\ \cline{2-3}
 & $1\text{ }\mu\text{s}$ (micro-giây) & $= 10^{-6}\text{ s}$ \\ \cline{2-3}
 & $1\text{ ns}$ (nano-giây) & $= 10^{-9}\text{ s}$ \\ \cline{2-3}
 & $1\text{ ps}$ (pico-giây) & $= 10^{-12}\text{ s}$ \\ \cline{2-3}
\multirow{2}{4.2cm}{Số lượng lệnh ($IC$)} & $1\text{ triệu lệnh}$ & $= 10^6\text{ lệnh}$ \\ \cline{2-3}
 & $1\text{ tỷ lệnh}$ (hoặc $1.0\text{E}9$) & $= 10^9\text{ lệnh}$ ($x\text{E}y = x \times 10^y$) \\ \hline
\end{tabularx}
\end{center}
\textbf{Quan hệ rút gọn giữa GHz và ns:} Vì $1\text{ GHz} = 10^9\text{ Hz}$ và $1\text{ ns} = 10^{-9}\text{ s}$ nên:
$$T_0\text{ (tính bằng ns)} = \frac{1}{f_0\text{ (tính bằng GHz)}}$$
\end{tcolorbox}

\textbf{1. Chu kỳ xung nhịp ($T_0$) và Tần số xung nhịp ($f_0$)}
\begin{itemize}
    \item \textbf{Chu kỳ xung nhịp ($T_0$ -- Clock period):} Là thời gian kéo dài của một chu kỳ xung nhịp (đơn vị: s, ms, ns, ps).
    \item \textbf{Tần số xung nhịp ($f_0$ -- Clock rate / Clock speed):} Là số chu kỳ xung nhịp diễn ra trong $1$ giây (đơn vị: Hz, MHz, GHz).
    \item \textbf{Công thức liên hệ:}
    $$T_0 = \frac{1}{f_0} \quad \text{và} \quad f_0 = \frac{1}{T_0}$$
\end{itemize}

\textbf{2. Tổng số chu kỳ xung nhịp ($n$) và Thời gian thực thi của CPU ($t_{CPU}$)}
\begin{itemize}
    \item Gọi $IC$ (Instruction Count) là tổng số lệnh máy được thực hiện của chương trình; $CPI$ (Cycles Per Instruction) là số chu kỳ xung nhịp trung bình để thực hiện một lệnh.
    \item \textbf{Tổng số chu kỳ xung nhịp ($n$) cần thiết để chạy xong chương trình:}
    $$n = IC \times CPI$$
    \item \textbf{Thời gian CPU thực hiện chương trình ($t_{CPU}$):}
    $$t_{CPU} = n \times T_0 = \frac{n}{f_0} = IC \times CPI \times T_0 = \frac{IC \times CPI}{f_0}$$
    \item \textbf{Các công thức suy diễn ngược (khi đề bài cho trước $t_{CPU}$):}
    $$n = t_{CPU} \times f_0 \quad ; \quad CPI = \frac{t_{CPU} \times f_0}{IC} \quad ; \quad IC = \frac{t_{CPU} \times f_0}{CPI}$$
\end{itemize}

\textbf{3. Công thức tính CPI trung bình ($CPI_{TB}$) khi có nhiều lớp lệnh}\\
Nếu chương trình gồm $K$ loại lệnh khác nhau, mỗi loại lệnh $i$ có số chu kỳ riêng là $CPI_i$:
\begin{itemize}
    \item \textbf{Trường hợp 1:} Đề bài cho số lượng lệnh cụ thể $IC_1, IC_2, \dots, IC_K$ của từng loại:
    $$IC = \sum_{i=1}^{K} IC_i \quad ; \quad n = \sum_{i=1}^{K} (CPI_i \times IC_i) \quad ; \quad CPI_{TB} = \frac{n}{IC} = \frac{\sum_{i=1}^{K} (CPI_i \times IC_i)}{\sum_{i=1}^{K} IC_i}$$
    \item \textbf{Trường hợp 2:} Đề bài cho tỷ lệ phần trăm $p_1\%, p_2\%, \dots, p_K\%$ của từng loại lệnh (với $\sum p_i\% = 100\% = 1$):
    $$CPI_{TB} = \sum_{i=1}^{K} (p_i\% \times CPI_i)$$
\end{itemize}

\textbf{4. Định nghĩa Hiệu năng ($P$) và So sánh tốc độ giữa hai máy tính}
\begin{itemize}
    \item Hiệu năng ($P$ -- Performance) là nghịch đảo của thời gian thực hiện chương trình:
    $$P = \frac{1}{t_{CPU}}$$
    Máy tính nào có thời gian thực hiện chương trình \textbf{nhỏ hơn} thì máy tính đó \textbf{nhanh hơn} (có hiệu năng cao hơn).
    \item Khi đề bài hỏi \textbf{``Máy A nhanh hơn máy B bao nhiêu lần?''}, ta lấy hiệu năng của máy A chia cho hiệu năng của máy B (tức là lấy \textbf{thời gian của máy B chia cho thời gian của máy A}):
    $$k = \frac{P_A}{P_B} = \frac{t_B}{t_A}$$
    \item Nếu hai máy A và B có \textbf{cùng kiến trúc tập lệnh (ISA)} và chạy cùng một chương trình thì số lệnh thực thi trên hai máy là bằng nhau ($IC_A = IC_B = IC$). Khi đó $IC$ tự triệt tiêu:
    $$\frac{P_A}{P_B} = \frac{t_B}{t_A} = \frac{CPI_B \times T_B}{CPI_A \times T_A} = \frac{CPI_B / f_B}{CPI_A / f_A}$$
\end{itemize}

\textbf{5. Chỉ số tốc độ thực thi MIPS và MFLOPS}
\begin{itemize}
    \item \textbf{MIPS (Millions of Instructions Per Second):} Số triệu lệnh máy thực hiện được trong $1$ giây:
    $$MIPS = \frac{IC}{t_{CPU} \times 10^6} = \frac{f_0}{CPI \times 10^6}$$
    \begin{itemize}
        \item Nếu thay $f_0$ theo đơn vị \textbf{MHz} ($10^6\text{ Hz}$): $MIPS = \dfrac{f_0\text{ (MHz)}}{CPI}$.
        \item Nếu thay $f_0$ theo đơn vị \textbf{GHz} ($10^9\text{ Hz}$): $MIPS = \dfrac{f_0\text{ (GHz)} \times 1000}{CPI}$.
        \item Suy ngược ra $CPI$ khi biết $MIPS$: $CPI = \dfrac{f_0}{MIPS \times 10^6} = \dfrac{f_0\text{ (GHz)} \times 1000}{MIPS}$.
    \end{itemize}
    \item \textbf{Hạn chế của chỉ số MIPS:} Chỉ đo tốc độ thực thi lệnh mà không phản ánh khả năng của từng lệnh; không dùng để so sánh hai máy khác kiến trúc tập lệnh (khác ISA); thay đổi theo từng chương trình trên cùng một máy; đôi khi một chương trình có nhiều lệnh đơn giản hơn lại cho chỉ số MIPS cao hơn nhưng tổng thời gian chạy thực tế lại lâu hơn.
    \item \textbf{MFLOPS (Millions of Floating-Point Operations Per Second):} Số triệu phép toán dấu phẩy động thực hiện được trong $1$ giây:
    $$MFLOPS = \frac{\text{Số phép toán dấu phẩy động thực thi}}{t_{CPU} \times 10^6}$$
    Các bội số lớn hơn: GFLOPS ($10^9$), TFLOPS ($10^{12}$), PFLOPS ($10^{15}$).
\end{itemize}

\newpage
\section{Phần 2: Toàn bộ bài tập Chương 1 và Hướng dẫn giải}

\subsection{Nhóm 1: Khái niệm cơ bản và Các thành phần của máy tính}

\textbf{Bài tập 1.1 (Đề Cuối kỳ 2020.2, 2021.1, 2024.2 \& Ôn tập Cuối kỳ 2025.1, 2025.2)}\\
Hãy chỉ ra tên của thành phần phần cứng trong cấu trúc máy tính thực hiện chức năng chính: ``Chứa toàn bộ các lệnh và dữ liệu của chương trình hiện đang được thực hiện''?
\begin{enumerate}[label=\Alph*.]
    \item Bộ nhớ chính (RAM)
    \item Tập thanh ghi (Register File)
    \item Bộ nhớ đệm nhanh (Cache memory)
    \item Bộ nhớ ngoài (Ổ đĩa SSD/HDD)
\end{enumerate}
\textbf{Hướng dẫn giải:}
\begin{itemize}
    \item \textbf{Kiến thức áp dụng:} Theo mô hình máy tính von Neumann (Slide Chương 1, trang 6 và trang 16), mọi chương trình khi đang được thực hiện đều phải nằm trong Bộ nhớ chính (Main Memory / RAM). Bộ nhớ ngoài (SSD/HDD) chỉ lưu trữ phần mềm và dữ liệu khi chưa chạy.
    \item \textbf{Kết luận:} Đáp án đúng là \textbf{A (Bộ nhớ chính)}.
\end{itemize}

\subsection{Nhóm 2: Tính chu kỳ xung nhịp ($T_0$), CPI và Thời gian thực thi ($t_{CPU}$)}

\begin{tcolorbox}[bwbox, title=Quy trình 3 bước giải bài toán tính thời gian thực thi của CPU]
\begin{enumerate}
    \item \textbf{Bước 1 (Đổi đơn vị):} Đổi số lệnh $IC$ ra đơn vị lệnh ($1\text{ triệu} = 10^6$, $1\text{ tỷ} = 10^9$); đổi tần số $f_0$ ra Hz ($1\text{ MHz} = 10^6\text{ Hz}$, $1\text{ GHz} = 10^9\text{ Hz}$).Lưu ý: Chỉ lấy số lượng \textbf{lệnh mã máy} sau khi biên dịch, không lấy số dòng lệnh ngôn ngữ bậc cao.
    \item \textbf{Bước 2 (Áp dụng công thức):} Tính thời gian theo đơn vị giây (s):
    $$t_{CPU} = \frac{IC \times CPI}{f_0}$$
    \item \textbf{Bước 3 (Đổi về đơn vị của đáp án):}
    \begin{itemize}
        \item Đổi từ giây (s) sang mili-giây (ms): Nhân kết quả với $10^3$ (vì $10^{-3}\text{ s} = 1\text{ ms}$).
        \item Đổi từ giây (s) sang nano-giây (ns): Nhân kết quả với $10^9$ (vì $10^{-9}\text{ s} = 1\text{ ns}$).
    \end{itemize}
\end{enumerate}
\end{tcolorbox}

\textbf{Bài tập 1.2 (Slide Chương 1 -- Trang 24)}\\
Một bộ xử lý có tần số xung nhịp $f_0 = 4\text{ GHz}$. Hãy tính chu kỳ xung nhịp $T_0$ của bộ xử lý đó theo đơn vị nano-giây (ns).

\textbf{Hướng dẫn giải:}
\begin{itemize}
    \item \textbf{Bước 1:} Đổi tần số về đơn vị Hz: $f_0 = 4\text{ GHz} = 4 \times 10^9\text{ Hz}$.
    \item \textbf{Bước 2:} Áp dụng công thức tính chu kỳ xung nhịp:
    $$T_0 = \frac{1}{f_0} = \frac{1}{4 \times 10^9} = 0.25 \times 10^{-9}\text{ s}$$
    \item \textbf{Bước 3:} Đổi sang đơn vị nano-giây ($1\text{ ns} = 10^{-9}\text{ s}$): $T_0 = 0.25\text{ ns}$.
    \item \textbf{Kết luận:} Đáp số là \textbf{0.25 ns}.
\end{itemize}

\textbf{Bài tập 1.3 (Đề Cuối kỳ 2020.2 \& Đề Cuối kỳ 2021.1)}\\
\textit{Biến thể 1 (Đề 2020.2):} Bộ xử lý trung tâm có tốc độ xung nhịp cố định là $2.9\text{ GHz}$. Hãy cho biết chu kỳ xung nhịp tương ứng của bộ xử lý đó là bao nhiêu?\\
A. $7.9\text{ ns}$ \quad B. $0.13\text{ ps}$ \quad C. $0.13\text{ ms}$ \quad D. $0.345\text{ ns}$\\[4pt]
\textit{Biến thể 2 (Ôn tập Cuối kỳ 2025.1 \& 2025.2 -- Câu 1):} Bộ xử lý có tốc độ xung nhịp là $7.9\text{ GHz}$. Cho biết chu kỳ xung nhịp tương ứng là bao nhiêu?\\
A. $7.9\text{ ns}$ \quad B. $0.13\text{ ps}$ \quad C. $0.13\text{ ms}$ \quad D. $0.13\text{ ns}$

\textbf{Hướng dẫn giải:}
\begin{itemize}
    \item \textbf{Xét Biến thể 1 ($f_0 = 2.9\text{ GHz} = 2.9 \times 10^9\text{ Hz}$):}
    $$T_0 = \frac{1}{f_0} = \frac{1}{2.9 \times 10^9}\text{ s} \approx 0.3448 \times 10^{-9}\text{ s} \approx 0.345\text{ ns}$$
    \item \textbf{Xét Biến thể 2 ($f_0 = 7.9\text{ GHz} = 7.9 \times 10^9\text{ Hz}$):}
    $$T_0 = \frac{1}{f_0} = \frac{1}{7.9 \times 10^9}\text{ s} \approx 0.1266 \times 10^{-9}\text{ s} \approx 0.13\text{ ns}$$
    \item \textbf{Kết luận:} Biến thể 1 chọn \textbf{D (0.345 ns)}; Biến thể 2 chọn \textbf{D (0.13 ns)}.
\end{itemize}

\textbf{Bài tập 1.4 (Đề Giữa kỳ 2018.2 -- Mã đề 198 Câu 2 \& Mã đề 199 Câu 17)}\\
Một phần mềm viết bằng ngôn ngữ hợp ngữ với $30$ lệnh, được biên dịch thành $40$ lệnh mã máy. Biết rằng $CPI$ trung bình là $2$, và CPU có tốc độ xung nhịp $4\text{ GHz}$. Hỏi thời gian thực hiện phần mềm là bao lâu?
\begin{enumerate}[label=\Alph*.]
    \item $60\text{ ns}$
    \item $80\text{ ns}$
    \item $20\text{ ns}$
    \item $15\text{ ns}$
\end{enumerate}
\textbf{Hướng dẫn giải:}
\begin{itemize}
    \item \textbf{Bước 1:} Xác định các thông số chuẩn. CPU trực tiếp thực thi các \textbf{lệnh mã máy}, do đó tổng số lệnh thực thi là $IC = 40\text{ lệnh}$ (con số $30$ lệnh hợp ngữ trước khi dịch là dữ kiện gây nhiễu). Ta có $CPI = 2$ và tần số $f_0 = 4\text{ GHz} = 4 \times 10^9\text{ Hz}$.
    \item \textbf{Bước 2:} Tính thời gian thực thi của CPU:
    $$t_{CPU} = \frac{IC \times CPI}{f_0} = \frac{40 \times 2}{4 \times 10^9} = \frac{80}{4} \times 10^{-9}\text{ s} = 20 \times 10^{-9}\text{ s}$$
    \item \textbf{Bước 3:} Đổi sang đơn vị nano-giây: $t_{CPU} = 20\text{ ns}$.
    \item \textbf{Kết luận:} Đáp án đúng là \textbf{C (20 ns)}.
\end{itemize}

\textbf{Bài tập 1.5 (Đề Giữa kỳ 2023.2 -- Đề 1, Câu 1)}\\
Máy tính có bộ xử lý với tốc độ xung nhịp là $f = 4\text{ GHz}$. Giả sử có chương trình được thực hiện trên máy tính đó với số lệnh $n = 5\text{ triệu lệnh}$ và $CPI$ trung bình bằng $2$. Hãy cho biết thời gian của CPU ($t_{CPU}$) để hoàn thành chương trình đó (tính theo ms)?

\textbf{Hướng dẫn giải:}
\begin{itemize}
    \item \textbf{Bước 1:} Đổi đơn vị:
    \begin{itemize}
        \item Số lượng lệnh: $IC = 5\text{ triệu lệnh} = 5 \times 10^6\text{ lệnh}$.
        \item Chỉ số $CPI = 2$.
        \item Tần số xung nhịp: $f_0 = 4\text{ GHz} = 4 \times 10^9\text{ Hz}$.
    \end{itemize}
    \item \textbf{Bước 2:} Tính thời gian thực thi của CPU:
    $$t_{CPU} = \frac{IC \times CPI}{f_0} = \frac{5 \times 10^6 \times 2}{4 \times 10^9} = \frac{10}{4} \times 10^{-3}\text{ s} = 2.5 \times 10^{-3}\text{ s}$$
    \item \textbf{Bước 3:} Vì $10^{-3}\text{ s} = 1\text{ ms}$, nên $t_{CPU} = 2.5\text{ ms}$.
    \item \textbf{Kết luận:} Đáp số là \textbf{2.5 ms}.
\end{itemize}

\textbf{Bài tập 1.6 (Đề Cuối kỳ 2024.1 -- \#54583d \& Ôn tập Cuối kỳ 2025.2 -- P1, Câu 4)}\\
Máy tính có bộ xử lý với tần số xung nhịp $5\text{ GHz}$ và $CPI$ trung bình bằng $2$. Giả sử có chương trình được thực hiện trên máy tính đó với $6\text{ triệu lệnh}$. Hãy cho biết thời gian của bộ xử lý để hoàn thành chương trình đó?
\begin{enumerate}[label=\Alph*.]
    \item $24\text{ ms}$
    \item $60\text{ ms}$
    \item $2.4\text{ ms}$
    \item $240\text{ ms}$
\end{enumerate}
\textbf{Hướng dẫn giải:}
\begin{itemize}
    \item \textbf{Bước 1:} Đổi đơn vị: $IC = 6 \times 10^6\text{ lệnh}$; $CPI = 2$; $f_0 = 5\text{ GHz} = 5 \times 10^9\text{ Hz}$.
    \item \textbf{Bước 2:} Tính thời gian thực thi của CPU:
    $$t_{CPU} = \frac{IC \times CPI}{f_0} = \frac{6 \times 10^6 \times 2}{5 \times 10^9} = \frac{12}{5} \times 10^{-3}\text{ s} = 2.4 \times 10^{-3}\text{ s} = 2.4\text{ ms}$$
    \item \textbf{Kết luận:} Đáp án đúng là \textbf{C (2.4 ms)}.
\end{itemize}

\textbf{Bài tập 1.7 (Đề Cuối kỳ 2020.2 \& 2021.1)}\\
Một hệ thống máy tính tích hợp lõi vi xử lý hoạt động ở tần số xung nhịp cố định là $1.8\text{ GHz}$ và sở hữu số chu kỳ lệnh trung bình trên mỗi câu lệnh là $CPI = 2.1$. Giả sử có một chương trình ứng dụng bao gồm tổng cộng $3\text{ triệu lệnh}$ cần thực thi trọn vẹn trên hệ thống này. Hãy tính tổng thời gian xử lý cần thiết để CPU hoàn thành chương trình trên?
\begin{enumerate}[label=\Alph*.]
    \item $3.5\text{ ms}$
    \item $11.3\text{ ms}$
    \item $35.0\text{ ms}$
    \item $3500.0\text{ ms}$
\end{enumerate}
\textbf{Hướng dẫn giải:}
\begin{itemize}
    \item \textbf{Bước 1:} Đổi đơn vị: $IC = 3 \times 10^6\text{ lệnh}$; $CPI = 2.1$; $f_0 = 1.8\text{ GHz} = 1.8 \times 10^9\text{ Hz}$.
    \item \textbf{Bước 2:} Tính thời gian thực thi:
    $$t_{CPU} = \frac{IC \times CPI}{f_0} = \frac{3 \times 10^6 \times 2.1}{1.8 \times 10^9} = \frac{6.3}{1.8} \times 10^{-3}\text{ s} = 3.5 \times 10^{-3}\text{ s} = 3.5\text{ ms}$$
    \item \textbf{Kết luận:} Đáp án đúng là \textbf{A (3.5 ms)}.
\end{itemize}

\textbf{Bài tập 1.8 (Đề Cuối kỳ 2022.2 \& Cuối kỳ 2024.2 -- Trang 12)}\\
Cho một chương trình có $6\text{ triệu lệnh}$ được thực hiện trên máy tính có bộ xử lý với tốc độ xung nhịp $2\text{ GHz}$ và $CPI$ trung bình là $1.5$. Hãy cho biết thời gian của bộ xử lý để hoàn thành chương trình đó?
\begin{enumerate}[label=\Alph*.]
    \item $72\text{ ms}$
    \item $450\text{ ms}$
    \item $4.5\text{ ms}$
    \item $45\text{ ms}$
\end{enumerate}
\textbf{Hướng dẫn giải:}
\begin{itemize}
    \item \textbf{Bước 1:} Đổi đơn vị: $IC = 6 \times 10^6\text{ lệnh}$; $CPI = 1.5$; $f_0 = 2\text{ GHz} = 2 \times 10^9\text{ Hz}$.
    \item \textbf{Bước 2:} Tính thời gian thực thi:
    $$t_{CPU} = \frac{IC \times CPI}{f_0} = \frac{6 \times 10^6 \times 1.5}{2 \times 10^9} = \frac{9}{2} \times 10^{-3}\text{ s} = 4.5 \times 10^{-3}\text{ s} = 4.5\text{ ms}$$
    \item \textbf{Kết luận:} Đáp án đúng là \textbf{C (4.5 ms)}.
\end{itemize}

\textbf{Bài tập 1.9 (Đề Giữa kỳ 2019.2 -- Mã đề 243 Câu 11, Mã đề 244 Câu 8 \& Giữa kỳ 2022.2 -- Đề 0008 Câu 1)}\\
Hãy tính chỉ số $CPI$ trung bình của phần mềm trong 3 đề thi dưới đây:
\begin{enumerate}[label=\alph*)]
    \item \textbf{(Mã đề 243 -- Câu 11):} Bộ xử lý có tần số xung nhịp $2\text{ GHz}$ thực hiện phần mềm gồm $400\,000\text{ lệnh máy}$ trong $0.0003\text{ giây}$.
    \item \textbf{(Mã đề 244 -- Câu 8):} Bộ xử lý có tần số xung nhịp $300\text{ MHz}$ thực hiện phần mềm gồm $1\,000\,000\text{ lệnh máy}$ trong $0.004\text{ giây}$.
    \item \textbf{(Giữa kỳ 2022.2 -- Đề 0008, Câu 1):} Bộ xử lý có tần số xung nhịp $50\text{ MHz}$ thực hiện phần mềm gồm $800\,000\text{ lệnh máy}$ trong $0.04\text{ giây}$.
\end{enumerate}
\textbf{Hướng dẫn giải:}\\
Từ công thức gốc $t_{CPU} = \dfrac{IC \times CPI}{f_0}$, ta rút ra công thức tính $CPI$:
$$CPI = \frac{t_{CPU} \times f_0}{IC}$$
\begin{itemize}
    \item \textbf{Giải câu a:} Với $f_0 = 2\text{ GHz} = 2 \times 10^9\text{ Hz}$, $IC = 400\,000 = 4 \times 10^5\text{ lệnh}$, $t_{CPU} = 0.0003\text{ s} = 3 \times 10^{-4}\text{ s}$:
    $$CPI = \frac{3 \times 10^{-4} \times 2 \times 10^9}{4 \times 10^5} = \frac{6 \times 10^5}{4 \times 10^5} = 1.5$$
    \item \textbf{Giải câu b:} Với $f_0 = 300\text{ MHz} = 300 \times 10^6\text{ Hz}$, $IC = 1\,000\,000 = 10^6\text{ lệnh}$, $t_{CPU} = 0.004\text{ s}$:
    $$CPI = \frac{0.004 \times 300 \times 10^6}{10^6} = 0.004 \times 300 = 1.2$$
    \item \textbf{Giải câu c:} Với $f_0 = 50\text{ MHz} = 50 \times 10^6\text{ Hz}$, $IC = 800\,000 = 0.8 \times 10^6\text{ lệnh}$, $t_{CPU} = 0.04\text{ s}$:
    $$CPI = \frac{0.04 \times 50 \times 10^6}{0.8 \times 10^6} = \frac{2}{0.8} = 2.5$$
    \item \textbf{Kết luận:} Câu a: $CPI = 1.5$; Câu b: $CPI = 1.2$; Câu c: $CPI = 2.5$.
\end{itemize}

\subsection{Nhóm 3: Bài tập về chỉ số tốc độ thực thi lệnh MIPS}

\textbf{Bài tập 1.10 (Slide Chương 1 -- Trang 35 \& Trang 36)}
\begin{enumerate}[label=\alph*)]
    \item Tính chỉ số $MIPS$ của bộ xử lý với $f_0 = 2\text{ GHz}$ và $CPI = 4$.
    \item Tính chỉ số $CPI$ của bộ xử lý với $f_0 = 1\text{ GHz}$ và tốc độ đạt $400\text{ MIPS}$.
\end{enumerate}
\textbf{Hướng dẫn giải:}
\begin{itemize}
    \item \textbf{Giải câu a:} Đổi $f_0 = 2\text{ GHz} = 2 \times 10^9\text{ Hz}$. Áp dụng công thức:
    $$MIPS = \frac{f_0}{CPI \times 10^6} = \frac{2 \times 10^9}{4 \times 10^6} = \frac{2000}{4} = 500\text{ MIPS}$$
    \item \textbf{Giải câu b:} Đổi $f_0 = 1\text{ GHz} = 10^9\text{ Hz}$. Áp dụng công thức suy ra $CPI$:
    $$CPI = \frac{f_0}{MIPS \times 10^6} = \frac{10^9}{400 \times 10^6} = \frac{1000}{400} = 2.5$$
    \item \textbf{Kết luận:} Câu a: \textbf{500 MIPS}; Câu b: \textbf{CPI = 2.5}.
\end{itemize}

\textbf{Bài tập 1.11 (Đề Cuối kỳ 2020.2, 2022.2 -- \#1fd28d, 2024.1 -- \#3e5442, 2024.2 \& Ôn tập 2025.2)}\\
Hãy tính số triệu lệnh thực hiện được trong một giây (đơn vị MIPS) cho 3 câu hỏi trắc nghiệm sau:
\begin{enumerate}[label=\alph*)]
    \item \textbf{(Đề CK 2020.2):} Tần số xung nhịp $f_0 = 1.2\text{ GHz}$ và $CPI = 1.5$.\\
    A. $1\text{ MIPS}$ \quad B. $1800\text{ MIPS}$ \quad C. $800\text{ MIPS}$ \quad D. $1250\text{ MIPS}$
    \item \textbf{(Đề CK 2022.2 \& 2024.2 -- \#1fd28d):} Tần số xung nhịp $f_0 = 2.5\text{ GHz}$ và $CPI = 2$.\\
    A. $1250\text{ MIPS}$ \quad B. $1.25\text{ MIPS}$ \quad C. $5\text{ MIPS}$ \quad D. $500\text{ MIPS}$
    \item \textbf{(Đề CK 2024.1 -- \#3e5442 \& Ôn tập 2025.2 Câu 21):} Tần số xung nhịp $f_0 = 3\text{ GHz}$ và $CPI = 1.5$.\\
    A. $2\text{ MIPS}$ \quad B. $4.5\text{ MIPS}$ \quad C. $2000\text{ MIPS}$ \quad D. $500\text{ MIPS}$
\end{enumerate}
\textbf{Hướng dẫn giải:}\\
Áp dụng công thức $MIPS = \dfrac{f_0}{CPI \times 10^6} = \dfrac{f_0\text{ (GHz)} \times 1000}{CPI}$:
\begin{itemize}
    \item \textbf{Câu a:} $MIPS = \dfrac{1.2 \times 1000}{1.5} = \dfrac{1200}{1.5} = 800\text{ MIPS}$ $\rightarrow$ Chọn \textbf{C}.
    \item \textbf{Câu b:} $MIPS = \dfrac{2.5 \times 1000}{2} = \dfrac{2500}{2} = 1250\text{ MIPS}$ $\rightarrow$ Chọn \textbf{A}.
    \item \textbf{Câu c:} $MIPS = \dfrac{3 \times 1000}{1.5} = \dfrac{3000}{1.5} = 2000\text{ MIPS}$ $\rightarrow$ Chọn \textbf{C}.
\end{itemize}

\textbf{Bài tập 1.12 (Đề Cuối kỳ 2021.2 -- Mã đề B1 Câu 1.1 \& Mã đề B5 Câu 1.2)}
\begin{enumerate}[label=\alph*)]
    \item \textbf{(Mã đề B1 -- Câu 1.1):} Máy tính A có bộ xử lý với tần số xung nhịp $1.5\text{ GHz}$, có khả năng thực hiện $600\text{ triệu lệnh}$ trên một giây. Tính $CPI$ của nó?
    \item \textbf{(Mã đề B5 -- Câu 1.2):} Cho máy tính A với bộ xử lý có $MIPS = 900$. Hãy tính $CPI$ của máy tính đó nếu biết chu kỳ xung nhịp của bộ xử lý là $T_A = 300\text{ ps}$.
\end{enumerate}
\textbf{Hướng dẫn giải:}
\begin{itemize}
    \item \textbf{Giải câu a:} Thực hiện $600\text{ triệu lệnh/giây}$ tức là $MIPS = 600$. Với $f_0 = 1.5\text{ GHz} = 1.5 \times 10^9\text{ Hz}$:
    $$CPI = \frac{f_0}{MIPS \times 10^6} = \frac{1.5 \times 10^9}{600 \times 10^6} = \frac{1500}{600} = 2.5$$
    \item \textbf{Giải câu b:}
    \begin{itemize}
        \item Bước 1: Đổi chu kỳ xung nhịp sang giây: $T_A = 300\text{ ps} = 300 \times 10^{-12}\text{ s}$.
        \item Bước 2: Tìm tần số xung nhịp $f_A = \dfrac{1}{T_A} = \dfrac{1}{300 \times 10^{-12}}\text{ Hz} = \dfrac{10^{12}}{300}\text{ Hz}$.
        \item Bước 3: Thay vào công thức tính $CPI$:
        $$CPI = \frac{f_A}{MIPS \times 10^6} = \frac{10^{12}}{300 \times 900 \times 10^6} = \frac{10^6}{270\,000} = \frac{100}{27} \approx 3.704$$
    \end{itemize}
    \item \textbf{Kết luận:} Câu a: \textbf{CPI = 2.5}; Câu b: $\mathbf{CPI = \frac{100}{27} \approx 3.704}$.
\end{itemize}

\subsection{Nhóm 4: Tính CPI trung bình và Thời gian thực thi với nhiều lớp lệnh}

\begin{tcolorbox}[bwbox, title=Quy trình giải bài toán có nhiều nhóm lệnh (Class A, B, C...)]
\begin{itemize}
    \item \textbf{Dạng 1 (Cho số lượng lệnh $IC_i$ của từng nhóm):}
    \begin{enumerate}
        \item Tính tổng số chu kỳ: $n = (IC_1 \times CPI_1) + (IC_2 \times CPI_2) + \dots$
        \item Tính tổng số lệnh: $IC = IC_1 + IC_2 + \dots$
        \item Lấy $CPI_{TB} = \dfrac{n}{IC}$ và thời gian $t_{CPU} = \dfrac{n}{f_0}$.
    \end{enumerate}
    \item \textbf{Dạng 2 (Cho tỷ lệ phần trăm $p_i\%$ của từng nhóm):}
    \begin{enumerate}
        \item Nếu đề bài ẩn tỷ lệ của một nhóm (ghi là ``còn lại''), lấy $100\%$ trừ đi các nhóm đã biết để tìm tỷ lệ phần trăm của nhóm đó.
        \item Tính $CPI_{TB} = (p_1\% \times CPI_1) + (p_2\% \times CPI_2) + \dots$ (nhớ đổi $p\%$ ra số thập phân, ví dụ $35\% = 0.35$).
    \end{enumerate}
\end{itemize}
\end{tcolorbox}

\textbf{Bài tập 1.13 (Slide Chương 1 -- Trang 32, Đề Giữa kỳ 2018.2 -- Mã đề 204 Câu 19 \& Cuối kỳ 2021.2 -- Mã đề B4 Câu 1)}\\
Tính chỉ số $CPI$ trung bình cho các dãy lệnh được cho theo bảng số lượng lệnh dưới đây:
\begin{enumerate}[label=\alph*)]
    \item \textbf{(Đề Giữa kỳ 2018.2 -- Mã đề 204, Câu 19):} Phần mềm X có 3 nhóm lệnh A, B, C với $CPI$ tương ứng là $1, 2, 3$ và số lệnh tương ứng là $4, 1, 1$.
    \item \textbf{(Slide Chương 1 -- Trang 32):} Ba loại lệnh A, B, C có $CPI$ lần lượt là $1, 2, 3$. Dãy lệnh 1 có số lệnh $(A, B, C) = (20, 10, 20)$; Dãy lệnh 2 có số lệnh $(A, B, C) = (40, 10, 10)$.
    \item \textbf{(Đề Cuối kỳ 2021.2 -- Mã đề B4, Câu 1):} Ba loại lệnh A, B, C có $CPI$ lần lượt là $1, 2, 4$.
    \begin{itemize}
        \item Dãy lệnh 1 có số lệnh $(A, B, C) = (1, 2, 2)$.
        \item Dãy lệnh 2 có số lệnh $(A, B, C) = (5, 2, 1)$.
        \item Dãy lệnh 3 có số lệnh $(A, B, C) = (4, 1, 1)$.
    \end{itemize}
\end{enumerate}
\textbf{Hướng dẫn giải:}\\
Áp dụng công thức $CPI_{TB} = \dfrac{\sum (CPI_i \times IC_i)}{\sum IC_i}$:
\begin{itemize}
    \item \textbf{Giải câu a:}
    $$CPI_{TB} = \frac{(1 \times 4) + (2 \times 1) + (3 \times 1)}{4 + 1 + 1} = \frac{4 + 2 + 3}{6} = \frac{9}{6} = 1.5$$
    \item \textbf{Giải câu b:}
    \begin{itemize}
        \item Dãy lệnh 1: $CPI_{TB1} = \dfrac{(1 \times 20) + (2 \times 10) + (3 \times 20)}{20 + 10 + 20} = \dfrac{20 + 20 + 60}{50} = \dfrac{100}{50} = 2.0$.
        \item Dãy lệnh 2: $CPI_{TB2} = \dfrac{(1 \times 40) + (2 \times 10) + (3 \times 10)}{40 + 10 + 10} = \dfrac{40 + 20 + 30}{60} = \dfrac{90}{60} = 1.5$.
    \end{itemize}
    \item \textbf{Giải câu c:}
    \begin{itemize}
        \item Dãy lệnh 1: $CPI_{TB1} = \dfrac{(1 \times 1) + (2 \times 2) + (4 \times 2)}{1 + 2 + 2} = \dfrac{1 + 4 + 8}{5} = \dfrac{13}{5} = 2.6$.
        \item Dãy lệnh 2: $CPI_{TB2} = \dfrac{(1 \times 5) + (2 \times 2) + (4 \times 1)}{5 + 2 + 1} = \dfrac{5 + 4 + 4}{8} = \dfrac{13}{8} = 1.625$.
        \item Dãy lệnh 3: $CPI_{TB3} = \dfrac{(1 \times 4) + (2 \times 1) + (4 \times 1)}{4 + 1 + 1} = \dfrac{4 + 2 + 4}{6} = \dfrac{10}{6} \approx 1.67$.
    \end{itemize}
\end{itemize}

\textbf{Bài tập 1.14 (Đề Giữa kỳ 2025.2 -- Đề 0013 Câu 75, Câu 944 \& Đề 0012 Câu 110)}
\begin{enumerate}[label=\alph*)]
    \item \textbf{(Đề 0013 -- Câu 75):} Một phần mềm gồm: $15\%$ lệnh tốn $3$ chu kỳ, còn lại là lệnh $4$ chu kỳ, $35\%$ lệnh $5$ chu kỳ và $10\%$ lệnh $2$ chu kỳ. $CPI$ trung bình của phần mềm là bao nhiêu?
    \item \textbf{(Đề 0013 -- Câu 944):} Một phần mềm gồm: $10\%$ lệnh tốn $4$ chu kỳ, $40\%$ lệnh $6$ chu kỳ, $30\%$ lệnh $8$ chu kỳ, còn lại là lệnh $10$ chu kỳ. $CPI$ trung bình của phần mềm là bao nhiêu?
    \item \textbf{(Đề 0012 -- Câu 110):} Bộ xử lý RISC-V có tần số xung nhịp $1\text{ GHz}$ thực hiện phần mềm gồm $200\,000$ lệnh máy với $CPI = 3$, $200\,000$ lệnh máy với $CPI = 4$ và $300\,000$ lệnh máy với $CPI = 2$. Hỏi phần mềm đó chạy trong bao lâu?
\end{enumerate}
\textbf{Hướng dẫn giải:}
\begin{itemize}
    \item \textbf{Giải câu a:}
    \begin{itemize}
        \item Bước 1: Tính tỷ lệ phần trăm của nhóm lệnh ``còn lại'' (nhóm $4$ chu kỳ):
        $$p_{\text{còn lại}} = 100\% - (15\% + 35\% + 10\%) = 100\% - 60\% = 40\% = 0.40$$
        \item Bước 2: Tính $CPI$ trung bình:
        $$CPI_{TB} = (0.15 \times 3) + (0.40 \times 4) + (0.35 \times 5) + (0.10 \times 2) = 0.45 + 1.60 + 1.75 + 0.20 = 4.0$$
    \end{itemize}
    \item \textbf{Giải câu b:}
    \begin{itemize}
        \item Bước 1: Tỷ lệ phần trăm của nhóm lệnh ``còn lại'' ($10$ chu kỳ):
        $$p_{\text{còn lại}} = 100\% - (10\% + 40\% + 30\%) = 20\% = 0.20$$
        \item Bước 2: Tính $CPI$ trung bình:
        $$CPI_{TB} = (0.10 \times 4) + (0.40 \times 6) + (0.30 \times 8) + (0.20 \times 10) = 0.4 + 2.4 + 2.4 + 2.0 = 7.2$$
    \end{itemize}
    \item \textbf{Giải câu c:}
    \begin{itemize}
        \item Bước 1: Tính tổng số chu kỳ xung nhịp ($n$) của cả phần mềm:
        $$n = (200\,000 \times 3) + (200\,000 \times 4) + (300\,000 \times 2) = 600\,000 + 800\,000 + 600\,000 = 2\,000\,000 = 2 \times 10^6\text{ chu kỳ}$$
        \item Bước 2: Với tần số $f_0 = 1\text{ GHz} = 10^9\text{ Hz}$, thời gian chạy phần mềm là:
        $$t_{CPU} = \frac{n}{f_0} = \frac{2 \times 10^6}{10^9} = 2 \times 10^{-3}\text{ s} = 0.002\text{ s} = 2\text{ ms}$$
    \end{itemize}
    \item \textbf{Kết luận:} Câu a: \textbf{4} (hoặc \textbf{4.0}); Câu b: \textbf{7.2}; Câu c: \textbf{2 ms} (hoặc \textbf{0.002 s}).
\end{itemize}

\textbf{Bài tập 1.15 (Đề Giữa kỳ 2024.1 -- \#a6b720 \& Ôn tập Cuối kỳ 2025.2 -- P1, Câu 10)}\\
Một bộ xử lý chạy một chương trình với tần số xung nhịp $4\text{ GHz}$. Chương trình gồm ba loại lệnh với các đặc điểm sau:
\begin{itemize}
    \item Loại A: chiếm $40\%$ tổng số lệnh với $CPI$ là $1$
    \item Loại B: chiếm $35\%$ tổng số lệnh với $CPI$ là $3$
    \item Loại C: chiếm $25\%$ tổng số lệnh với $CPI$ là $5$
\end{itemize}
Tổng số lệnh cho chương trình là $2\text{ tỷ lệnh}$. Tổng thời gian thực thi là bao nhiêu?
\begin{enumerate}[label=\Alph*.]
    \item $1.1\text{ giây}$
    \item $1.25\text{ giây}$
    \item $1.35\text{ giây}$
    \item $1.4\text{ giây}$
\end{enumerate}
\textbf{Hướng dẫn giải:}
\begin{itemize}
    \item \textbf{Bước 1:} Tính $CPI$ trung bình của chương trình:
    $$CPI_{TB} = (0.40 \times 1) + (0.35 \times 3) + (0.25 \times 5) = 0.40 + 1.05 + 1.25 = 2.7$$
    \item \textbf{Bước 2:} Đổi đơn vị: Tổng số lệnh $IC = 2\text{ tỷ} = 2 \times 10^9\text{ lệnh}$; tần số $f_0 = 4\text{ GHz} = 4 \times 10^9\text{ Hz}$.
    \item \textbf{Bước 3:} Tính tổng thời gian thực thi:
    $$t_{CPU} = \frac{IC \times CPI_{TB}}{f_0} = \frac{2 \times 10^9 \times 2.7}{4 \times 10^9} = \frac{5.4}{4} = 1.35\text{ giây}$$
    \item \textbf{Kết luận:} Đáp án đúng là \textbf{C (1.35 giây)}.
\end{itemize}

\textbf{Bài tập 1.16 (Ôn tập Cuối kỳ 2025.2 -- P1, Câu 55 \& Câu 70)}\\
Cho một bộ xử lý $2\text{ GHz}$ thực thi tập lệnh gồm năm lớp lệnh A, B, C, D, E với $CPI$ lần lượt là $2, 3, 4, 6$ và $10$. Máy tính đang thực hiện một chương trình trong đó các lớp lệnh A, B, C, D, E chiếm lần lượt $30\%, 30\%, 15\%, 15\%$ và $10\%$ tổng số lệnh.
\begin{enumerate}[label=\alph*)]
    \item \textbf{(Câu 70):} Tìm $CPI$ trung bình của bộ xử lý này.
    \item \textbf{(Câu 55):} Nếu thời gian thực thi của chương trình này là $20\text{ giây}$, hãy tìm tổng số lệnh ($IC$). (Viết đáp án không chứa dấu phân tách chữ số).
\end{enumerate}
\textbf{Hướng dẫn giải:}
\begin{itemize}
    \item \textbf{Giải câu a:} Tính $CPI$ trung bình:
    $$CPI_{TB} = (0.30 \times 2) + (0.30 \times 3) + (0.15 \times 4) + (0.15 \times 6) + (0.10 \times 10)$$
    $$CPI_{TB} = 0.6 + 0.9 + 0.6 + 0.9 + 1.0 = 4.0$$
    \item \textbf{Giải câu b:} Ta có $t_{CPU} = 20\text{ s}$ và $f_0 = 2\text{ GHz} = 2 \times 10^9\text{ Hz}$. Từ công thức $t_{CPU} = \dfrac{IC \times CPI_{TB}}{f_0}$, rút ra tổng số lệnh $IC$:
    $$IC = \frac{t_{CPU} \times f_0}{CPI_{TB}} = \frac{20 \times 2 \times 10^9}{4} = 10 \times 10^9 = 10\,000\,000\,000\text{ lệnh}$$
    \item \textbf{Kết luận:} Câu a điền \textbf{4}; Câu b điền \textbf{10000000000}.
\end{itemize}

\textbf{Bài tập 1.17 (Đề Giữa kỳ 2021.1 -- IT3030E, Problem 3)}\\
Cho một chương trình gồm $2000\text{ lệnh}$ thuộc 4 lớp lệnh: $750$ lệnh số học (arithmetics), $250$ lệnh ghi bộ nhớ (stores), $500$ lệnh đọc bộ nhớ (loads) và $500$ lệnh rẽ nhánh (branches). Giả sử số chu kỳ xung nhịp ($CPI$) cho các lệnh số học, load, store và branch lần lượt là $1, 5, 5$ và $2$.
\begin{enumerate}[label=\alph*)]
    \item Tính thời gian thực thi của chương trình trên bộ xử lý có tốc độ $2\text{ GHz}$.
    \item Tìm $CPI$ trung bình của chương trình này.
    \item Nếu số lượng lệnh load trong chương trình giảm đi $50\%$, hãy tính độ tăng tốc (speedup) và $CPI$ mới.
\end{enumerate}
\textbf{Hướng dẫn giải:}
\begin{itemize}
    \item \textbf{Giải câu a:}
    \begin{itemize}
        \item Bước 1: Tính tổng số chu kỳ xung nhịp ban đầu ($n_{\text{cũ}}$):
        $$n_{\text{cũ}} = (750 \times 1) + (250 \times 5) + (500 \times 5) + (500 \times 2) = 750 + 1250 + 2500 + 1000 = 5500\text{ chu kỳ}$$
        \item Bước 2: Thời gian thực thi trên bộ xử lý $f_0 = 2\text{ GHz} = 2 \times 10^9\text{ Hz}$:
        $$t_{\text{cũ}} = \frac{n_{\text{cũ}}}{f_0} = \frac{5500}{2 \times 10^9} = 2750 \times 10^{-9}\text{ s} = 2.75 \times 10^{-6}\text{ s} = 2.75\text{ }\mu\text{s}$$
    \end{itemize}
    \item \textbf{Giải câu b:} Chỉ số $CPI$ trung bình ban đầu:
    $$CPI_{\text{cũ}} = \frac{n_{\text{cũ}}}{IC_{\text{cũ}}} = \frac{5500}{2000} = 2.75$$
    \item \textbf{Giải câu c:}
    \begin{itemize}
        \item Khi số lệnh load giảm $50\%$, số lệnh load chỉ còn: $500 \times 50\% = 250\text{ lệnh}$ (giảm đi $250$ lệnh).
        \item Tổng số lệnh mới: $IC_{\text{mới}} = 2000 - 250 = 1750\text{ lệnh}$.
        \item Tổng số chu kỳ mới (giảm đi $250 \times 5 = 1250\text{ chu kỳ}$):
        $$n_{\text{mới}} = 5500 - 1250 = 4250\text{ chu kỳ}$$
        \item Chỉ số $CPI$ trung bình mới:
        $$CPI_{\text{mới}} = \frac{n_{\text{mới}}}{IC_{\text{mới}}} = \frac{4250}{1750} = \frac{17}{7} \approx 2.43$$
        \item Độ tăng tốc (Speedup) do tần số $f_0$ không đổi nên bằng tỷ số chu kỳ cũ chia cho chu kỳ mới:
        $$\text{Speedup} = \frac{t_{\text{cũ}}}{t_{\text{mới}}} = \frac{n_{\text{cũ}}}{n_{\text{mới}}} = \frac{5500}{4250} = \frac{22}{17} \approx 1.294\text{ lần}$$
    \end{itemize}
\end{itemize}

\subsection{Nhóm 5: So sánh hiệu năng và cân bằng thời gian giữa các bộ xử lý}

\textbf{Bài tập 1.18 (Slide Chương 1 -- Trang 29, Đề Giữa kỳ 2018.2 -- Mã đề 198 Câu 12 \& Đề Cuối kỳ 2021.2 -- Mã đề B1 Câu 1.2)}\\
Hai máy tính A và B có cùng kiến trúc tập lệnh (ISA). Hãy xác định máy nào nhanh hơn và nhanh hơn bao nhiêu lần trong 3 trường hợp thông số dưới đây:
\begin{enumerate}[label=\alph*)]
    \item \textbf{(Slide Chương 1 -- Trang 29):} Máy A có $T_A = 250\text{ ps}, CPI_A = 2.0$; Máy B có $T_B = 500\text{ ps}, CPI_B = 1.2$.
    \item \textbf{(Đề Cuối kỳ 2021.2 -- Mã đề B1, Câu 1.2):} Máy A có $T_A = 250\text{ ps}, CPI_A = 2.0$; Máy B có $T_B = 750\text{ ps}, CPI_B = 1.2$.
    \item \textbf{(Đề Giữa kỳ 2018.2 -- Mã đề 198 Câu 12 \& Mã đề 199 Câu 4):} Máy A có $T_A = 300\text{ ps}, CPI_A = 2.0$; Máy B có $T_B = 500\text{ ps}, CPI_B = 1.1$.
\end{enumerate}
\textbf{Hướng dẫn giải:}\\
Vì hai máy A và B có cùng kiến trúc tập lệnh (ISA), nên khi chạy cùng một chương trình thì số lệnh thực thi là như nhau: $IC_A = IC_B = IC$. Thời gian thực thi trên mỗi máy là $t = IC \times CPI \times T_0$:
\begin{itemize}
    \item \textbf{Giải câu a:}
    \begin{itemize}
        \item $t_A = IC \times 2.0 \times 250\text{ ps} = 500 \times IC\text{ (ps)}$.
        \item $t_B = IC \times 1.2 \times 500\text{ ps} = 600 \times IC\text{ (ps)}$.
        \item Vì $t_A < t_B$ ($500 < 600$) nên \textbf{máy A nhanh hơn máy B}, và nhanh hơn số lần là: $\dfrac{t_B}{t_A} = \dfrac{600}{500} = \mathbf{1.2\text{ lần}}$.
    \end{itemize}
    \item \textbf{Giải câu b:}
    \begin{itemize}
        \item $t_A = IC \times 2.0 \times 250\text{ ps} = 500 \times IC\text{ (ps)}$.
        \item $t_B = IC \times 1.2 \times 750\text{ ps} = 900 \times IC\text{ (ps)}$.
        \item Vì $t_A < t_B$ ($500 < 900$) nên \textbf{máy A nhanh hơn máy B}, và nhanh hơn số lần là: $\dfrac{t_B}{t_A} = \dfrac{900}{500} = \mathbf{1.8\text{ lần}}$.
    \end{itemize}
    \item \textbf{Giải câu c:}
    \begin{itemize}
        \item $t_A = IC \times 2.0 \times 300\text{ ps} = 600 \times IC\text{ (ps)}$.
        \item $t_B = IC \times 1.1 \times 500\text{ ps} = 550 \times IC\text{ (ps)}$.
        \item Vì $t_B < t_A$ ($550 < 600$) nên \textbf{máy B nhanh hơn máy A} (chọn đáp án ``B nhanh hơn A'').
    \end{itemize}
\end{itemize}

\textbf{Bài tập 1.19 (Đề Cuối kỳ 2023.1, 2023.2 -- \#2262f, Giữa kỳ 2024.1 -- \#8cf496 \& Ôn tập 2025.2 -- Câu 12)}
\begin{enumerate}[label=\alph*)]
    \item \textbf{(Đề CK 2023.1 -- \#2262f \& Ôn tập 2025.2 Câu 12):} Xét hai bộ xử lý cùng thực hiện một công việc. P1 có tốc độ xung nhịp $4\text{ GHz}$, $CPI$ trung bình là $0.9$ và yêu cầu thực thi $5.0\text{E}9$ lệnh. P2 có tốc độ xung nhịp $3\text{ GHz}$, $CPI$ trung bình là $0.75$ và yêu cầu thực thi $1.0\text{E}9$ lệnh. Hỏi P2 nhanh hơn P1 bao nhiêu lần khi thực hiện công việc này?
    \item \textbf{(Đề GK 2024.1 -- \#8cf496):} Xét hai bộ xử lý cùng thực hiện một công việc. P1 có tốc độ xung nhịp $1.5\text{ GHz}$, $CPI$ trung bình là $1.8$ và yêu cầu thực thi $3.0\text{E}9$ lệnh. P2 có tốc độ xung nhịp $2\text{ GHz}$, $CPI$ trung bình là $1.2$ và yêu cầu thực thi $4.0\text{E}9$ lệnh. Hỏi P2 nhanh hơn P1 bao nhiêu lần?
\end{enumerate}
\textbf{Hướng dẫn giải:}
\begin{itemize}
    \item \textbf{Giải câu a:}
    \begin{itemize}
        \item Bước 1: Tính thời gian thực thi của P1 ($t_1$):
        $$t_1 = \frac{IC_1 \times CPI_1}{f_1} = \frac{5.0 \times 10^9 \times 0.9}{4 \times 10^9} = \frac{4.5}{4} = 1.125\text{ giây}$$
        \item Bước 2: Tính thời gian thực thi của P2 ($t_2$):
        $$t_2 = \frac{IC_2 \times CPI_2}{f_2} = \frac{1.0 \times 10^9 \times 0.75}{3 \times 10^9} = \frac{0.75}{3} = 0.25\text{ giây}$$
        \item Bước 3: Đề hỏi ``P2 nhanh hơn P1 bao nhiêu lần'', ta lấy thời gian của P1 chia cho thời gian của P2:
        $$\frac{P_2}{P_1} = \frac{t_1}{t_2} = \frac{1.125}{0.25} = 4.5\text{ lần}$$
    \end{itemize}
    \item \textbf{Giải câu b:}
    \begin{itemize}
        \item Bước 1: Tính thời gian thực thi của P1 ($t_1$):
        $$t_1 = \frac{IC_1 \times CPI_1}{f_1} = \frac{3.0 \times 10^9 \times 1.8}{1.5 \times 10^9} = \frac{5.4}{1.5} = 3.6\text{ giây}$$
        \item Bước 2: Tính thời gian thực thi của P2 ($t_2$):
        $$t_2 = \frac{IC_2 \times CPI_2}{f_2} = \frac{4.0 \times 10^9 \times 1.2}{2 \times 10^9} = \frac{4.8}{2} = 2.4\text{ giây}$$
        \item Bước 3: Tỷ số tốc độ của P2 so với P1:
        $$\frac{P_2}{P_1} = \frac{t_1}{t_2} = \frac{3.6}{2.4} = 1.5\text{ lần}$$
    \end{itemize}
    \item \textbf{Kết luận:} Câu a điền \textbf{4.5}; Câu b điền \textbf{1.5}.
\end{itemize}

\textbf{Bài tập 1.20 (Đề Giữa kỳ 2021.1 -- IT3030E, Problem 2)}\\
Giả sử có hai bộ xử lý P1 và P2 cùng thực hiện một kiến trúc tập lệnh (ISA) gồm 4 lớp lệnh A, B, C, D. Tốc độ xung nhịp và $CPI$ của từng lớp lệnh được cho trong bảng sau:
\begin{center}
\begin{tabularx}{0.85\textwidth}{|C{2cm}|C{2.5cm}|Z|Z|Z|Z|}
\hline
\textbf{Bộ xử lý} & \textbf{Clock Rate} & \textbf{Class A} & \textbf{Class B} & \textbf{Class C} & \textbf{Class D} \\ \hline
P1 & $2.5\text{ GHz}$ & $2$ & $1.5$ & $2$ & $1$ \\ \hline
P2 & $3\text{ GHz}$ & $1$ & $2$ & $1$ & $1$ \\ \hline
\end{tabularx}
\end{center}
Cho một chương trình có $1\,000\,000\text{ lệnh}$ được chia thành các lớp như sau: $10\%$ lớp A, $20\%$ lớp B, $50\%$ lớp C và $20\%$ lớp D.
\begin{enumerate}[label=\alph*)]
    \item Tính $CPI$ trung bình (global CPI) cho mỗi bộ xử lý.
    \item Tìm tổng số chu kỳ xung nhịp yêu cầu trong cả hai trường hợp.
    \item Bộ xử lý nào nhanh hơn?
\end{enumerate}
\textbf{Hướng dẫn giải:}
\begin{itemize}
    \item \textbf{Giải câu a (Tính $CPI$ trung bình):}
    \begin{itemize}
        \item Đối với P1:
        $$CPI_1 = (0.10 \times 2) + (0.20 \times 1.5) + (0.50 \times 2) + (0.20 \times 1) = 0.2 + 0.3 + 1.0 + 0.2 = 1.7$$
        \item Đối với P2:
        $$CPI_2 = (0.10 \times 1) + (0.20 \times 2) + (0.50 \times 1) + (0.20 \times 1) = 0.1 + 0.4 + 0.5 + 0.2 = 1.2$$
    \end{itemize}
    \item \textbf{Giải câu b (Tính tổng số chu kỳ $n = IC \times CPI$ với $IC = 10^6\text{ lệnh}$):}
    \begin{itemize}
        \item Số chu kỳ của P1: $n_1 = 10^6 \times 1.7 = 1.7 \times 10^6\text{ chu kỳ}$.
        \item Số chu kỳ của P2: $n_2 = 10^6 \times 1.2 = 1.2 \times 10^6\text{ chu kỳ}$.
    \end{itemize}
    \item \textbf{Giải câu c (So sánh thời gian thực thi):}
    \begin{itemize}
        \item Thời gian thực thi trên P1: $t_1 = \dfrac{n_1}{f_1} = \dfrac{1.7 \times 10^6}{2.5 \times 10^9} = 0.68 \times 10^{-3}\text{ s} = 0.68\text{ ms}$.
        \item Thời gian thực thi trên P2: $t_2 = \dfrac{n_2}{f_2} = \dfrac{1.2 \times 10^6}{3 \times 10^9} = 0.40 \times 10^{-3}\text{ s} = 0.40\text{ ms}$.
        \item Vì $t_2 < t_1$ ($0.40\text{ ms} < 0.68\text{ ms}$) nên \textbf{bộ xử lý P2 nhanh hơn P1} (nhanh hơn $\dfrac{0.68}{0.40} = 1.7\text{ lần}$).
    \end{itemize}
\end{itemize}

\textbf{Bài tập 1.21 (Slide Chương 1 -- Trang 26)}\\
Hai máy tính A và B cùng chạy một chương trình:
\begin{itemize}
    \item Máy tính A có tần số xung nhịp $f_A = 2\text{ GHz}$, thời gian CPU thực hiện chương trình là $t_A = 10\text{ s}$.
    \item Máy tính B cần đạt thời gian thực hiện chương trình là $t_B = 6\text{ s}$. Biết rằng số chu kỳ xung nhịp khi chạy chương trình trên máy B ($n_B$) bằng $1.2\text{ lần}$ số chu kỳ xung nhịp khi chạy chương trình trên máy A ($n_A$).
\end{itemize}
Hãy xác định tần số xung nhịp cần thiết cho máy B ($f_B$)?

\textbf{Hướng dẫn giải:}
\begin{itemize}
    \item \textbf{Bước 1:} Tính tổng số chu kỳ xung nhịp khi chạy chương trình trên máy A:
    $$n_A = t_A \times f_A = 10\text{ s} \times (2 \times 10^9\text{ Hz}) = 20 \times 10^9\text{ chu kỳ}$$
    \item \textbf{Bước 2:} Tính tổng số chu kỳ xung nhịp khi chạy chương trình trên máy B:
    $$n_B = 1.2 \times n_A = 1.2 \times 20 \times 10^9 = 24 \times 10^9\text{ chu kỳ}$$
    \item \textbf{Bước 3:} Tính tần số xung nhịp $f_B$ cần thiết để máy B hoàn thành trong $t_B = 6\text{ s}$:
    $$f_B = \frac{n_B}{t_B} = \frac{24 \times 10^9}{6} = 4 \times 10^9\text{ Hz} = 4\text{ GHz}$$
    \item \textbf{Kết luận:} Đáp số là $\mathbf{f_B = 4\text{ GHz}}$.
\end{itemize}

\textbf{Bài tập 1.22 (Tổng hợp 4 mã đề tự luận Đề Cuối kỳ 2021.2 -- Mã đề B2, B3, B6, B7 Câu 1)}\\
Giả sử có một chương trình chạy trên hai máy tính A và B. Hãy xác định khi dịch trên máy B thì chương trình đó được dịch thành bao nhiêu lệnh ($IC_B$) để cả hai máy có cùng thời gian thực hiện của CPU ($t_A = t_B$) trong 4 mã đề sau:
\begin{enumerate}[label=\alph*)]
    \item \textbf{(Mã đề B2):} Máy A có $CPI_A = 1.3, f_A = 600\text{ MHz}, IC_A = 100\,000\text{ lệnh}$. Máy B có $CPI_B = 2.5, f_B = 750\text{ MHz}$.
    \item \textbf{(Mã đề B3):} Máy A có $CPI_A = 1.2, f_A = 500\text{ MHz}, IC_A = 90\,000\text{ lệnh}$. Máy B có $CPI_B = 2.4, f_B = 600\text{ MHz}$.
    \item \textbf{(Mã đề B6):} Máy A có $CPI_A = 1.4, f_A = 750\text{ MHz}, IC_A = 80\,000\text{ lệnh}$. Máy B có $CPI_B = 2.8, f_B = 1.5\text{ GHz}$.
    \item \textbf{(Mã đề B7):} Máy A có $CPI_A = 2.4, f_A = 800\text{ MHz}, IC_A = 70\,000\text{ lệnh}$. Máy B có $CPI_B = 3.6, f_B = 1\text{ GHz}$.
\end{enumerate}
\textbf{Hướng dẫn giải:}\\
Phương trình cân bằng thời gian thực thi giữa hai máy ($t_A = t_B$):
$$\frac{IC_A \times CPI_A}{f_A} = \frac{IC_B \times CPI_B}{f_B} \iff IC_B = \frac{IC_A \times CPI_A \times f_B}{f_A \times CPI_B}$$
(Lưu ý: Phải đưa $f_A$ và $f_B$ về cùng một đơn vị là MHz trước khi thay số; với $1\text{ GHz} = 1000\text{ MHz}$ và $1.5\text{ GHz} = 1500\text{ MHz}$):
\begin{itemize}
    \item \textbf{Giải câu a (Đề B2):}
    $$IC_B = \frac{100\,000 \times 1.3 \times 750}{600 \times 2.5} = \frac{97\,500\,000}{1500} = 65\,000\text{ lệnh}$$
    \item \textbf{Giải câu b (Đề B3):}
    $$IC_B = \frac{90\,000 \times 1.2 \times 600}{500 \times 2.4} = \frac{64\,800\,000}{1200} = 54\,000\text{ lệnh}$$
    \item \textbf{Giải câu c (Đề B6 -- đổi $f_B = 1.5\text{ GHz} = 1500\text{ MHz}$):}
    $$IC_B = \frac{80\,000 \times 1.4 \times 1500}{750 \times 2.8} = \frac{168\,000\,000}{2100} = 80\,000\text{ lệnh}$$
    \item \textbf{Giải câu d (Đề B7 -- đổi $f_B = 1\text{ GHz} = 1000\text{ MHz}$):}
    $$IC_B = \frac{70\,000 \times 2.4 \times 1000}{800 \times 3.6} = \frac{168\,000\,000}{2880} = \frac{175\,000}{3} \approx 58\,333\text{ lệnh}$$
\end{itemize}

\textbf{Bài tập 1.23 (Đề Cuối kỳ 2021.2 -- Mã đề B8, Câu 1)}\\
Giả sử có 3 bộ xử lý P1, P2, P3 có cùng kiến trúc tập lệnh (ISA). Tần số xung nhịp và $CPI$ của mỗi bộ xử lý được cho trong bảng sau:
\begin{center}
\begin{tabularx}{0.65\textwidth}{|C{2.5cm}|Z|Z|}
\hline
\textbf{Bộ xử lý} & \textbf{Tần số xung nhịp ($f_0$)} & \textbf{CPI} \\ \hline
P1 & $2\text{ GHz}$ & $1.5$ \\ \hline
P2 & $1.5\text{ GHz}$ & $1.0$ \\ \hline
P3 & $3\text{ GHz}$ & $2.5$ \\ \hline
\end{tabularx}
\end{center}
\begin{enumerate}[label=1.\arabic*.]
    \item Bộ xử lý nào có hiệu suất cao nhất?
    \item Mỗi bộ xử lý thực thi một chương trình trong $10\text{ s}$. Tính số chu kỳ và số lệnh của mỗi bộ xử lý.
    \item Người ta muốn giảm thời gian thực thi đi $30\%$, nhưng điều này lại làm tăng $CPI$ lên $20\%$. Tính tần số xung nhịp mới của mỗi bộ xử lý lúc này.
\end{enumerate}
\textbf{Hướng dẫn giải:}
\begin{itemize}
    \item \textbf{Giải ý 1.1 (Tìm bộ xử lý có hiệu suất cao nhất):}\\
    Vì cả 3 bộ xử lý có cùng kiến trúc tập lệnh (ISA), bộ xử lý nào thực hiện được nhiều lệnh nhất trong $1\text{ giây}$ (tức tỷ số $\dfrac{f_0}{CPI}$ lớn nhất) sẽ có hiệu suất cao nhất:
    \begin{itemize}
        \item P1: $\dfrac{f_1}{CPI_1} = \dfrac{2\text{ GHz}}{1.5} \approx 1.33 \times 10^9\text{ lệnh/giây}$ ($1333\text{ MIPS}$).
        \item P2: $\dfrac{f_2}{CPI_2} = \dfrac{1.5\text{ GHz}}{1.0} = 1.50 \times 10^9\text{ lệnh/giây}$ ($1500\text{ MIPS}$).
        \item P3: $\dfrac{f_3}{CPI_3} = \dfrac{3\text{ GHz}}{2.5} = 1.20 \times 10^9\text{ lệnh/giây}$ ($1200\text{ MIPS}$).
    \end{itemize}
    Vậy \textbf{bộ xử lý P2 có hiệu suất cao nhất}.

    \item \textbf{Giải ý 1.2 (Tính số chu kỳ $n = t \times f_0$ và số lệnh $IC = \dfrac{n}{CPI}$ trong $t = 10\text{ s}$):}
    \begin{itemize}
        \item Bộ xử lý P1: $n_1 = 10 \times 2 \times 10^9 = 20 \times 10^9\text{ chu kỳ}$; \quad $IC_1 = \dfrac{20 \times 10^9}{1.5} \approx 1.33 \times 10^{10}\text{ lệnh}$.
        \item Bộ xử lý P2: $n_2 = 10 \times 1.5 \times 10^9 = 15 \times 10^9\text{ chu kỳ}$; \quad $IC_2 = \dfrac{15 \times 10^9}{1.0} = 1.50 \times 10^{10}\text{ lệnh}$.
        \item Bộ xử lý P3: $n_3 = 10 \times 3 \times 10^9 = 30 \times 10^9\text{ chu kỳ}$; \quad $IC_3 = \dfrac{30 \times 10^9}{2.5} = 1.20 \times 10^{10}\text{ lệnh}$.
    \end{itemize}

    \item \textbf{Giải ý 1.3 (Tính tần số xung nhịp mới $f_{\text{mới}}$):}
    \begin{itemize}
        \item Giảm thời gian đi $30\%$ nghĩa là thời gian mới bằng $70\%$ thời gian cũ: $t_{\text{mới}} = (1 - 0.30) \times t_{\text{cũ}} = 0.7 \times 10\text{ s} = 7\text{ s}$.
        \item Tăng $CPI$ lên $20\%$ nghĩa là $CPI$ mới bằng $120\%$ $CPI$ cũ: $CPI_{\text{mới}} = (1 + 0.20) \times CPI_{\text{cũ}} = 1.2 \times CPI_{\text{cũ}}$.
        \item Do số lệnh $IC$ không đổi, tần số xung nhịp mới liên hệ với tần số cũ theo công thức:
        $$f_{\text{mới}} = \frac{IC \times CPI_{\text{mới}}}{t_{\text{mới}}} = \frac{IC \times 1.2 \times CPI_{\text{cũ}}}{0.7 \times t_{\text{cũ}}} = \frac{1.2}{0.7} \times f_{\text{cũ}} = \frac{12}{7} \times f_{\text{cũ}}$$
        \item Thay số cho từng bộ xử lý:
        \begin{itemize}
            \item $f_{1\text{ mới}} = \dfrac{12}{7} \times 2\text{ GHz} = \dfrac{24}{7} \approx \mathbf{3.43\text{ GHz}}$.
            \item $f_{2\text{ mới}} = \dfrac{12}{7} \times 1.5\text{ GHz} = \dfrac{18}{7} \approx \mathbf{2.57\text{ GHz}}$.
            \item $f_{3\text{ mới}} = \dfrac{12}{7} \times 3\text{ GHz} = \dfrac{36}{7} \approx \mathbf{5.14\text{ GHz}}$.
        \end{itemize}
    \end{itemize}
\end{itemize}
\chapter{HỆ THỐNG MÁY TÍNH}

\section{Phần 1: Lý thuyết trọng tâm}

\subsection{Các thành phần cơ bản của máy tính}

Mọi hệ thống máy tính đều hoạt động dựa trên sự phối hợp của các mô-đun phần cứng chính: Bộ xử lý trung tâm (CPU), Bộ nhớ máy tính và Hệ thống vào-ra, được kết nối với nhau thông qua Liên kết hệ thống (Bus).

\textbf{1. Bộ xử lý trung tâm (CPU -- Central Processing Unit)}
\begin{itemize}
    \item \textbf{Chức năng:} Điều khiển hoạt động của toàn bộ máy tính và trực tiếp xử lý dữ liệu.
    \item \textbf{Nguyên tắc hoạt động:} CPU hoạt động theo chương trình nằm trong bộ nhớ chính và là thành phần có tốc độ xử lý nhanh nhất trong hệ thống.
    \item \textbf{Ba thành phần cơ bản bên trong CPU:}
    \begin{enumerate}
        \item \textbf{Đơn vị điều khiển (CU -- Control Unit):} Giải mã lệnh và phát các tín hiệu điều khiển hoạt động của máy tính theo chương trình đã định sẵn.
        \item \textbf{Đơn vị số học và logic (ALU -- Arithmetic and Logic Unit):} Thực hiện các phép toán số học (cộng, trừ, nhân, chia, so sánh) và các phép toán logic (AND, OR, XOR, NOT, dịch bit).
        \item \textbf{Tập thanh ghi (RF -- Register File):} Gồm các thanh ghi nằm ngay bên trong CPU, có tốc độ truy cập nhanh nhất, chịu trách nhiệm lưu giữ các thông tin tạm thời phục vụ cho hoạt động tức thời của CPU.
    \end{enumerate}
\end{itemize}

\textbf{2. Bộ nhớ máy tính (Memory System)}
\begin{itemize}
    \item \textbf{Chức năng và thao tác:} Lưu trữ chương trình và dữ liệu dưới dạng nhị phân. Có hai thao tác cơ bản với bộ nhớ là thao tác \textbf{Ghi (Write)} và thao tác \textbf{Đọc (Read)}.
    \item \textbf{Phân biệt Bộ nhớ trong và Bộ nhớ ngoài:}
    \begin{itemize}
        \item \textit{Bộ nhớ trong:} Gồm \textbf{Bộ nhớ chính} (RAM, ROM) và \textbf{Bộ nhớ đệm} (Cache).
        \item \textit{Bộ nhớ ngoài (Thiết bị lưu trữ -- Storage):} Gồm ổ đĩa cứng từ tính (HDD), ổ thể rắn bán dẫn (SSD), ổ nhớ Flash, thẻ nhớ, đĩa quang (CD, DVD).
    \end{itemize}
    \item \textbf{Đặc điểm của 3 thành phần bộ nhớ chính:}
    \begin{enumerate}
        \item \textbf{Bộ nhớ chính (Main Memory):} Tồn tại trên mọi máy tính, chứa các lệnh và dữ liệu của chương trình đang được thực hiện. Bộ nhớ chính sử dụng bộ nhớ bán dẫn, được tổ chức thành các \textbf{ngăn nhớ} được đánh địa chỉ (thường đánh địa chỉ theo từng byte).
        \begin{itemize}
            \item \textit{Tính chất quan trọng của ngăn nhớ:} Mỗi ngăn nhớ trong máy tính có \textbf{nội dung} (giá trị dữ liệu) có thể thay đổi, nhưng \textbf{địa chỉ vật lý} của ngăn nhớ đó luôn cố định. CPU muốn đọc hoặc ghi một ngăn nhớ thì bắt buộc phải biết địa chỉ của ngăn nhớ đó.
        \end{itemize}
        \item \textbf{Bộ nhớ đệm (Cache Memory):} Là bộ nhớ bán dẫn có tốc độ nhanh được đặt đệm giữa CPU và bộ nhớ chính nhằm \textbf{tăng tốc độ CPU truy cập bộ nhớ}. Cache có dung lượng nhỏ hơn bộ nhớ chính, thường được chia thành các mức (L1, L2, L3) và tích hợp trên cùng chip bộ xử lý.
        \item \textbf{Thiết bị lưu trữ (Storage / Bộ nhớ ngoài):} Dùng để lưu giữ lâu dài tài nguyên phần mềm và dữ liệu; có dung lượng rất lớn nhưng tốc độ chậm. Về mặt liên kết hệ thống, thiết bị lưu trữ được kết nối với máy tính dưới dạng các \textbf{thiết bị vào-ra}.
    \end{enumerate}
\end{itemize}

\textbf{3. Hệ thống vào-ra (Input/Output System)}
\begin{itemize}
    \item \textbf{Chức năng và thao tác:} Trao đổi thông tin giữa máy tính với thế giới bên ngoài thông qua hai thao tác cơ bản là \textbf{Vào dữ liệu (Input)} và \textbf{Ra dữ liệu (Output)}.
    \item \textbf{Hai thành phần chính của hệ thống vào-ra:}
    \begin{enumerate}
        \item \textbf{Các thiết bị vào-ra (I/O Devices / Thiết bị ngoại vi -- Peripherals):} Làm nhiệm vụ chuyển đổi dữ liệu giữa bên trong và bên ngoài máy tính. Gồm 4 nhóm: Thiết bị vào (bàn phím, chuột), Thiết bị ra (màn hình, máy in), Thiết bị lưu trữ (HDD, SSD), và Thiết bị truyền thông (card mạng, modem). \textit{Lưu ý:} Các thiết bị ngoại vi (như bàn phím, màn hình) \textbf{không nối trực tiếp} với CPU hay Bus hệ thống mà phải kết nối thông qua Mô-đun vào-ra.
        \item \textbf{Các mô-đun vào-ra (I/O Modules):} Làm nhiệm vụ nối ghép các thiết bị vào-ra với máy tính (qua Bus hệ thống). Mỗi mô-đun vào-ra có một hoặc một vài \textbf{cổng vào-ra (I/O Ports)}, mỗi cổng được đánh một địa chỉ xác định. CPU muốn trao đổi dữ liệu với thiết bị ngoại vi nào thì phải thông qua địa chỉ của cổng vào-ra tương ứng.
    \end{enumerate}
\end{itemize}

\begin{center}
\begin{center}
\begin{tikzpicture}[
    box/.style={draw=black, thick, minimum width=3.4cm, minimum height=1.3cm, align=center},
    smallbox/.style={draw=black, thick, minimum width=2.8cm, minimum height=0.85cm, align=center},
    arrow/.style={<->, >=Stealth, thick}
]
    % Khối CPU bên trái (cao trải đều từ Bộ nhớ chính xuống Mô-đun vào-ra)
    \node[box, minimum width=3.8cm, minimum height=3.2cm] (cpu) at (0, 0) {\textbf{CPU}\\[4pt]\footnotesize (PC, IR, CU, ALU, RF)};
    
    % Khối Bộ nhớ chính (ở trên) và Mô-đun vào-ra (ở dưới)
    \node[box] (mem) at (5.6, 0.95) {\textbf{Bộ nhớ chính}\\[2pt]\footnotesize (Các ngăn nhớ)};
    \node[box] (iom) at (5.6, -0.95) {\textbf{Mô-đun vào-ra}\\[2pt]\footnotesize (Các cổng vào-ra)};
    
    % Các thiết bị ngoại vi bên phải Mô-đun vào-ra
    \node[smallbox] (dev1) at (9.8, -0.45) {Thiết bị ngoại vi 1};
    \node[smallbox] (dev2) at (9.8, -1.45) {Thiết bị ngoại vi 2};

    % Thanh Bus hệ thống nằm chính giữa CPU và (Bộ nhớ / Mô-đun vào-ra)
    \coordinate (bus_top) at (2.8, 1.85);
    \coordinate (bus_bot) at (2.8, -1.85);
    \draw[thick, double distance=3.5pt] (bus_top) -- (bus_bot);
    \node[above=3pt, font=\footnotesize\bfseries] at (bus_top) {Bus hệ thống};
    
    % Các mũi tên kết nối nằm ngang hoàn toàn (dùng phép chiếu vuông góc |-)
    \draw[arrow] (cpu.east) -- ($(bus_top |- cpu.east) + (-0.08, 0)$);
    \draw[arrow] ($(bus_top |- mem.west) + (0.08, 0)$) -- (mem.west);
    \draw[arrow] ($(bus_top |- iom.west) + (0.08, 0)$) -- (iom.west);
    
    % Mũi tên nối Mô-đun vào-ra với Thiết bị ngoại vi
    \draw[arrow] (iom.east |- dev1.west) -- (dev1.west);
    \draw[arrow] (iom.east |- dev2.west) -- (dev2.west);
\end{tikzpicture}
\end{center}
\end{center}

\subsection{Ba hoạt động cơ bản của máy tính}

\textbf{1. Hoạt động thực hiện chương trình (Chu trình lệnh)}
\begin{itemize}
    \item Máy tính thực hiện chương trình bằng cách lặp đi lặp lại \textbf{chu trình lệnh} gồm hai bước chính: \textbf{Tìm nạp lệnh (Fetch instruction)} và \textbf{Thực hiện lệnh (Execute instruction)}. Hoạt động này chỉ dừng lại khi gặp lệnh thực hiện bị lỗi, gặp lệnh dừng hoặc tắt máy.
    \item \textbf{Quá trình Tìm nạp lệnh (Fetch):}
    \begin{itemize}
        \item \textbf{Bộ đếm chương trình PC (Program Counter):} Là thanh ghi của CPU luôn chứa \textbf{địa chỉ của lệnh tiếp theo} sẽ được nạp vào để thực hiện (trên kiến trúc Intel x86, thanh ghi này được gọi là \textbf{IP -- Instruction Pointer}).
        \item Bắt đầu mỗi chu trình lệnh, CPU phát địa chỉ từ thanh ghi \textbf{PC} ra bộ nhớ chính để tìm ngăn nhớ chứa lệnh.
        \item Lệnh đọc từ bộ nhớ chính được nạp vào \textbf{Thanh ghi lệnh IR (Instruction Register)} của CPU.
        \item Sau khi lệnh được nạp vào IR, nếu chương trình thực hiện tuần tự, CPU sẽ \textbf{tự động tăng nội dung thanh ghi PC} để trỏ sang lệnh kế tiếp (với tập lệnh 32-bit cố định của RISC-V, mỗi lệnh chiếm $4\text{ byte}$ nên $PC \leftarrow PC + 4$).
    \end{itemize}
    \item \textbf{Quá trình Giải mã và Thực hiện lệnh (Decode \& Execute):}
    \begin{itemize}
        \item Ngay sau khi lấy lệnh vào thanh ghi \textbf{IR}, CPU tiến hành \textbf{giải mã lệnh} đang nằm trong thanh ghi \textbf{IR} và phát các tín hiệu điều khiển thực hiện thao tác mà lệnh yêu cầu.
        \item \textbf{Ba kiểu thao tác cơ bản của lệnh:}
        \begin{enumerate}
            \item Trao đổi dữ liệu giữa CPU với bộ nhớ chính hoặc với mô-đun vào-ra.
            \item Thực hiện các phép toán số học hoặc phép toán logic với các dữ liệu (trên ALU).
            \item Chuyển điều khiển trong chương trình (rẽ nhánh hoặc nhảy đến vị trí khác, bản chất là thay đổi giá trị của thanh ghi PC bằng địa chỉ đích).
        \end{enumerate}
    \end{itemize}
\end{itemize}

\textbf{2. Hoạt động Ngắt (Interrupt)}
\begin{itemize}
    \item \textbf{Định nghĩa:} Ngắt là cơ chế cho phép CPU \textbf{tạm dừng chương trình đang thực hiện} để chuyển sang thực hiện một chương trình con có sẵn trong bộ nhớ gọi là \textbf{chương trình con xử lý ngắt (Interrupt Handler)}.
    \item \textbf{Các loại nguyên nhân gây ngắt:}
    \begin{itemize}
        \item Ngắt do lỗi khi thực hiện lệnh (ví dụ: tràn số học, chia cho 0, mã lệnh sai).
        \item Ngắt do lỗi phần cứng.
        \item Ngắt do bộ định thời (Timer).
        \item Ngắt do bộ điều khiển vào-ra gửi tín hiệu ngắt đến CPU để yêu cầu trao đổi dữ liệu.
    \end{itemize}
    \item \textbf{Quy trình CPU xử lý ngắt từ mô-đun vào-ra:}
    \begin{enumerate}
        \item Sau khi hoàn thành mỗi một lệnh máy, CPU kiểm tra tín hiệu ngắt. Nếu không có ngắt, CPU tìm nạp lệnh tiếp theo.
        \item Nếu có tín hiệu ngắt, CPU \textbf{tạm dừng (suspend)} chương trình đang thực hiện.
        \item \textbf{Cất ngữ cảnh} (lưu địa chỉ trở về trên PC và các thông tin liên quan đến chương trình bị ngắt vào ngăn xếp).
        \item Thiết lập bộ đếm chương trình \textbf{PC} trỏ đến địa chỉ bắt đầu của chương trình con xử lý ngắt tương ứng.
        \item Thực hiện chương trình con xử lý ngắt, sau đó \textbf{khôi phục ngữ cảnh} và trở về tiếp tục thực hiện chương trình đang bị tạm dừng.
    \end{enumerate}
    \item \textbf{Xử lý khi có nhiều yêu cầu ngắt:} Gồm hai phương pháp là \textbf{Xử lý ngắt tuần tự} (thực hiện lần lượt hết chương trình xử lý ngắt X rồi mới thực hiện chương trình xử lý ngắt Y) và \textbf{Xử lý ngắt ưu tiên} (chương trình xử lý ngắt Y có mức ưu tiên cao hơn được phép ngắt tạm thời chương trình xử lý ngắt X đang chạy).
\end{itemize}

\textbf{3. Hoạt động Vào-Ra (Input/Output Operations)}
\begin{itemize}
    \item Là hoạt động trao đổi dữ liệu giữa mô-đun vào-ra với bên trong máy tính.
    \item \textbf{Hai kiểu hoạt động vào-ra cơ bản:}
    \begin{enumerate}
        \item CPU trực tiếp trao đổi dữ liệu với mô-đun vào-ra thông qua các lệnh vào-ra trong chương trình (Vào-ra bằng chương trình hoặc Vào-ra điều khiển bằng ngắt).
        \item CPU trao quyền điều khiển hệ thống bus cho phép mô-đun vào-ra trao đổi dữ liệu trực tiếp với bộ nhớ chính mà không đi qua CPU (\textbf{DMA -- Direct Memory Access}).
    \end{enumerate}
\end{itemize}

\subsection{Liên kết trong máy tính (Hệ thống Bus và Liên kết điểm-điểm)}

\textbf{1. Luồng thông tin kết nối tới từng mô-đun}
\begin{itemize}
    \item \textbf{Mô-đun nhớ:} Nhận địa chỉ để xác định ngăn nhớ; nhận dữ liệu khi ghi; đưa dữ liệu hoặc lệnh ra khi đọc (bộ nhớ không phân biệt lệnh và dữ liệu); nhận tín hiệu điều khiển Đọc (Read) và Ghi (Write).
    \item \textbf{Mô-đun vào-ra:} Nhận địa chỉ để xác định cổng vào-ra; chuyển dữ liệu hai chiều giữa bên trong và bên ngoài; nhận tín hiệu điều khiển Đọc/Ghi từ CPU; phát tín hiệu điều khiển đến thiết bị ngoại vi và phát tín hiệu yêu cầu ngắt đến CPU.
    \item \textbf{CPU:} Phát địa chỉ đến mô-đun nhớ hoặc mô-đun vào-ra; đọc lệnh từ bộ nhớ; đọc/ghi dữ liệu với bộ nhớ hoặc mô-đun vào-ra; phát tín hiệu điều khiển đọc/ghi và nhận tín hiệu yêu cầu ngắt.
\end{itemize}

\textbf{2. Liên kết Bus trong máy tính}\\
\textbf{Bus} là tập hợp các đường dây kết nối dùng để vận chuyển thông tin giữa các mô-đun của máy tính với nhau. \textbf{Độ rộng bus} là số đường dây của bus có thể truyền các bit thông tin đồng thời (khái niệm độ rộng bus \textbf{chỉ dùng cho Bus địa chỉ và Bus dữ liệu}, không cố định cho Bus điều khiển).

\begin{tcolorbox}[bwbox, title=Bảng tổng hợp 3 loại Bus chức năng trong Bus hệ thống]
\begin{center}
\renewcommand{\arraystretch}{1.25}
\begin{tabularx}{\linewidth}{|L{3.2cm}|L{2.8cm}|Y|}
\hline
\textbf{Loại Bus} & \textbf{Hướng truyền} & \textbf{Chức năng và Công thức tính toán} \\ \hline
\textbf{Bus địa chỉ} \newline (Address Bus) & \textbf{Một chiều} \newline (Từ CPU/DMA đến Bộ nhớ và I/O) & Vận chuyển địa chỉ để xác định vị trí ngăn nhớ hay cổng vào-ra. \newline Nếu có độ rộng $N\text{ bit}$ ($A_{N-1} \dots A_0$), số địa chỉ tối đa (không gian địa chỉ) là: \textbf{$2^N\text{ địa chỉ}$} (từ $00\dots0_{(2)}$ đến $11\dots1_{(2)}$). \\ \hline
\textbf{Bus dữ liệu} \newline (Data Bus) & \textbf{Hai chiều} \newline (Giữa CPU, Bộ nhớ và I/O) & Vận chuyển \textbf{lệnh} từ bộ nhớ đến CPU và vận chuyển \textbf{dữ liệu} giữa các thành phần. \newline Nếu có độ rộng $M\text{ bit}$ ($D_{M-1} \dots D_0$, thường là $8, 16, 32, 64\text{ bit}$), mỗi thời điểm truyền đồng thời được $M\text{ bit}$ ($= M/8\text{ byte}$). \\ \hline
\textbf{Bus điều khiển} \newline (Control Bus) & \textbf{Tùy từng đường tín hiệu} & Vận chuyển các tín hiệu điều khiển đọc/ghi, điều khiển ngắt và điều khiển xin nhượng bus. Độ rộng phụ thuộc vào thiết kế cụ thể của từng bộ xử lý (không xác định cố định theo bus dữ liệu). \\ \hline
\end{tabularx}
\end{center}
\end{tcolorbox}

\textbf{3. Các tín hiệu điều khiển điển hình trên Bus điều khiển}
\begin{itemize}
    \item \textbf{Nhóm tín hiệu điều khiển đọc/ghi (do CPU phát ra):}
    \begin{itemize}
        \item \texttt{MEMR} (Memory Read): Điều khiển đọc dữ liệu/lệnh từ ngăn nhớ có địa chỉ xác định đưa lên bus dữ liệu chuyển vào CPU.
        \item \texttt{MEMW} (Memory Write): Điều khiển ghi dữ liệu trên bus dữ liệu (từ CPU) vào ngăn nhớ có địa chỉ xác định.
        \item \texttt{IOR} (I/O Read): Điều khiển đọc dữ liệu từ cổng vào-ra có địa chỉ xác định đưa lên bus dữ liệu chuyển vào CPU.
        \item \texttt{IOW} (I/O Write): Điều khiển ghi dữ liệu trên bus dữ liệu (từ CPU) ra cổng vào-ra có địa chỉ xác định.
    \end{itemize}
    \item \textbf{Nhóm tín hiệu điều khiển ngắt:}
    \begin{itemize}
        \item \texttt{INTR} (Interrupt Request): Từ bộ điều khiển vào-ra gửi đến yêu cầu ngắt CPU để trao đổi vào-ra (tín hiệu này có thể bị che/cấm).
        \item \texttt{INTA} (Interrupt Acknowledge): Từ CPU phát ra báo cho bộ điều khiển vào-ra biết CPU chấp nhận ngắt.
        \item \texttt{NMI} (Non-Maskable Interrupt): Tín hiệu ngắt không thể che gửi đến CPU.
        \item \texttt{Reset}: Tín hiệu từ bên ngoài gửi đến để khởi động lại máy tính.
    \end{itemize}
    \item \textbf{Nhóm tín hiệu điều khiển bus (phục vụ DMA):}
    \begin{itemize}
        \item \texttt{BRQ} (Bus Request): Từ mô-đun vào-ra (DMAC) gửi đến yêu cầu CPU chuyển nhượng quyền sử dụng bus.
        \item \texttt{BGT} (Bus Grant): Từ CPU phát ra chấp nhận chuyển nhượng quyền sử dụng bus.
        \item \texttt{Lock/Unlock}: Tín hiệu cấm hoặc cho phép xin chuyển nhượng bus.
    \end{itemize}
\end{itemize}

\textbf{4. Hạn chế của Bus dùng chung (Shared Bus) và Liên kết điểm-điểm}
\begin{itemize}
    \item \textbf{Hạn chế của Bus dùng chung:} Cần có bộ phân xử bus; tại một thời điểm chỉ phục vụ được một yêu cầu trao đổi dữ liệu nên độ trễ lớn; khó đồng bộ hóa khi tốc độ truyền cao.
    \item \textbf{Giải pháp khắc phục:} Sử dụng cấu trúc \textbf{Đa bus (Multiple bus)} hoặc \textbf{Liên kết điểm-điểm (Point-to-point interconnection)}.
    \item \textbf{Các liên kết điểm-điểm và bus ngoại vi phổ biến:}
    \begin{itemize}
        \item \textbf{QPI (Quick Path Interconnect):} Đa kết nối trực tiếp từng cặp giữa các lõi xử lý (Cores) và Chipset, không cần phân xử bus, truyền dữ liệu dạng gói theo kiến trúc phân lớp.
        \item \textbf{PCI Express (PCIe):} Liên kết điểm-điểm tốc độ cao thay thế cho bus dùng chung PCI truyền thống để nối ghép các thiết bị vào-ra băng thông lớn (card đồ họa, ổ SSD NVMe).
        \item \textbf{SATA:} Chuẩn kết nối chuyên dụng chỉ dùng để kết nối với các ổ lưu trữ vật lý (HDD, SSD SATA).
        \item \textbf{USB (Universal Serial Bus):} Chuẩn bus ngoại vi nối tiếp đa dụng kết nối giữa mô-đun vào-ra và các thiết bị ngoại vi.
    \end{itemize}
\end{itemize}

\newpage
\section{Phần 2: Toàn bộ bài tập Chương 2 và Hướng dẫn giải}

\subsection{Nhóm 1: Cấu trúc CPU, Bộ nhớ và Hệ thống vào-ra}

\textbf{Bài tập 2.1 (Đề Giữa kỳ 2018.2 -- Mã đề 204, Câu 5, Câu 6 \& Câu 13)}\\
Hãy chọn phương án đúng cho 3 câu hỏi lý thuyết cơ bản sau:
\begin{enumerate}[label=\alph*)]
    \item \textbf{(Câu 5):} Đơn vị nào bên trong CPU chịu trách nhiệm lưu giữ các thông tin tạm thời phục vụ cho hoạt động của CPU?\\
    A. Cache \quad B. Bộ nhớ chính \quad C. Tập thanh ghi (RF) \quad D. Biến nhớ
    \item \textbf{(Câu 6):} Loại bộ nhớ nào sau đây \textbf{không} thuộc về bộ nhớ trong?\\
    A. Bộ nhớ ROM \quad B. Bộ nhớ ổ cứng HDD \quad C. Bộ nhớ chính \quad D. Bộ nhớ cache
    \item \textbf{(Câu 13):} Ngăn nhớ nằm ở chỗ nào trong hệ thống máy tính?\\
    A. CPU \quad B. Thiết bị ngoại vi \quad C. Module vào-ra \quad D. Bộ nhớ chính
\end{enumerate}
\textbf{Hướng dẫn giải:}
\begin{itemize}
    \item \textbf{Giải câu a:} Theo Slide Chương 2 (trang 7), bên trong CPU có 3 thành phần là CU, ALU và Tập thanh ghi (RF). Trong đó, \textbf{Tập thanh ghi (RF)} chứa các thông tin tạm thời phục vụ trực tiếp cho hoạt động của CPU $\rightarrow$ Chọn \textbf{C}.
    \item \textbf{Giải câu b:} ROM, RAM (bộ nhớ chính) và Cache đều là bộ nhớ bán dẫn thuộc bộ nhớ trong. \textbf{Ổ cứng HDD} là thiết bị lưu trữ từ tính thuộc \textbf{bộ nhớ ngoài} $\rightarrow$ Chọn \textbf{B}.
    \item \textbf{Giải câu c:} Theo Slide Chương 2 (trang 9), \textbf{Bộ nhớ chính} được tổ chức thành các ngăn nhớ được đánh địa chỉ $\rightarrow$ Chọn \textbf{D}.
\end{itemize}

\textbf{Bài tập 2.2 (Đề Giữa kỳ 2022.2 -- Đề 0009 Câu 11, Đề 0008 Câu 11 \& Ôn tập Cuối kỳ 2025.2 -- Câu 41)}
\begin{enumerate}[label=\alph*)]
    \item \textbf{(Đề 0009 -- Câu 11):} Điền từ còn thiếu vào 2 chỗ trống trong phát biểu sau: ``Mỗi ngăn nhớ trong máy tính có \dots\dots\dots thay đổi được, nhưng \dots\dots\dots của nó luôn cố định.''
    \item \textbf{(Đề 0008 -- Câu 11 \& Ôn tập 2025.2 -- Câu 41):} Loại bộ nhớ giúp CPU tăng tốc độ truy cập bộ nhớ chính có tên là gì? (Hoặc: Tác dụng chính của bộ nhớ Cache là gì?)\\
    A. Giảm chi phí cho hệ thống máy tính \quad B. Tăng dung lượng bộ nhớ chính\\
    C. Tăng tốc độ truy cập đến bộ nhớ chính \quad D. Tăng tốc độ truy cập đến ổ đĩa cứng
\end{enumerate}
\textbf{Hướng dẫn giải:}
\begin{itemize}
    \item \textbf{Giải câu a:} Theo Slide Chương 2 (trang 9), mỗi ngăn nhớ có \textbf{nội dung} (hoặc giá trị) có thể thay đổi khi ghi dữ liệu mới, nhưng \textbf{địa chỉ} (hoặc địa chỉ vật lý) của ngăn nhớ đó luôn cố định.
    \item \textbf{Giải câu b:} Bộ nhớ đặt đệm giữa CPU và bộ nhớ chính để tăng tốc độ CPU truy cập bộ nhớ chính là \textbf{Bộ nhớ đệm (Cache)} $\rightarrow$ Điền \textbf{Cache} (hoặc chọn đáp án \textbf{C}).
\end{itemize}

\textbf{Bài tập 2.3 (Đề Giữa kỳ 2018.2 -- Mã đề 199, Câu 20)}\\
Bàn phím được kết nối trực tiếp với bộ phận chức năng nào trong hệ thống máy tính?
\begin{enumerate}[label=\Alph*.]
    \item Với bộ nhớ chính
    \item Với mô-đun vào-ra
    \item Với CPU
    \item Với bus hệ thống
\end{enumerate}
\textbf{Hướng dẫn giải:}
\begin{itemize}
    \item \textbf{Kiến thức áp dụng:} Theo sơ đồ Hệ thống vào-ra (Slide Chương 2, trang 12--14), bàn phím là một thiết bị vào-ra (thiết bị ngoại vi). Các thiết bị ngoại vi không thể nối trực tiếp vào CPU hay Bus hệ thống mà bắt buộc phải kết nối thông qua các cổng của \textbf{mô-đun vào-ra}.
    \item \textbf{Kết luận:} Đáp án đúng là \textbf{B (Với mô-đun vào-ra)}.
\end{itemize}

\subsection{Nhóm 2: Chu trình lệnh, các thanh ghi PC / IR và Hoạt động Ngắt}

\textbf{Bài tập 2.4 (Đề Giữa kỳ 2018.2 -- Mã đề 198 Câu 17, Mã đề 204 Câu 4 \& Câu 18, Giữa kỳ 2022.2 -- Đề 0008 Câu 10 \& Ôn tập 2025.2 -- Câu 56, 76)}\\
Hãy trả lời chuỗi câu hỏi về các thanh ghi hệ thống trong quá trình Tìm nạp và Giải mã lệnh:
\begin{enumerate}[label=\alph*)]
    \item \textbf{(Đề GK 2022.2 -- Đề 0008 \& Mã đề 204 Câu 4):} Trong chu trình hoạt động cơ bản của máy tính, ở giai đoạn tìm nạp lệnh (nhận lệnh từ bộ nhớ chính), CPU sẽ tác động tới 2 thanh ghi hệ thống nào? Viết tên đầy đủ hoặc viết tắt.
    \item \textbf{(Ôn tập 2025.2 Câu 76 \& Mã đề 198 Câu 17):} Thanh ghi nào trong CPU luôn giữ địa chỉ của lệnh tiếp theo sẽ được lấy từ bộ nhớ và được tự động tăng mỗi khi CPU nhận xong một lệnh?\\
    A. Thanh ghi Stack Pointer (SP) \quad B. Thanh ghi Program Counter (PC)\\
    C. Thanh ghi Instruction Register (IR) \quad D. Thanh ghi Memory Address Register (MAR)
    \item \textbf{(Mã đề 204 -- Câu 18):} Trong quá trình giải mã lệnh, CPU sẽ tiến hành giải mã lệnh đang nằm trên thanh ghi nào?\\
    A. Thanh ghi lệnh IR \quad B. Thanh ghi con trỏ lệnh PC/IP \quad C. Thanh ghi dữ liệu \quad D. Thanh ghi trạng thái
    \item \textbf{(Ôn tập 2025.2 -- Câu 56):} Trong chu trình thực hiện lệnh cơ bản (Instruction Cycle), giai đoạn nào diễn ra ngay sau khi CPU lấy lệnh (Fetch) từ bộ nhớ?\\
    A. Giai đoạn Ghi lại kết quả (Write Back) \quad B. Giai đoạn Thực thi (Execute)\\
    C. Giai đoạn Giải mã (Decode), nơi CPU xác định loại lệnh và các toán hạng \quad D. Giai đoạn Điều khiển ngắt
\end{enumerate}
\textbf{Hướng dẫn giải:}
\begin{itemize}
    \item \textbf{Giải câu a:} Hai thanh ghi phục vụ quá trình tìm nạp lệnh là \textbf{PC (Program Counter -- Bộ đếm chương trình)} và \textbf{IR (Instruction Register -- Thanh ghi lệnh)}.
    \item \textbf{Giải câu b:} Thanh ghi giữ địa chỉ của lệnh tiếp theo và tự động tăng sau mỗi lần nạp lệnh là \textbf{Thanh ghi Bộ đếm chương trình (Program Counter -- PC)} $\rightarrow$ Chọn \textbf{B}.
    \item \textbf{Giải câu c:} Lệnh sau khi đọc từ bộ nhớ lên được chứa trong thanh ghi lệnh \textbf{IR}, do đó Đơn vị điều khiển (CU) sẽ giải mã lệnh đang nằm trên \textbf{thanh ghi lệnh IR} $\rightarrow$ Chọn \textbf{A}.
    \item \textbf{Giải câu d:} Ngay sau khi tìm nạp lệnh (Fetch), CPU phải tiến hành \textbf{Giai đoạn Giải mã (Decode)} để biết lệnh đó yêu cầu thực hiện thao tác gì và lấy toán hạng ở đâu $\rightarrow$ Chọn \textbf{C}.
\end{itemize}

\textbf{Bài tập 2.5 (Đề Giữa kỳ 2018.2 -- Mã đề 198 Câu 14, Mã đề 199 Câu 8 \& Ôn tập 2025.2 -- Câu 15)}
\begin{enumerate}[label=\alph*)]
    \item \textbf{(Mã đề 198 Câu 14):} Đáp án nào mô tả đúng về cơ chế ngắt (Interrupt)?\\
    A. Tạm dừng chương trình đang thực hiện để chuyển đến thực hiện chương trình con khác (chương trình con xử lý ngắt).\\
    B. Kết thúc phần mềm đang chạy.\\
    C. Tạm dừng phần mềm đang hoạt động vĩnh viễn.\\
    D. Hủy bỏ lệnh máy đang thực hiện dở mà không quay lại.
    \item \textbf{(Ôn tập 2025.2 -- Câu 15):} Với phương pháp vào-ra bằng chương trình, thành phần nào của máy tính trực tiếp điều khiển vào-ra?\\
    A. DMAC \quad B. Thiết bị vào-ra \quad C. CPU \quad D. RAM
\end{enumerate}
\textbf{Hướng dẫn giải:}
\begin{itemize}
    \item \textbf{Giải câu a:} Theo định nghĩa tại Slide Chương 2 (trang 22), ngắt là cơ chế cho phép CPU tạm dừng chương trình đang thực hiện để chuyển sang thực hiện chương trình con xử lý ngắt có sẵn trong bộ nhớ, sau đó quay trở về chương trình cũ $\rightarrow$ Chọn \textbf{A}.
    \item \textbf{Giải câu b:} Theo Slide Chương 2 (trang 26), khi vào-ra bằng chương trình, \textbf{CPU} trực tiếp điều khiển trao đổi dữ liệu với mô-đun vào-ra bằng các lệnh trong chương trình $\rightarrow$ Chọn \textbf{C (CPU)}.
\end{itemize}

\subsection{Nhóm 3: Liên kết hệ thống, Các loại Bus và Tính toán không gian địa chỉ}

\begin{tcolorbox}[bwbox, title=Quy trình giải các bài toán về Độ rộng Bus địa chỉ và Dung lượng RAM]
\begin{enumerate}
    \item \textbf{Công thức cơ bản:} Nếu Bus địa chỉ có độ rộng $N\text{ bit}$ và bộ nhớ đánh địa chỉ theo từng byte:
    $$\text{Dung lượng không gian địa chỉ tối đa} = 2^N\text{ byte}$$
    \item \textbf{Đổi nhanh giữa đơn vị dung lượng và lũy thừa của 2:}
    \begin{itemize}
        \item $1\text{ KiB} = 2^{10}\text{ byte} \implies x\text{ KiB} = 2^k \times 2^{10} = 2^{10+k}\text{ byte}$.
        \item $1\text{ MiB} = 2^{20}\text{ byte} \implies x\text{ MiB} = 2^k \times 2^{20} = 2^{20+k}\text{ byte}$.
        \item $1\text{ GiB} = 2^{30}\text{ byte} \implies x\text{ GiB} = 2^k \times 2^{30} = 2^{30+k}\text{ byte}$.
    \end{itemize}
    (Với $2 = 2^1$; $4 = 2^2$; $8 = 2^3$; $16 = 2^4$; $32 = 2^5$; $64 = 2^6$; $128 = 2^7$; $256 = 2^8$; $512 = 2^9$).
    \item \textbf{Bài toán ``Lắp thêm 1 thanh RAM lớn nhất và không thừa'':}
    \begin{itemize}
        \item Bước 1: Tính tổng không gian địa chỉ tối đa mà CPU quản lý được: $S_{\max} = 2^N\text{ (MiB hoặc GiB)}$.
        \item Bước 2: Tính phần dung lượng còn trống: $S_{\text{trống}} = S_{\max} - S_{\text{đã cắm}}$.
        \item Bước 3: Vì dung lượng của \textbf{1 thanh RAM vật lý} luôn phải là một lũy thừa của $2$ ($1, 2, 4, 8, 16, 32, 64\text{ MiB}...$), ta chọn thanh RAM có dung lượng là \textbf{lũy thừa của 2 lớn nhất mà nhỏ hơn hoặc bằng $S_{\text{trống}}$}.
    \end{itemize}
\end{enumerate}
\end{tcolorbox}

\textbf{Bài tập 2.6 (Đề Giữa kỳ 2018.2 -- Mã đề 198 Câu 1 \& Câu 10, Ôn tập 2025.2 -- Câu 13 \& Câu 26)}\\
Hãy chọn phương án đúng cho 4 câu hỏi về đặc điểm của Bus hệ thống:
\begin{enumerate}[label=\alph*)]
    \item \textbf{(Mã đề 198 -- Câu 1):} Một CPU có bus dữ liệu và bus địa chỉ có độ rộng đều bằng $32\text{ bit}$. Hỏi độ rộng của bus điều khiển bằng bao nhiêu?\\
    A. $64\text{ bit}$ \quad B. $2^{32}\text{ bit}$ \quad C. Không xác định \quad D. $32\text{ bit}$
    \item \textbf{(Mã đề 198 -- Câu 10):} Loại bus nào vận chuyển nội dung lệnh từ bộ nhớ tới CPU?\\
    A. Bus điều khiển \quad B. Bus vào-ra \quad C. Bus dữ liệu \quad D. Bus địa chỉ
    \item \textbf{(Ôn tập 2025.2 -- Câu 13):} Bus dữ liệu (Data Bus) trong kiến trúc máy tính chịu trách nhiệm truyền tải loại thông tin nào?\\
    A. Chỉ các dữ liệu giữa các thành phần trong hệ thống máy tính.\\
    B. Các lệnh và dữ liệu giữa các thành phần trong hệ thống máy tính.\\
    C. Địa chỉ bộ nhớ của dữ liệu cần truy cập.\\
    D. Tín hiệu điều khiển để đồng bộ hóa hoạt động.
    \item \textbf{(Ôn tập 2025.2 -- Câu 26):} Bus Địa chỉ (Address Bus) có đặc điểm nào sau đây, đặc biệt là về hướng truyền?\\
    A. Là bus một chiều (unidirectional), truyền địa chỉ từ CPU/DMA đến bộ nhớ/I/O.\\
    B. Là bus hai chiều (bidirectional), truyền cả địa chỉ và dữ liệu.\\
    C. Là bus hai chiều, truyền tín hiệu điều khiển và dữ liệu.\\
    D. Là bus một chiều, truyền tín hiệu điều khiển từ bộ nhớ đến CPU.
\end{enumerate}
\textbf{Hướng dẫn giải:}
\begin{itemize}
    \item \textbf{Giải câu a:} Theo Slide Chương 2 (trang 31), khái niệm độ rộng bus chỉ áp dụng cho bus địa chỉ và bus dữ liệu. Số lượng đường dây điều khiển phụ thuộc vào thiết kế riêng của từng hệ thống chứ không phụ thuộc vào độ rộng bus dữ liệu hay địa chỉ $\rightarrow$ Chọn \textbf{C (Không xác định)}.
    \item \textbf{Giải câu b và câu c:} Theo Slide Chương 2 (trang 34), Bus dữ liệu có 2 chức năng: (1) Vận chuyển lệnh từ bộ nhớ đến CPU, và (2) Vận chuyển dữ liệu giữa các thành phần của máy tính với nhau $\rightarrow$ Câu b chọn \textbf{C (Bus dữ liệu)}; Câu c chọn \textbf{B (Các lệnh và dữ liệu)}.
    \item \textbf{Giải câu d:} Chỉ có CPU (hoặc DMAC) phát địa chỉ ra để định vị ngăn nhớ hoặc cổng vào-ra, do đó Bus địa chỉ là bus một chiều $\rightarrow$ Chọn \textbf{A}.
\end{itemize}

\textbf{Bài tập 2.7 (Đề Giữa kỳ 2019.2 -- Mã đề 243 Câu 5, 6, 13 \& Mã đề 244 Câu 7, 9, 13)}
\begin{enumerate}[label=\alph*)]
    \item \textbf{(Mã đề 243 Câu 6 \& Mã đề 244 Câu 9):} Với lệnh hợp ngữ ghi bộ nhớ \texttt{sw t2, 4(t4)} (hoặc \texttt{sw s2, 8(s4)}), giá trị tổng \texttt{4 + t4} (hoặc \texttt{8 + s4}) sẽ được truyền đi từ CPU tới bộ nhớ chính trên đường bus nào?\\
    A. Bus dữ liệu \quad B. Bus địa chỉ \quad C. Bus điều khiển \quad D. Bus ngoại vi
    \item \textbf{(Mã đề 243 Câu 5 \& Mã đề 244 Câu 7):} Loại bus ngoại vi nào dưới đây \textbf{chỉ dùng} để kết nối với ổ lưu trữ vật lý (HDD, SSD)?\\
    A. HDMI \quad B. PCIe \quad C. SATA \quad D. USB
    \item \textbf{(Mã đề 243 Câu 13 \& Mã đề 244 Câu 13):} Trong kiến trúc máy tính, khái niệm USB nói về điều gì?\\
    A. Ổ lưu trữ \quad B. Bus ngoại vi giữa mô-đun I/O và thiết bị \quad C. Khe cắm mở rộng RAM \quad D. Nhà sản xuất linh kiện
\end{enumerate}
\textbf{Hướng dẫn giải:}
\begin{itemize}
    \item \textbf{Giải câu a:} Trong lệnh \texttt{sw t2, 4(t4)}, giá trị \texttt{4 + t4} là \textbf{địa chỉ hiệu dụng} của ngăn nhớ cần ghi dữ liệu, nên giá trị này được CPU phát ra trên \textbf{Bus địa chỉ} (còn nội dung của thanh ghi \texttt{t2} mới truyền trên Bus dữ liệu) $\rightarrow$ Chọn \textbf{B (Bus địa chỉ)}.
    \item \textbf{Giải câu b:} Chuẩn \textbf{SATA} (Serial ATA) là chuẩn kết nối chuyên dụng chỉ dùng cho các ổ đĩa lưu trữ (ổ cứng HDD, SSD SATA, ổ đĩa quang) $\rightarrow$ Chọn \textbf{C (SATA)}.
    \item \textbf{Giải câu c:} USB (Universal Serial Bus) là chuẩn \textbf{bus ngoại vi giữa mô-đun vào-ra và thiết bị} $\rightarrow$ Chọn \textbf{B}.
\end{itemize}

\textbf{Bài tập 2.8 (Tổng hợp các câu tính độ rộng Bus địa chỉ từ Đề Giữa kỳ 2019.2, 2022.2 \& Cuối kỳ 2020.2, 2022.2, 2025.2)}\\
Biết bộ nhớ chính được đánh địa chỉ theo từng byte. Hãy thực hiện các tính toán sau:
\begin{enumerate}[label=\alph*)]
    \item \textbf{(Đề GK 2022.2 -- Đề 0008, Câu 3):} Máy tính có bus dữ liệu $32\text{ bit}$ và bus địa chỉ $23\text{ bit}$. Hỏi không gian địa chỉ bộ nhớ có kích thước bao nhiêu byte (hoặc MiB)?
    \item \textbf{(Đề GK 2022.2 -- Đề 0006, Câu 3):} Nếu \texttt{0x3FF} là địa chỉ bộ nhớ lớn nhất mà CPU có thể quản lý được, thì bus địa chỉ có độ rộng bằng bao nhiêu bit?
    \item \textbf{(Đề CK 2020.2 \& Ôn tập 2025.2 Câu 11):} Không gian địa chỉ bộ nhớ là $2\text{ GiB}$, độ rộng bus địa chỉ là bao nhiêu bit?
    \item \textbf{(Đề GK 2019.2 -- Mã đề 243 Câu 16 \& Mã đề 244 Câu 4):} Không gian địa chỉ bộ nhớ là $64\text{ GiB}$ (Mã đề 243) hoặc $128\text{ GiB}$ (Mã đề 244), độ rộng bus địa chỉ là bao nhiêu bit?
    \item \textbf{(Đề CK 2022.2 -- \#bb55ef \& CK 2024.2):} Một hệ thống máy tính gồm có $4\text{ MiB}$ bộ nhớ cache, $8\text{ GiB}$ bộ nhớ chính, $1\text{ TiB}$ bộ nhớ thể rắn SSD. Hỏi bus địa chỉ của hệ thống cần có độ rộng tối thiểu bao nhiêu bit?
\end{enumerate}
\textbf{Hướng dẫn giải:}
\begin{itemize}
    \item \textbf{Giải câu a:} Không gian địa chỉ chỉ phụ thuộc vào độ rộng bus địa chỉ ($N = 23\text{ bit}$):
    $$\text{Kích thước} = 2^{23}\text{ byte} = 2^3 \times 2^{20}\text{ byte} = \mathbf{8\text{ MiB}}\text{ (hoặc } 8\,388\,608\text{ byte)}$$
    \item \textbf{Giải câu b:} Đổi địa chỉ lớn nhất \texttt{0x3FF} từ hệ Hexa sang hệ nhị phân:
    $$\texttt{0x3FF} = 0011\,1111\,1111_{(2)} \rightarrow \text{gồm } 10\text{ bit } 1\text{ liên tiếp } (= 2^{10} - 1)$$
    Địa chỉ chạy từ \texttt{0x000} đến \texttt{0x3FF} gồm $2^{10} = 1024$ địa chỉ $\implies$ Độ rộng bus địa chỉ là \textbf{10 bit}.
    \item \textbf{Giải câu c:} Ta có $2\text{ GiB} = 2^1 \times 2^{30}\text{ byte} = 2^{31}\text{ byte} \implies$ Độ rộng bus địa chỉ là \textbf{31 bit}.
    \item \textbf{Giải câu d:}
    \begin{itemize}
        \item Với $64\text{ GiB} = 2^6 \times 2^{30}\text{ byte} = 2^{36}\text{ byte} \implies$ Độ rộng bus địa chỉ là \textbf{36 bit}.
        \item Với $128\text{ GiB} = 2^7 \times 2^{30}\text{ byte} = 2^{37}\text{ byte} \implies$ Độ rộng bus địa chỉ là \textbf{37 bit}.
    \end{itemize}
    \item \textbf{Giải câu e:}
    \begin{itemize}
        \item Bộ nhớ Cache tự động ánh xạ từ bộ nhớ chính; ổ cứng SSD là thiết bị lưu trữ ngoài giao tiếp qua mô-đun vào-ra theo khối chứ không đánh địa chỉ từng byte trên bus địa chỉ CPU.
        \item Bus địa chỉ hệ thống chỉ cần định vị trực tiếp cho $8\text{ GiB}$ bộ nhớ chính (RAM):
        $$8\text{ GiB} = 2^3 \times 2^{30}\text{ byte} = 2^{33}\text{ byte} \implies \text{Độ rộng tối thiểu là }\mathbf{33\text{ bit}}$$
    \end{itemize}
\end{itemize}

\textbf{Bài tập 2.9 (Đề Giữa kỳ 2018.2 -- Mã đề 204 Câu 16 \& Đề Giữa kỳ 2019.2 -- Mã đề 243 Câu 14, Mã đề 244 Câu 6)}
\begin{enumerate}[label=\alph*)]
    \item \textbf{(Mã đề 204 -- Câu 16):} Một hệ thống máy tính cần sử dụng lượng bộ nhớ chính lên tới $128\text{ KiB}$, cổng vào-ra $64\text{ KiB}$ (dùng chung không gian địa chỉ), và ổ cứng SSD $256\text{ KiB}$. Hỏi bus địa chỉ phải có độ rộng tối thiểu bằng bao nhiêu bit?\\
    A. $17\text{ bit}$ \quad B. $18\text{ bit}$ \quad C. $19\text{ bit}$ \quad D. $20\text{ bit}$
    \item \textbf{(Mã đề 243 -- Câu 14):} Bus hệ thống của một máy tính có đường bus địa chỉ $24\text{ bit}$. Biết rằng máy tính đã cắm 1 thanh RAM $8\text{ MiB}$. Hỏi máy tính có thể lắp thêm \textbf{1 thanh RAM} kích thước \textbf{lớn nhất} là bao nhiêu và \textbf{không thừa}? (Giả định máy tính không có cổng vào-ra).\\
    A. $2\text{ MiB}$ \quad B. $4\text{ MiB}$ \quad C. $16\text{ MiB}$ \quad D. $8\text{ MiB}$
    \item \textbf{(Mã đề 244 -- Câu 6):} Tương tự câu b, nhưng máy tính có đường bus địa chỉ $26\text{ bit}$ và đã cắm 1 thanh RAM $8\text{ MiB}$. Hỏi có thể lắp thêm \textbf{1 thanh RAM} kích thước lớn nhất là bao nhiêu và không thừa?\\
    A. $8\text{ MiB}$ \quad B. $64\text{ MiB}$ \quad C. $32\text{ MiB}$ \quad D. $16\text{ MiB}$
\end{enumerate}
\textbf{Hướng dẫn giải:}
\begin{itemize}
    \item \textbf{Giải câu a:}
    \begin{itemize}
        \item Ổ cứng SSD là bộ nhớ ngoài nên không tính vào không gian địa chỉ trực tiếp của CPU.
        \item Tổng không gian địa chỉ cần đánh cho Bộ nhớ chính và Cổng vào-ra là: $128\text{ KiB} + 64\text{ KiB} = 192\text{ KiB}$.
        \item Với $17\text{ bit}$ chỉ đánh địa chỉ được $2^{17} = 128\text{ KiB}$ (thiếu). Để đánh địa chỉ cho $192\text{ KiB}$ (nằm giữa $128\text{ KiB}$ và $256\text{ KiB} = 2^{18}$), độ rộng bus địa chỉ tối thiểu phải là \textbf{18 bit} $\rightarrow$ Chọn \textbf{B (18 bit)}.
    \end{itemize}
    \item \textbf{Giải câu b (Bus địa chỉ 24 bit):}
    \begin{itemize}
        \item Bước 1: Tổng không gian bộ nhớ tối đa mà $24\text{ bit}$ quản lý được là: $2^{24}\text{ byte} = 2^4 \times 2^{20}\text{ byte} = 16\text{ MiB}$.
        \item Bước 2: Máy đã cắm $8\text{ MiB}$, vậy dung lượng địa chỉ còn trống là: $16\text{ MiB} - 8\text{ MiB} = 8\text{ MiB}$.
        \item Bước 3: Lắp thêm 1 thanh RAM có kích thước \textbf{8 MiB} là vừa khít và không thừa $\rightarrow$ Chọn \textbf{D (8 MiB)}.
    \end{itemize}
    \item \textbf{Giải câu c (Bus địa chỉ 26 bit):}
    \begin{itemize}
        \item Bước 1: Tổng không gian bộ nhớ tối đa mà $26\text{ bit}$ quản lý được là: $2^{26}\text{ byte} = 2^6 \times 2^{20}\text{ byte} = 64\text{ MiB}$.
        \item Bước 2: Máy đã cắm $8\text{ MiB}$, vậy phần không gian địa chỉ còn trống là: $64\text{ MiB} - 8\text{ MiB} = 56\text{ MiB}$.
        \item Bước 3: Vì dung lượng của \textbf{1 thanh RAM} thương mại phải là lũy thừa của $2$ ($8, 16, 32, 64\text{ MiB}$): nếu lắp thanh $64\text{ MiB}$ thì tổng là $8 + 64 = 72\text{ MiB} > 64\text{ MiB}$ (bị thừa $8\text{ MiB}$ không nhận diện được). Do đó thanh RAM đơn lẻ lớn nhất lắp vào mà không bị thừa là thanh \textbf{32 MiB} (tổng $8 + 32 = 40\text{ MiB} \le 64\text{ MiB}$) $\rightarrow$ Chọn \textbf{C (32 MiB)}.
    \end{itemize}
\end{itemize}

\textbf{Bài tập 2.10 (Đề Cuối kỳ 2021.2 -- Mã đề B3, Câu 4.1 \& 4.2)}\\
Máy tính dùng $36\text{ bit}$ địa chỉ để đánh địa chỉ cho bộ nhớ theo byte, $16\text{ bit}$ để đánh địa chỉ cho cổng vào-ra riêng biệt. Hãy cho biết:
\begin{enumerate}[label=4.\arabic*.]
    \item Số byte nhớ tối đa được đánh địa chỉ? Địa chỉ đầu và địa chỉ cuối dưới dạng Hexa?
    \item Số cổng vào-ra tối đa được đánh địa chỉ vào-ra riêng biệt? Địa chỉ đầu và địa chỉ cuối dưới dạng Hexa?
\end{enumerate}
\textbf{Hướng dẫn giải:}
\begin{itemize}
    \item \textbf{Giải ý 4.1 (Không gian địa chỉ bộ nhớ $36\text{ bit}$):}
    \begin{itemize}
        \item Số byte nhớ tối đa được đánh địa chỉ: $2^{36}\text{ byte} = 2^6 \times 2^{30}\text{ byte} = \mathbf{64\text{ GiB}}$.
        \item Vì mỗi chữ số Hexa biểu diễn cho $4\text{ bit}$ nhị phân nên $36\text{ bit}$ ứng với $36 \div 4 = 9\text{ chữ số Hexa}$:
        \begin{itemize}
            \item Địa chỉ đầu (tất cả $36\text{ bit}$ bằng 0): \texttt{0x000000000} (9 chữ số 0).
            \item Địa chỉ cuối (tất cả $36\text{ bit}$ bằng 1): \texttt{0xFFFFFFFFF} (9 chữ số F).
        \end{itemize}
    \end{itemize}
    \item \textbf{Giải ý 4.2 (Không gian địa chỉ cổng vào-ra $16\text{ bit}$):}
    \begin{itemize}
        \item Số cổng vào-ra tối đa: $2^{16} = \mathbf{65\,536\text{ cổng}}$ (hay $64\text{K cổng}$).
        \item Vì $16\text{ bit}$ ứng với $16 \div 4 = 4\text{ chữ số Hexa}$:
        \begin{itemize}
            \item Địa chỉ cổng đầu: \texttt{0x0000}.
            \item Địa chỉ cổng cuối: \texttt{0xFFFF}.
        \end{itemize}
    \end{itemize}
\end{itemize}
\chapter{KIẾN TRÚC TẬP LỆNH RISC-V}

\section{Phần 1: Lý thuyết trọng tâm}

\subsection{Tổng quan về Kiến trúc tập lệnh và Tập 32 thanh ghi}

\textbf{1. Đặc trưng của kiến trúc tập lệnh thu gọn RISC và RISC-V}
\begin{itemize}
    \item \textbf{Phân biệt CISC và RISC:}
    \begin{itemize}
        \item \textit{CISC (Complex Instruction Set Computer):} Máy tính với tập lệnh phức tạp (điển hình là Intel x86, Motorola 680x0). Độ dài lệnh thay đổi từ $1$ đến $15\text{ byte}$, cho phép lệnh tính toán trực tiếp với toán hạng nằm trong bộ nhớ.
        \item \textit{RISC (Reduced Instruction Set Computer):} Máy tính với tập lệnh thu gọn (điển hình là RISC-V, MIPS, ARM).
    \end{itemize}
    \item \textbf{Các đặc trưng cốt lõi của kiến trúc RISC (và RISC-V 32-bit):}
    \begin{enumerate}
        \item Các lệnh chuẩn có độ dài cố định bằng $32\text{ bit}$ ($= 4\text{ byte}$). Do đó, địa chỉ của mọi lệnh máy chuẩn trong bộ nhớ luôn là \textbf{bội số của 4} (chữ số Hexa tận cùng là \texttt{0, 4, 8, C}).
        \item Các lệnh số học và logic chỉ thao tác với các toán hạng nằm trong \textbf{thanh ghi} (hoặc hằng số nằm ngay trong lệnh).
        \item Chỉ truy cập bộ nhớ dữ liệu thông qua các lệnh chuyên biệt là \textbf{Load} (đọc từ bộ nhớ vào thanh ghi) và \textbf{Store} (ghi từ thanh ghi ra bộ nhớ).
    \end{enumerate}
\end{itemize}

\textbf{2. Tập 32 thanh ghi đa dụng của RISC-V 32-bit (RV32I)}\\
Bộ xử lý RISC-V 32-bit có $32$ thanh ghi số nguyên, mỗi thanh ghi rộng $32\text{ bit}$, đánh số từ \texttt{x0} đến \texttt{x31}. Vì $32 = 2^5$, mỗi thanh ghi được mã hóa trong lệnh máy bằng đúng \textbf{5 bit} (từ \texttt{00000} đến \texttt{11111}).

\begin{tcolorbox}[bwbox, title={Bảng tra cứu 32 thanh ghi của RISC-V và Quy ước bảo toàn khi gọi hàm}]
\begin{center}
\renewcommand{\arraystretch}{1.2}
\small
\begin{tabularx}{\linewidth}{|L{2.0cm}|C{1.8cm}|Y|L{2.8cm}|}
\hline
\textbf{Tên gọi} & \textbf{Số hiệu} & \textbf{Chức năng sử dụng trong chương trình} & \textbf{Quy ước bảo toàn} \\ \hline
\texttt{zero} & \texttt{x0} & Luôn chứa hằng số $0$ (thanh ghi chỉ đọc, mọi thao tác ghi vào \texttt{x0} đều bị bỏ qua) & Không áp dụng \\ \hline
\texttt{ra} & \texttt{x1} & Return Address: Chứa địa chỉ trở về khi gọi chương trình con & Caller-saved \\ \hline
\texttt{sp} & \texttt{x2} & Stack Pointer: Con trỏ ngăn xếp (luôn trỏ vào đỉnh ngăn xếp) & Callee-saved \\ \hline
\texttt{gp} & \texttt{x3} & Global Pointer: Con trỏ vùng dữ liệu tĩnh toàn cục & Không áp dụng \\ \hline
\texttt{tp} & \texttt{x4} & Thread Pointer: Con trỏ luồng thực thi & Không áp dụng \\ \hline
\texttt{t0 -- t2} & \texttt{x5 -- x7} & Temporaries: Các thanh ghi chứa giá trị trung gian tạm thời & Caller-saved \\ \hline
\texttt{s0 / fp} & \texttt{x8} & Saved register / Frame pointer: Lưu biến hoặc làm con trỏ khung & Callee-saved \\ \hline
\texttt{s1} & \texttt{x9} & Saved register: Thanh ghi lưu biến của chương trình & Callee-saved \\ \hline
\texttt{a0 -- a1} & \texttt{x10 -- x11} & Arguments / Return values: Truyền tham số đầu vào và nhận kết quả trả về của hàm & Caller-saved \\ \hline
\texttt{a2 -- a7} & \texttt{x12 -- x17} & Arguments: Truyền các tham số đầu vào thứ 3 đến thứ 8 của hàm & Caller-saved \\ \hline
\texttt{s2 -- s11} & \texttt{x18 -- x27} & Saved registers: Các thanh ghi lưu biến bền vững của chương trình & Callee-saved \\ \hline
\texttt{t3 -- t6} & \texttt{x28 -- x31} & Temporaries: Các thanh ghi chứa giá trị trung gian tạm thời & Caller-saved \\ \hline
\end{tabularx}
\end{center}
\textbf{Quy tắc bảo toàn thanh ghi khi gọi chương trình con:}
\begin{itemize}
    \item \textbf{Nhóm tạm thời (\texttt{t0 -- t6}, \texttt{a0 -- a7}, \texttt{ra}):} Thuộc nhóm \textbf{Caller-saved} (hàm gọi tự bảo vệ trên ngăn xếp nếu muốn dùng lại sau lời gọi hàm; hàm được gọi có quyền thay đổi tự do).
    \item \textbf{Nhóm lưu trữ (\texttt{s0 -- s11}, \texttt{sp}):} Thuộc nhóm \textbf{Callee-saved} (hàm được gọi bắt buộc phải cất giá trị cũ vào ngăn xếp trước khi sử dụng và khôi phục lại nguyên trạng trước khi thoát hàm).
\end{itemize}
\end{tcolorbox}

\subsection{Tổ chức dữ liệu trong bộ nhớ và Quy tắc Đầu nhỏ (Little-Endian)}

\textbf{1. Các đơn vị dữ liệu và Quy tắc căn chỉnh biên địa chỉ (Alignment)}
\begin{itemize}
    \item \textbf{Byte ($8\text{ bit}$):} Chiếm $1$ ngăn nhớ. Địa chỉ có thể là số chẵn hoặc số lẻ bất kỳ.
    \item \textbf{Halfword -- Nửa từ ($16\text{ bit} = 2\text{ byte}$):} Chiếm $2$ ngăn nhớ liên tiếp. Địa chỉ hiệu dụng bắt buộc phải là \textbf{số chẵn} (chia hết cho $2$, chữ số Hexa cuối là \texttt{0, 2, 4, 6, 8, A, C, E}).
    \item \textbf{Word -- Từ nhớ ($32\text{ bit} = 4\text{ byte}$):} Chiếm $4$ ngăn nhớ liên tiếp. Địa chỉ của một Word là địa chỉ của byte đầu tiên (byte có địa chỉ nhỏ nhất) thuộc Word đó và bắt buộc phải \textbf{chia hết cho 4} (chữ số Hexa cuối cùng chỉ có thể là \texttt{0, 4, 8} hoặc \texttt{C}).
\end{itemize}

\textbf{2. Thứ tự lưu trữ các byte trong bộ nhớ: Đầu nhỏ (Little-Endian) vs. Đầu to (Big-Endian)}
\begin{itemize}
    \item \textbf{Đầu nhỏ (Little-Endian):} Byte có ý nghĩa thấp nhất (LSB -- 2 chữ số Hexa ngoài cùng bên phải) được lưu ở ngăn nhớ có \textbf{địa chỉ nhỏ nhất}; byte có ý nghĩa cao nhất (MSB -- 2 chữ số Hexa ngoài cùng bên trái) được lưu ở ngăn nhớ có \textbf{địa chỉ lớn nhất}.
    \begin{itemize}
        \item \textit{Các kiến trúc sử dụng:} \textbf{RISC-V} và \textbf{Intel x86} sử dụng Little-Endian.
    \end{itemize}
    \item \textbf{Đầu to (Big-Endian):} Byte có ý nghĩa cao nhất (MSB) được lưu ở ngăn nhớ có địa chỉ nhỏ nhất; byte có ý nghĩa thấp nhất (LSB) được lưu ở ngăn nhớ có địa chỉ lớn nhất.
    \begin{itemize}
        \item \textit{Các kiến trúc sử dụng:} Motorola 680x0, SunSPARC dùng Big-Endian. Riêng kiến trúc \textbf{ARM} hỗ trợ cả Little-Endian hoặc Big-Endian tùy phiên bản cấu hình.
    \end{itemize}
\end{itemize}

\begin{center}
\begin{tikzpicture}[scale=0.95]
    \node[font=\bfseries] at (0, 2.6) {Word 32-bit = \texttt{0x1A2B3C4D} lưu tại địa chỉ \texttt{0x10010008}};
    
    % Little-Endian
    \node[font=\bfseries\small] at (-3.5, 1.9) {Thứ tự Đầu nhỏ (Little-Endian -- RISC-V)};
    \draw[thick] (-5.5, -0.5) rectangle (-3.5, 1.5);
    \foreach \y/\byte/\addr in {1.0/0x4D (LSB)/0x10010008, 0.5/0x3C/0x10010009, 0.0/0x2B/0x1001000A, -0.5/0x1A (MSB)/0x1001000B} {
        \draw (-5.5, \y) -- (-3.5, \y);
        \node[font=\footnotesize\ttfamily] at (-4.5, \y+0.25) {\byte};
        \node[right, font=\footnotesize\ttfamily] at (-3.4, \y+0.25) {\addr};
    }

    % Big-Endian
    \node[font=\bfseries\small] at (3.5, 1.9) {Thứ tự Đầu to (Big-Endian)};
    \draw[thick] (1.5, -0.5) rectangle (3.5, 1.5);
    \foreach \y/\byte/\addr in {1.0/0x1A (MSB)/0x10010008, 0.5/0x2B/0x10010009, 0.0/0x3C/0x1001000A, -0.5/0x4D (LSB)/0x1001000B} {
        \draw (1.5, \y) -- (3.5, \y);
        \node[font=\footnotesize\ttfamily] at (2.5, \y+0.25) {\byte};
        \node[right, font=\footnotesize\ttfamily] at (3.6, \y+0.25) {\addr};
    }
\end{tikzpicture}
\end{center}

\subsection{Các nhóm lệnh hợp ngữ RISC-V cốt lõi}

\textbf{1. Nhóm lệnh số học và Nạp hằng số vào thanh ghi}
\begin{itemize}
    \item \texttt{add rd, rs1, rs2}: Cộng hai thanh ghi ($rd = rs1 + rs2$).
    \item \texttt{sub rd, rs1, rs2}: Trừ hai thanh ghi ($rd = rs1 - rs2$).
    \item \texttt{addi rd, rs1, imm}: Cộng thanh ghi với hằng số 12-bit có dấu ($rd = rs1 + \text{SignExt}(imm)$).
    \begin{itemize}
        \item \textit{Giới hạn của hằng số \texttt{imm}:} Vì \texttt{imm} là số nguyên có dấu $12\text{ bit}$, giá trị hợp lệ bắt buộc phải nằm trong đoạn $[\mathbf{-2048, +2047}]$ (viết ở hệ Hexa 12-bit là từ \texttt{0x800} đến \texttt{0x7FF}, hoặc mở rộng dấu 32-bit là từ \texttt{0xFFFFF800} đến \texttt{0x000007FF}). Nếu truyền số $\ge 2048$ (như \texttt{0x800}) vào lệnh \texttt{addi}, trình dịch sẽ báo lỗi.
        \item RISC-V \textbf{không có lệnh \texttt{subi}}; để trừ một hằng số, ta dùng lệnh \texttt{addi} với hằng số âm.
    \end{itemize}
    \item \texttt{lui rd, upimm} (Load Upper Immediate): Đưa hằng số 20-bit (\texttt{upimm}) vào $20\text{ bit}$ cao của thanh ghi đích \texttt{rd} (từ bit 31 đến bit 12) và xóa $12\text{ bit}$ thấp (bit 11 đến bit 0) về $0$.
\end{itemize}

\begin{tcolorbox}[bwbox, title={Quy tắc nạp hằng số 32-bit bất kỳ vào thanh ghi (Bản chất của lệnh giả \texttt{li})}]
Để nạp hằng số 32-bit dạng \texttt{0xABCDE123} vào thanh ghi \texttt{rd} mà không dùng lệnh giả \texttt{li}, ta tách thành $20\text{ bit}$ cao (5 chữ số Hexa đầu \texttt{0xABCDE}) và $12\text{ bit}$ thấp (3 chữ số Hexa cuối \texttt{0x123}), rồi kết hợp \texttt{lui} với \texttt{addi}:
\begin{enumerate}
    \item \textbf{Trường hợp 1 (Chữ số Hexa thứ 3 từ phải sang nhỏ hơn 8 -- từ \texttt{0} đến \texttt{7}):} Bit số 11 bằng $0$ nên phần 12-bit thấp là số dương. Giữ nguyên 5 chữ số Hexa đầu cho lệnh \texttt{lui}:
    \begin{itemize}
        \item Ví dụ nạp \texttt{0x20241120}: \texttt{lui t0, 0x20241} $\rightarrow$ \texttt{addi t0, t0, 0x120}.
    \end{itemize}
    \item \textbf{Trường hợp 2 (Chữ số Hexa thứ 3 từ phải sang từ \texttt{8} đến \texttt{F}):} Bit số 11 bằng $1$ nên lệnh \texttt{addi} sẽ mở rộng dấu và hiểu 12-bit thấp là một \textbf{số âm} (bị trừ lùi đi $2^{12} = \texttt{0x1000}$). Để bù lại, ta bắt buộc phải \textbf{cộng thêm 1 vào 5 chữ số Hexa đầu} trong lệnh \texttt{lui}:
    \begin{itemize}
        \item Ví dụ nạp \texttt{0x12345EDC} (chữ thứ 3 từ phải sang là \texttt{E} $\ge 8$, với \texttt{0xEDC} $= -292$):\\
        Dùng \texttt{lui s1, 0x12346} (tăng \texttt{12345} lên \texttt{12346}), sau đó \texttt{addi s1, s1, -292}.
    \end{itemize}
\end{enumerate}
\end{tcolorbox}

\textbf{2. Nhóm lệnh Đọc và Ghi bộ nhớ (Load / Store)}
\begin{itemize}
    \item \textbf{Thao tác với Word ($32\text{ bit} = 4\text{ byte}$):}
    \begin{itemize}
        \item \texttt{lw rd, imm(rs1)}: Đọc 1 word (4 byte) tại địa chỉ $\texttt{rs1} + \text{SignExt}(\texttt{imm})$ nạp vào thanh ghi \texttt{rd}.
        \item \texttt{sw rs2, imm(rs1)}: Ghi 1 word (4 byte) từ thanh ghi nguồn \texttt{rs2} ra bộ nhớ tại địa chỉ $\texttt{rs1} + \text{SignExt}(\texttt{imm})$.
    \end{itemize}
    \item \textbf{Thao tác với Halfword ($16\text{ bit} = 2\text{ byte}$):}
    \begin{itemize}
        \item \texttt{lh rd, imm(rs1)}: Đọc 2 byte từ bộ nhớ vào 16 bit thấp của \texttt{rd}, mở rộng \textbf{có dấu} (Sign-extend: lấy bit số 15 điền vào toàn bộ 16 bit cao của \texttt{rd}).
        \item \texttt{lhu rd, imm(rs1)}: Đọc 2 byte từ bộ nhớ vào 16 bit thấp của \texttt{rd}, mở rộng \textbf{không dấu} (Zero-extend: luôn điền \texttt{0x0000} vào 16 bit cao của \texttt{rd}).
        \item \texttt{sh rs2, imm(rs1)}: Ghi 2 byte bên phải (16 bit thấp) của thanh ghi \texttt{rs2} ra bộ nhớ.
    \end{itemize}
    \item \textbf{Thao tác với Byte ($8\text{ bit} = 1\text{ byte}$):}
    \begin{itemize}
        \item \texttt{lb rd, imm(rs1)}: Đọc 1 byte từ bộ nhớ vào 8 bit thấp của \texttt{rd}, mở rộng \textbf{có dấu} (Sign-extend: nếu byte đọc lên có chữ số Hexa đầu từ \texttt{8} đến \texttt{F} thì 24 bit cao của \texttt{rd} điền toàn \texttt{0xFFFFFF}; nếu từ \texttt{0} đến \texttt{7} thì điền \texttt{0x000000}).
        \item \texttt{lbu rd, imm(rs1)}: Đọc 1 byte từ bộ nhớ vào 8 bit thấp của \texttt{rd}, mở rộng \textbf{không dấu} (Zero-extend: luôn điền \texttt{0x000000} vào 24 bit cao của \texttt{rd}).
        \item \texttt{sb rs2, imm(rs1)}: Chỉ ghi 1 byte bên phải (8 bit thấp) của thanh ghi \texttt{rs2} ra bộ nhớ.
    \end{itemize}
\end{itemize}

\textbf{3. Nhóm lệnh Logic và Dịch bit}
\begin{itemize}
    \item \textbf{Các phép toán logic từng bit (Bitwise):}
    \begin{itemize}
        \item \texttt{and rd, rs1, rs2} và \texttt{andi rd, rs1, imm}: Phép AND ($1\text{ AND }1 = 1$, còn lại bằng $0$). Dùng để xóa một số bit về $0$ và trích xuất các bit cần giữ.
        \item \texttt{or rd, rs1, rs2} và \texttt{ori rd, rs1, imm}: Phép OR ($0\text{ OR }0 = 0$, còn lại bằng $1$). Dùng để thiết lập một số bit lên $1$.
        \item \texttt{xor rd, rs1, rs2} và \texttt{xori rd, rs1, imm}: Phép XOR (hai bit khác nhau thì bằng $1$, giống nhau thì bằng $0$). Dùng để đảo giá trị các bit (thay thế cho lệnh \texttt{not rd, rs1} bằng \texttt{xori rd, rs1, -1}) hoặc xóa thanh ghi về $0$ (\texttt{xor s0, s0, s0}).
        \item \textit{Lưu ý đặc biệt của RISC-V:} Hằng số 12-bit \texttt{imm} trong các lệnh \texttt{andi, ori, xori} của RISC-V luôn được \textbf{mở rộng có dấu (Sign-extended)} thành 32-bit trước khi thực hiện phép toán.
    \end{itemize}
    \item \textbf{Các phép dịch bit (Shift):} Với \texttt{uimm} là hằng số không dấu $5\text{ bit}$ (từ $0$ đến $31$), hoặc dịch theo 5 bit thấp của thanh ghi \texttt{rs2}:
    \begin{itemize}
        \item \textit{Dịch trái logic (\texttt{sll, slli}):} Đẩy các bit sang trái $k$ vị trí, chèn $k$ bit \texttt{0} vào bên phải. Tương đương \textbf{nhân với $2^k$}.
        \item \textit{Dịch phải logic (\texttt{srl, srli}):} Đẩy các bit sang phải $k$ vị trí, luôn chèn $k$ bit \texttt{0} vào bên trái. Tương đương \textbf{chia số không dấu cho $2^k$} lấy phần nguyên.
        \item \textit{Dịch phải số học (\texttt{sra, srai}):} Đẩy các bit sang phải $k$ vị trí, nhưng \textbf{giữ nguyên bit dấu} (nếu bit 31 đang là \texttt{1} thì chèn toàn bit \texttt{1} vào bên trái; nếu bit 31 là \texttt{0} thì chèn bit \texttt{0}). Tương đương \textbf{chia số có dấu cho $2^k$} lấy phần nguyên.
        \item \textit{Quy đổi nhanh khi dịch theo bội số của 4 bit:} Vì $1\text{ chữ số Hexa} = 4\text{ bit}$, nên dịch trái/phải $4\text{ bit}$ là dịch $1$ chữ số Hexa; dịch $8\text{ bit}$ là dịch $2$ chữ số Hexa; dịch $20\text{ bit}$ là dịch $5$ chữ số Hexa; dịch $24\text{ bit}$ là dịch $6$ chữ số Hexa.
    \end{itemize}
\end{itemize}

\textbf{4. Nhóm lệnh So sánh, Rẽ nhánh, Nhảy và Quản lý Ngăn xếp (Stack)}
\begin{itemize}
    \item \textbf{Lệnh thiết lập nếu nhỏ hơn (Set Less Than):}
    \begin{itemize}
        \item \texttt{slt rd, rs1, rs2} và \texttt{slti rd, rs1, imm}: So sánh hai số nguyên \textbf{có dấu}. Nếu $\texttt{rs1} < \texttt{rs2}$ (hoặc $\text{SignExt}(\texttt{imm})$) thì $\texttt{rd} = 1$, ngược lại $\texttt{rd} = 0$.
        \item \texttt{sltu rd, rs1, rs2} và \texttt{sltiu rd, rs1, imm}: So sánh hai số nguyên \textbf{không dấu}.
    \end{itemize}
    \item \textbf{Các lệnh rẽ nhánh có điều kiện (Branch):}
    \begin{itemize}
        \item \texttt{beq rs1, rs2, L1} (rẽ nhánh nếu $\texttt{rs1} == \texttt{rs2}$); \quad \texttt{bne rs1, rs2, L1} (rẽ nhánh nếu $\texttt{rs1} \neq \texttt{rs2}$).
        \item \texttt{blt rs1, rs2, L1} (rẽ nhánh nếu $\texttt{rs1} < \texttt{rs2}$ theo số có dấu); \quad \texttt{bge rs1, rs2, L1} (nếu $\texttt{rs1} \ge \texttt{rs2}$ có dấu).
        \item \texttt{bltu rs1, rs2, L1} (rẽ nhánh nếu $\texttt{rs1} < \texttt{rs2}$ không dấu); \quad \texttt{bgeu rs1, rs2, L1} (nếu $\texttt{rs1} \ge \texttt{rs2}$ không dấu).
        \item \textit{Lưu ý khi dịch câu lệnh \texttt{if / while} từ C sang hợp ngữ:} Điều kiện kiểm tra trong lệnh rẽ nhánh nhảy tới nhãn thoát (\texttt{done} hoặc \texttt{else}) thường là \textbf{điều kiện ngược} với điều kiện viết trong câu lệnh C (ví dụ: trong C viết \texttt{if (i == j)} thì hợp ngữ dùng \texttt{bne} để nhảy qua khối lệnh nếu $i \neq j$).
    \end{itemize}
    \item \textbf{Các lệnh gọi và trở về từ chương trình con (Thủ tục):}
    \begin{itemize}
        \item \texttt{jal rd, Label} (Jump and Link): Thực hiện đồng thời 2 việc: (1) Lưu địa chỉ của lệnh kế tiếp vào thanh ghi đích $\texttt{rd} = PC + 4$, và (2) Nhảy tới nhãn đích $PC = PC + \text{imm}$.
        \begin{itemize}
            \item Khi gọi chương trình con: Dùng \texttt{jal ra, ProcLabel} (lưu địa chỉ trở về vào \texttt{ra}).
            \item Khi nhảy không điều kiện (tương đương lệnh giả \texttt{j Label}): Dùng \texttt{jal x0, Label} (hủy địa chỉ trở về vào \texttt{x0}).
        \end{itemize}
        \item \texttt{jalr rd, rs1, imm} (Jump and Link Register): Lưu $\texttt{rd} = PC + 4$ và nhảy tới địa chỉ $PC = \texttt{rs1} + \text{SignExt}(\texttt{imm})$.
        \begin{itemize}
            \item Khi trở về từ chương trình con (tương đương lệnh giả \texttt{ret} hoặc \texttt{jr ra}): Dùng \texttt{jalr x0, ra, 0} (hoặc viết dạng \texttt{jalr x0, 0(ra)}).
        \end{itemize}
    \end{itemize}
    \item \textbf{Thao tác với Ngăn xếp (Stack):}
    \begin{itemize}
        \item Ngăn xếp là vùng nhớ hoạt động theo cơ chế \textbf{LIFO (Last In -- First Out: Vào sau -- Ra trước)} nằm trong bộ nhớ chính (RAM). Thanh ghi \texttt{sp} (\texttt{x2}) luôn chứa địa chỉ của \textbf{đỉnh ngăn xếp}.
        \item Trong RISC-V, ngăn xếp phát triển theo chiều \textbf{giảm địa chỉ bộ nhớ}:
        \begin{enumerate}
            \item \textbf{Khi cất (Push) $k$ thanh ghi (mỗi thanh ghi $4\text{ byte}$) vào ngăn xếp:} Phải giảm \texttt{sp} trước bằng lệnh \texttt{addi sp, sp, -4*k}, sau đó dùng lệnh \texttt{sw} lưu các thanh ghi vào \texttt{0(sp)}, \texttt{4(sp)}, \dots
            \item \textbf{Khi lấy (Pop / Khôi phục) $k$ thanh ghi ra khỏi ngăn xếp:} Phải dùng lệnh \texttt{lw} đọc các thanh ghi từ \texttt{0(sp)}, \texttt{4(sp)}, \dots ra đúng vị trí cũ trước, sau đó mới tăng \texttt{sp} trả lại bằng lệnh \texttt{addi sp, sp, 4*k}.
        \end{enumerate}
    \end{itemize}
\end{itemize}

\subsection{Mã máy 32-bit của RISC-V, Lệnh nén 16-bit và 4 Phương thức định địa chỉ}

\textbf{1. Sáu khuôn dạng lệnh máy 32-bit của RISC-V (RV32I)}\\
Mọi câu lệnh chuẩn 32-bit của RISC-V đều thuộc 1 trong 6 khuôn dạng: \textbf{R, I, S, B, U, J}. Trường \texttt{opcode} luôn cố định ở 7 bit cuối cùng (\texttt{bit 6:0}).

\begin{tcolorbox}[bwbox, title={Bảng cấu trúc 6 khuôn dạng lệnh máy 32-bit của RISC-V}]
\begin{center}
\renewcommand{\arraystretch}{1.3}
\footnotesize
\begin{tabularx}{\linewidth}{|C{1.4cm}|C{2.4cm}|C{2.0cm}|C{2.0cm}|C{1.8cm}|C{2.2cm}|Z|}
\hline
\textbf{Kiểu} & \textbf{Bit 31:25 (7b)} & \textbf{Bit 24:20 (5b)} & \textbf{Bit 19:15 (5b)} & \textbf{Bit 14:12 (3b)} & \textbf{Bit 11:7 (5b)} & \textbf{Bit 6:0 (7b)} \\ \hline
\textbf{R-type} & \texttt{funct7} & \texttt{rs2} & \texttt{rs1} & \texttt{funct3} & \texttt{rd} & \texttt{opcode} \\ \hline
\textbf{I-type} & \multicolumn{2}{c|}{\texttt{imm[11:0]} (12 bit)} & \texttt{rs1} & \texttt{funct3} & \texttt{rd} & \texttt{opcode} \\ \hline
\textbf{S-type} & \texttt{imm[11:5]} & \texttt{rs2} & \texttt{rs1} & \texttt{funct3} & \texttt{imm[4:0]} & \texttt{opcode} \\ \hline
\textbf{B-type} & \texttt{imm[12,10:5]} & \texttt{rs2} & \texttt{rs1} & \texttt{funct3} & \texttt{imm[4:1,11]} & \texttt{opcode} \\ \hline
\textbf{U-type} & \multicolumn{4}{c|}{\texttt{imm[31:12]} (20 bit)} & \texttt{rd} & \texttt{opcode} \\ \hline
\textbf{J-type} & \multicolumn{4}{c|}{\texttt{imm[20,10:1,11,19:12]} (20 bit)} & \texttt{rd} & \texttt{opcode} \\ \hline
\end{tabularx}
\end{center}
\textbf{Các đặc điểm nhận dạng nhanh trong đề thi:}
\begin{itemize}
    \item Khuôn dạng duy nhất \textbf{không chứa hằng số (\texttt{imm})} là kiểu \textbf{R-type} (chỉ dùng cho các phép toán thanh ghi -- thanh ghi).
    \item Khuôn dạng bắt buộc sử dụng đầy đủ cả hai trường \textbf{\texttt{funct3} và \texttt{funct7}} là kiểu \textbf{R-type} (ngoài ra các lệnh dịch bit tức thì \texttt{slli, srli, srai} thuộc kiểu I nhưng dùng 7 bit đầu \texttt{imm[11:5]} đóng vai trò như \texttt{funct7} và 5 bit \texttt{imm[4:0]} làm số lượng bit cần dịch \texttt{uimm}).
    \item Hai khuôn dạng có trường hằng số dài \textbf{20 bit} là kiểu \textbf{U-type} (\texttt{lui, auipc}) và kiểu \textbf{J-type} (\texttt{jal}). Lưu ý: Lệnh \texttt{jalr} thuộc kiểu \textbf{I-type} (12-bit \texttt{imm}), không phải J-type.
    \item Hai khuôn dạng \textbf{không có thanh ghi đích \texttt{rd}} là kiểu \textbf{S-type} (Store) và kiểu \textbf{B-type} (Branch).
\end{itemize}
\end{tcolorbox}

\begin{tcolorbox}[bwbox, title={Bảng tra cứu nhanh Opcode, Funct3 và Funct7 của tập lệnh RV32I}]
\begin{center}
\renewcommand{\arraystretch}{1.2}
\small
\begin{tabularx}{\linewidth}{|L{2.5cm}|C{2.5cm}|C{2.1cm}|C{1.9cm}|Y|}
\hline
\textbf{Nhóm / Kiểu} & \textbf{Opcode (6:0)} & \textbf{funct3} & \textbf{funct7} & \textbf{Tên lệnh hợp ngữ tương ứng} \\ \hline
\multirow{4}{2.5cm}{\textbf{R-type} \newline (Số học / Logic)} & \multirow{4}{2.5cm}{\centering\texttt{0110011} (51)} & \texttt{000} & \texttt{0000000} & \texttt{add rd, rs1, rs2} \\ \cline{3-5}
 & & \texttt{000} & \texttt{0100000} & \texttt{sub rd, rs1, rs2} \\ \cline{3-5}
 & & \texttt{001 / 010} & \texttt{0000000} & \texttt{sll} (\texttt{001}), \texttt{slt} (\texttt{010}), \texttt{sltu} (\texttt{011}) \\ \cline{3-5}
 & & \texttt{100 / 101} & \texttt{000/010..} & \texttt{xor} (\texttt{100}), \texttt{srl/sra} (\texttt{101}), \texttt{or} (\texttt{110}), \texttt{and} (\texttt{111}) \\ \hline
\multirow{2}{2.5cm}{\textbf{I-type} \newline (Tính toán)} & \multirow{2}{2.5cm}{\centering\texttt{0010011} (19)} & \texttt{000 / 010} & -- & \texttt{addi} (\texttt{000}), \texttt{slti} (\texttt{010}), \texttt{sltiu} (\texttt{011}) \\ \cline{3-5}
 & & \texttt{100 / 110} & -- & \texttt{xori} (\texttt{100}), \texttt{ori} (\texttt{110}), \texttt{andi} (\texttt{111}), \texttt{slli} (\texttt{001}), \texttt{srli/srai} (\texttt{101}) \\ \hline
\textbf{I-type} (Load) & \texttt{0000011} (3) & \texttt{000..101} & -- & \texttt{lb} (\texttt{000}), \texttt{lh} (\texttt{001}), \texttt{lw} (\texttt{010}), \texttt{lbu} (\texttt{100}), \texttt{lhu} (\texttt{101}) \\ \hline
\textbf{S-type} (Store) & \texttt{0100011} (35) & \texttt{000..010} & -- & \texttt{sb} (\texttt{000}), \texttt{sh} (\texttt{001}), \texttt{sw} (\texttt{010}) \\ \hline
\textbf{B-type} (Branch) & \texttt{1100011} (99) & \texttt{000..111} & -- & \texttt{beq} (\texttt{000}), \texttt{bne} (\texttt{001}), \texttt{blt} (\texttt{100}), \texttt{bge} (\texttt{101}), \texttt{bltu} (\texttt{110}), \texttt{bgeu} (\texttt{111}) \\ \hline
\textbf{U-type} & \texttt{0110111} (55) & -- & -- & \texttt{lui rd, upimm} (\texttt{auipc} có opcode \texttt{0010111}) \\ \hline
\textbf{J-type} \& \textbf{I-jalr} & \texttt{1101111} (111) & -- & -- & \texttt{jal rd, label} (Riêng \texttt{jalr} là kiểu I, opcode \texttt{1100111}, \texttt{funct3=000}) \\ \hline
\end{tabularx}
\end{center}
\end{tcolorbox}

\textbf{2. Tập lệnh nén 16-bit (Compressed Instructions -- RVC):}
\begin{itemize}
    \item Để giảm dung lượng bộ nhớ lưu trữ mã chương trình, RISC-V (cũng như ARM Thumb) cung cấp tập lệnh nén có độ dài \textbf{16 bit} ($2\text{ byte}$) cho các thao tác thường gặp, ký hiệu bằng tiền tố \texttt{c.} (ví dụ: \texttt{c.add, c.addi, c.lw}).
    \item \textbf{Các hạn chế của lệnh nén 16-bit so với lệnh 32-bit:}
    \begin{enumerate}
        \item Giảm số lượng thanh ghi có thể truy cập trực tiếp (nhiều lệnh nén chỉ dùng 3 bit để định danh 8 thanh ghi phổ biến từ \texttt{x8} đến \texttt{x15}).
        \item Giảm số lượng toán hạng trong lệnh (thường gộp thanh ghi đích trùng với thanh ghi nguồn thứ nhất: $\texttt{rd} = \texttt{rd} + \texttt{rs2}$).
        \item Phạm vi biểu diễn của hằng số tức thời (\texttt{imm}) và khoảng cách nhảy rẽ nhánh bị thu hẹp đáng kể.
    \end{enumerate}
\end{itemize}

\textbf{3. Bốn phương thức định địa chỉ của RISC-V}
\begin{enumerate}
    \item \textbf{Định địa chỉ thanh ghi (Register Addressing):} Toán hạng nằm hoàn toàn trong các thanh ghi (áp dụng cho các lệnh kiểu R).
    \item \textbf{Định địa chỉ tức thì (Immediate Addressing):} Toán hạng là hằng số nằm ngay bên trong lệnh máy (áp dụng cho toán hạng hằng số của \texttt{addi, andi, ori, lui...}).
    \item \textbf{Định địa chỉ cơ sở (Base Addressing):} Toán hạng nằm trong bộ nhớ, địa chỉ được tính bằng tổng nội dung thanh ghi cơ sở và hằng số độ lệch: $\text{Address} = \texttt{rs1} + \text{SignExt}(\texttt{imm})$ (áp dụng cho các lệnh Load và Store).
    \item \textbf{Định địa chỉ tương đối với PC (PC-relative Addressing):} Địa chỉ của lệnh đích cần nhảy tới bằng tổng giá trị thanh ghi \texttt{PC} (địa chỉ của chính lệnh rẽ nhánh/nhảy hiện tại) cộng với hằng số độ lệch:
    $$\text{Địa chỉ đích (BTA hoặc JTA)} = PC_{\text{hiện tại}} + \text{imm} \iff \text{imm} = \text{Địa chỉ đích} - PC_{\text{hiện tại}}$$
\end{enumerate}

\newpage
\section{Phần 2: Toàn bộ bài tập Chương 3 và Hướng dẫn giải}

\subsection{Nhóm 1: Thanh ghi, Khuôn dạng lệnh, Phương thức định địa chỉ và Kiểm tra hợp lệ}

\textbf{Bài tập 3.1 (Đề Cuối kỳ 2023.1 -- \#08de08, 2023.2 \& Ôn tập Cuối kỳ 2025.1, 2025.2 -- Câu 53)}\\
Phát biểu nào sau đây là đúng về các thanh ghi tạm thời (\texttt{t0 -- t6}) và các thanh ghi lưu trữ (\texttt{s0 -- s11}) trong RISC-V?
\begin{enumerate}[label=\Alph*.]
    \item Giá trị của các thanh ghi lưu trữ luôn phải được lưu bởi bên được gọi (callee).
    \item Giá trị của các thanh ghi tạm thời luôn phải được lưu bởi bên gọi (caller).
    \item Trong một số quy ước gọi hàm, thanh ghi tạm thời là ``caller-saved'', còn thanh ghi lưu trữ là ``callee-saved''.
    \item Trong một số quy ước gọi hàm, thanh ghi tạm thời là ``callee-saved'', còn thanh ghi lưu trữ là ``caller-saved''.
\end{enumerate}
\textbf{Hướng dẫn giải:}
\begin{itemize}
    \item \textbf{Kiến thức áp dụng:} Theo quy ước gọi hàm của RISC-V, các thanh ghi tạm thời (\texttt{t0 -- t6}) thuộc nhóm \textbf{caller-saved} (bên gọi chỉ cất vào ngăn xếp nếu cần dùng lại sau lời gọi hàm, không phải ``luôn luôn'' phải cất), còn các thanh ghi lưu trữ (\texttt{s0 -- s11}) thuộc nhóm \textbf{callee-saved} (bên được gọi chỉ cất vào ngăn xếp nếu nó làm thay đổi các thanh ghi này).
    \item \textbf{Kết luận:} Đáp án đúng là \textbf{C}.
\end{itemize}

\textbf{Bài tập 3.2 (Đề Giữa kỳ 2022.2 -- Đề 0008 Câu 10 \& Đề 0006 Câu 5)}
\begin{enumerate}[label=\alph*)]
    \item Trong 32 thanh ghi được đánh địa chỉ của bộ xử lý 32-bit, tên của 1 thanh ghi chỉ đọc (read-only) luôn chứa hằng số 0 là gì?
    \item Trong khuôn dạng lệnh R của bộ xử lý 32-bit, tại sao độ rộng của vùng thông tin thanh ghi đích \texttt{rd} phải là 5 bit?
\end{enumerate}
\textbf{Hướng dẫn giải:}
\begin{itemize}
    \item \textbf{Giải câu a:} Thanh ghi chỉ đọc luôn có giá trị bằng 0 là thanh ghi \textbf{\texttt{zero}} (hoặc \textbf{\texttt{x0}}).
    \item \textbf{Giải câu b:} Vì bộ xử lý có tất cả $32$ thanh ghi đa dụng. Để mã hóa phân biệt $32$ thanh ghi khác nhau cần số bit là $\log_2(32) = \mathbf{5\text{ bit}}$ (vì $2^5 = 32$).
\end{itemize}

\textbf{Bài tập 3.3 (Đề Giữa kỳ 2022.2 -- Đề 0006, Câu 8)}\\
Một phần mềm viết cho bộ xử lý RISC-V 32-bit có $150$ lệnh hợp ngữ chuẩn (không dùng lệnh giả và không dùng lệnh nén) thì chiếm bao nhiêu byte lệnh mã máy trong bộ nhớ?

\textbf{Hướng dẫn giải:}
\begin{itemize}
    \item Mỗi lệnh mã máy chuẩn của RISC-V 32-bit có độ dài cố định là $32\text{ bit} = 4\text{ byte}$.
    \item Vậy $150$ lệnh chiếm dung lượng là: $150 \times 4\text{ byte} = \mathbf{600\text{ byte}}$.
\end{itemize}

\textbf{Bài tập 3.4 (Đề Giữa kỳ 2019.2 -- Mã đề 243 Câu 17 \& Mã đề 244 Câu 10)}
\begin{enumerate}[label=\alph*)]
    \item \textbf{(Mã đề 243):} Một chương trình hợp ngữ có $8$ lệnh \texttt{addi}, $5$ lệnh \texttt{sub}, $1$ lệnh \texttt{lui}, $3$ lệnh \texttt{xori}. Hỏi mã máy của phần mềm trên chiếm bao nhiêu byte dung lượng?\\
    A. $18$ \quad B. $72$ \quad C. $17$ \quad D. $68$
    \item \textbf{(Mã đề 244):} Một chương trình hợp ngữ có $5$ lệnh \texttt{addi}, $2$ lệnh \texttt{sub}, $3$ lệnh \texttt{lui}, $8$ lệnh \texttt{xori}. Hỏi mã máy của phần mềm trên chiếm bao nhiêu byte dung lượng?\\
    A. $72$ \quad B. $68$ \quad C. $18$ \quad D. $17$
\end{enumerate}
\textbf{Hướng dẫn giải:}
\begin{itemize}
    \item \textbf{Giải câu a:} Tổng số lệnh máy $= 8 + 5 + 1 + 3 = 17\text{ lệnh}$. Mỗi lệnh dài $4\text{ byte}$ nên dung lượng mã máy là: $17 \times 4 = 68\text{ byte}$ $\rightarrow$ Chọn \textbf{D (68)}.
    \item \textbf{Giải câu b:} Tổng số lệnh máy $= 5 + 2 + 3 + 8 = 18\text{ lệnh}$. Dung lượng mã máy là: $18 \times 4 = 72\text{ byte}$ $\rightarrow$ Chọn \textbf{A (72)}.
\end{itemize}

\textbf{Bài tập 3.5 (Đề Giữa kỳ 2019.2 -- Mã đề 243 Câu 20 \& Mã đề 244 Câu 2)}\\
Đoạn hợp ngữ gồm 3 lệnh chuẩn liên tiếp (\texttt{lui s0, ...} $\rightarrow$ \texttt{addi s0, s0, ...} $\rightarrow$ \texttt{lw t0, 4(s0)}).
\begin{enumerate}[label=\alph*)]
    \item \textbf{(Mã đề 243):} Nếu lệnh đầu tiên ở địa chỉ \texttt{0xBCDE8004}, hỏi lệnh cuối cùng (lệnh thứ 3) ở địa chỉ nào?\\
    A. Không xác định \quad B. \texttt{0xBCDE8006} \quad C. \texttt{0x20201238} \quad D. \texttt{0xBCDE800C}
    \item \textbf{(Mã đề 244):} Nếu lệnh đầu tiên ở địa chỉ \texttt{0x8004FEDC}, hỏi lệnh cuối cùng (lệnh thứ 3) ở địa chỉ nào?\\
    A. Không xác định \quad B. \texttt{0x8004FEE4} \quad C. \texttt{0x8004FEDA} \quad D. \texttt{0x80004000}
\end{enumerate}
\textbf{Hướng dẫn giải:}
\begin{itemize}
    \item Mỗi lệnh dài $4\text{ byte}$. Lệnh 1 ở địa chỉ $A$, lệnh 2 ở $A + 4$, lệnh 3 ở địa chỉ $A + 8$:
    \item \textbf{Câu a:} $\texttt{0xBCDE8004} + 8 = \texttt{0xBCDE800C}$ $\rightarrow$ Chọn \textbf{D}.
    \item \textbf{Câu b:} $\texttt{0x8004FEDC} + 8 = \texttt{0x8004FEE4}$ (vì $\texttt{0xDC} + 8 = 220 + 8 = 228 = \texttt{0xE4}$) $\rightarrow$ Chọn \textbf{B}.
\end{itemize}

\textbf{Bài tập 3.6 (Đề Cuối kỳ 2022.2 -- \#17e1a7, 2024.1 -- \#366b03, 2024.2 \& Ôn tập 2025.2 -- Câu 34)}\\
Phương pháp định địa chỉ nào được sử dụng mã hóa cho toán hạng hằng số của lệnh hợp ngữ RISC-V sau: \texttt{addi s0, s3, -5} (hoặc \texttt{addi x18, x19, -5})?
\begin{enumerate}[label=\Alph*.]
    \item Định địa chỉ cơ sở
    \item Định địa chỉ tương đối với PC
    \item Định địa chỉ tức thì
    \item Định địa chỉ thanh ghi
\end{enumerate}
\textbf{Hướng dẫn giải:}
\begin{itemize}
    \item Toán hạng hằng số \texttt{-5} nằm trực tiếp ngay trong mã máy của câu lệnh \texttt{addi}, do đó sử dụng phương pháp \textbf{Định địa chỉ tức thì (Immediate Addressing)}.
    \item \textbf{Kết luận:} Đáp án đúng là \textbf{C}.
\end{itemize}

\textbf{Bài tập 3.7 (Đề Giữa kỳ 2018.2 -- Mã đề 198 Câu 19, Mã đề 199 Câu 13 \& Mã đề 204 Câu 15)}\\
Trong các phương pháp định địa chỉ toán hạng dưới đây, phương pháp nào có tốc độ thực thi lệnh chậm hơn?
\begin{enumerate}[label=\Alph*.]
    \item Định địa chỉ thanh ghi
    \item Định địa chỉ trực tiếp / Định địa chỉ cơ sở (truy cập bộ nhớ)
    \item Định địa chỉ tức thì
    \item Định địa chỉ trực tiếp và định địa chỉ tức thì
\end{enumerate}
\textbf{Hướng dẫn giải:}
\begin{itemize}
    \item Định địa chỉ tức thì (lấy hằng số có sẵn trong lệnh tại IR) và định địa chỉ thanh ghi (đọc thanh ghi trong CPU) không cần truy cập ra ngoài CPU nên tốc độ rất nhanh.
    \item Định địa chỉ trực tiếp / cơ sở phải truy xuất toán hạng nằm trong bộ nhớ chính (RAM) qua Bus hệ thống nên có tốc độ thực thi chậm hơn.
    \item \textbf{Kết luận:} Đáp án đúng là \textbf{B}.
\end{itemize}

\textbf{Bài tập 3.8 (Đề Cuối kỳ 2020.2 -- Trang 5 \& Cuối kỳ 2022.2, 2024.2 -- \#3ace1f)}
\begin{enumerate}[label=\alph*)]
    \item \textbf{(Đề CK 2020.2):} Trong các giá trị địa chỉ bộ nhớ hệ thống dưới đây, giá trị nào là địa chỉ hợp lệ và căn chỉnh đúng biên (aligned address) của một từ nhớ (word) dữ liệu 32-bit?\\
    A. \texttt{0x1001A00C} \quad B. \texttt{0x10800016} \quad C. \texttt{0x100020CE} \quad D. \texttt{0x1050009A}
    \item \textbf{(Đề CK 2022.2 \& 2024.2 -- \#3ace1f):} Trong các giá trị địa chỉ dưới đây, giá trị nào có thể là địa chỉ của word nhớ trong bộ nhớ chính của máy tính RISC-V?\\
    A. \texttt{0x00404508} \quad B. \texttt{0x0050601A} \quad C. \texttt{0x00601002} \quad D. \texttt{0x00402346}
    \item \textbf{(Đề GK 2018.2 -- Mã đề 198 Câu 9):} Địa chỉ nào sau đây \textbf{không phải} là địa chỉ hợp lệ của một lệnh máy 32-bit?\\
    A. \texttt{0x188A} \quad B. \texttt{0xFB019AC} \quad C. \texttt{0x18D88} \quad D. \texttt{0x198C124}
\end{enumerate}
\textbf{Hướng dẫn giải:}
\begin{itemize}
    \item \textbf{Quy tắc:} Địa chỉ của một word 32-bit hoặc một lệnh máy 32-bit bắt buộc phải \textbf{chia hết cho 4}, nghĩa là chữ số Hexa cuối cùng bên phải chỉ được phép là \textbf{\texttt{0, 4, 8} hoặc \texttt{C} ($=12$)}. Các số tận cùng bằng \texttt{2, 6, A} ($=10$), \texttt{E} ($=14$) là số chẵn nhưng không chia hết cho 4.
    \item \textbf{Câu a:} Chỉ có \texttt{0x1001A00C} tận cùng bằng \texttt{C} chia hết cho 4 $\rightarrow$ Chọn \textbf{A}.
    \item \textbf{Câu b:} Chỉ có \texttt{0x00404508} tận cùng bằng \texttt{8} chia hết cho 4 $\rightarrow$ Chọn \textbf{A}.
    \item \textbf{Câu c:} Chỉ có \texttt{0x188A} tận cùng bằng \texttt{A} ($10$) không chia hết cho 4 $\rightarrow$ Chọn \textbf{A}.
\end{itemize}

\textbf{Bài tập 3.9 (Đề Cuối kỳ 2023.1 -- Trang 6 \& Ôn tập Cuối kỳ 2025.2 -- P1, Câu 28)}\\
Toán hạng địa chỉ đích trong câu lệnh lưu trữ từ nhớ \texttt{sw} dưới đây của hệ RISC-V có hợp lệ để thực thi hay không?
\begin{lstlisting}
    lui  s0, 0x10010
    addi t0, zero, 100
    sw   t0, 2(s0)
\end{lstlisting}
\begin{enumerate}[label=\Alph*.]
    \item Có, hoàn toàn hợp lệ.
    \item Có, nhưng đoạn mã này yêu cầu quyền đặc biệt để thực thi.
    \item Không, nó không có một địa chỉ hợp lệ (do vi phạm căn lề ranh giới word nhớ -- alignment error).
    \item Không, nó không nằm trong phân đoạn dữ liệu hợp lệ.
\end{enumerate}
\textbf{Hướng dẫn giải:}
\begin{itemize}
    \item Lệnh \texttt{lui s0, 0x10010} nạp địa chỉ \texttt{0x10010000} vào \texttt{s0}.
    \item Lệnh \texttt{sw t0, 2(s0)} ghi một word 32-bit (4 byte) vào địa chỉ $\texttt{0x10010000} + 2 = \texttt{0x10010002}$.
    \item Địa chỉ \texttt{0x10010002} không chia hết cho 4, vi phạm quy tắc căn chỉnh biên word nhớ.
    \item \textbf{Kết luận:} Đáp án đúng là \textbf{C}.
\end{itemize}

\textbf{Bài tập 3.10 (Đề Cuối kỳ 2020.2 -- Trang 2, Trang 8 \& Đề Giữa kỳ 2018.2 -- Mã đề 204 Câu 7)}\\
Hãy xác định tổng số lần CPU trực tiếp truy nhập bộ nhớ dữ liệu (Data Memory) trong mỗi đoạn chương trình sau:
\begin{enumerate}[label=\alph*)]
    \item \textbf{(Đề CK 2020.2 -- Trang 2):} \texttt{lw t0, 0(s0)} $\rightarrow$ \texttt{addi t1, t0, 1} $\rightarrow$ \texttt{sw t1, 4(s0)} $\rightarrow$ \texttt{add t2, t1, t0}.
    \item \textbf{(Đề CK 2020.2 -- Trang 8):} \texttt{lw a1, 0(a2)} $\rightarrow$ \texttt{addi a1, a1, -4} $\rightarrow$ \texttt{slli a1, a1, 0} $\rightarrow$ \texttt{sw zero, 8(a1)}.
    \item \textbf{(Đề GK 2018.2 -- Mã đề 204, Câu 7):} \texttt{lui t0, 0x2019} $\rightarrow$ \texttt{addi t0, t0, 0x104} $\rightarrow$ \texttt{slli t1, t0, 3} $\rightarrow$ \texttt{sub t2, t1, t3}.
\end{enumerate}
\textbf{Hướng dẫn giải:}
\begin{itemize}
    \item Chỉ có các lệnh Load (\texttt{lw, lh, lb...}) và Store (\texttt{sw, sh, sb}) mới truy nhập vào bộ nhớ dữ liệu (Data Memory).
    \item \textbf{Câu a:} Có 1 lệnh \texttt{lw} và 1 lệnh \texttt{sw} $\implies$ Truy nhập bộ nhớ dữ liệu \textbf{2 lần}.
    \item \textbf{Câu b:} Có 1 lệnh \texttt{lw} và 1 lệnh \texttt{sw} $\implies$ Truy nhập bộ nhớ dữ liệu \textbf{2 lần}.
    \item \textbf{Câu c:} Không có lệnh Load/Store nào $\implies$ Truy nhập bộ nhớ dữ liệu \textbf{0 lần}.
\end{itemize}

\textbf{Bài tập 3.11 (Ôn tập Cuối kỳ 2025.2 -- P1, Câu 18)}\\
Phát biểu nào sau đây là đúng về thứ tự lưu trữ các byte trong bộ nhớ chính?
\begin{enumerate}[label=\Alph*.]
    \item Intel x86 dùng big-endian
    \item SunSPARC dùng middle-endian
    \item ARM dùng little-endian hoặc big-endian tùy phiên bản
    \item RISC-V dùng big-endian
\end{enumerate}
\textbf{Hướng dẫn giải:}
\begin{itemize}
    \item Theo Slide Chương 3 (trang 16): Intel x86 và RISC-V dùng little-endian; Motorola 680x0 và SunSPARC dùng big-endian; còn kiến trúc \textbf{ARM} dùng little-endian hoặc big-endian tùy phiên bản.
    \item \textbf{Kết luận:} Đáp án đúng là \textbf{C}.
\end{itemize}

\textbf{Bài tập 3.12 (Đề Cuối kỳ 2025.1 -- \#9853fd \& Ôn tập Cuối kỳ 2025.2 -- P1, Câu 42)}
\begin{enumerate}[label=\alph*)]
    \item \textbf{(Đề CK 2025.1 -- \#9853fd):} Phạm vi giá trị của toán hạng hằng số (immediate) trong lệnh \texttt{addi} của RISC-V là gì?\\
    A. $[-4095, +4095]$ \quad B. $[-2^{31}, +2^{31}-1]$ \quad C. $[-2048, +2048]$ \quad D. $[-2048, +2047]$
    \item \textbf{(Ôn tập 2025.2 -- P1, Câu 42):} Những câu lệnh nào sau đây bị lỗi biên dịch trong phần mềm RARS?\\
    1. \texttt{lui x1, 0x1} $\rightarrow$ \texttt{addi x1, x0, -2048} \qquad 2. \texttt{lui x1, 0x1} $\rightarrow$ \texttt{addi x1, x0, 0x800}\\
    3. \texttt{lui x1, 0x1} $\rightarrow$ \texttt{addi x1, x0, 0xFFFF800} \quad 4. \texttt{addi x1, x0, 0x801}\\
    5. \texttt{li x1, 0x800} \qquad 6. \texttt{lui x1, 0x1} $\rightarrow$ \texttt{addi x1, x0, 2048}\\
    7. \texttt{lui x1, 0x1} $\rightarrow$ \texttt{addi x1, x0, 0xFFFFF800}
\end{enumerate}
\textbf{Hướng dẫn giải:}
\begin{itemize}
    \item \textbf{Giải câu a:} Trường immediate của lệnh \texttt{addi} là số nguyên có dấu $12\text{ bit}$, có dải biểu diễn từ $-2^{11}$ đến $+2^{11}-1$, tức là $[\mathbf{-2048, +2047}]$ $\rightarrow$ Chọn \textbf{D}.
    \item \textbf{Giải câu b:} Mọi giá trị truyền vào \texttt{addi} nằm ngoài đoạn $[-2048, +2047]$ (tức ngoài đoạn \texttt{0xFFFFF800} đến \texttt{0x7FF}) đều báo lỗi biên dịch:
    \begin{itemize}
        \item Mục \textbf{2} (\texttt{0x800} $= 2048$), mục \textbf{3} (\texttt{0xFFFF800} chỉ có 4 chữ \texttt{F} nên là số dương lớn), mục \textbf{4} (\texttt{0x801} $= 2049$) và mục \textbf{6} ($2048$) đều vượt quá $+2047$ $\implies$ \textbf{Bị lỗi biên dịch (Chọn 2, 3, 4, 6)}.
    \end{itemize}
\end{itemize}

\textbf{Bài tập 3.13 (Đề Giữa kỳ 2021.1 -- IT3030E, Problem 6)}\\
Các lệnh chuẩn của RISC-V/MIPS dài 32 bit. Một cách để tiết kiệm bộ nhớ chương trình là giảm độ dài lệnh xuống 16 bit như tập lệnh nén (RISC-V Compressed RVC hoặc ARM Thumb). Hãy nêu ngắn gọn những hạn chế (limitations) tất yếu của các lệnh 16-bit này?

\textbf{Hướng dẫn giải:}
\begin{itemize}
    \item Khi thu gọn độ dài mã lệnh từ $32\text{ bit}$ xuống còn $16\text{ bit}$, số bit dành cho các trường bị cắt giảm mạnh dẫn đến 4 hạn chế chính:
    \begin{enumerate}
        \item \textbf{Hạn chế số lượng thanh ghi truy cập được:} Thay vì dùng 5 bit để truy cập cả 32 thanh ghi, nhiều lệnh 16-bit chỉ dùng 3 bit nên chỉ truy cập được 8 thanh ghi thông dụng (\texttt{x8 -- x15}).
        \item \textbf{Giảm số địa chỉ toán hạng trong lệnh:} Chuyển từ cấu trúc 3 toán hạng (\texttt{add rd, rs1, rs2}) sang 2 toán hạng (\texttt{c.add rd, rs2}, trong đó \texttt{rd} vừa là nguồn vừa là đích), làm tăng số lệnh cần viết nếu muốn giữ nguyên giá trị thanh ghi nguồn.
        \item \textbf{Thu hẹp phạm vi hằng số tức thời (Immediate):} Hằng số chỉ còn từ 5 đến 6 bit thay vì 12 bit hay 20 bit.
        \item \textbf{Thu hẹp khoảng cách nhảy và rẽ nhánh:} Độ lệch nhảy tương đối với PC ngắn hơn nhiều so với lệnh 32-bit.
    \end{enumerate}
\end{itemize}

\subsection{Nhóm 2: Thực thi lệnh Số học, Logic, Dịch bit và Khởi tạo thanh ghi}

\textbf{Bài tập 3.14 (Đề Cuối kỳ 2020.2 -- Trang 4 \& Đề Cuối kỳ 2021.1 -- Trang 3)}\\
Xác định nội dung giá trị số Hexa lưu trữ tại thanh ghi \texttt{t1} sau khi thực hiện đoạn lệnh hợp ngữ tính toán cơ bản dưới đây:
\begin{lstlisting}
    addi s5, zero, 199    # 0xC7 tuong ung gia tri thap phan 199
    addi s6, zero, 229
    add  t1, s6, s5
\end{lstlisting}
\textbf{Hướng dẫn giải:}
\begin{itemize}
    \item Thực hiện phép cộng trong hệ thập phân: $\texttt{t1} = 229 + 199 = 428_{(10)}$.
    \item Đổi $428_{(10)}$ sang hệ Hexa:
    $$428 = 256 + 172 = 1 \times 16^2 + 10 \times 16^1 + 12 \times 16^0 = \texttt{0x1AC}$$
    \item Viết đầy đủ 8 chữ số Hexa (32-bit): \textbf{\texttt{0x000001AC}} (hoặc \textbf{\texttt{0x1AC}}).
\end{itemize}

\textbf{Bài tập 3.15 (Đề Cuối kỳ 2022.2 \& Cuối kỳ 2024.2 -- \#014753)}\\
Cho đoạn chương trình hợp ngữ RISC-V sau (giả sử ban đầu $\texttt{s0} = 5$ và $\texttt{s1} = 4$, hoặc $\texttt{s0} = 7, \texttt{s1} = 2$):
\begin{lstlisting}
    add  t0, s0, s1
    sub  t1, s0, s1
    addi t2, t0, 10
    addi t3, t1, -5
\end{lstlisting}
Hãy xác định giá trị của \texttt{t2} và \texttt{t3} sau khi thực hiện đoạn chương trình:
\begin{enumerate}[label=\Alph*.]
    \item $\texttt{t2} = 19, \texttt{t3} = -3$ (khi $\texttt{s0} = 5.5$? Không, ứng với $\texttt{s0} + \texttt{s1} = 9$ và $\texttt{s0} - \texttt{s1} = 2$: không nguyên) $\rightarrow$ Kiểm tra với $\texttt{s0} = 7, \texttt{s1} = 2$: $\texttt{t0} = 9 \implies \texttt{t2} = 19$; $\texttt{t1} = 5 \implies \texttt{t3} = 0$.
    \item $\texttt{t2} = 19, \texttt{t3} = 0$
    \item $\texttt{t2} = 15, \texttt{t3} = -5$
    \item $\texttt{t2} = 24, \texttt{t3} = -8$
\end{enumerate}
\textbf{Hướng dẫn giải:}
\begin{itemize}
    \item Ta có hệ phương trình giữa \texttt{t2} và \texttt{t3}:
    $$\begin{cases} \texttt{t2} = (\texttt{s0} + \texttt{s1}) + 10 \\ \texttt{t3} = (\texttt{s0} - \texttt{s1}) - 5 \end{cases} \implies \texttt{t2} + \texttt{t3} = 2 \times \texttt{s0} + 5$$
    \item Vì \texttt{s0} là số nguyên nên $2 \times \texttt{s0} + 5$ bắt buộc phải là một \textbf{số lẻ} (tức tổng $\texttt{t2} + \texttt{t3}$ phải là số lẻ).
    \item Kiểm tra 4 đáp án:
    \begin{itemize}
        \item Đáp án A: $19 + (-3) = 16$ (chẵn -- loại).
        \item Đáp án C: $15 + (-5) = 10$ (chẵn -- loại).
        \item Đáp án D: $24 + (-8) = 16$ (chẵn -- loại).
        \item Đáp án B: $19 + 0 = 19$ (lẻ, ứng với $\texttt{s0} = 7, \texttt{s1} = 2$: $\texttt{t0} = 7 + 2 = 9 \implies \texttt{t2} = 19$; $\texttt{t1} = 7 - 2 = 5 \implies \texttt{t3} = 5 - 5 = 0$).
    \end{itemize}
    \item \textbf{Kết luận:} Đáp án đúng là \textbf{B ($\texttt{t2} = 19, \texttt{t3} = 0$)}.
\end{itemize}

\textbf{Bài tập 3.16 (Đề Giữa kỳ 2024.1 -- \#4610c2, Ôn tập 2025.2 -- Câu 75 \& Đề Giữa kỳ 2018.2 -- Mã đề 198 Câu 18)}
\begin{enumerate}[label=\alph*)]
    \item \textbf{(Đề GK 2024.1 -- \#4610c2 \& Ôn tập 2025.2 Câu 75):} Những lệnh nào sau đây có thể dùng để xóa nội dung của thanh ghi \texttt{s0} về 0?\\
    1. \texttt{and s0, s0, zero} \quad 2. \texttt{xor s0, s0, s0} \quad 3. \texttt{xor s0, s0, zero} \quad 4. \texttt{and s0, s0, s0}
    \item \textbf{(Đề GK 2018.2 -- Mã đề 198, Câu 18):} Trong hai lệnh \texttt{add s4, s4, s4} và \texttt{slli s4, s4, 1}, chúng cùng thực hiện phép toán nào? Hai lệnh \texttt{mul s4, s4, s4} và \texttt{srl s4, s4, s4} có tương đương với hai lệnh trên không?
\end{enumerate}
\textbf{Hướng dẫn giải:}
\begin{itemize}
    \item \textbf{Giải câu a:}
    \begin{itemize}
        \item \texttt{and s0, s0, zero}: $\texttt{s0} \ \& \ 0 = 0$ (Đúng).
        \item \texttt{xor s0, s0, s0}: Hai số giống nhau XOR với nhau luôn bằng $0$ (Đúng).
        \item \texttt{xor s0, s0, zero} và \texttt{and s0, s0, s0} đều giữ nguyên giá trị \texttt{s0}. Vậy chọn \textbf{1 và 2}.
    \end{itemize}
    \item \textbf{Giải câu b:} Lệnh \texttt{add s4, s4, s4} ($\texttt{s4} + \texttt{s4} = 2 \times \texttt{s4}$) và lệnh \texttt{slli s4, s4, 1} (dịch trái 1 bit $= 2 \times \texttt{s4}$) tương đương nhau; còn lệnh nhân bình phương \texttt{mul s4, s4, s4} và dịch phải \texttt{srl} không tương đương.
\end{itemize}

\textbf{Bài tập 3.17 (Đề Giữa kỳ 2024.1 -- \#99ef6, \#e10492 \& Ôn tập 2025.2 -- Câu 6)}
\begin{enumerate}[label=\alph*)]
    \item \textbf{(Câu \#99ef6 \& Câu 6):} Đoạn chương trình nào tương đương với lệnh giả \texttt{li t0, 0x20241120}?\\
    A. \texttt{lui t0, 0x202411} $\rightarrow$ \texttt{ori t0, t0, 0x20} \quad B. \texttt{lui t0, 0x2024} $\rightarrow$ \texttt{addi t0, t0, 0x1120}\\
    C. \texttt{lui t0, 0x2024} $\rightarrow$ \texttt{ori t0, t0, 0x1120} \quad D. \texttt{lui t0, 0x20241} $\rightarrow$ \texttt{addi t0, t0, 0x120}
    \item \textbf{(Câu \#e10492):} Chọn đoạn lệnh RISC-V thực hiện câu lệnh C: \texttt{a += 0x12345678;} (với biến \texttt{a} nằm trong thanh ghi \texttt{s0}).
\end{enumerate}
\textbf{Hướng dẫn giải:}
\begin{itemize}
    \item \textbf{Giải câu a:} Số 32-bit \texttt{0x20241120} có 20 bit cao (5 chữ Hexa đầu) là \texttt{0x20241} và 12 bit thấp (3 chữ Hexa cuối) là \texttt{0x120} (dương). Do đó dùng: \texttt{lui t0, 0x20241} $\rightarrow$ \texttt{addi t0, t0, 0x120} $\rightarrow$ Chọn \textbf{D}.
    \item \textbf{Giải câu b:} Đầu tiên nạp hằng số \texttt{0x12345678} vào thanh ghi tạm \texttt{t0} bằng 2 lệnh \texttt{lui t0, 0x12345} và \texttt{addi t0, t0, 0x678}, sau đó cộng dồn vào \texttt{s0} bằng lệnh \texttt{add s0, s0, t0}.
\end{itemize}

\textbf{Bài tập 3.18 (Đề Giữa kỳ 2022.2 -- Đề 0009, 0008, 0007, 0006 Câu 9)}\\
Hãy lập trình hợp ngữ chuẩn RISC-V (không dùng lệnh giả \texttt{li}) để thực hiện các phép gán hằng số vào thanh ghi dưới đây:
\begin{enumerate}[label=\alph*)]
    \item \textbf{(Đề 0008):} $\texttt{t1} = \texttt{0x234}$ và $\texttt{t2} = \texttt{0x56789}$.
    \item \textbf{(Đề 0009):} $\texttt{t1} = \texttt{0xE7F}$ và $\texttt{t2} = \texttt{0x11062023}$.
    \item \textbf{(Đề 0006):} $\texttt{s5} = \texttt{0xCA03CA06}$ và $\texttt{s6} = \texttt{0x2B1A}$.
\end{enumerate}
\textbf{Hướng dẫn giải:}
\begin{itemize}
    \item \textbf{Giải câu a:}
    \begin{itemize}
        \item Với $\texttt{t1} = \texttt{0x234}$ (nhỏ hơn \texttt{0x800}, nằm gọn trong 12-bit dương): \texttt{ori t1, x0, 0x234} (hoặc \texttt{addi t1, x0, 0x234}).
        \item Với $\texttt{t2} = \texttt{0x00056789}$ (5 chữ đầu là \texttt{0x56}, 3 chữ cuối là $\texttt{0x789} < \texttt{0x800}$):
        \texttt{lui t2, 0x56} $\rightarrow$ \texttt{addi t2, t2, 0x789}.
    \end{itemize}
    \item \textbf{Giải câu b:}
    \begin{itemize}
        \item Với $\texttt{t1} = \texttt{0x00000E7F}$: Vì chữ Hexa thứ 3 là $\texttt{E} \ge 8$ (nếu dùng trực tiếp \texttt{addi/ori 0xE7F} sẽ bị mở rộng dấu âm thành \texttt{0xFFFFFE7F}), nên ta phải nạp \texttt{lui t1, 1} (\texttt{0x00001000}) rồi cộng số âm $\texttt{0xE7F} - \texttt{0x1000} = -385$: \texttt{lui t1, 1} $\rightarrow$ \texttt{addi t1, t1, -385}.
        \item Với $\texttt{t2} = \texttt{0x11062023}$ (3 chữ cuối là $\texttt{0x023} < \texttt{0x800}$):
        \texttt{lui t2, 0x11062} $\rightarrow$ \texttt{addi t2, t2, 0x023}.
    \end{itemize}
    \item \textbf{Giải câu c:}
    \begin{itemize}
        \item Với $\texttt{s5} = \texttt{0xCA03CA06}$: Chữ số thứ 3 từ phải sang là $\texttt{A} \ge 8$ ($\texttt{0xA06} = 2566 - 4096 = -1530$). Phải tăng 5 chữ đầu từ \texttt{0xCA03C} lên \texttt{0xCA03D}:
        \texttt{lui s5, 0xCA03D} $\rightarrow$ \texttt{addi s5, s5, -1530}.
        \item Với $\texttt{s6} = \texttt{0x00002B1A}$: Chữ thứ 3 từ phải sang là $\texttt{B} \ge 8$ ($\texttt{0xB1A} = 2842 - 4096 = -1254$). Tăng 5 chữ đầu từ \texttt{0x2} lên \texttt{0x3}:
        \texttt{lui s6, 0x3} $\rightarrow$ \texttt{addi s6, s6, -1254}.
    \end{itemize}
\end{itemize}

\textbf{Bài tập 3.19 (Đề Giữa kỳ 2024.1 -- \#1b55cf \& Ôn tập Cuối kỳ 2025.2 -- P1, Câu 24)}\\
Cho đoạn mã hợp ngữ RISC-V dưới đây. Hãy xác định giá trị của thanh ghi \texttt{s2} ở dạng số Hexa sau khi thực thi:
\begin{lstlisting}
    li   s0, 1
    lui  s1, 0xFFFFF
    ori  s1, s1, 2047
    add  s2, s0, s1
\end{lstlisting}
\textbf{Hướng dẫn giải:}
\begin{itemize}
    \item Lệnh 1: $\texttt{s0} = \texttt{0x00000001}$.
    \item Lệnh 2: $\texttt{s1} = \texttt{0xFFFFF000}$.
    \item Lệnh 3: Vì $2047 = \texttt{0x7FF}$ (bit 11 bằng 0, mở rộng 32-bit là \texttt{0x000007FF}), nên sau phép OR ta có $\texttt{s1} = \texttt{0xFFFFF7FF}$.
    \item Lệnh 4: $\texttt{s2} = \texttt{0xFFFFF7FF} + 1 = \mathbf{\texttt{0xFFFFF800}}$ (Nếu đề bài viết \texttt{lui s1, 0xFFFF} gồm 4 chữ \texttt{F} thì kết quả là \textbf{\texttt{0x0FFFF800}}).
\end{itemize}

\textbf{Bài tập 3.20 (Đề Giữa kỳ 2024.1 -- \#5e7591, Cuối kỳ 2025.1 -- \#48b64c \& Ôn tập 2025.2 -- Câu 64)}
\begin{enumerate}[label=\alph*)]
    \item \textbf{(Đề GK 2024.1 -- \#5e7591 \& Ôn tập 2025.2 Câu 64):} Tính giá trị thập phân của \texttt{s4} sau đoạn lệnh: \texttt{addi s1, zero, 87} $\rightarrow$ \texttt{addi s2, zero, 53} $\rightarrow$ \texttt{xor s3, s2, zero} $\rightarrow$ \texttt{addi s3, s3, 1} $\rightarrow$ \texttt{add s4, s1, s3}.
    \item \textbf{(Đề CK 2025.1 -- \#48b64c):} Tính giá trị thập phân của \texttt{s4} sau đoạn lệnh: \texttt{addi s1, zero, 80} $\rightarrow$ \texttt{addi s2, zero, 56} $\rightarrow$ \texttt{xor s3, s2, zero} $\rightarrow$ \texttt{addi s3, s3, 1} $\rightarrow$ \texttt{add s4, s1, s3}.
\end{enumerate}
\textbf{Hướng dẫn giải:}
\begin{itemize}
    \item Phép toán \texttt{xor s3, s2, zero} giữ nguyên giá trị của \texttt{s2} ($\texttt{s3} = \texttt{s2}$).
    \item \textbf{Câu a:} $\texttt{s3} = 53 + 1 = 54 \implies \texttt{s4} = 87 + 54 = \mathbf{141}$.
    \item \textbf{Câu b:} $\texttt{s3} = 56 + 1 = 57 \implies \texttt{s4} = 80 + 57 = \mathbf{137}$.
\end{itemize}

\textbf{Bài tập 3.21 (Đề Giữa kỳ 2018.2 -- Mã đề 204, Câu 9)}\\
Cho biết giá trị của thanh ghi \texttt{s1} sau khi thực hiện xong đoạn lệnh sau:
\begin{lstlisting}
    addi t1, x0, -1
    addi t2, zero, 5
    or   s1, t1, t2
\end{lstlisting}
\begin{enumerate}[label=\Alph*.]
    \item $5$
    \item \texttt{0xFFFFFFFF}
    \item $-5$
    \item $1$
\end{enumerate}
\textbf{Hướng dẫn giải:}
\begin{itemize}
    \item Số nguyên âm $-1$ trong biểu diễn bù 2 32-bit gồm toàn bộ 32 bit bằng $1$: $\texttt{t1} = \texttt{0xFFFFFFFF}$.
    \item Khi thực hiện phép OR với \texttt{0xFFFFFFFF}, tất cả các bit của kết quả đều bằng $1$: $\texttt{s1} = \texttt{0xFFFFFFFF}$.
    \item \textbf{Kết luận:} Đáp án đúng là \textbf{B (\texttt{0xFFFFFFFF})}.
\end{itemize}

\textbf{Bài tập 3.22 (Đề Cuối kỳ 2022.2 -- \#585834, 2024.1 -- \#d9f3db, 2024.2 \& Ôn tập 2025.2 -- Câu 54)}
\begin{enumerate}[label=\alph*)]
    \item \textbf{(Đề CK 2024.1 -- \#d9f3db \& Ôn tập 2025.2 Câu 54):} Cho biết nội dung thanh ghi \texttt{s1} theo số Hexa sau đoạn lệnh: \texttt{lui s0, 0x11223} $\rightarrow$ \texttt{ori s0, s0, 0x4AB} $\rightarrow$ \texttt{ori s1, s0, 0xFF}.
    \item \textbf{(Đề CK 2022.2 \& 2024.2 -- \#585834):} Nếu sau hai lệnh đầu thanh ghi $\texttt{s0} = \texttt{0x223389CD}$, thực hiện lệnh thứ ba là \texttt{andi s1, s0, 0xFF} thì nội dung thanh ghi \texttt{s1} bằng bao nhiêu?
\end{enumerate}
\textbf{Hướng dẫn giải:}
\begin{itemize}
    \item \textbf{Giải câu a:}
    \begin{itemize}
        \item \texttt{lui s0, 0x11223} $\implies \texttt{s0} = \texttt{0x11223000}$.
        \item \texttt{ori s0, s0, 0x4AB} $\implies \texttt{s0} = \texttt{0x112234AB}$.
        \item \texttt{ori s1, s0, 0xFF}: Phép OR với \texttt{0x000000FF} giữ nguyên 6 chữ Hexa đầu và bật 2 chữ cuối thành \texttt{FF} $\implies \mathbf{\texttt{s1 = 0x112234FF}}$.
    \end{itemize}
    \item \textbf{Giải câu b:} Lệnh \texttt{andi s1, s0, 0xFF} thực hiện phép AND với \texttt{0x000000FF}, xóa 6 chữ số Hexa đầu về \texttt{0} và giữ lại 2 chữ số cuối: $\mathbf{\texttt{s1 = 0x000000CD}}$.
\end{itemize}

\textbf{Bài tập 3.23 (Đề Cuối kỳ 2020.2, 2022.2 -- \#5070f9, 2024.2, Giữa kỳ 2023.2 \& Giữa kỳ 2025.2 -- Đề 0012 Câu 108)}\\
Hãy xác định kết quả dạng Hexa cho các bài tập phép toán logic dưới đây:
\begin{enumerate}[label=\alph*)]
    \item \textbf{(Đề CK 2020.2 -- Trang 3):} Cho $\texttt{s0} = \texttt{0x1234ABCD}$ và $\texttt{s1} = \texttt{0xFFFF0000}$. Tìm \texttt{s2} sau lệnh \texttt{xor s2, s0, s1}.
    \item \textbf{(Đề CK 2020.2 -- Trang 7):} Cho $\texttt{t0} = \texttt{0xAAAAAAAA}$ và $\texttt{t1} = \texttt{0x0000FFC0}$. Tìm \texttt{t3} sau lệnh \texttt{and t3, t0, t1}.
    \item \textbf{(Đề CK 2022.2, 2024.2 -- \#5070f9 \& Ôn tập 2025.2 Câu 64):} Cho $\texttt{t0} = \texttt{0x5A6B7C8D}$ và $\texttt{t1} = \texttt{0x000FFF00}$. Tìm \texttt{t2, t3} sau 2 lệnh: \texttt{xor t2, t0, t1} và \texttt{xor t3, t2, t1}.
    \item \textbf{(Đề GK 2025.2 -- Đề 0012, Câu 108):} Cho $\texttt{t0} = \texttt{0x35B9D2E4}$ và $\texttt{t1} = \texttt{0x0F0F0F0F}$. Tìm \texttt{t2} sau lệnh \texttt{xor t2, t0, t1}.
    \item \textbf{(Đề GK 2023.2 -- Đề 2, Câu 4):} Cho $\texttt{t0} = \texttt{0xA1B2C3D4}$. Viết một lệnh logic để biến đổi \texttt{t0} thành \texttt{0x286CB0F5} (giả sử thanh ghi mặt nạ \texttt{t1} đã được nạp sẵn hằng số cần thiết).
\end{enumerate}
\textbf{Hướng dẫn giải:}
\begin{itemize}
    \item \textbf{Giải câu a:} XOR với \texttt{0xFFFF0000} sẽ đảo 4 chữ số đầu \texttt{1234} (lấy $15$ trừ từng chữ số: $15-1=\texttt{E}, 15-2=\texttt{D}, 15-3=\texttt{C}, 15-4=\texttt{B}$) và giữ nguyên 4 chữ cuối \texttt{ABCD} $\implies \mathbf{\texttt{s2 = 0xEDCBABCD}}$.
    \item \textbf{Giải câu b:} Trong $\texttt{t0} = \texttt{0xAAAAAAAA}$ ($\texttt{A} = 1010_2$) và $\texttt{t1} = \texttt{0x0000FFC0}$ ($\texttt{C} = 1100_2$), lấy từng chữ số AND với nhau: $1010_2 \ \& \ 1100_2 = 1000_2 = \texttt{8}$ $\implies \mathbf{\texttt{t3 = 0x0000AA80}}$.
    \item \textbf{Giải câu c:} Đảo 3 chữ giữa \texttt{B7C} của \texttt{0x5A6B7C8D} được $\mathbf{\texttt{t2 = 0x5A64838D}}$. Khi lấy \texttt{t2} XOR tiếp với \texttt{t1} sẽ hoàn nguyên về giá trị của \texttt{t0}: $\mathbf{\texttt{t3 = 0x5A6B7C8D}}$.
    \item \textbf{Giải câu d:} Đảo các chữ số ở vị trí chẵn (\texttt{5, 9, 2, 4} $\rightarrow \texttt{A, 6, D, B}$) và giữ nguyên các chữ số ở vị trí lẻ (\texttt{3, B, D, E}) $\implies \mathbf{\texttt{t2 = 0x3AB6DDEB}}$.
    \item \textbf{Giải câu e:} Lấy $\texttt{0xA1B2C3D4} \oplus \texttt{0x286CB0F5} = \texttt{0x89DE7321}$. Với $\texttt{t1} = \texttt{0x89DE7321}$, lệnh cần viết là: \textbf{\texttt{xor t0, t0, t1}}.
\end{itemize}

\textbf{Bài tập 3.24 (Đề Giữa kỳ 2022.2 -- Đề 0008 Câu 5 \& Đề 0007 Câu 5)}\\
Thực hiện các phép toán logic 32-bit sau và biểu diễn kết quả ở hệ 16 (Hexa):
\begin{enumerate}[label=\alph*)]
    \item \texttt{0xB8C32AC9 and 0x00FFFF00} \quad và \quad \texttt{0xCF072D76 or 0xF0F0F0F0}.
    \item \texttt{0x0000FFFF and 0x079CC0D1} \quad và \quad \texttt{0xAC4F0130 or 0xFFFF0000}.
\end{enumerate}
\textbf{Hướng dẫn giải:}
\begin{itemize}
    \item \textbf{Giải câu a:}
    \begin{itemize}
        \item \texttt{0xB8C32AC9 and 0x00FFFF00} $= \mathbf{\texttt{0x00C32A00}}$.
        \item \texttt{0xCF072D76 or 0xF0F0F0F0} $= \mathbf{\texttt{0xFFF7FDF6}}$.
    \end{itemize}
    \item \textbf{Giải câu b:}
    \begin{itemize}
        \item \texttt{0x0000FFFF and 0x079CC0D1} $= \mathbf{\texttt{0x0000C0D1}}$.
        \item \texttt{0xAC4F0130 or 0xFFFF0000} $= \mathbf{\texttt{0xFFFF0130}}$.
    \end{itemize}
\end{itemize}

\textbf{Bài tập 3.25 (Đề Cuối kỳ 2020.2 -- Trang 6 \& Đề Cuối kỳ 2021.1 -- Trang 5)}\\
Cho đoạn mã nguồn C thực hiện thao tác xử lý trên bit: \texttt{X = (Y \& 0x0F) << 4;}. Giả thiết hai biến số nguyên \texttt{X} và \texttt{Y} lần lượt được phân bổ tại các thanh ghi \texttt{s2} và \texttt{s1}. Hãy chỉ ra chuỗi lệnh hợp ngữ RISC-V dịch chính xác câu lệnh C trên?
\begin{enumerate}[label=\Alph*.]
    \item \texttt{and t1, s1, 0xF} $\rightarrow$ \texttt{slli s2, t1, 16}
    \item \texttt{andi t1, s1, 15} $\rightarrow$ \texttt{slli s2, t1, 4}
    \item \texttt{andi t1, s1, 0xFFFF} $\rightarrow$ \texttt{slli s2, t1, 2}
    \item \texttt{slli t1, s1, 15} $\rightarrow$ \texttt{andi s2, t1, 4}
\end{enumerate}
\textbf{Hướng dẫn giải:}
\begin{itemize}
    \item Bước 1: Tính biểu thức trong ngoặc \texttt{Y \& 0x0F}. Vì $\texttt{0x0F} = 15_{(10)}$ là hằng số tức thời nên phải dùng lệnh \texttt{andi}: \texttt{andi t1, s1, 15}.
    \item Bước 2: Dịch trái kết quả đi $4\text{ bit}$ (\texttt{<< 4}) và lưu vào \texttt{s2} (\texttt{X}) bằng lệnh \texttt{slli}: \texttt{slli s2, t1, 4}.
    \item \textbf{Kết luận:} Đáp án đúng là \textbf{B}.
\end{itemize}

\textbf{Bài tập 3.26 (Đề Giữa kỳ 2023.2 -- Đề 1 Câu 4 \& Đề Cuối kỳ 2021.2 -- Mã đề B8 Câu 2.2)}
\begin{enumerate}[label=\alph*)]
    \item \textbf{(Đề GK 2023.2):} Cho $\texttt{t0} = \texttt{0xAAAAAAAA}, \texttt{t1} = \texttt{0x00FF0000}$. Tìm giá trị Hexa của \texttt{t2} và \texttt{t3} sau hai lệnh: \texttt{srli t2, t0, 3} và \texttt{or t3, t2, t1}.
    \item \textbf{(Đề CK 2021.2 -- Mã đề B8):} Cho $\texttt{t0} = \texttt{0xBEADFEED}, \texttt{t1} = \texttt{0xDEADFADE}$. Tìm giá trị Hexa của \texttt{t2} sau hai lệnh: \texttt{slli t2, t0, 4} và \texttt{or t2, t2, t1}.
\end{enumerate}
\textbf{Hướng dẫn giải:}
\begin{itemize}
    \item \textbf{Giải câu a:}
    \begin{itemize}
        \item Đổi $\texttt{t0} = \texttt{0xAAAAAAAA} = 1010\,1010\,\dots\,1010_{(2)}$.
        \item Khi dịch phải logic $3\text{ bit}$ (\texttt{srli t2, t0, 3}), 3 bit \texttt{000} chèn vào bên trái ghép với bit \texttt{1} đầu tiên thành \texttt{0001} (\texttt{1}), các nhóm 4 bit tiếp theo đều trở thành \texttt{0101} (\texttt{5}) $\implies \mathbf{\texttt{t2 = 0x15555555}}$.
        \item Thực hiện \texttt{or t3, t2, t1} với $\texttt{t1} = \texttt{0x00FF0000}$: $\mathbf{\texttt{t3 = 0x15FF5555}}$.
    \end{itemize}
    \item \textbf{Giải câu b:}
    \begin{itemize}
        \item \texttt{slli t2, t0, 4} (dịch trái 1 chữ số Hexa): $\texttt{t2} = \texttt{0xEADFEED0}$.
        \item \texttt{or t2, t2, t1}: Lấy \texttt{0xEADFEED0} OR từng chữ số với \texttt{0xDEADFADE} được: $\mathbf{\texttt{t2 = 0xFEFFFEFE}}$.
    \end{itemize}
\end{itemize}

\subsection{Nhóm 3: Lệnh So sánh (\texttt{slt, slti, sltu, blt, bltu}) và Rẽ nhánh}

\textbf{Bài tập 3.27 (Đề Cuối kỳ 2021.2 -- Mã đề B8, Câu 2.1)}\\
Cho câu lệnh hợp ngữ: \texttt{slti t2, t0, X}. Giả sử thanh ghi \texttt{t0} có giá trị \texttt{0xAD100002}. Với những giá trị nào của hằng số \texttt{X} (viết dưới dạng thập phân) thì thanh ghi \texttt{t2} sẽ có giá trị bằng $1$?

\textbf{Hướng dẫn giải:}
\begin{itemize}
    \item Lệnh \texttt{slti} thực hiện so sánh số nguyên \textbf{có dấu} giữa thanh ghi \texttt{t0} (32-bit) và hằng số tức thời \texttt{X} (12-bit có dấu trong RISC-V, có dải giá trị hợp lệ là $[-2048, +2047]$).
    \item Xét $\texttt{t0} = \texttt{0xAD100002}$: Chữ số Hexa đầu tiên là $\texttt{A} = 1010_{(2)}$ có bit dấu bằng $1$, nên \texttt{t0} là một số nguyên âm rất nhỏ ($\approx -1.39 \times 10^9$), luôn nhỏ hơn $-2048$.
    \item Do đó, với \textbf{mọi giá trị hợp lệ của hằng số \texttt{X} thuộc đoạn $[-2048, +2047]$} (hoặc $[-32768, +32767]$ nếu xét chuẩn 16-bit), điều kiện $\texttt{t0} < \texttt{X}$ luôn luôn đúng và làm cho $\texttt{t2} = 1$.
\end{itemize}

\textbf{Bài tập 3.28 (Đề Cuối kỳ 2025.1 -- \#a91b41, Giữa kỳ 2024.1 -- \#8d33e2, \#484325 \& Ôn tập 2025.2 -- Câu 23, 61)}
\begin{enumerate}[label=\alph*)]
    \item \textbf{(Đề CK 2025.1 -- \#a91b41):} Tìm giá trị thập phân của \texttt{t1} sau khi chạy đoạn mã:
\begin{lstlisting}
    li   s1, -15
    li   s2, 17
    slt  t0, s1, s2
    bne  t0, zero, else
    addi t1, s1, 14
    j    endif
else:
    addi t1, s2, -1
endif: nop
\end{lstlisting}
    \item \textbf{(Ôn tập 2025.2 -- Câu 23):} Cho $\texttt{x22} = 10, \texttt{x23} = 5, \texttt{x24} = 0$. Chạy: \texttt{slt x24, x22, x23} $\rightarrow$ \texttt{bne x24, x0, L1} $\rightarrow$ \texttt{addi x24, x24, 10} $\rightarrow$ \texttt{L1: addi x24, x24, 1}. Tìm giá trị cuối của \texttt{x24}?
    \item \textbf{(Ôn tập 2025.2 -- Câu 61):} Cho $\texttt{x5} = \texttt{0x0000000A}$ ($10$) và $\texttt{x6} = \texttt{0x0000000A}$. Chạy: \texttt{addi x7, x0, 100} $\rightarrow$ \texttt{bne x5, x6, Skip} $\rightarrow$ \texttt{addi x7, x7, 50} $\rightarrow$ \texttt{Skip: addi x7, x7, 1}. Tìm \texttt{x7}?
    \item \textbf{(Đề GK 2024.1 -- \#8d33e2 \& \#484325):}
    \begin{itemize}
        \item \textit{Bài 1:} $\texttt{t0} = -2, \texttt{t1} = 2$. Chạy \texttt{blt t0, t1, else}. Nhánh không nhảy tính $\texttt{a0} = 1 + 5$; nhánh \texttt{else} tính $\texttt{a0} = 2 + 7$. Tìm \texttt{a0}?
        \item \textit{Bài 2:} $\texttt{t0} = 10, \texttt{t1} = 10$. Chạy \texttt{bltu t0, t1, else}. Nhánh không nhảy tính $\texttt{a0} = 10 + 5$; nhánh \texttt{else} tính $\texttt{a0} = 20 + 7$. Tìm \texttt{a0}?
    \end{itemize}
\end{enumerate}
\textbf{Hướng dẫn giải:}
\begin{itemize}
    \item \textbf{Giải câu a:} Vì $-15 < 17$ là Đúng nên $\texttt{t0} = 1$. Lệnh \texttt{bne t0, zero, else} thỏa mãn ($\texttt{t0} \neq 0$) nên nhảy tới nhãn \texttt{else} thực hiện: $\texttt{t1} = \texttt{s2} - 1 = 17 - 1 = \mathbf{16}$.
    \item \textbf{Giải câu b:} Vì $10 < 5$ là Sai nên $\texttt{x24} = 0$. Lệnh \texttt{bne} không nhảy, thực hiện cộng $10$ ($\texttt{x24} = 10$) rồi trôi tiếp qua nhãn \texttt{L1} cộng thêm $1$ $\implies \mathbf{\texttt{x24 = 11}}$.
    \item \textbf{Giải câu c:} Vì $\texttt{x5} == \texttt{x6}$ nên lệnh \texttt{bne} không nhảy. Thanh ghi \texttt{x7} nhận giá trị: $100 + 50 + 1 = \mathbf{151}$.
    \item \textbf{Giải câu d:}
    \begin{itemize}
        \item \textit{Bài 1:} $-2 < 2$ là Đúng $\implies$ Nhảy tới \texttt{else}: $\texttt{a0} = 2 + 7 = \mathbf{9}$.
        \item \textit{Bài 2:} $10 < 10$ là Sai $\implies$ Không nhảy: $\texttt{a0} = 10 + 5 = \mathbf{15}$.
    \end{itemize}
\end{itemize}

\subsection{Nhóm 4: Truy cập bộ nhớ (\texttt{lw, lh, lb, lbu, sw}) và Thứ tự Little-Endian}

\textbf{Bài tập 3.29 (Đề Cuối kỳ 2022.2 \& Cuối kỳ 2024.2 -- \#56e53f)}\\
Giả sử thanh ghi \texttt{s0} đang chứa giá trị địa chỉ \texttt{0x10021000}. Lựa chọn chức năng đúng của câu lệnh hợp ngữ \texttt{lw t0, 8(s0)} (hoặc \texttt{sw t0, 8(s0)}):
\begin{enumerate}[label=\Alph*.]
    \item Ghi giá trị hằng số \texttt{0x10021008} vào thanh ghi \texttt{t0}.
    \item Nạp nội dung word nhớ ở địa chỉ \texttt{0x10021008} đưa vào thanh ghi \texttt{t0}.
    \item Lấy nội dung thanh ghi \texttt{t0} cộng thêm 8 rồi lưu ra word nhớ ở địa chỉ \texttt{0x10021000}.
    \item Lưu nội dung thanh ghi \texttt{t0} ra word nhớ ở địa chỉ \texttt{0x10021008}.
\end{enumerate}
\textbf{Hướng dẫn giải:}
\begin{itemize}
    \item Địa chỉ hiệu dụng của toán hạng bộ nhớ là $\texttt{s0} + 8 = \texttt{0x10021000} + 8 = \texttt{0x10021008}$.
    \item Nếu đề bài hỏi lệnh \textbf{\texttt{lw t0, 8(s0)}} (load word): Chức năng là \textbf{B (Nạp nội dung word nhớ ở địa chỉ \texttt{0x10021008} đưa vào thanh ghi \texttt{t0})}.
    \item Nếu đề bài hỏi lệnh \textbf{\texttt{sw t0, 8(s0)}} (store word): Chức năng là \textbf{D (Lưu nội dung thanh ghi \texttt{t0} ra word nhớ ở địa chỉ \texttt{0x10021008})}.
\end{itemize}

\textbf{Bài tập 3.30 (Đề Cuối kỳ 2025.1 -- \#1f833e)}\\
Đoạn lệnh nào dưới đây dùng để nạp một word từ địa chỉ bộ nhớ \texttt{0x10010004} vào thanh ghi \texttt{t1} (mà không dùng lệnh giả)?
\begin{enumerate}[label=\Alph*.]
    \item \texttt{li t0, 0x10010004} $\rightarrow$ \texttt{lw t1, zero(t0)}
    \item \texttt{lui t0, 0x1001} $\rightarrow$ \texttt{addi t0, t0, 0x0004} $\rightarrow$ \texttt{lw t1, 0(t0)}
    \item \texttt{lw t1, 0x10010004(zero)}
    \item \texttt{lui t0, 0x10010} $\rightarrow$ \texttt{addi t0, t0, 0x004} $\rightarrow$ \texttt{lw t1, 0(t0)}
\end{enumerate}
\textbf{Hướng dẫn giải:}
\begin{itemize}
    \item Địa chỉ 32-bit \texttt{0x10010004} có 20 bit cao (5 chữ số Hexa đầu) là \texttt{0x10010} và 12 bit thấp (3 chữ số Hexa cuối) là \texttt{0x004}.
    \item Lệnh \texttt{lui t0, 0x10010} tạo ra \texttt{0x10010000}; lệnh \texttt{addi t0, t0, 0x004} cộng thêm $4$ tạo thành địa chỉ \texttt{0x10010004}; sau đó dùng \texttt{lw t1, 0(t0)} đọc word nhớ lên \texttt{t1}. (Phương án B sai vì \texttt{0x1001} thiếu một chữ số 0).
    \item \textbf{Kết luận:} Đáp án đúng là \textbf{D}.
\end{itemize}

\textbf{Bài tập 3.31 (Đề Giữa kỳ 2018.2 -- Mã đề 198 Câu 20, Mã đề 199 Câu 6 \& Mã đề 204 Câu 1)}\\
Biến \texttt{X} thuộc kiểu word 32-bit nằm ở địa chỉ \texttt{0x0FB028C4} (hoặc \texttt{0xAFB028C4}). Đoạn chương trình nào sau đây tương đương với câu lệnh trong ngôn ngữ C: \texttt{X = X + 1;}?
\begin{enumerate}[label=\Alph*.]
    \item \texttt{addi t3, x0, 0x0FB028C4} $\rightarrow$ \texttt{lw s3, 0(t3)} $\rightarrow$ \texttt{addi s3, s3, 1} $\rightarrow$ \texttt{sw s3, 0(t3)}
    \item \texttt{lui t3, 0x0FB03} $\rightarrow$ \texttt{addi t3, t3, -1852} (\textit{hoặc \texttt{lui} + \texttt{ori}}) $\rightarrow$ \texttt{lw s3, 0(t3)} $\rightarrow$ \texttt{addi s3, s3, 1} $\rightarrow$ \texttt{sw s3, 0(t3)}
    \item \texttt{lui t3, 0x28C4} $\rightarrow$ \texttt{ori t3, t3, 0x0FB0} $\rightarrow$ \texttt{lw s3, 0(t3)} $\rightarrow$ \texttt{addi s3, s3, 1} $\rightarrow$ \texttt{sw s3, 0(t3)}
    \item \texttt{lui t3, 0x0FB03} $\rightarrow$ \texttt{addi t3, t3, -1852} $\rightarrow$ \texttt{sw s3, 0(t3)} $\rightarrow$ \texttt{addi s3, s3, 1} $\rightarrow$ \texttt{lw s3, 0(t3)}
\end{enumerate}
\textbf{Hướng dẫn giải:}
\begin{itemize}
    \item Quy trình chuẩn để tăng biến \texttt{X} trong bộ nhớ lên 1 đơn vị gồm 3 bước:
    \begin{enumerate}
        \item Nạp phần cao của địa chỉ bằng \texttt{lui}, ghép phần thấp để tạo địa chỉ \texttt{0x0FB028C4} trong thanh ghi con trỏ \texttt{t3} (Loại A vì \texttt{addi} không nạp được hằng số 32-bit; loại C vì nạp ngược phần thấp lên phần cao).
        \item Đọc (\texttt{lw}) giá trị của \texttt{X} từ bộ nhớ vào thanh ghi \texttt{s3}, cộng thêm 1 (\texttt{addi s3, s3, 1}).
        \item Ghi (\texttt{sw}) kết quả từ \texttt{s3} trả lại bộ nhớ (Loại D vì thực hiện \texttt{sw} trước khi \texttt{lw}).
    \end{enumerate}
    \item \textbf{Kết luận:} Đáp án đúng là \textbf{B}.
\end{itemize}

\textbf{Bài tập 3.32 (Đề Cuối kỳ 2022.2 -- \#757298, 2024.1 -- \#c99f87, 2024.2 \& Ôn tập 2025.2 -- Câu 20)}\\
Biến \texttt{x} thuộc kiểu word 32-bit nằm ở địa chỉ \texttt{0x10010200}. Đoạn chương trình RISC-V nào dưới đây tương đương với câu lệnh ngôn ngữ C: \texttt{x = x \& 15;}?
\begin{enumerate}[label=\Alph*.]
    \item \texttt{andi t3, t3, 15}
    \item \texttt{lui t3, 0x10200} $\rightarrow$ \texttt{addi t3, t3, 0x100} $\rightarrow$ \texttt{lw s3, 0(t3)} $\rightarrow$ \texttt{andi s3, s3, 15} $\rightarrow$ \texttt{sw s3, 0(t3)}
    \item \texttt{lui t3, 0x10010} $\rightarrow$ \texttt{addi t3, t3, 0x200} $\rightarrow$ \texttt{lw s3, 0(t3)} $\rightarrow$ \texttt{andi s3, s3, 15} $\rightarrow$ \texttt{sw s3, 0(t3)}
    \item \texttt{lui t3, 0x10010} $\rightarrow$ \texttt{addi t3, t3, 0x200} $\rightarrow$ \texttt{sw s3, 0(t3)} $\rightarrow$ \texttt{andi s3, s3, 15} $\rightarrow$ \texttt{lw s3, 0(t3)}
\end{enumerate}
\textbf{Hướng dẫn giải:}
\begin{itemize}
    \item Tạo địa chỉ \texttt{0x10010200} trong \texttt{t3} bằng cặp lệnh \texttt{lui t3, 0x10010} và \texttt{addi t3, t3, 0x200}.
    \item Đọc \texttt{x} từ bộ nhớ lên \texttt{s3} (\texttt{lw s3, 0(t3)}), thực hiện AND với 15 (\texttt{andi s3, s3, 15}), rồi ghi trả lại bộ nhớ (\texttt{sw s3, 0(t3)}).
    \item \textbf{Kết luận:} Đáp án đúng là \textbf{C}.
\end{itemize}

\textbf{Bài tập 3.33 (Đề Giữa kỳ 2018.2 -- Mã đề 198 Câu 11, Mã đề 204 Câu 12 \& Đề Giữa kỳ 2019.2 -- Mã đề 243 Câu 19, Mã đề 244 Câu 3)}
\begin{enumerate}[label=\alph*)]
    \item \textbf{(Mã đề 243 -- Câu 19):} Trong bộ nhớ chính có 4 byte theo thứ tự địa chỉ tăng dần là \texttt{1A, 2B, 3C, 4D} được lưu theo thứ tự đầu nhỏ (Little-Endian). Hỏi số nguyên 32-bit tương ứng bằng bao nhiêu?\\
    A. \texttt{0x4321DCBA} \quad B. \texttt{0x4D3C2B1A} \quad C. \texttt{0x1A2B3C4D} \quad D. \texttt{0xDCBA4321}
    \item \textbf{(Mã đề 244 -- Câu 3):} Trong bộ nhớ chính có 4 byte theo thứ tự địa chỉ tăng dần là \texttt{4D, 3C, 2B, 1A} được lưu theo thứ tự đầu nhỏ. Hỏi số nguyên 32-bit tương ứng bằng bao nhiêu?\\
    A. \texttt{0x4D3C2B1A} \quad B. \texttt{0x1A2B3C4D} \quad C. \texttt{0x4321DCBA} \quad D. \texttt{0xDCBA4321}
    \item \textbf{(Mã đề 204 -- Câu 12):} Số 32-bit \texttt{0xC81A56} (tức \texttt{0x00C81A56}) khi được lưu trữ vào bộ nhớ chính theo thứ tự đầu nhỏ thì trình tự các byte theo địa chỉ tăng dần như thế nào?\\
    A. \texttt{56, 1A, C8} \quad B. \texttt{56, 1A, C8, 00} \quad C. \texttt{00, C8, 1A, 56} \quad D. \texttt{C8, 1A, 56}
    \item \textbf{(Mã đề 198 -- Câu 11):} Một số 32-bit được lưu ở 4 ngăn nhớ liên tiếp từ địa chỉ 612 đến 615 lần lượt chứa các giá trị Hexa: \texttt{E} (tức \texttt{0E}), \texttt{48}, \texttt{21}, \texttt{A} (tức \texttt{0A}). Giá trị của số đó ở hệ 16 bằng bao nhiêu?\\
    A. \texttt{E4821A} nếu đầu to \quad B. \texttt{A4821E} nếu đầu to \quad C. \texttt{0A48210E} nếu đầu nhỏ \quad D. \texttt{0E48210A} nếu đầu to
\end{enumerate}
\textbf{Hướng dẫn giải:}
\begin{itemize}
    \item \textbf{Giải câu a:} Theo Little-Endian, ghép ngược từ byte ở địa chỉ lớn nhất về địa chỉ nhỏ nhất: \textbf{\texttt{0x4D3C2B1A}} $\rightarrow$ Chọn \textbf{B}.
    \item \textbf{Giải câu b:} Ghép ngược từ phải sang trái: \textbf{\texttt{0x1A2B3C4D}} $\rightarrow$ Chọn \textbf{B}.
    \item \textbf{Giải câu c:} Số 32-bit \texttt{0xC81A56} viết đủ 8 chữ số Hexa (4 byte) là \texttt{0x00C81A56}. Lưu theo Little-Endian từ byte thấp nhất đến byte cao nhất là: \textbf{\texttt{56, 1A, C8, 00}} $\rightarrow$ Chọn \textbf{B}.
    \item \textbf{Giải câu d:} Mỗi ngăn nhớ 1 byte phải biểu diễn bằng 2 chữ số Hexa: tại 612 là \texttt{0E}, 613 là \texttt{48}, 614 là \texttt{21}, 615 là \texttt{0A}. Nếu lưu theo thứ tự đầu to (Big-Endian) thì ghép xuôi từ 612 đến 615 được \textbf{\texttt{0E48210A}} (nếu đầu nhỏ thì ghép ngược là \texttt{0A21480E}) $\rightarrow$ Chọn \textbf{D}.
\end{itemize}

\textbf{Bài tập 3.34 (Đề Giữa kỳ 2023.2 -- Đề 1 Câu 3, Đề 2 Câu 1 \& Giữa kỳ 2022.2 -- Đề 0008, 0007 Câu 2)}
\begin{enumerate}[label=\alph*)]
    \item \textbf{(Đề GK 2023.2 -- Đề 1, Câu 3):} Bộ nhớ chính lưu theo kiểu đầu nhỏ (Little-Endian). Thanh ghi \texttt{s6} chứa địa chỉ một ngăn nhớ dữ liệu. Sau khi thực hiện: \texttt{lui s0, 0x1B32C} $\rightarrow$ \texttt{addi s0, s0, 0x2E5} (để $\texttt{s0} = \texttt{0x1B32C2E5}$) $\rightarrow$ \texttt{sw s0, 0(s6)} $\rightarrow$ \texttt{lb s1, 1(s6)}. Hãy chỉ ra nội dung của \texttt{s0} và \texttt{s1} theo số Hexa.
    \item \textbf{(Đề GK 2023.2 -- Đề 2, Câu 1):} Thực hiện đoạn lệnh: \texttt{lui s0, 0x4567C} $\rightarrow$ \texttt{addi s0, s0, -802} (để $\texttt{s0} = \texttt{0x4567BCDE}$) $\rightarrow$ \texttt{addi s1, zero, -16} $\rightarrow$ \texttt{and s2, s0, s1}. Hãy chỉ ra nội dung của \texttt{s1} và \texttt{s2} theo số Hexa.
    \item \textbf{(Đề GK 2022.2 -- Đề 0008 \& 0007, Câu 2):}
    \begin{itemize}
        \item \textit{Đề 0008:} Từ địa chỉ \texttt{0x100110F8} đến \texttt{0x100110FB} lưu 4 byte theo Little-Endian là: \texttt{0x5C, 0x52, 0x0E, 0xEF}. Giá trị biến 32-bit này bằng bao nhiêu? Viết lệnh nạp biến này vào \texttt{s1}.
        \item \textit{Đề 0007:} Biến 32-bit có giá trị \texttt{0x4CDC001A} lưu ở địa chỉ \texttt{0x1001272C}. Hãy điền các byte vào 4 địa chỉ từ \texttt{0x1001272C} đến \texttt{0x1001272F}.
    \end{itemize}
\end{enumerate}
\textbf{Hướng dẫn giải:}
\begin{itemize}
    \item \textbf{Giải câu a:}
    \begin{itemize}
        \item Thanh ghi $\mathbf{\texttt{s0 = 0x1B32C2E5}}$.
        \item Khi \texttt{sw s0, 0(s6)} ra bộ nhớ Little-Endian: tại \texttt{0(s6)} là \texttt{0xE5}, tại \texttt{1(s6)} là \texttt{0xC2}, tại \texttt{2(s6)} là \texttt{0x32}, tại \texttt{3(s6)} là \texttt{0x1B}.
        \item Lệnh \texttt{lb s1, 1(s6)} đọc byte \texttt{0xC2}. Vì $\texttt{C} \ge 8$ (số âm), mở rộng dấu thành: $\mathbf{\texttt{s1 = 0xFFFFFFC2}}$.
    \end{itemize}
    \item \textbf{Giải câu b:}
    \begin{itemize}
        \item Số $-16$ đổi sang mã bù 2 32-bit là: $\mathbf{\texttt{s1 = 0xFFFFFFF0}}$.
        \item Lệnh \texttt{and s2, s0, s1} xóa chữ số Hexa cuối cùng của \texttt{0x4567BCDE} về 0: $\mathbf{\texttt{s2 = 0x4567BCD0}}$.
    \end{itemize}
    \item \textbf{Giải câu c:}
    \begin{itemize}
        \item \textit{Đề 0008:} Ghép ngược từ \texttt{0x100110FB} về \texttt{0x100110F8} được \textbf{\texttt{0xEF0E525C}}. Lệnh nạp vào \texttt{s1}: \texttt{lui t0, 0x10011} $\rightarrow$ \texttt{lw s1, 0x0F8(t0)}.
        \item \textit{Đề 0007:} Tại \texttt{0x1001272C} là \textbf{\texttt{0x1A}}; tại \texttt{0x1001272D} là \textbf{\texttt{0x00}}; tại \texttt{0x1001272E} là \textbf{\texttt{0xDC}}; tại \texttt{0x1001272F} là \textbf{\texttt{0x4C}}.
    \end{itemize}
\end{itemize}

\textbf{Bài tập 3.35 (Đề Giữa kỳ 2024.1 -- \#6a94a2 \& Ôn tập Cuối kỳ 2025.2 -- P1, Câu 49)}\\
Địa chỉ của toán hạng đích trong lệnh \texttt{sw} bên dưới là gì (viết ở dạng số Hexa)?
\begin{lstlisting}
    li   s0, 0x1001BCDE
    addi s1, s0, 0x321
    sw   s0, 1(s1)
\end{lstlisting}
\textbf{Hướng dẫn giải:}
\begin{itemize}
    \item Lệnh 1: $\texttt{s0} = \texttt{0x1001BCDE}$.
    \item Lệnh 2: $\texttt{s1} = \texttt{0x1001BCDE} + \texttt{0x321} = \texttt{0x1001BFFF}$ (vì $\texttt{0xCDE} + \texttt{0x321} = \texttt{0xFFF}$).
    \item Lệnh 3 (\texttt{sw s0, 1(s1)}): Địa chỉ của toán hạng đích là:
    $$\text{Địa chỉ} = \texttt{s1} + 1 = \texttt{0x1001BFFF} + 1 = \mathbf{\texttt{0x1001C000}}$$
    (Nếu đề bài cắt ngắn hằng số thành \texttt{li s0, 0x1001BC} và \texttt{addi s1, s0, 0x32} như ở ảnh chụp bị che của Câu 49 Ôn tập 2025.2 thì địa chỉ là $\texttt{0x1001BC} + \texttt{0x32} + 1 = \mathbf{\texttt{0x1001EF}}$).
\end{itemize}
\subsection{Nhóm 5: Chuyển đổi giữa Ngôn ngữ C và Hợp ngữ RISC-V (Biểu thức, Mảng, Con trỏ, Rẽ nhánh)}

\textbf{Bài tập 3.36 (Đề Giữa kỳ 2018.2 -- Mã đề 198 Câu 13, Mã đề 199 Câu 5 \& Mã đề 204 Câu 2)}\\
Biến \texttt{X, Y, Z} nằm ở các thanh ghi \texttt{s1, s2, s3} tương ứng. Hỏi mỗi đoạn mã hợp ngữ dưới đây thực hiện phép toán nào?
\begin{enumerate}[label=\alph*)]
    \item \textbf{(Mã đề 198 \& 199):} \texttt{addi t1, s1, 2} $\rightarrow$ \texttt{slli s2, s2, 2} $\rightarrow$ \texttt{sub s3, t1, s2}.\\
    A. $Z = 2 - Y \times 2$ \qquad B. $Z = X + 2 - Y \times 4$\\
    C. $Z = X \times 2 + Y / 2$ \qquad D. $Z = X - Y \times 4$
    \item \textbf{(Mã đề 204):} \texttt{add t1, s1, s2} $\rightarrow$ \texttt{srli t2, s2, 1} $\rightarrow$ \texttt{add s3, t1, t2}.\\
    A. $Z = X + Y \times 2$ \qquad B. $Z = X + Y / 2$\\
    C. $Z = X / 2 + Y$ \qquad D. $Z = X + Y + Y / 2$
\end{enumerate}
\textbf{Hướng dẫn giải:}
\begin{itemize}
    \item \textbf{Giải câu a:}
    \begin{itemize}
        \item \texttt{addi t1, s1, 2}: $\texttt{t1} = X + 2$.
        \item \texttt{slli s2, s2, 2}: Dịch trái $2\text{ bit}$ tương đương nhân với $2^2 = 4 \implies \texttt{s2} = Y \times 4$.
        \item \texttt{sub s3, t1, s2}: $Z = \texttt{t1} - \texttt{s2} = X + 2 - Y \times 4$ $\rightarrow$ Chọn \textbf{B}.
    \end{itemize}
    \item \textbf{Giải câu b:}
    \begin{itemize}
        \item \texttt{add t1, s1, s2}: $\texttt{t1} = X + Y$.
        \item \texttt{srli t2, s2, 1}: Dịch phải $1\text{ bit}$ tương đương chia nguyên cho $2^1 = 2 \implies \texttt{t2} = Y / 2$.
        \item \texttt{add s3, t1, t2}: $Z = X + Y + Y / 2$ $\rightarrow$ Chọn \textbf{D}.
    \end{itemize}
\end{itemize}

\textbf{Bài tập 3.37 (Đề Cuối kỳ 2020.2 -- Trang 9, 2022.2 -- Trang 14, 2024.1 -- \#f4, 2024.2 \& Ôn tập 2025.2 -- Câu 68)}\\
Đoạn lệnh hợp ngữ RISC-V dưới đây tương ứng với câu lệnh C nào (biết rằng các thanh ghi \texttt{s0, s1, s2, s3} được sử dụng tương ứng cho các biến \texttt{f, g, h, i})?

\noindent\begin{minipage}{\linewidth}
\begin{lstlisting}
    slli t0, s1, 2
    sub  s3, s0, s2
    add  s3, s3, t0
    addi s3, s3, 5
\end{lstlisting}
\end{minipage}
\begin{enumerate}[label=\Alph*.]
    \item $i = f + h + 4 * g - 5$
    \item $i = f - h + 4 * g + 5$
    \item $i = f - h + 4 * g - 5$
    \item $i = f + h - 4 * g - 5$
\end{enumerate}
\textbf{Hướng dẫn giải:}
\begin{itemize}
    \item Lệnh 1 (\texttt{slli t0, s1, 2}): Dịch trái \texttt{s1} ($g$) đi $2\text{ bit}$, tương đương $\texttt{t0} = 4 * g$.
    \item Lệnh 2 (\texttt{sub s3, s0, s2}): $\texttt{s3} = f - h$.
    \item Lệnh 3 (\texttt{add s3, s3, t0}): $\texttt{s3} = f - h + 4 * g$.
    \item Lệnh 4 (\texttt{addi s3, s3, 5}): $i = \texttt{s3} = f - h + 4 * g + 5$.
    \item \textbf{Kết luận:} Đáp án đúng là \textbf{B ($i = f - h + 4 * g + 5$)}.
\end{itemize}

\textbf{Bài tập 3.38 (Đề Giữa kỳ 2021.1 -- IT3030E, Problem 1)}\\
Dịch câu lệnh C sau đây sang hợp ngữ RISC-V mà không sử dụng lệnh giả (pseudo instructions):
$$\texttt{x = (4*x - y) \& (x + 5*z);}$$
Biết rằng các biến \texttt{x, y, z} được lưu lần lượt trong các thanh ghi \texttt{s1, s2, s3}.

\textbf{Hướng dẫn giải:}
\begin{itemize}
    \item \textbf{Phân tích:} Vì bộ xử lý cơ bản thực hiện phép nhân với hằng số qua phép dịch bit và phép cộng, ta tách $4*x = (x \ll 2)$ và $5*z = 4*z + z = (z \ll 2) + z$:

\noindent\begin{minipage}{\linewidth}
\begin{lstlisting}
    slli t0, s1, 2      # t0 = 4 * x
    sub  t0, t0, s2     # t0 = 4 * x - y
    slli t1, s3, 2      # t1 = 4 * z
    add  t1, t1, s3     # t1 = 4 * z + z = 5 * z
    add  t1, s1, t1     # t1 = x + 5 * z
    and  s1, t0, t1     # x = (4 * x - y) & (x + 5 * z)
\end{lstlisting}
\end{minipage}
\end{itemize}

\textbf{Bài tập 3.39 (Đề Giữa kỳ 2022.2 -- Đề 0009, 0008, 0007, 0006 \& Giữa kỳ 2025.2 -- Đề 0012 Câu 107)}
\begin{enumerate}[label=\alph*)]
    \item \textbf{(Đề 0009, 0008, 0012):} Viết 1 lệnh hợp ngữ còn thiếu vào chỗ trống giữa 2 lệnh \texttt{add t2, s3, s4} và \texttt{sub s3, t3, s4} để thực hiện biểu thức (với \texttt{s3} là biến ở vế trái):
    \begin{itemize}
        \item TH1 (Đề 0009): $y = (x + y) \times 16 - x$ (với $y$ ở \texttt{s3}, $x$ ở \texttt{s4}).
        \item TH2 (Đề 0008): $x = (x + y) \times 4 - y$ (với $x$ ở \texttt{s3}, $y$ ở \texttt{s4}).
        \item TH3 (Đề 0012): $a = (a + b) \times 64 - a$ (với $a$ ở \texttt{s3} ở vế trái, $\texttt{s4} = a$ bị trừ).
    \end{itemize}
    \item \textbf{(Đề 0008 -- Câu 6):} Thay thế chuỗi 3 lệnh liên tiếp: \texttt{add s0, t2, t2} $\rightarrow$ \texttt{slli s0, s0, 1} $\rightarrow$ \texttt{slli s0, s0, 2} bằng 1 lệnh hợp ngữ duy nhất.
    \item \textbf{(Đề 0008 Câu 7, Đề 0007 Câu 6 \& Đề 0006 Câu 10):} Lập trình hợp ngữ RISC-V cho 3 biểu thức:
    \begin{itemize}
        \item Biểu thức 1: $\texttt{t0} = 4 * (\texttt{t1} + \texttt{t2}) + 5$.
        \item Biểu thức 2: $\texttt{t0} = \texttt{t1} + 8 * (\texttt{t2} - 7)$.
        \item Biểu thức 3: $\texttt{s2} = 8 * (\texttt{t3} + \texttt{t7}) + 6$.
    \end{itemize}
\end{enumerate}
\textbf{Hướng dẫn giải:}
\begin{itemize}
    \item \textbf{Giải câu a:} Thanh ghi \texttt{t2} đã giữ tổng $(x + y)$, lệnh thứ 3 lấy $\texttt{t3} - \texttt{s4}$, nên lệnh ở giữa cần dịch trái \texttt{t2} và cất vào \texttt{t3}:
    \begin{itemize}
        \item TH1 (nhân $16 = 2^4$): \textbf{\texttt{slli t3, t2, 4}}.
        \item TH2 (nhân $4 = 2^2$): \textbf{\texttt{slli t3, t2, 2}}.
        \item TH3 (nhân $64 = 2^6$): \textbf{\texttt{slli t3, t2, 6}}.
    \end{itemize}
    \item \textbf{Giải câu b:} Lệnh \texttt{add s0, t2, t2} nhân $2$, \texttt{slli 1} nhân $2$, \texttt{slli 2} nhân $4$. Tổng cộng nhân $2 \times 2 \times 4 = 16 = 2^4$. Thay bằng 1 lệnh duy nhất là: \textbf{\texttt{slli s0, t2, 4}}.
    \item \textbf{Giải câu c:}
    \begin{itemize}
        \item Biểu thức 1: \texttt{add t0, t1, t2} $\rightarrow$ \texttt{slli t0, t0, 2} $\rightarrow$ \texttt{addi t0, t0, 5}.
        \item Biểu thức 2: \texttt{addi t0, t2, -7} $\rightarrow$ \texttt{slli t0, t0, 3} $\rightarrow$ \texttt{add t0, t1, t0}.
        \item Biểu thức 3: \texttt{add s2, t3, t7} $\rightarrow$ \texttt{slli s2, s2, 3} $\rightarrow$ \texttt{addi s2, s2, 6}.
    \end{itemize}
\end{itemize}

\textbf{Bài tập 3.40 (Đề Giữa kỳ 2024.1 -- \#f96f43, \#12416e \& Ôn tập Cuối kỳ 2025.2 -- P1, Câu 32)}
\begin{enumerate}[label=\alph*)]
    \item \textbf{(Câu \#f96f43):} Chọn đoạn lệnh RISC-V thực hiện câu lệnh C: \texttt{*a = b + c;} (giả sử \texttt{a} là con trỏ trong thanh ghi \texttt{a1}, còn \texttt{b, c} là các biến nguyên trong \texttt{a2, a3}).\\
    A. \texttt{lw t0, 0(a1)} $\rightarrow$ \texttt{add t0, a2, a3} \qquad B. \texttt{add t0, a2, a3} $\rightarrow$ \texttt{lw t0, 0(a1)}\\
    C. \texttt{sw t0, 0(a1)} $\rightarrow$ \texttt{add t0, a2, a3} \qquad D. \texttt{add t0, a2, a3} $\rightarrow$ \texttt{sw t0, 0(a1)}
    \item \textbf{(Câu \#12416e \& Câu 32):} Chọn đoạn lệnh RISC-V thực hiện câu lệnh C: \texttt{a = b + *c;} (giả sử \texttt{a, b} nằm trong \texttt{a1, a2}, còn \texttt{c} là một con trỏ nằm trong \texttt{a3}).\\
    A. \texttt{li t0, 0(a3)} $\rightarrow$ \texttt{add a1, a2, t0} \qquad B. \texttt{sw t0, 0(a3)} $\rightarrow$ \texttt{add a1, a2, t0}\\
    C. \texttt{la t0, 0(a3)} $\rightarrow$ \texttt{add a1, a2, t0} \qquad D. \texttt{lw t0, 0(a3)} $\rightarrow$ \texttt{add a1, a2, t0}
\end{enumerate}
\textbf{Hướng dẫn giải:}
\begin{itemize}
    \item \textbf{Giải câu a:} Toán tử \texttt{*a} ở vế trái nghĩa là lưu kết quả vào ô nhớ được trỏ bởi \texttt{a1}. Ta tính tổng \texttt{add t0, a2, a3} trước, rồi ghi ra bộ nhớ bằng \texttt{sw t0, 0(a1)} $\rightarrow$ Chọn \textbf{D}.
    \item \textbf{Giải câu b:} Toán tử \texttt{*c} ở vế phải nghĩa là đọc giá trị từ ô nhớ được trỏ bởi \texttt{a3}. Ta đọc bằng \texttt{lw t0, 0(a3)} trước, rồi cộng với \texttt{a2} bằng \texttt{add a1, a2, t0} $\rightarrow$ Chọn \textbf{D}.
\end{itemize}

\textbf{Bài tập 3.41 (Đề Cuối kỳ 2022.2 -- \#13504d, 2024.2 \& Đề Giữa kỳ 2024.1 -- \#d0125d)}
\begin{enumerate}[label=\alph*)]
    \item \textbf{(Đề CK 2022.2 \& 2024.2 -- \#13504d):} Biến \texttt{k} ở thanh ghi \texttt{s1}; \texttt{P, Q} là hai mảng số nguyên 32-bit có địa chỉ cơ sở ở \texttt{s6} và \texttt{s7}. Đoạn lệnh: \texttt{lw t0, 160(s7)} $\rightarrow$ \texttt{add t0, s1, t0} $\rightarrow$ \texttt{sw t0, 80(s6)} tương ứng với câu lệnh C nào?\\
    A. \texttt{P[80] = k + Q[160]} \qquad B. \texttt{Q[40] = k + P[20]}\\
    C. \texttt{P[20] = k + Q[40]} \qquad D. \texttt{Q[160] = k + P[80]}
    \item \textbf{(Đề GK 2024.1 -- \#d0125d):} Cho \texttt{s1, s2} giữ địa chỉ cơ sở của mảng số nguyên 32-bit \texttt{A} và \texttt{B}. Đoạn lệnh: \texttt{lw t0, 24(s1)} $\rightarrow$ \texttt{addi s0, t0, 1} $\rightarrow$ \texttt{sw s0, 20(s2)} tương ứng với câu lệnh C nào?\\
    A. \texttt{A[24] = B[20] + 1} \qquad B. \texttt{B[20] = A[24] + 1}\\
    C. \texttt{A[5] = B[6] + 1} \qquad D. \texttt{B[5] = A[6] + 1}
\end{enumerate}
\textbf{Hướng dẫn giải:}
\begin{itemize}
    \item Mỗi phần tử số nguyên 32-bit chiếm $4\text{ byte}$, do đó $\text{Chỉ số mảng} = \text{Độ lệch byte (Offset)} \div 4$:
    \item \textbf{Câu a:} $\texttt{160(s7)} \rightarrow Q[160 \div 4] = Q[40]$; $\texttt{80(s6)} \rightarrow P[80 \div 4] = P[20]$. Vậy biểu thức là \texttt{P[20] = k + Q[40]} $\rightarrow$ Chọn \textbf{C}.
    \item \textbf{Câu b:} $\texttt{24(s1)} \rightarrow A[24 \div 4] = A[6]$; $\texttt{20(s2)} \rightarrow B[20 \div 4] = B[5]$. Vậy biểu thức là \texttt{B[5] = A[6] + 1} $\rightarrow$ Chọn \textbf{D}.
\end{itemize}

\textbf{Bài tập 3.42 (Đề Giữa kỳ 2023.2 -- Đề 1 Câu 5 \& Đề 2 Câu 3)}\\
Giả thiết \texttt{f, g, h, i, j} là các biến số nguyên 32-bit ở \texttt{s0, s1, s2, s3, s4}; \texttt{A} và \texttt{B} là hai mảng số nguyên 32-bit có địa chỉ cơ sở ở \texttt{s6} và \texttt{s7}.
\begin{enumerate}[label=\alph*)]
    \item \textbf{(Đề 1 -- Câu 5):} Viết đoạn hợp ngữ RISC-V (có giải thích) cho câu lệnh C: \texttt{f = g + h - A[20];} (với địa chỉ cơ sở của \texttt{A} ở \texttt{s7}).
    \item \textbf{(Đề 2 -- Câu 3):} Viết câu lệnh C tương ứng với 2 đoạn hợp ngữ sau:
    \begin{itemize}
        \item Đoạn 1: \texttt{lw t1, 20(s7)} $\rightarrow$ \texttt{sub s2, s3, t1}.
        \item Đoạn 2: \texttt{add t0, s1, s0} $\rightarrow$ \texttt{slli t1, s4, 2} $\rightarrow$ \texttt{add t1, s6, t1} $\rightarrow$ \texttt{sw t0, 0(t1)}.
    \end{itemize}
\end{enumerate}
\textbf{Hướng dẫn giải:}
\begin{itemize}
    \item \textbf{Giải câu a:} Phần tử \texttt{A[20]} có độ lệch là $20 \times 4 = 80\text{ byte}$:

\noindent\begin{minipage}{\linewidth}
\begin{lstlisting}
    lw   t0, 80(s7)     # t0 = A[20] (do lech 20 * 4 = 80 byte)
    add  t1, s1, s2     # t1 = g + h
    sub  s0, t1, t0     # f = (g + h) - A[20]
\end{lstlisting}
\end{minipage}

    \item \textbf{Giải câu b:}
    \begin{itemize}
        \item Đoạn 1: \texttt{20(s7)} là phần tử $\texttt{B[20/4]} = \texttt{B[5]}$, \texttt{s2} là $h$, \texttt{s3} là $i \implies$ Câu lệnh C: \textbf{\texttt{h = i - B[5];}}.
        \item Đoạn 2: $\texttt{t0} = g + f$; \texttt{slli t1, s4, 2} tính $4 \times j$, cộng với \texttt{s6} ra địa chỉ của \texttt{A[j]}, rồi lưu \texttt{t0} vào \texttt{A[j]} $\implies$ Câu lệnh C: \textbf{\texttt{A[j] = g + f;}}.
    \end{itemize}
\end{itemize}

\textbf{Bài tập 3.43 (Đề Cuối kỳ 2021.2 -- Mã đề B2 Câu 2 \& Mã đề B3 Câu 2)}\\
Viết đoạn chương trình hợp ngữ tối thiểu của RISC-V (không dùng lệnh giả, có lời giải thích) để thực hiện các câu lệnh C sau (biết biến \texttt{a} nằm trong thanh ghi \texttt{s2}):
\begin{enumerate}[label=\alph*)]
    \item \textbf{(Mã đề B2):} \texttt{x[2] = x[7] + a;} với \texttt{x} là mảng 32-bit có địa chỉ cơ sở là \texttt{0x00283A00}.
    \item \textbf{(Mã đề B3):} \texttt{x[3] = x[9] + a;} với \texttt{x} là mảng 32-bit có địa chỉ cơ sở là \texttt{0x00D28900}.
\end{enumerate}
\textbf{Hướng dẫn giải:}
\begin{itemize}
    \item \textbf{Giải câu a (Địa chỉ cơ sở \texttt{0x00283A00}):}\\
    Chữ số Hexa thứ 3 từ phải sang là $\texttt{A} \ge 8$ ($\texttt{0xA00} = 2560 - 4096 = -1536$), nên khi dùng \texttt{lui} phải tăng \texttt{0x00283} lên \texttt{0x00284}:

\noindent\begin{minipage}{\linewidth}
\begin{lstlisting}
    lui  t0, 0x00284      # t0 = 0x00284000
    addi t0, t0, -1536    # t0 = 0x00283A00 (dia chi co so mang x)
    lw   t1, 28(t0)       # t1 = x[7] (do lech 7 * 4 = 28 byte)
    add  t1, t1, s2       # t1 = x[7] + a
    sw   t1, 8(t0)        # x[2] = x[7] + a (do lech 2 * 4 = 8 byte)
\end{lstlisting}
\end{minipage}

    \item \textbf{Giải câu b (Địa chỉ cơ sở \texttt{0x00D28900}):}\\
    Chữ số thứ 3 từ phải sang là $\texttt{9} \ge 8$ ($\texttt{0x900} = 2304 - 4096 = -1792$), tăng \texttt{0x00D28} lên \texttt{0x00D29}:

\noindent\begin{minipage}{\linewidth}
\begin{lstlisting}
    lui  t0, 0x00D29      # t0 = 0x00D29000
    addi t0, t0, -1792    # t0 = 0x00D28900 (dia chi co so mang x)
    lw   t1, 36(t0)       # t1 = x[9] (do lech 9 * 4 = 36 byte)
    add  t1, t1, s2       # t1 = x[9] + a
    sw   t1, 12(t0)       # x[3] = x[9] + a (do lech 3 * 4 = 12 byte)
\end{lstlisting}
\end{minipage}
\end{itemize}

\textbf{Bài tập 3.44 (Đề Giữa kỳ 2019.2 -- Mã đề 243 Câu 7, 8, Mã đề 244 Câu 16, 18, Giữa kỳ 2022.2 -- Đề 0008 Câu 8, Đề 0007 Câu 7, Giữa kỳ 2024.1 -- \#3602ec \& Cuối kỳ 2025.1 -- \#da417b)}\\
Hãy xác định địa chỉ hoặc chỉ số phần tử mảng được truy cập trong các bài toán sau:
\begin{enumerate}[label=\alph*)]
    \item \textbf{(Mã đề 243 Câu 7 \& Mã đề 244 Câu 18):} Mảng 32-bit \texttt{A} có \texttt{A[6]} ở \texttt{0x20200508}, hỏi \texttt{A[9]} ở đâu? Nếu \texttt{A[2]} ở \texttt{0x20200008}, hỏi \texttt{A[5]} ở đâu?
    \item \textbf{(Mã đề 243 Câu 8 \& Mã đề 244 Câu 16):} Mảng 32-bit \texttt{A} có địa chỉ cơ sở \texttt{0x20201234}.
    \begin{itemize}
        \item TH1 (Mã đề 243): Chạy \texttt{lui s0, 0x20201} $\rightarrow$ \texttt{addi s0, s0, 0x230} $\rightarrow$ \texttt{lw t0, 8(s0)}. Hỏi \texttt{t0} nhận phần tử nào?
        \item TH2 (Mã đề 244): Chạy \texttt{lui s0, 0x20201} $\rightarrow$ \texttt{addi s0, s0, 0x238} $\rightarrow$ \texttt{lw t0, 4(s0)}. Hỏi \texttt{t0} nhận phần tử nào?
    \end{itemize}
    \item \textbf{(Đề GK 2022.2 -- Đề 0008 Câu 8 \& Đề 0007 Câu 7):} Mảng 32-bit \texttt{A} ở địa chỉ cơ sở \texttt{0x1000CDC0}.
    \begin{itemize}
        \item TH1 (Đề 0008): Khi $\texttt{s1} = \texttt{0x1000CDC4}$, thực hiện \texttt{lw t1, 4(s1)}. Hỏi các địa chỉ byte nào được nạp vào \texttt{t1} và đó là phần tử nào của \texttt{A}?
        \item TH2 (Đề 0007): Khi $\texttt{s1} = \texttt{0x1000CDC8}$, thực hiện \texttt{lw t1, -4(s1)}. Hỏi đó là phần tử nào của \texttt{A}?
    \end{itemize}
    \item \textbf{(Đề GK 2024.1 -- \#3602ec, Ôn tập 2025.2 Câu 72 \& CK 2025.1 -- \#da417b):}
    \begin{itemize}
        \item TH1: Mảng \texttt{A[]} ở địa chỉ cơ sở \texttt{0x20230200}. Chạy \texttt{lui s0, 0x20230} $\rightarrow$ \texttt{addi s0, s0, 0x20C} $\rightarrow$ \texttt{lw t0, 4(s0)}. Hỏi \texttt{t0} giữ phần tử nào?
        \item TH2: Mảng \texttt{A[]} ở địa chỉ cơ sở \texttt{0x20240214}. Chạy \texttt{lui t1, 0x20240} $\rightarrow$ \texttt{addi t1, t1, 0x218} $\rightarrow$ \texttt{lw t0, 4(t1)}. Hỏi \texttt{t0} giữ phần tử nào?
    \end{itemize}
\end{enumerate}
\textbf{Hướng dẫn giải:}
\begin{itemize}
    \item \textbf{Giải câu a:} Cách nhau 3 phần tử $= 3 \times 4 = 12\text{ byte} = \texttt{0x0C}$. Vậy $\texttt{A[9]} = \texttt{0x20200508} + \texttt{0x0C} = \mathbf{\texttt{0x20200514}}$; và $\texttt{A[5]} = \texttt{0x20200008} + \texttt{0x0C} = \mathbf{\texttt{0x20200014}}$.
    \item \textbf{Giải câu b:}
    \begin{itemize}
        \item TH1: Địa chỉ đọc $= \texttt{0x20201230} + 8 = \texttt{0x20201238}$. Trừ địa chỉ cơ sở \texttt{0x20201234} được $4\text{ byte} \implies 4 \div 4 = 1 \rightarrow \mathbf{\texttt{A[1]}}$.
        \item TH2: Địa chỉ đọc $= \texttt{0x20201238} + 4 = \texttt{0x2020123C}$. Trừ \texttt{0x20201234} được $8\text{ byte} \implies 8 \div 4 = 2 \rightarrow \mathbf{\texttt{A[2]}}$.
    \end{itemize}
    \item \textbf{Giải câu c:}
    \begin{itemize}
        \item TH1: Địa chỉ đọc $= \texttt{0x1000CDC4} + 4 = \texttt{0x1000CDC8}$ (gồm 4 byte từ \texttt{0x1000CDC8} đến \texttt{0x1000CDCB}). Độ lệch so với \texttt{0x1000CDC0} là $8\text{ byte} \implies \mathbf{\texttt{A[2]}}$.
        \item TH2: Địa chỉ đọc $= \texttt{0x1000CDC8} - 4 = \texttt{0x1000CDC4}$. Độ lệch là $4\text{ byte} \implies \mathbf{\texttt{A[1]}}$.
    \end{itemize}
    \item \textbf{Giải câu d:}
    \begin{itemize}
        \item TH1: Địa chỉ đọc $= \texttt{0x2023020C} + 4 = \texttt{0x20230210}$. Trừ \texttt{0x20230200} được $\texttt{0x10} = 16\text{ byte} \implies 16 \div 4 = 4 \rightarrow \mathbf{\texttt{A[4]}}$.
        \item TH2: Địa chỉ đọc $= \texttt{0x20240218} + 4 = \texttt{0x2024021C}$. Trừ \texttt{0x20240214} được $8\text{ byte} \implies 8 \div 4 = 2 \rightarrow \mathbf{\texttt{A[2]}}$.
    \end{itemize}
\end{itemize}

\textbf{Bài tập 3.45 (Đề Cuối kỳ 2022.2 \& Cuối kỳ 2024.2 -- \#0542f \& \#70e201)}
\begin{enumerate}[label=\alph*)]
    \item \textbf{(Câu \#0542f):} Mảng \texttt{M[]} gồm 100 phần tử số nguyên 32-bit bắt đầu từ địa chỉ \texttt{0x10020400}. Đoạn lệnh hoán đổi \texttt{M[i]} và \texttt{M[j]} dùng \texttt{lw t0, 0(s6)} và \texttt{lw t1, 0(s7)}. Nếu $i = 9, j = 6$, nội dung các thanh ghi \texttt{s6} và \texttt{s7} cần có giá trị bao nhiêu?
    \item \textbf{(Câu \#70e201):} Mảng 2 chiều \texttt{M} kích thước $(4 \times 5)$ gồm các phần tử 32-bit, lưu trữ theo hàng bắt đầu tại địa chỉ \texttt{0x10024000}. Để lệnh \texttt{lw t0, 0(s6)} đọc phần tử \texttt{M[2, 2]} (với $i = 0..3, j = 0..4$), nội dung thanh ghi \texttt{s6} phải bằng bao nhiêu?
\end{enumerate}
\textbf{Hướng dẫn giải:}
\begin{itemize}
    \item \textbf{Giải câu a:}
    \begin{itemize}
        \item Với $i = 9$: Độ lệch $= 9 \times 4 = 36\text{ byte} = \texttt{0x24} \implies \texttt{s6} = \texttt{0x10020400} + \texttt{0x24} = \mathbf{\texttt{0x10020424}}$.
        \item Với $j = 6$: Độ lệch $= 6 \times 4 = 24\text{ byte} = \texttt{0x18} \implies \texttt{s7} = \texttt{0x10020400} + \texttt{0x18} = \mathbf{\texttt{0x10020418}}$.
    \end{itemize}
    \item \textbf{Giải câu b:} Mảng có $5$ cột. Với $i = 2, j = 2$, số phần tử đứng trước \texttt{M[2, 2]} là $2 \times 5 + 2 = 12\text{ phần tử}$. Độ lệch byte là $12 \times 4 = 48\text{ byte} = \texttt{0x30}$. Vậy $\texttt{s6} = \texttt{0x10024000} + \texttt{0x30} = \mathbf{\texttt{0x10024030}}$.
\end{itemize}

\textbf{Bài tập 3.46 (Đề Cuối kỳ 2021.2 -- Mã đề B1 Câu 4.1 \& Ôn tập Cuối kỳ 2025.2 -- P2, Câu 67)}\\
Trên máy tính có các thanh ghi 32-bit, giả sử có một mảng dữ liệu hai chiều \texttt{B} kích thước $5 \times 7$ (5 hàng, 7 cột) của các phần tử 32-bit. Mảng \texttt{B} được lưu trữ theo hàng trong bộ nhớ chính, cấp phát ban đầu từ địa chỉ \texttt{0x000A5020}. Để trỏ vào phần tử \texttt{B[i, j]}, yêu cầu kết hợp thanh ghi \texttt{R1} để trỏ vào đầu hàng thứ $i$, thanh ghi \texttt{R2} để trỏ tương đối vào phần tử thứ $j$ trong hàng $i$ tính từ đầu hàng.
\begin{enumerate}[label=\alph*)]
    \item \textbf{(Đề CK 2021.2 -- B1):} Với chỉ số tính từ $1$ ($i = 1..5, j = 1..7$), các thanh ghi \texttt{R1} và \texttt{R2} phải được gán giá trị bằng bao nhiêu (dạng Hexa) khi trỏ vào phần tử \texttt{B[2, 5]}?
    \item \textbf{(Ôn tập 2025.2 -- P2, Câu 67):} Với chỉ số tính từ $0$ ($i = 0..4, j = 0..6$), cho biết địa chỉ của phần tử \texttt{B[1, 4]}?
\end{enumerate}
\textbf{Hướng dẫn giải:}
\begin{itemize}
    \item Mỗi hàng có $7$ phần tử 32-bit ($4\text{ byte}$), nên kích thước 1 hàng là $7 \times 4 = 28\text{ byte} = \texttt{0x1C}\text{ byte}$.
    \item \textbf{Giải câu a:} Phần tử \texttt{B[2, 5]} (khi đếm từ 1) nằm ở hàng thứ 2, cột thứ 5:
    \begin{itemize}
        \item Thanh ghi \texttt{R1} trỏ vào đầu hàng 2 (cách đầu mảng 1 hàng $= \texttt{0x1C}$): $\texttt{R1} = \texttt{0x000A5020} + \texttt{0x1C} = \mathbf{\texttt{0x000A503C}}$.
        \item Thanh ghi \texttt{R2} trỏ tương đối vào phần tử thứ 5 trong hàng (cách đầu hàng $5 - 1 = 4\text{ phần tử} = 16\text{ byte} = \texttt{0x10}$): $\mathbf{\texttt{R2 = 0x00000010}}$.
    \end{itemize}
    \item \textbf{Giải câu b:} Phần tử \texttt{B[1, 4]} (khi đếm từ 0) cũng chính là phần tử ở hàng thứ 2, cột thứ 5 nói trên. Địa chỉ của nó là: $\texttt{R1} + \texttt{R2} = \texttt{0x000A503C} + \texttt{0x10} = \mathbf{\texttt{0x000A504C}}$.
\end{itemize}

\textbf{Bài tập 3.47 (Đề Cuối kỳ 2022.2 -- \#bcfd00, 2024.2 \& Đề Cuối kỳ 2025.1 -- \#31d5bf)}
\begin{enumerate}[label=\alph*)]
    \item \textbf{(Đề CK 2022.2 \& 2024.2 -- \#bcfd00):} Cho các biến \texttt{a, b, c} được gán tương ứng trong \texttt{s0, s1, s2}. Hãy chỉ ra câu lệnh C tương ứng với đoạn hợp ngữ sau:

\noindent\begin{minipage}{\linewidth}
\begin{lstlisting}
    slt  t0, s2, s1
    bne  t0, zero, L1
    add  s0, s0, s1
    j    L2
L1: add  s0, s0, s2
L2:
\end{lstlisting}
\end{minipage}

    A. \texttt{if (b > c) a = a + b; else a = a + c;}\\
    B. \texttt{if (b <= c) a = a + b; else a = a + c;}\\
    C. \texttt{if (b < c) a = a + b; else a = a + c;}\\
    D. \texttt{if (b >= c) a = a + b; else a = a + c;}
    \item \textbf{(Đề CK 2025.1 -- \#31d5bf \& Ôn tập 2025.2 -- P2, Câu 46):} Cho các biến \texttt{f, g, h, i, j} lần lượt nằm trong \texttt{s0, s1, s2, s3, s4}. Đoạn mã hợp ngữ sau tương đương với đoạn mã C nào?

\noindent\begin{minipage}{\linewidth}
\begin{lstlisting}
    blt  s2, s1, L1
    beq  s3, s4, L1
    add  s0, s0, s1
    jal  x0, L2
L1: sub  s0, s0, s1
L2:
\end{lstlisting}
\end{minipage}

    A. \texttt{if ((g > h) \&\& (i == j)) f = f + g; else f = f - g;}\\
    B. \texttt{if ((h >= g) \&\& (i == j)) f = f + g; else f = f - g;}\\
    C. \texttt{if ((g > h) \&\& (i != j)) f = f + g; else f = f - g;}\\
    D. \texttt{if ((h >= g) \&\& (i != j)) f = f + g; else f = f - g;}
\end{enumerate}
\textbf{Hướng dẫn giải:}
\begin{itemize}
    \item \textbf{Giải câu a:} Lệnh \texttt{slt t0, s2, s1} kiểm tra $\texttt{s2} < \texttt{s1}$ ($c < b$, tức $b > c$). Nếu $b > c$ ($\texttt{t0} = 1 \neq 0$) thì nhảy xuống \texttt{L1} thực hiện $a = a + c$. Nếu không nhảy (tức \textbf{$b \le c$}) thì thực hiện lệnh ngay bên dưới là $a = a + b$ $\rightarrow$ Chọn \textbf{B}.
    \item \textbf{Giải câu b:} Để thực hiện lệnh $f = f + g$ (\texttt{add s0, s0, s1}), chương trình phải \textbf{không nhảy} ở cả hai lệnh \texttt{blt} và \texttt{beq}:
    \begin{itemize}
        \item Không nhảy ở \texttt{blt s2, s1, L1} (nhảy khi $h < g$) $\implies$ Điều kiện là \textbf{$h \ge g$}.
        \item Không nhảy ở \texttt{beq s3, s4, L1} (nhảy khi $i == j$) $\implies$ Điều kiện là \textbf{$i \neq j$}.
        \item Kết hợp lại: \texttt{if ((h >= g) \&\& (i != j)) f = f + g; else f = f - g;} $\rightarrow$ Chọn \textbf{D}.
    \end{itemize}
\end{itemize}

\textbf{Bài tập 3.48 (Đề Giữa kỳ 2023.2 -- Đề 2 Câu 5, Giữa kỳ 2022.2 -- Đề 0007 Câu 8 \& Cuối kỳ 2021.2 -- Mã đề B5 Câu 3, Mã đề B8 Câu 3)}\\
Viết đoạn chương trình hợp ngữ tối thiểu của RISC-V (không dùng lệnh giả, có lời giải thích) cho 4 câu lệnh rẽ nhánh \texttt{if/else} sau:
\begin{enumerate}[label=\alph*)]
    \item \textbf{(Đề GK 2023.2 -- Đề 2, Câu 5):} \texttt{if (i <= j) f = f + g; else f = f + h;} (với \texttt{f, g, h, i, j} ở \texttt{s0, s1, s2, s3, s4}).
    \item \textbf{(Đề GK 2022.2 -- Đề 0007, Câu 8):} \texttt{if (s1 == 0) t2 = 1; else t2 = 2;}.
    \item \textbf{(Đề CK 2021.2 -- Mã đề B5, Câu 3):} \texttt{if (a == b) i += 1; else i += 3;} (với \texttt{a, b, i} ở \texttt{s0, s1, s2}).
    \item \textbf{(Đề CK 2021.2 -- Mã đề B8, Câu 3):} \texttt{if (a == 1) m[3] = b - a; else m[4] = b + a*2;} (với \texttt{a, b} ở \texttt{s0, s1}; \texttt{m} là mảng 32-bit có địa chỉ cơ sở \texttt{0x13090134}).
\end{enumerate}
\textbf{Hướng dẫn giải:}
\begin{itemize}
    \item \textbf{Giải câu a:} Điều kiện ngược của $i \le j$ là $j < i$ (\texttt{blt s4, s3, Else}):

\noindent\begin{minipage}{\linewidth}
\begin{lstlisting}
    blt  s4, s3, Else     # Neu j < i (tuc i > j) thi nhay xuong Else
    add  s0, s0, s1       # f = f + g
    jal  x0, Exit         # Nhay khong dieu kien toi Exit
Else:
    add  s0, s0, s2       # f = f + h
Exit:
\end{lstlisting}
\end{minipage}

    \item \textbf{Giải câu b:}

\noindent\begin{minipage}{\linewidth}
\begin{lstlisting}
    bne  s1, x0, Else     # Neu s1 != 0 thi nhay xuong Else
    addi t2, x0, 1        # t2 = 1
    jal  x0, Exit         # Nhay qua nhanh Else
Else:
    addi t2, x0, 2        # t2 = 2
Exit:
\end{lstlisting}
\end{minipage}

    \item \textbf{Giải câu c:}

\noindent\begin{minipage}{\linewidth}
\begin{lstlisting}
    bne  s0, s1, Else     # Neu a != b thi nhay xuong Else
    addi s2, s2, 1        # i = i + 1
    jal  x0, Exit         # Nhay qua nhanh Else
Else:
    addi s2, s2, 3        # i = i + 3
Exit:
\end{lstlisting}
\end{minipage}

    \item \textbf{Giải câu d:}

\noindent\begin{minipage}{\linewidth}
\begin{lstlisting}
    lui  t0, 0x13090      # Nap 20 bit cao cua dia chi mang m
    addi t0, t0, 0x134    # t0 = 0x13090134 (dia chi co so mang m)
    addi t1, x0, 1        # t1 = 1
    bne  s0, t1, Else     # Neu a != 1 thi nhay xuong Else
    sub  t2, s1, s0       # t2 = b - a
    sw   t2, 12(t0)       # m[3] = b - a (do lech 3 * 4 = 12 byte)
    jal  x0, Exit         # Nhay toi Exit
Else:
    slli t2, s0, 1        # t2 = a * 2
    add  t2, s1, t2       # t2 = b + a * 2
    sw   t2, 16(t0)       # m[4] = b + a * 2 (do lech 4 * 4 = 16B)
Exit:
\end{lstlisting}
\end{minipage}
\end{itemize}

\textbf{Bài tập 3.49 (Đề Giữa kỳ 2021.1 -- IT3030E, Problem 5)}\\
Cho đoạn mã C sau:
$$\texttt{if (i < j) f = A[2*i] + B[4*j];}$$
Giả sử các biến \texttt{f, i, j} được gán cho các thanh ghi \texttt{s0, s1, s2}; địa chỉ cơ sở của mảng 32-bit \texttt{A} và \texttt{B} nằm trong \texttt{s6} và \texttt{s7}. Hãy dịch câu lệnh trên sang hợp ngữ RISC-V và viết mã máy cho từng lệnh.

\textbf{Hướng dẫn giải:}
\begin{itemize}
    \item \textbf{Phân tích độ lệch byte:} Mỗi phần tử rộng $4\text{ byte}$. Phần tử \texttt{A[2*i]} có độ lệch byte là $4 \times (2*i) = 8*i = (i \ll 3)$; phần tử \texttt{B[4*j]} có độ lệch byte là $4 \times (4*j) = 16*j = (j \ll 4)$:

\noindent\begin{minipage}{\linewidth}
\begin{lstlisting}
    bge  s1, s2, Exit     # Ma may: 0x0124D063 (offset = +32)
    slli t0, s1, 3        # Ma may: 0x00349293 (t0 = 8 * i)
    add  t0, s6, t0       # Ma may: 0x005B02B3 (dia chi A[2*i])
    lw   t0, 0(t0)        # Ma may: 0x0002A283 (t0 = A[2*i])
    slli t1, s2, 4        # Ma may: 0x00491313 (t1 = 16 * j)
    add  t1, s7, t1       # Ma may: 0x006B8333 (dia chi B[4*j])
    lw   t1, 0(t1)        # Ma may: 0x00032303 (t1 = B[4*j])
    add  s0, t0, t1       # Ma may: 0x00628433 (f = A[2*i] + B[4*j])
Exit:
\end{lstlisting}
\end{minipage}
\end{itemize}

\textbf{Bài tập 3.50 (Tổng hợp 4 bài tự luận viết hợp ngữ vòng lặp \texttt{for} trong Đề Cuối kỳ 2021.2 -- Mã đề B1, B4, B6, B7)}\\
Viết đoạn chương trình hợp ngữ tối thiểu của RISC-V (không dùng lệnh giả, có lời giải thích) để thực hiện các vòng lặp ngôn ngữ C sau đây:
\begin{enumerate}[label=\alph*)]
    \item \textbf{(Mã đề B1, Câu 3):} \texttt{for (i = 0; i < 50; i += 2) a = a + m[5];} (với \texttt{m} là mảng các phần tử 32-bit có địa chỉ cơ sở \texttt{0x07E1100A}; giả sử \texttt{i, a} ở \texttt{s0, s1}).
    \item \textbf{(Mã đề B4, Câu 2):} \texttt{for (i = 1; i != 20; i += 3) a[i] = a[5] + 1;} (với \texttt{a} là mảng 32-bit có địa chỉ cơ sở \texttt{0x0FE3B128}; \texttt{i} ở \texttt{s0}).
    \item \textbf{(Mã đề B6, Câu 3):} \texttt{for (i = 0; i < 100; i++) N[i] = N[i] * 16;} (với \texttt{N} là mảng 32-bit có địa chỉ cơ sở \texttt{0x100A4023}; \texttt{i} ở \texttt{s0}).
    \item \textbf{(Mã đề B7, Câu 3):} \texttt{for (i = 0; i < 100; i++) M[i] = M[i] * 8;} (với \texttt{M} là mảng 32-bit có địa chỉ cơ sở \texttt{0x07AD1309}; \texttt{i} ở \texttt{s0}).
\end{enumerate}
\textbf{Hướng dẫn giải:}
\begin{itemize}
    \item \textbf{Giải câu a (Mã đề B1):} Phần tử \texttt{m[5]} cố định nên đọc ra ngoài vòng lặp để tối ưu:

\noindent\begin{minipage}{\linewidth}
\begin{lstlisting}
    lui  t0, 0x07E11      # Nap 20 bit cao cua dia chi co so m
    addi t0, t0, 0x00A    # t0 = 0x07E1100A (dia chi co so mang m)
    lw   t1, 20(t0)       # t1 = m[5] (do lech 5 * 4 = 20 byte)
    addi s0, x0, 0        # i = 0
    addi t2, x0, 50       # t2 = 50 (gia tri gioi han)
Loop:
    bge  s0, t2, Done     # Neu i >= 50 thi thoat vong lap
    add  s1, s1, t1       # a = a + m[5]
    addi s0, s0, 2        # i = i + 2
    jal  x0, Loop         # Quay lai Loop
Done:
\end{lstlisting}
\end{minipage}

    \item \textbf{Giải câu b (Mã đề B4):}

\noindent\begin{minipage}{\linewidth}
\begin{lstlisting}
    lui  t0, 0x0FE3B      # Nap 20 bit cao cua dia chi co so a
    addi t0, t0, 0x128    # t0 = 0x0FE3B128 (dia chi co so mang a)
    lw   t1, 20(t0)       # t1 = a[5]
    addi t1, t1, 1        # t1 = a[5] + 1
    addi s0, x0, 1        # i = 1
    addi t2, x0, 20       # t2 = 20
Loop:
    beq  s0, t2, Done     # Neu i == 20 thi thoat
    slli t3, s0, 2        # t3 = i * 4 (do lech byte cua a[i])
    add  t3, t0, t3       # t3 = dia chi cua a[i]
    sw   t1, 0(t3)        # a[i] = a[5] + 1
    addi s0, s0, 3        # i = i + 3
    jal  x0, Loop         # Quay lai Loop
Done:
\end{lstlisting}
\end{minipage}

    \item \textbf{Giải câu c và câu d (Mã đề B6 \& B7):} Nhân với $16 = 2^4$ dùng \texttt{slli ..., 4}; nhân với $8 = 2^3$ dùng \texttt{slli ..., 3}. (Đoạn mã dưới viết cho Đề B6, với Đề B7 chỉ cần thay địa chỉ \texttt{0x07AD1} / \texttt{0x309} và dịch trái \texttt{3}):

\noindent\begin{minipage}{\linewidth}
\begin{lstlisting}
    lui  t0, 0x100A4      # Nap 20 bit cao dia chi mang N
    addi t0, t0, 0x023    # t0 = 0x100A4023 (dia chi co so N[0])
    addi s0, x0, 0        # i = 0
    addi t1, x0, 100      # t1 = 100
Loop:
    bge  s0, t1, Done     # Neu i >= 100 thi thoat
    lw   t2, 0(t0)        # t2 = N[i]
    slli t2, t2, 4        # t2 = N[i] * 16 (De B7 thay 4 bang 3)
    sw   t2, 0(t0)        # N[i] = N[i] * 16
    addi t0, t0, 4        # Tang con tro t0 sang N[i+1]
    addi s0, s0, 1        # i++
    jal  x0, Loop         # Lap lai
Done:
\end{lstlisting}
\end{minipage}
\end{itemize}

\subsection{Nhóm 6: Phân tích Vòng lặp và Đếm tổng số lệnh thực thi}

\textbf{Bài tập 3.51 (Đề Cuối kỳ 2020.2 -- Trang 4 \& Đề Cuối kỳ 2021.1 -- Trang 3)}\\
Hãy phân tích đoạn mã nguồn chương trình vòng lặp hợp ngữ RISC-V dưới đây và xác định giá trị hằng số hệ thập phân thu được tại thanh ghi kết quả \texttt{a0} sau khi thực thi hoàn tất chương trình:

\noindent\begin{minipage}{\linewidth}
\begin{lstlisting}
    addi a1, zero, 14
    addi s0, zero, 0
    add  s1, zero, zero
loop:
    slt  t0, a1, s1
    bne  t0, zero, finish
    add  s0, s0, s1
    addi s1, s1, 2
    j    loop
finish:
    add  a0, s0, zero
\end{lstlisting}
\end{minipage}
\begin{enumerate}[label=\Alph*.]
    \item $54$
    \item $14$
    \item $56$
    \item $28$
\end{enumerate}
\textbf{Hướng dẫn giải:}
\begin{itemize}
    \item Khởi tạo: $\texttt{a1} = 14$, biến tổng $\texttt{s0} = 0$, biến chạy $\texttt{s1} = 0$.
    \item Trong mỗi vòng lặp, lệnh \texttt{slt t0, a1, s1} kiểm tra $14 < \texttt{s1}$. Chừng nào $\texttt{s1} \le 14$ thì $\texttt{t0} = 0$ (không nhảy tới \texttt{finish}), chương trình cộng \texttt{s1} vào \texttt{s0} và tăng \texttt{s1} thêm $2$. Khi $\texttt{s1} = 16$ ($14 < 16$ đúng) thì thoát vòng lặp.
    \item Vậy chương trình tính tổng các số chẵn từ $0$ đến $14$:
    $$\texttt{a0} = \texttt{s0} = 0 + 2 + 4 + 6 + 8 + 10 + 12 + 14 = 56$$
    \item \textbf{Kết luận:} Đáp án đúng là \textbf{C (56)}.
\end{itemize}

\textbf{Bài tập 3.52 (Đề Giữa kỳ 2024.1 -- \#f4fd9f \& Ôn tập Cuối kỳ 2025.2 -- P1, Câu 35)}\\
Sau khi chạy chương trình dưới đây, giá trị thập phân của thanh ghi \texttt{s5} là bao nhiêu?

\noindent\begin{minipage}{\linewidth}
\begin{lstlisting}
.data
A: .word 1, 2, 3, 4, 5, 6, 7, 8, 9, 10, 11, 12, 13, 14, 15
.text
    li   s1, 0
    li   s3, 10
    la   s2, A
    li   s5, 0
loop:
    add  t1, s1, s1
    add  t1, t1, t1
    add  t1, t1, s2
    lw   t0, 0(t1)
    add  s5, s5, t0
    addi s1, s1, 2
    bne  s1, s3, loop
\end{lstlisting}
\end{minipage}
\textbf{Hướng dẫn giải:}
\begin{itemize}
    \item Hai lệnh \texttt{add t1, s1, s1} ($\texttt{t1} = 2 \times \texttt{s1}$) và \texttt{add t1, t1, t1} ($\texttt{t1} = 4 \times \texttt{s1}$) tính độ lệch byte của phần tử \texttt{A[s1]}.
    \item Biến chỉ số \texttt{s1} bắt đầu từ $0$, mỗi vòng tăng thêm $2$, và dừng khi $\texttt{s1} == 10$. Vậy \texttt{s1} lần lượt nhận các giá trị: $0, 2, 4, 6, 8$.
    \item Tổng cộng dồn vào \texttt{s5} là (chú ý $\texttt{A[0]} = 1$):
    $$\texttt{s5} = \texttt{A[0]} + \texttt{A[2]} + \texttt{A[4]} + \texttt{A[6]} + \texttt{A[8]} = 1 + 3 + 5 + 7 + 9 = \mathbf{25}$$
\end{itemize}

\textbf{Bài tập 3.53 (Ôn tập Cuối kỳ 2025.2 -- P1, Câu 33)}\\
Sau khi thực thi đoạn mã sau, giá trị cuối cùng của thanh ghi \texttt{x15} là bao nhiêu (thập phân)? Giả sử ban đầu \texttt{x14} chứa $4$ và \texttt{x15} chứa $0$:

\noindent\begin{minipage}{\linewidth}
\begin{lstlisting}
Loop:
    addi x15, x15, 1
    addi x14, x14, -1
    bge  x14, x0, Loop
\end{lstlisting}
\end{minipage}
\begin{enumerate}[label=\Alph*.]
    \item $4$
    \item $0$
    \item $5$
    \item $6$
\end{enumerate}
\textbf{Hướng dẫn giải:}
\begin{itemize}
    \item Lệnh \texttt{bge x14, x0, Loop} quay lại vòng lặp chừng nào $\texttt{x14} \ge 0$.
    \item Lần 1: $\texttt{x15} = 1, \texttt{x14} = 3 \ge 0$ (lặp lại).
    \item Lần 2: $\texttt{x15} = 2, \texttt{x14} = 2 \ge 0$ (lặp lại).
    \item Lần 3: $\texttt{x15} = 3, \texttt{x14} = 1 \ge 0$ (lặp lại).
    \item Lần 4: $\texttt{x15} = 4, \texttt{x14} = 0 \ge 0$ (vẫn thỏa mãn $\ge 0$ nên lặp lại lần nữa).
    \item Lần 5: $\texttt{x15} = 5, \texttt{x14} = -1 < 0$ (thoát vòng lặp).
    \item \textbf{Kết luận:} Đáp án đúng là \textbf{C (5)}.
\end{itemize}

\textbf{Bài tập 3.54 (Đề Cuối kỳ 2021.2 -- Mã đề B1 Câu 2 \& Mã đề B5 Câu 2)}
\begin{enumerate}[label=\alph*)]
    \item \textbf{(Mã đề B1, Câu 2):} Cho đoạn chương trình hợp ngữ:

\noindent\begin{minipage}{\linewidth}
\begin{lstlisting}
    addi t1, zero, 100
LOOP:
    lw   s1, 0(s0)
    add  s2, s2, s1
    addi s0, s0, 4
    addi t1, t1, -1
    bne  t1, x0, LOOP
\end{lstlisting}
\end{minipage}

    2.1. Khi thực hiện đoạn chương trình trên, có bao nhiêu lệnh được thực hiện?\\
    2.2. Biết \texttt{t1} chứa biến \texttt{i}, \texttt{s2} chứa biến \texttt{a}, \texttt{s0} chứa địa chỉ cơ sở mảng \texttt{A}. Hãy chuyển sang ngôn ngữ C.
    \item \textbf{(Mã đề B5, Câu 2):} Cho đoạn chương trình hợp ngữ (giả sử ban đầu \texttt{s0} giữ địa chỉ cơ sở mảng \texttt{B}, và lệnh đầu gán $\texttt{t1} = \texttt{s0} + 600$):

\noindent\begin{minipage}{\linewidth}
\begin{lstlisting}
    addi t1, s0, 600
LOOP:
    lw   s1, 0(s0)
    add  s2, s2, s1
    lw   s1, 4(s0)
    add  s2, s2, s1
    addi s0, s0, 8
    bne  t1, s0, LOOP
\end{lstlisting}
\end{minipage}

    2.1. Có bao nhiêu lệnh được thực hiện? \quad 2.2. Chuyển đoạn lệnh trên sang ngôn ngữ C.
\end{enumerate}
\textbf{Hướng dẫn giải:}
\begin{itemize}
    \item \textbf{Giải câu a (Mã đề B1):}
    \begin{itemize}
        \item \textit{Ý 2.1:} Có $1$ lệnh khởi tạo bên ngoài (\texttt{addi t1, zero, 100}) và $100$ lần lặp, mỗi lần lặp thực hiện $5$ lệnh trong thân \texttt{LOOP}. Tổng số lệnh được thực hiện là: $1 + 100 \times 5 = \mathbf{501\text{ lệnh}}$.
        \item \textit{Ý 2.2:} Đoạn mã C tương ứng: \texttt{for (i = 100; i != 0; i--) \{ a = a + *A; A++; \}} (hoặc viết theo chỉ số mảng: \texttt{for (i = 0; i < 100; i++) a = a + A[i];}).
    \end{itemize}
    \item \textbf{Giải câu b (Mã đề B5):}
    \begin{itemize}
        \item \textit{Ý 2.1:} Mỗi vòng lặp địa chỉ trong \texttt{s0} tăng thêm $8\text{ byte}$ (tương ứng $2$ phần tử 32-bit). Để tăng hết $600\text{ byte}$, số lần thực hiện vòng lặp \texttt{LOOP} là $600 \div 8 = 75\text{ vòng}$. Mỗi vòng có $6$ lệnh, cộng $1$ lệnh khởi tạo ở ngoài là: $1 + 75 \times 6 = \mathbf{451\text{ lệnh}}$.
        \item \textit{Ý 2.2:} Đoạn mã C tương ứng (duyệt $150$ phần tử của mảng \texttt{B}, mỗi bước cộng 2 phần tử liên tiếp vào \texttt{b}): \texttt{for (i = 0; i < 150; i += 2) b = b + B[i] + B[i+1];}.
    \end{itemize}
\end{itemize}

\textbf{Bài tập 3.55 (Đề Cuối kỳ 2021.2 -- Mã đề B2 Câu 3, Mã đề B3 Câu 3 \& Mã đề B4 Câu 3)}
\begin{enumerate}[label=\alph*)]
    \item \textbf{(Mã đề B2 \& B3, Câu 3):} Cho đoạn chương trình vòng lặp sau:

\noindent\begin{minipage}{\linewidth}
\begin{lstlisting}
LOOP: slt  t2, zero, t1
      bne  t2, zero, ELSE
      jal  x0, DONE
ELSE: addi s2, s2, 2
      addi t1, t1, -1
      jal  x0, LOOP
DONE:
\end{lstlisting}
\end{minipage}

    3.1. Xác định giá trị thanh ghi \texttt{s2} sau khi chạy xong nếu khởi tạo: $\texttt{t1} = 20, \texttt{s2} = 0$ (Đề B2) hoặc $\texttt{t1} = 18, \texttt{s2} = 0$ (Đề B3).\\
    3.2. Nếu \texttt{t1} khởi tạo bằng $N$ ($N$ nguyên dương), có bao nhiêu lệnh được thực hiện?
    \item \textbf{(Mã đề B4, Câu 3):} Cho đoạn chương trình 2 vòng lặp lồng nhau:

\noindent\begin{minipage}{\linewidth}
\begin{lstlisting}
LOOP:  addi t2, zero, 0xA
LOOP2: addi s2, s2, 14
       addi t2, t2, -1
       bne  t2, x0, LOOP2
       addi t1, t1, -1
       bne  t1, x0, LOOP
DONE:
\end{lstlisting}
\end{minipage}

    3.1. Giả sử ban đầu $\texttt{t1} = 18, \texttt{s2} = 0$. Xác định giá trị của \texttt{s2} sau khi thực hiện xong.\\
    3.2. Nếu \texttt{t1} khởi tạo bằng $N$ ($N$ nguyên dương), có bao nhiêu lệnh được thực hiện?
\end{enumerate}
\textbf{Hướng dẫn giải:}
\begin{itemize}
    \item \textbf{Giải câu a (Mã đề B2 \& B3):}
    \begin{itemize}
        \item \textit{Ý 3.1:} Chừng nào $0 < \texttt{t1}$ ($\texttt{t1} > 0$), chương trình nhảy vào \texttt{ELSE} cộng thêm $2$ vào \texttt{s2} và giảm \texttt{t1} đi $1$.
        \begin{itemize}
            \item Với $\texttt{t1} = 20$ (Đề B2): $\texttt{s2} = 20 \times 2 = \mathbf{40}$.
            \item Với $\texttt{t1} = 18$ (Đề B3): $\texttt{s2} = 18 \times 2 = \mathbf{36}$.
        \end{itemize}
        \item \textit{Ý 3.2:} Trong $N$ lần lặp đầu tiên (khi $\texttt{t1} = N, N-1, \dots, 1$), mỗi lần thực hiện $5$ lệnh (\texttt{slt, bne, addi s2, addi t1, jal LOOP}). Ở lần cuối cùng khi $\texttt{t1} = 0$, chương trình thực hiện thêm $3$ lệnh (\texttt{slt, bne, jal DONE}) để thoát. Vậy tổng số lệnh thực hiện là: $\mathbf{5N + 3\text{ lệnh}}$.
    \end{itemize}
    \item \textbf{Giải câu b (Mã đề B4):}
    \begin{itemize}
        \item \textit{Ý 3.1:} Vòng lặp trong (\texttt{LOOP2}) lặp $\texttt{0xA} = 10\text{ lần}$, mỗi lần cộng $14$ vào \texttt{s2} (tức mỗi vòng ngoài tăng $10 \times 14 = 140$). Với $\texttt{t1} = 18$ vòng ngoài, giá trị cuối của \texttt{s2} là: $18 \times 140 = \mathbf{2520}$.
        \item \textit{Ý 3.2:} Trong mỗi vòng lặp ngoài (\texttt{LOOP}), số lệnh thực hiện gồm: $1$ lệnh (\texttt{addi t2}) $+ 10 \times 3\text{ lệnh}$ (của \texttt{LOOP2}) $+ 2\text{ lệnh}$ cuối (\texttt{addi t1, bne t1}) $= 33\text{ lệnh}$. Lặp $N$ vòng ngoài sẽ thực hiện tổng cộng: $\mathbf{33N\text{ lệnh}}$.
    \end{itemize}
\end{itemize}

\textbf{Bài tập 3.56 (Đề Cuối kỳ 2021.2 -- Mã đề B6 Câu 2 \& Mã đề B7 Câu 2)}\\
Cho đoạn chương trình hợp ngữ dưới đây:

\noindent\begin{minipage}{\linewidth}
\begin{lstlisting}
    addi t1, zero, K
    add  s2, x0, x0
LOOP:
    slt  t2, x0, t1
    beq  t2, x0, DONE
    addi s2, s2, M
    slli s2, s2, L
    addi t1, t1, -1
    jal  x0, LOOP
DONE:
\end{lstlisting}
\end{minipage}

Hãy xác định (2.1) Tổng số lệnh được thực hiện, và (2.2) Giá trị của thanh ghi \texttt{s2} sau khi kết thúc đoạn lệnh trên trong 2 mã đề:
\begin{enumerate}[label=\alph*)]
    \item \textbf{(Mã đề B6):} $K = 5, M = 3, L = 2$.
    \item \textbf{(Mã đề B7):} $K = 7, M = 1, L = 1$.
\end{enumerate}
\textbf{Hướng dẫn giải:}
\begin{itemize}
    \item \textbf{Công thức đếm số lệnh tổng quát với $K$ vòng lặp:} Có $2$ lệnh khởi tạo ở ngoài $+ K \times 6\text{ lệnh}$ trong thân \texttt{LOOP} $+ 2\text{ lệnh}$ kiểm tra thoát (\texttt{slt, beq}) khi $\texttt{t1} = 0 \implies \text{Tổng số lệnh} = 6K + 4$.
    \item \textbf{Giải câu a (Mã đề B6: $K = 5, M = 3, L = 2$):}
    \begin{itemize}
        \item \textit{Ý 2.1:} Tổng số lệnh thực hiện $= 6 \times 5 + 4 = \mathbf{34\text{ lệnh}}$.
        \item \textit{Ý 2.2:} Trong mỗi vòng lặp, \texttt{s2} được cộng thêm $3$ rồi dịch trái $2\text{ bit}$ (nhân $4$): $\texttt{s2}_{\text{mới}} = (\texttt{s2}_{\text{cũ}} + 3) \times 4$:
        \begin{itemize}
            \item Vòng 1 ($\texttt{t1}=5$): $\texttt{s2} = (0 + 3) \times 4 = 12$.
            \item Vòng 2 ($\texttt{t1}=4$): $\texttt{s2} = (12 + 3) \times 4 = 60$.
            \item Vòng 3 ($\texttt{t1}=3$): $\texttt{s2} = (60 + 3) \times 4 = 252$.
            \item Vòng 4 ($\texttt{t1}=2$): $\texttt{s2} = (252 + 3) \times 4 = 1020$.
            \item Vòng 5 ($\texttt{t1}=1$): $\texttt{s2} = (1020 + 3) \times 4 = \mathbf{4092}$.
        \end{itemize}
    \end{itemize}
    \item \textbf{Giải câu b (Mã đề B7: $K = 7, M = 1, L = 1$):}
    \begin{itemize}
        \item \textit{Ý 2.1:} Tổng số lệnh thực hiện $= 6 \times 7 + 4 = \mathbf{46\text{ lệnh}}$.
        \item \textit{Ý 2.2:} Trong mỗi vòng lặp, $\texttt{s2}_{\text{mới}} = (\texttt{s2}_{\text{cũ}} + 1) \times 2$:
        \begin{itemize}
            \item Vòng 1: $2$; \quad Vòng 2: $6$; \quad Vòng 3: $14$; \quad Vòng 4: $30$; \quad Vòng 5: $62$; \quad Vòng 6: $126$; \quad Vòng 7: $(126 + 1) \times 2 = \mathbf{254}$.
        \end{itemize}
    \end{itemize}
\end{itemize}

\subsection{Nhóm 7: Chương trình con (Thủ tục), Ngăn xếp (Stack) và Lệnh \texttt{jal / jalr}}

\textbf{Bài tập 3.57 (Đề Giữa kỳ 2024.1 -- \#5ebf59 \& Ôn tập Cuối kỳ 2025.2 -- P1, Câu 50)}\\
Những lệnh nào lưu địa chỉ của $PC + 4$ vào một thanh ghi trước khi thực hiện rẽ nhánh / nhảy trong RISC-V?
\begin{enumerate}[label=\Alph*.]
    \item \texttt{jal}
    \item \texttt{beq}
    \item \texttt{j}
    \item \texttt{jr}
    \item \texttt{jalr}
\end{enumerate}
\textbf{Hướng dẫn giải:}
\begin{itemize}
    \item Hai lệnh nhảy và liên kết (Jump and Link) trong RISC-V thực hiện lưu địa chỉ trở về $PC + 4$ vào thanh ghi đích trước khi nhảy là \textbf{\texttt{jal}} và \textbf{\texttt{jalr}}.
    \item \textbf{Kết luận:} Chọn cả \textbf{A (\texttt{jal})} và \textbf{E (\texttt{jalr})}.
\end{itemize}

\textbf{Bài tập 3.58 (Đề Cuối kỳ 2022.2 -- \#a3c257, 2024.1 -- \#e79a4d, 2024.2 \& Ôn tập 2025.2 -- P2, Câu 31)}\\
Cho thân chương trình con hợp ngữ RISC-V gồm có 3 đoạn lệnh như dưới đây:

\noindent\begin{minipage}{\linewidth}
\begin{lstlisting}
# Doan lenh cat noi dung s2 (x18) va s3 (x19) ra ngan xep:
    addi sp, sp, -8
    sw   s3, 4(sp)
    sw   s2, 0(sp)
# Doan lenh khong lien quan den sp va ngan xep:
    ...
# Doan lenh khoi phuc lai noi dung s2, s3 va sp:
\end{lstlisting}
\end{minipage}

Hãy chỉ ra đoạn lệnh phù hợp để khôi phục lại nội dung \texttt{s2, s3} và \texttt{sp}?
\begin{enumerate}[label=\Alph*.]
    \item \texttt{lw s2, 0(sp)} $\rightarrow$ \texttt{lw s3, 4(sp)} $\rightarrow$ \texttt{addi sp, sp, 8}
    \item \texttt{addi sp, sp, 8} $\rightarrow$ \texttt{lw s3, 4(sp)} $\rightarrow$ \texttt{lw s2, 0(sp)}
    \item \texttt{lw s2, 4(sp)} $\rightarrow$ \texttt{lw s3, 0(sp)} $\rightarrow$ \texttt{addi sp, sp, 8}
    \item \texttt{addi sp, sp, 8} $\rightarrow$ \texttt{lw s3, 0(sp)} $\rightarrow$ \texttt{lw s2, 4(sp)}
\end{enumerate}
\textbf{Hướng dẫn giải:}
\begin{itemize}
    \item Lúc cất vào ngăn xếp, \texttt{s2} được lưu tại \texttt{0(sp)} và \texttt{s3} được lưu tại \texttt{4(sp)}.
    \item Khi khôi phục, phải đọc (\texttt{lw}) \texttt{s2} từ \texttt{0(sp)} và \texttt{s3} từ \texttt{4(sp)} ra trước, sau đó mới cộng trả lại $8\text{ byte}$ cho con trỏ \texttt{sp} (\texttt{addi sp, sp, 8}).
    \item \textbf{Kết luận:} Đáp án đúng là \textbf{A (Đoạn 1)}.
\end{itemize}

\textbf{Bài tập 3.59 (Đề Cuối kỳ 2022.2 -- \#9aec2f, 2024.1 -- \#f023de, 2024.2 \& Ôn tập 2025.2 -- P2, Câu 65)}\\
Cho đoạn lệnh như dưới đây, bên trái là địa chỉ của lệnh tương ứng:

\noindent\begin{minipage}{\linewidth}
\begin{lstlisting}
0x00402034        jal  ra, XYZ
0x00402038        ...
0x00402060  XYZ:  addi sp, sp, -8
    ...           ...
0x00402090        ret               # (tuc jalr x0, 0(ra))
\end{lstlisting}
\end{minipage}

Hãy cho biết các giá trị sau (viết theo số Hexa 8 chữ số, ví dụ \texttt{0x01234567}):
\begin{enumerate}[label=\alph*)]
    \item Nội dung các thanh ghi \texttt{PC} và \texttt{ra} (\texttt{x1}) ngay sau khi thực hiện lệnh \texttt{jal ra, XYZ}.
    \item Nội dung thanh ghi \texttt{PC} ngay sau khi thực hiện lệnh \texttt{ret} (\texttt{jalr x0, 0(ra)}).
\end{enumerate}
\textbf{Hướng dẫn giải:}
\begin{itemize}
    \item \textbf{Giải câu a:} Khi thực hiện lệnh \texttt{jal ra, XYZ} tại địa chỉ \texttt{0x00402034}:
    \begin{itemize}
        \item Thanh ghi \texttt{PC} nhảy tới địa chỉ của nhãn \texttt{XYZ}: $\mathbf{\texttt{PC = 0x00402060}}$.
        \item Thanh ghi \texttt{ra} (\texttt{x1}) lưu địa chỉ của lệnh kế tiếp ($PC + 4$): $\mathbf{\texttt{ra = 0x00402038}}$.
    \end{itemize}
    \item \textbf{Giải câu b:} Khi thực hiện lệnh \texttt{ret} (\texttt{jalr x0, 0(ra)}), giá trị trong thanh ghi \texttt{ra} được nạp trả lại cho \texttt{PC}: $\mathbf{\texttt{PC = 0x00402038}}$.
\end{itemize}

\textbf{Bài tập 3.60 (Đề Cuối kỳ 2022.2 -- \#526d03, 2024.1 -- \#115355 \& 2024.2)}\\
Cho chương trình con \texttt{count} sau:

\noindent\begin{minipage}{\linewidth}
\begin{lstlisting}
count:  addi s0, zero, 0
L1:     beq  a0, zero, done
        andi t0, a0, 1
        beq  t0, zero, L2
        addi s0, s0, 1
L2:     srli a0, a0, 1
        j    L1
done:   add  a0, s0, zero
        ret
\end{lstlisting}
\end{minipage}
\begin{enumerate}[label=\alph*)]
    \item Giả sử chương trình con được dịch ra mã máy và nạp vào bộ nhớ. Khi thực hiện lệnh \texttt{beq t0, zero, L2} mà nội dung thanh ghi $\texttt{t0} = 0$ thì sau khi hoàn thành lệnh đó, nội dung thanh ghi \texttt{PC} nhận giá trị bao nhiêu? (Xét cả 2 trường hợp đề bài: TH1 đặt lệnh đầu tiên \texttt{count} ở \texttt{0x00400100}; TH2 cắt từ dòng \texttt{beq t0, zero, L2} làm dòng đầu tiên ở \texttt{0x00400100}).
    \item Giả thiết trước khi thực hiện chương trình con \texttt{count}, thanh ghi $\texttt{a0} = \texttt{0x5678ABCD}$. Hãy cho biết nội dung kết quả trả về tại \texttt{a0} (hoặc \texttt{\$v0}) sau khi thực hiện chương trình con đó (viết theo số thập phân)?
\end{enumerate}
\textbf{Hướng dẫn giải:}
\begin{itemize}
    \item \textbf{Giải câu a:} Khi $\texttt{t0} = 0$, điều kiện của lệnh \texttt{beq t0, zero, L2} thỏa mãn nên chương trình rẽ nhánh đến nhãn \texttt{L2} (bỏ qua lệnh \texttt{addi s0, s0, 1}), tức là \texttt{PC} tăng thêm $2\text{ lệnh} = 8\text{ byte}$ so với địa chỉ của lệnh \texttt{beq t0, zero, L2}:
    \begin{itemize}
        \item Nếu đề bài cắt từ dòng \texttt{beq t0, zero, L2} là dòng đầu tiên ở \texttt{0x00400100} (như bản in trang 10 đề 2024.2): Nhãn \texttt{L2} nằm ở $\texttt{0x00400100} + 8 = \mathbf{\texttt{0x00400108}}$.
        \item Nếu đề bài in đầy đủ từ dòng \texttt{count: addi s0, zero, 0} ở \texttt{0x00400100}: Lệnh \texttt{beq t0, zero, L2} là lệnh thứ 4 (ở \texttt{0x0040010C}), nên nhãn \texttt{L2} (lệnh thứ 6) nằm ở \textbf{\texttt{0x00400114}}.
    \end{itemize}
    \item \textbf{Giải câu b:} Thuật toán dùng \texttt{andi t0, a0, 1} để kiểm tra từng bit của \texttt{a0} từ phải sang trái; nếu bit bằng 1 thì tăng bộ đếm \texttt{s0} lên 1. Vậy hàm trả về \textbf{tổng số bit 1} trong biểu diễn nhị phân của $\texttt{a0} = \texttt{0x5678ABCD}$:
    \begin{itemize}
        \item $\texttt{5} (0101_2) \rightarrow 2\text{ bit }1$; \quad $\texttt{6} (0110_2) \rightarrow 2\text{ bit }1$; \quad $\texttt{7} (0111_2) \rightarrow 3\text{ bit }1$; \quad $\texttt{8} (1000_2) \rightarrow 1\text{ bit }1$.
        \item $\texttt{A} (1010_2) \rightarrow 2\text{ bit }1$; \quad $\texttt{B} (1011_2) \rightarrow 3\text{ bit }1$; \quad $\texttt{C} (1100_2) \rightarrow 2\text{ bit }1$; \quad $\texttt{D} (1101_2) \rightarrow 3\text{ bit }1$.
        \item Tổng cộng: $2 + 2 + 3 + 1 + 2 + 3 + 2 + 3 = \mathbf{18}$.
    \end{itemize}
\end{itemize}

\subsection{Nhóm 8: Dịch Mã máy -- Hợp ngữ, Tính địa chỉ đích và Giải mã Bản đồ bộ nhớ (Memory Map)}

\textbf{Bài tập 3.61 (Ôn tập Cuối kỳ 2025.2 -- P1, Câu 65, 19, 78, 22, 11, 14, 5)}\\
Hãy dịch 7 câu lệnh hợp ngữ RISC-V dưới đây sang mã máy biểu diễn ở hệ Hexa (8 chữ số, bắt buộc có \texttt{0x} ở đầu):
\begin{enumerate}[label=\alph*)]
    \item \textbf{(Câu 65 -- Kiểu R):} \texttt{sub t0, t2, t1}
    \item \textbf{(Câu 19 -- Kiểu I tính toán):} \texttt{addi s2, s1, -14}
    \item \textbf{(Câu 78 -- Kiểu I đọc bộ nhớ):} \texttt{lh t5, -14(s1)} (và biến thể \texttt{lh t1, -13(s2)})
    \item \textbf{(Câu 22 -- Kiểu S ghi bộ nhớ):} \texttt{sw s11, -16(s4)}
    \item \textbf{(Câu 11 -- Kiểu B rẽ nhánh xuôi):} \texttt{beq s0, t5, L1} (biết địa chỉ lệnh là \texttt{0x00401050}, địa chỉ \texttt{L1} là \texttt{0x00401060}).
    \item \textbf{(Câu 14 -- Kiểu B rẽ nhánh ngược):} \texttt{bne t0, s0, L2} (biết địa chỉ lệnh là \texttt{0x0040008C}, địa chỉ \texttt{L2} là \texttt{0x00400070}).
    \item \textbf{(Câu 5 -- Kiểu J nhảy):} \texttt{jal x0, L1} (biết địa chỉ lệnh là \texttt{0x004030AC}, địa chỉ \texttt{L1} là \texttt{0x00403074}).
\end{enumerate}
\textbf{Hướng dẫn giải:}
\begin{itemize}
    \item \textbf{Giải câu a (\texttt{sub t0, t2, t1} -- Kiểu R):}
    \begin{itemize}
        \item Số hiệu thanh ghi: $\texttt{rd} = \texttt{t0} = \texttt{x5}$ (\texttt{00101}); $\texttt{rs1} = \texttt{t2} = \texttt{x7}$ (\texttt{00111}); $\texttt{rs2} = \texttt{t1} = \texttt{x6}$ (\texttt{00110}).
        \item Lệnh \texttt{sub} có $\texttt{funct7} = \texttt{0100000}$, $\texttt{funct3} = \texttt{000}$, $\texttt{opcode} = \texttt{0110011}$.
        \item Ghép 32 bit (\texttt{funct7 | rs2 | rs1 | funct3 | rd | op}):
        $$\texttt{0100000\_00110\_00111\_000\_00101\_0110011}_{(2)} = \mathbf{\texttt{0x406382B3}}$$
    \end{itemize}
    \item \textbf{Giải câu b (\texttt{addi s2, s1, -14} -- Kiểu I):}
    \begin{itemize}
        \item Số hiệu: $\texttt{rd} = \texttt{s2} = \texttt{x18}$ (\texttt{10010}); $\texttt{rs1} = \texttt{s1} = \texttt{x9}$ (\texttt{01001}); $\texttt{funct3} = \texttt{000}$; $\texttt{op} = \texttt{0010011}$.
        \item Đổi $\texttt{imm} = -14$ sang bù 2 12-bit: $+14 = \texttt{0000\,0000\,1110}_2 \implies -14 = \texttt{1111\,1111\,0010}_2 = \texttt{0xFF2}$.
        \item Ghép 32 bit (\texttt{imm[11:0] | rs1 | funct3 | rd | op}):
        $$\texttt{1111\,1111\,0010\_01001\_000\_10010\_0010011}_{(2)} = \mathbf{\texttt{0xFF248913}}$$
    \end{itemize}
    \item \textbf{Giải câu c (\texttt{lh t5, -14(s1)} -- Kiểu I Load):}
    \begin{itemize}
        \item Số hiệu: $\texttt{rd} = \texttt{t5} = \texttt{x30}$ (\texttt{11110}); $\texttt{rs1} = \texttt{s1} = \texttt{x9}$ (\texttt{01001}); $\texttt{funct3} = \texttt{001}$ (\texttt{lh}); $\texttt{op} = \texttt{0000011}$; $\texttt{imm} = -14 = \texttt{1111\,1111\,0010}_2$.
        \item Ghép 32 bit: $\texttt{1111\,1111\,0010\_01001\_001\_11110\_0000011}_{(2)} = \mathbf{\texttt{0xFF249F03}}$.
        \item (Với biến thể \texttt{lh t1, -13(s2)}: $\texttt{imm} = -13 = \texttt{1111\,1111\,0011}_2$, $\texttt{rs1} = 18 = \texttt{10010}_2$, $\texttt{rd} = 6 = \texttt{00110}_2 \implies \mathbf{\texttt{0xFF391303}}$).
    \end{itemize}
    \item \textbf{Giải câu d (\texttt{sw s11, -16(s4)} -- Kiểu S Store):}
    \begin{itemize}
        \item Số hiệu: $\texttt{rs2} = \texttt{s11} = \texttt{x27}$ (\texttt{11011}); $\texttt{rs1} = \texttt{s4} = \texttt{x20}$ (\texttt{10100}); $\texttt{funct3} = \texttt{010}$; $\texttt{op} = \texttt{0100011}$.
        \item Đổi $\texttt{imm} = -16$ sang bù 2 12-bit: $\texttt{1111\,1111\,0000}_2 \implies \texttt{imm[11:5]} = \texttt{1111111}$, $\texttt{imm[4:0]} = \texttt{10000}$.
        \item Ghép 32 bit (\texttt{imm[11:5] | rs2 | rs1 | funct3 | imm[4:0] | op}):
        $$\texttt{1111111\_11011\_10100\_010\_10000\_0100011}_{(2)} = \mathbf{\texttt{0xFFBA2823}}$$
    \end{itemize}
    \item \textbf{Giải câu e (\texttt{beq s0, t5, L1} -- Kiểu B):}
    \begin{itemize}
        \item Độ lệch $\texttt{imm} = \texttt{0x00401060} - \texttt{0x00401050} = \texttt{0x10} = +16 = \texttt{0\_0\_000000\_1000\_0}_{(2)}$ (từ bit 12 đến bit 0).
        \item Số hiệu: $\texttt{rs1} = \texttt{s0} = 8$ (\texttt{01000}); $\texttt{rs2} = \texttt{t5} = 30$ (\texttt{11110}); $\texttt{funct3} = \texttt{000}$; $\texttt{op} = \texttt{1100011}$.
        \item Ghép theo B-type (\texttt{imm[12,10:5] | rs2 | rs1 | funct3 | imm[4:1,11] | op}):
        $$\texttt{0\_000000\_11110\_01000\_000\_1000\_0\_1100011}_{(2)} = \mathbf{\texttt{0x01E40863}}$$
    \end{itemize}
    \item \textbf{Giải câu f (\texttt{bne t0, s0, L2} -- Kiểu B):}
    \begin{itemize}
        \item Độ lệch $\texttt{imm} = \texttt{0x00400070} - \texttt{0x0040008C} = -\texttt{0x1C} = -28 = \texttt{1\_1\_111111\_0010\_0}_{(2)}$ (từ bit 12 đến bit 0).
        \item Số hiệu: $\texttt{rs1} = \texttt{t0} = 5$ (\texttt{00101}); $\texttt{rs2} = \texttt{s0} = 8$ (\texttt{01000}); $\texttt{funct3} = \texttt{001}$; $\texttt{op} = \texttt{1100011}$.
        \item Ghép theo B-type: $\texttt{1\_111111\_01000\_00101\_001\_0010\_1\_1100011}_{(2)} = \mathbf{\texttt{0xFE8292E3}}$.
    \end{itemize}
    \item \textbf{Giải câu g (\texttt{jal x0, L1} -- Kiểu J):}
    \begin{itemize}
        \item Độ lệch $\texttt{imm} = \texttt{0x00403074} - \texttt{0x004030AC} = -\texttt{0x38} = -56 = \texttt{1\_11111111\_1\_1111100100\_0}_{(2)}$ (từ bit 20 đến bit 0).
        \item Ghép theo J-type (\texttt{imm[20,10:1,11,19:12] | rd | op}):
        $$\texttt{1\_1111100100\_1\_11111111\_00000\_1101111}_{(2)} = \mathbf{\texttt{0xFC9FF06F}}$$
    \end{itemize}
\end{itemize}

\textbf{Bài tập 3.62 (Đề Cuối kỳ 2024.1 -- \#b9e098 \& Ôn tập Cuối kỳ 2025.2 -- P1, Câu 8)}\\
Cho lệnh máy RISC-V được biểu diễn theo số Hexa: \texttt{0x012B2423}. Hãy chỉ ra lệnh hợp ngữ tương ứng:
\begin{enumerate}[label=\Alph*.]
    \item \texttt{add t0, s4, s5}
    \item \texttt{lw t2, 4(s6)}
    \item \texttt{sw s2, 8(s6)}
    \item \texttt{addi t0, s3, -4}
\end{enumerate}
\textbf{Hướng dẫn giải:}
\begin{itemize}
    \item Đổi 2 chữ số Hexa cuối $\texttt{23}_{16} = \texttt{0010\,0011}_2 \implies 7\text{ bit opcode} = \texttt{0100011}_2 = 35_{10}$ (Nhóm lệnh Store kiểu S).
    \item Tách 32 bit theo khuôn dạng S-type (\texttt{imm[11:5] | rs2 | rs1 | funct3 | imm[4:0] | op}):
    $$\texttt{0x012B2423} = \underbrace{\texttt{0000000}}_{\texttt{imm[11:5]}=0} \mid \underbrace{\texttt{10010}}_{\texttt{rs2}=18 (\texttt{s2})} \mid \underbrace{\texttt{10110}}_{\texttt{rs1}=22 (\texttt{s6})} \mid \underbrace{\texttt{010}}_{\texttt{funct3}=2 (\texttt{sw})} \mid \underbrace{\texttt{01000}}_{\texttt{imm[4:0]}=8} \mid \underbrace{\texttt{0100011}}_{\text{opcode}}$$
    \item Vậy câu lệnh hợp ngữ là: \textbf{\texttt{sw s2, 8(s6)}} $\rightarrow$ Chọn \textbf{C}.
\end{itemize}

\textbf{Bài tập 3.63 (Đề Cuối kỳ 2022.2 \& Cuối kỳ 2024.2 -- \#b7c77a)}\\
Cho lệnh máy được biểu diễn theo số Hexa: \texttt{0xFFC98293} (trong đề cũ chuyển từ MIPS là \texttt{0x2268FFFC}). Hãy chỉ ra lệnh hợp ngữ tương ứng:
\begin{enumerate}[label=\Alph*.]
    \item \texttt{lw t2, 4(s6)}
    \item \texttt{add t0, s4, s5}
    \item \texttt{sw s2, 8(s6)}
    \item \texttt{addi t0, s3, -4}
\end{enumerate}
\textbf{Hướng dẫn giải:}
\begin{itemize}
    \item Trong mã máy \texttt{0xFFC98293} (hoặc \texttt{0x2268FFFC}): hằng số ở đầu/cuối là $\texttt{0xFFC} = -4$, thanh ghi nguồn là $\texttt{s3}$ ($\texttt{x19} = 10011_2$), thanh ghi đích là $\texttt{t0}$ ($\texttt{x5} = 00101_2$), phép toán là cộng hằng số \texttt{addi}.
    \item \textbf{Kết luận:} Đáp án đúng là \textbf{D (\texttt{addi t0, s3, -4})}.
\end{itemize}

\textbf{Bài tập 3.64 (Đề Cuối kỳ 2024.1 -- \#3ad10f \& Ôn tập Cuối kỳ 2025.2 -- P1, Câu 62)}\\
Giả thiết lệnh \texttt{bge s0, s1, L1} của RISC-V có mã máy là \texttt{0x02945063}, được đặt tại word nhớ có địa chỉ \texttt{0x00404014}. Hãy chỉ ra địa chỉ đích của lệnh đó (nếu điều kiện thỏa mãn). Biết $\text{BTA} = PC + \text{SignExt}(\{\texttt{imm}_{12:1}, \texttt{1'b0}\})$.
\begin{enumerate}[label=\Alph*.]
    \item \texttt{0x0040401F}
    \item \texttt{0x00404034}
    \item \texttt{0x00404030}
    \item \texttt{0x00404018}
\end{enumerate}
\textbf{Hướng dẫn giải:}
\begin{itemize}
    \item Đổi mã máy \texttt{0x02945063} sang nhị phân và tách theo B-type:
    $$\underbrace{\texttt{0\_000001}}_{\texttt{imm[12,10:5]}} \mid \underbrace{\texttt{01001}}_{\texttt{rs2}=9} \mid \underbrace{\texttt{01000}}_{\texttt{rs1}=8} \mid \underbrace{\texttt{101}}_{\texttt{funct3}} \mid \underbrace{\texttt{0000\_0}}_{\texttt{imm[4:1,11]}} \mid \underbrace{\texttt{1100011}}_{\text{opcode}}$$
    \item Ghép lại hằng số 13-bit $\texttt{imm}_{12:0} = \{\texttt{imm}_{12}, \texttt{imm}_{11}, \texttt{imm}_{10:5}, \texttt{imm}_{4:1}, 0\}$:
    $$\texttt{imm} = \texttt{0\_0\_000001\_0000\_0}_{(2)} = 0000\,0010\,0000_{(2)} = \texttt{0x20} = 32_{(10)}$$
    \item Địa chỉ đích $\text{BTA} = PC + \texttt{imm} = \texttt{0x00404014} + \texttt{0x20} = \mathbf{\texttt{0x00404034}}$.
    \item \textbf{Kết luận:} Đáp án đúng là \textbf{B (\texttt{0x00404034})}.
\end{itemize}

\textbf{Bài tập 3.65 (Đề Cuối kỳ 2020.2 -- Trang 4, Giữa kỳ 2024.1 -- \#a2a36c \& Ôn tập 2025.2 -- Câu 47, 57)}
\begin{enumerate}[label=\alph*)]
    \item \textbf{(Đề CK 2020.2 -- Trang 4):} Lệnh \texttt{beq t3, s4, L1} trong RISC-V được đặt tại địa chỉ \texttt{0x00401000}, lệnh kế tiếp tại \texttt{0x00401004}. Biết nhãn đích \texttt{L1} nằm tại địa chỉ \texttt{0x0040104C}. Hãy xác định giá trị độ lệch byte (offset) và giá trị 12-bit \texttt{imm[12:1]} mã hóa trong lệnh.
    \item \textbf{(Đề GK 2024.1 -- \#a2a36c \& Ôn tập 2025.2 Câu 47):} Cho chương trình RISC-V bắt đầu từ địa chỉ \texttt{0x00400000}: (1) \texttt{addi s0, zero, 100}; (2) \texttt{slti t0, s0, -100}; (3) \texttt{beq t0, zero, else}; (4) \texttt{addi t2, t2, 1}; (5) \texttt{j endif}; (6) \texttt{else: addi t2, t2, -1}; (7) \texttt{endif: nop}. Giá trị thập phân của trường 12-bit immediate trong lệnh \texttt{beq t0, zero, else} là bao nhiêu?\\
    A. $8$ \qquad B. $12$ \qquad C. $3$ \qquad D. $6$
    \item \textbf{(Ôn tập 2025.2 -- Câu 57):} Giả sử \texttt{PC} hiện tại (địa chỉ của lệnh \texttt{beq}) là \texttt{0x2000}. Với câu lệnh \texttt{beq x1, x2, -12}, nếu $\texttt{x1} == \texttt{x2}$ thì địa chỉ lệnh tiếp theo (Target PC) là bao nhiêu?\\
    A. \texttt{0x1FE8} \qquad B. \texttt{0x1FF8} \qquad C. \texttt{0x1FF4} \qquad D. \texttt{0x200C}
\end{enumerate}
\textbf{Hướng dẫn giải:}
\begin{itemize}
    \item \textbf{Giải câu a:} Độ lệch tính theo byte là: $\text{Offset} = \texttt{0x0040104C} - \texttt{0x00401000} = \mathbf{\texttt{0x4C}} = \mathbf{76\text{ byte}}$. Nếu tính trường 12-bit \texttt{imm[12:1]} lưu trong mã máy (bỏ bit 0, tức chia 2) thì giá trị là $76 \div 2 = \mathbf{38}$ (\texttt{0x26}); nếu tính số lệnh nhảy qua (mỗi lệnh 4 byte) là $76 \div 4 = 19$ (hoặc $18$ nếu trừ $PC+4$ theo đề cũ).
    \item \textbf{Giải câu b:} Lệnh (3) \texttt{beq} nằm ở địa chỉ \texttt{0x00400008}, nhãn \texttt{else} ở lệnh (6) nằm ở địa chỉ \texttt{0x00400014}. Độ lệch byte là $\texttt{0x14} - \texttt{0x08} = 12\text{ byte}$. Trường 12-bit \texttt{imm[12:1]} được mã hóa trong lệnh có giá trị thập phân là: $12 \div 2 = \mathbf{6}$ $\rightarrow$ Chọn \textbf{D (6)}.
    \item \textbf{Giải câu c:} $\text{Target PC} = \texttt{0x2000} - 12 = \texttt{0x2000} - \texttt{0x0C} = \mathbf{\texttt{0x1FF4}}$ $\rightarrow$ Chọn \textbf{C}.
\end{itemize}

\textbf{Bài tập 3.66 (Đề Cuối kỳ 2023.1 -- \#5a6d48, \#e08504, 2025.1 -- \#794d71, \#0c9d, Giữa kỳ 2024.1 -- \#1c2f2d, \#65dce1 \& Ôn tập 2025.2 -- Câu 16, 60, 77)}\\
Hãy xác định khuôn dạng lệnh, thanh ghi đích (\texttt{rd}) và các thanh ghi nguồn (\texttt{rs1, rs2}) cho các mã máy RISC-V dưới đây:
\begin{enumerate}[label=\alph*)]
    \item \textbf{(CK 2023.1 -- \#e08504):} Lệnh mã máy \texttt{0x00628863} (\texttt{bne t0, t1, 16}) thuộc định dạng nào (\texttt{R, I, B, J})?
    \item \textbf{(CK 2023.1 -- \#5a6d48):} Lệnh mã máy \texttt{0x00008067} (\texttt{jalr zero, ra, 0}) có toán hạng đích lưu tại đâu?
    \item \textbf{(Ôn tập 2025.2 -- P2, Câu 16):} Cho mã máy \texttt{0000000 00111 00000 000 11111 0110011}. Toán hạng đích của nó là gì?
    \item \textbf{(GK 2024.1 -- \#1c2f2d \& CK 2025.1 -- \#794d71):} Tìm toán hạng đích (\texttt{rd}) của lệnh R-format có mã máy: (1) \texttt{0x007302b3}; và (2) \texttt{0x01c303b3}.
    \item \textbf{(GK 2024.1 -- \#65dce1 \& CK 2025.1 -- \#0c9d):} Tìm các thanh ghi nguồn (\texttt{rs1, rs2}) của lệnh \texttt{beq} có mã máy: (1) \texttt{0x01ee8463}; và (2) \texttt{0x01de0463}.
\end{enumerate}
\textbf{Hướng dẫn giải:}
\begin{itemize}
    \item \textbf{Câu a:} Lệnh \texttt{bne} có opcode $\texttt{0x63} = \texttt{1100011}_2$, thuộc định dạng \textbf{B-format (B-type)}.
    \item \textbf{Câu b:} Trong \texttt{jalr zero, ra, 0}, toán hạng đích là \textbf{thanh ghi \texttt{zero} (\texttt{x0})}.
    \item \textbf{Câu c:} Trường \texttt{rd} (5 bit đứng ngay trước 7 bit opcode) là $\texttt{11111}_2 = 31 \implies$ Thanh ghi \texttt{x31}, tức là \textbf{Thanh ghi \texttt{t6}}.
    \item \textbf{Câu d:} Đổi 3 chữ số Hexa cuối sang 12 bit nhị phân (\texttt{rd} 5 bit + \texttt{op} 7 bit):
    \begin{itemize}
        \item Với \texttt{0x007302b3}: $\texttt{2b3}_{16} = \underbrace{\texttt{00101}}_{\texttt{rd}=5}\,\underbrace{\texttt{0110011}}_{\text{op}} \implies \mathbf{\texttt{x5}}$ (\texttt{t0}).
        \item Với \texttt{0x01c303b3}: $\texttt{3b3}_{16} = \underbrace{\texttt{00111}}_{\texttt{rd}=7}\,\underbrace{\texttt{0110011}}_{\text{op}} \implies \mathbf{\texttt{x7}}$ (\texttt{t2}).
    \end{itemize}
    \item \textbf{Câu e:} Trong lệnh \texttt{beq rs1, rs2, label}, bit 24:20 là \texttt{rs2} và bit 19:15 là \texttt{rs1}:
    \begin{itemize}
        \item Với \texttt{0x01ee8463}: $\texttt{01ee8}_{16} = \texttt{0000000\_}\underbrace{\texttt{11110}}_{\texttt{rs2}=30}\_\underbrace{\texttt{11101}}_{\texttt{rs1}=29}\_\texttt{000} \implies$ Nguồn \texttt{rs1, rs2} là \textbf{\texttt{x29, x30}}.
        \item Với \texttt{0x01de0463}: $\texttt{01de0}_{16} = \texttt{0000000\_}\underbrace{\texttt{11101}}_{\texttt{rs2}=29}\_\underbrace{\texttt{11100}}_{\texttt{rs1}=28}\_\texttt{000} \implies$ Nguồn \texttt{rs1, rs2} là \textbf{\texttt{x28, x29}}.
    \end{itemize}
\end{itemize}

\textbf{Bài tập 3.67 (Đề Cuối kỳ 2025.1 -- \#15fd1c \& Ôn tập Cuối kỳ 2025.2 -- P2, Câu 48)}\\
Cho chương trình hợp ngữ RISC-V dưới đây đã được nạp vào bộ nhớ chính. Giả sử lệnh đầu tiên của chương trình nằm tại địa chỉ \texttt{0x00400000}:

\noindent\begin{minipage}{\linewidth}
\begin{lstlisting}
    lui  s0, 0xFFFF0          # 0x00400000
    addi s0, zero, 0          # 0x00400004
    slti t1, s0, 100          # 0x00400008
    beq  t1, zero, else       # 0x0040000C
    addi t2, zero, 1          # 0x00400010
    addi t3, zero, 20         # 0x00400014
    j    endif                # 0x00400018
else:
    addi t2, zero, -1         # 0x0040001C
    addi t3, zero, -20        # 0x00400020
endif: nop                    # 0x00400024
\end{lstlisting}
\end{minipage}

Giá trị hệ thập lục phân (Hexa) của thanh ghi \texttt{PC} ngay sau khi thực hiện lệnh \texttt{beq} là bao nhiêu?

\textbf{Hướng dẫn giải:}
\begin{itemize}
    \item Lệnh thứ 2 (\texttt{addi s0, zero, 0}) gán lại $\texttt{s0} = 0 + 0 = 0$ (ghi đè giá trị của lệnh \texttt{lui} đầu tiên).
    \item Lệnh thứ 3 (\texttt{slti t1, s0, 100}) kiểm tra $0 < 100 \rightarrow$ Đúng, nên $\texttt{t1} = 1$.
    \item Lệnh thứ 4 (\texttt{beq t1, zero, else} tại địa chỉ \texttt{0x0040000C}) kiểm tra $\texttt{t1} == 0$ ($1 == 0 \rightarrow$ Sai), nên \textbf{không rẽ nhánh}.
    \item Do không rẽ nhánh, thanh ghi \texttt{PC} tăng thêm $4$ để trỏ sang lệnh kế tiếp thứ 5:
    $$PC = \texttt{0x0040000C} + 4 = \mathbf{\texttt{0x00400010}}$$
\end{itemize}

\textbf{Bài tập 3.68 (Giải chi tiết Đề Giữa kỳ 2025.2 -- Đề 0013, Bản 1: Câu 67 đến Câu 77)}\\
Cho bản đồ bộ nhớ (Memory Map) của hệ thống RISC-V 32-bit lưu trữ theo thứ tự đầu nhỏ (Little-Endian) từ địa chỉ \texttt{0x400000} đến \texttt{0x40001B} như bảng dưới đây:
\begin{center}
\renewcommand{\arraystretch}{1.15}
\footnotesize
\begin{tabularx}{\linewidth}{|C{1.8cm}|C{3.1cm}|C{2.4cm}|Y|}
\hline
\textbf{Địa chỉ} & \textbf{4 byte (+0..+3)} & \textbf{Mã máy 32-bit} & \textbf{Giải mã sang lệnh hợp ngữ RISC-V} \\ \hline
\texttt{0x400000} & \texttt{B3, 02, 73, 00} & \texttt{0x007302B3} & \texttt{add t0, t1, t2} ($\texttt{rd}=\texttt{t0}, \texttt{rs1}=\texttt{t1}, \texttt{rs2}=\texttt{t2}$) \\ \hline
\texttt{0x400004} & \texttt{13, 0E, 54, 00} & \texttt{0x00540E13} & \texttt{addi t3, s0, 5} ($\texttt{rd}=\texttt{t3}, \texttt{rs1}=\texttt{s0}, \texttt{imm}=5$) \\ \hline
\texttt{0x400008} & \texttt{93, DE, 74, 00} & \texttt{0x0074DE93} & \texttt{srli t4, s1, 7} (dịch phải logic $7\text{ bit}$) \\ \hline
\texttt{0x40000C} & \texttt{23, 24, E9, 01} & \texttt{0x01E92423} & \texttt{sw t5, 8(s2)} ($\texttt{rs2}=\texttt{t5}, \texttt{rs1}=\texttt{s2}, \texttt{imm}=8$) \\ \hline
\texttt{0x400010} & \texttt{E3, 88, 0F, FE} & \texttt{0xFE0F88E3} & \texttt{beq t6, x0, -16} ($\texttt{rs1}=\texttt{t6}, \texttt{rs2}=\texttt{x0}$) \\ \hline
\texttt{0x400014} & \texttt{B7, 59, 34, 12} & \texttt{0x123459B7} & \texttt{lui s3, 0x12345} ($\texttt{rd}=\texttt{x19}=\texttt{s3}$) \\ \hline
\texttt{0x400018} & \texttt{EF, F0, 9F, FE} & \texttt{0xFE9FF0EF} & \texttt{jal ra, -24} ($\texttt{rd}=\texttt{x1}=\texttt{ra}$) \\ \hline
\end{tabularx}
\end{center}
Hãy trả lời các câu hỏi trong Đề 0013 (Bản 1):
\begin{enumerate}[label=\textbf{Câu \arabic*.}, start=67]
    \item Hãy cho biết giá trị số 32-bit ở địa chỉ \texttt{0x400000} ở hệ 16.
    \item Sau đoạn lệnh: \texttt{lui s0, 0x400} $\rightarrow$ \texttt{addi s0, s0, 0x18} $\rightarrow$ \texttt{lh t1, 0(s0)}, hai thanh ghi \texttt{s0} và \texttt{t1} bằng bao nhiêu?
    \item Lệnh ở địa chỉ \texttt{0x400014} có thanh ghi đích \texttt{rd} là gì?
    \item Lệnh thứ 3 trong bản đồ (tại \texttt{0x400008}) là lệnh tương đương với phép chia. Đó là phép chia với hằng số bằng mấy?
    \item Hằng số của lệnh ở địa chỉ \texttt{0x400008} bằng mấy (viết ở hệ 16)?
    \item Hãy cho biết thao tác của lệnh ở địa chỉ \texttt{0x400010} là gì (ví dụ: \texttt{add, bne, lui})?
\end{enumerate}
\textbf{Hướng dẫn giải:}
\begin{itemize}
    \item \textbf{Câu 67:} Ghép ngược 4 byte từ \texttt{0x400003} về \texttt{0x400000} theo Little-Endian: \textbf{\texttt{0x007302B3}}.
    \item \textbf{Câu 68:} Sau 2 lệnh đầu, $\mathbf{\texttt{s0 = 0x00400018}}$. Lệnh \texttt{lh t1, 0(s0)} đọc 2 byte tại \texttt{0x400018} (\texttt{0xEF}) và \texttt{0x400019} (\texttt{0xF0}), ghép lại là \texttt{0xF0EF}, mở rộng dấu âm 32-bit thành: $\mathbf{\texttt{t1 = 0xFFFFF0EF}}$.
    \item \textbf{Câu 69:} Lệnh ở \texttt{0x400014} là \texttt{0x123459B7}. Đổi 3 chữ cuối $\texttt{9B7}_{16} = \underbrace{\texttt{10011}}_{\texttt{rd}=19}\,\underbrace{\texttt{0110111}}_{\text{lui}}$. Thanh ghi $\texttt{x19}$ có tên là \textbf{\texttt{s3}}.
    \item \textbf{Câu 70:} Lệnh thứ 3 ở \texttt{0x400008} là \texttt{0x0074DE93} (\texttt{srli t4, s1, 7}). Dịch phải $7\text{ bit}$ tương đương chia cho $2^7 = \mathbf{128}$.
    \item \textbf{Câu 71:} Hằng số dịch của lệnh \texttt{0x0074DE93} là 12 bit đầu \texttt{0x007} $\implies$ Trả lời: \textbf{\texttt{0x7}} (hoặc \texttt{0x007}).
    \item \textbf{Câu 72:} Lệnh ở \texttt{0x400010} (\texttt{0xFE0F88E3}) có opcode $\texttt{1100011}_2$ (Branch) và $\texttt{funct3} = \texttt{000}_2$ $\implies$ Thao tác là \textbf{\texttt{beq}}.
\end{itemize}

\textbf{Bài tập 3.69 (Giải chi tiết Đề Giữa kỳ 2025.2 -- Đề 0013, Bản 2: Câu 936 đến Câu 946)}\\
Sử dụng cùng bản đồ bộ nhớ của Đề 0013 ở Bài tập 3.68 (nhưng tại địa chỉ \texttt{0x40000A} đổi byte \texttt{0x74} thành \texttt{0x34}, tức mã máy tại \texttt{0x400008} là \texttt{0x0034DE93}), hãy trả lời các câu hỏi sau:
\begin{enumerate}[label=\textbf{Câu \arabic*.}, start=936]
    \item Giá trị số 32-bit ở địa chỉ \texttt{0x400018} ở hệ 16 là bao nhiêu?
    \item Sau đoạn lệnh: \texttt{lui s0, 0x400} $\rightarrow$ \texttt{addi s0, s0, 0x16} $\rightarrow$ \texttt{lh t1, 0(s0)}, hai thanh ghi \texttt{s0} và \texttt{t1} bằng bao nhiêu?
    \item Lệnh ở địa chỉ \texttt{0x400000} có thanh ghi nguồn \texttt{rs1} là gì?
    \item Lệnh thứ 3 (tại \texttt{0x400008}: \texttt{0x0034DE93}) tương đương với phép chia cho hằng số bằng mấy?
    \item Trong bản đồ, lệnh \texttt{beq} duy nhất nằm ở địa chỉ nào?
    \item Hãy cho biết thao tác của lệnh ở địa chỉ \texttt{0x400008} là gì?
\end{enumerate}
\textbf{Hướng dẫn giải:}
\begin{itemize}
    \item \textbf{Câu 936:} Ghép ngược 4 byte từ \texttt{0x40001B} về \texttt{0x400018} được: \textbf{\texttt{0xFE9FF0EF}}.
    \item \textbf{Câu 937:} Sau 2 lệnh đầu, $\mathbf{\texttt{s0 = 0x00400016}}$. Lệnh \texttt{lh t1, 0(s0)} đọc 2 byte tại \texttt{0x400016} (\texttt{0x34}) và \texttt{0x400017} (\texttt{0x12}), ghép lại thành \texttt{0x1234}, mở rộng dấu dương được: $\mathbf{\texttt{t1 = 0x00001234}}$.
    \item \textbf{Câu 938:} Lệnh ở \texttt{0x400000} (\texttt{0x007302B3}) có $\texttt{rs1}$ ở bit 19:15 là $\texttt{00110}_2 = 6$, tức là thanh ghi \textbf{\texttt{t1}} (\texttt{x6}).
    \item \textbf{Câu 939:} Mã máy \texttt{0x0034DE93} là lệnh \texttt{srli t4, s1, 3}. Dịch phải $3\text{ bit}$ tương đương chia cho $2^3 = \mathbf{8}$.
    \item \textbf{Câu 940:} Lệnh \texttt{beq} duy nhất nằm ở địa chỉ \textbf{\texttt{0x400010}}.
    \item \textbf{Câu 941:} Lệnh ở \texttt{0x400008} (\texttt{0x0034DE93}) có opcode $\texttt{0010011}_2$, $\texttt{funct3} = \texttt{101}_2$, 7 bit cao bằng $0$ $\implies$ Thao tác là \textbf{\texttt{srli}}.
\end{itemize}

\textbf{Bài tập 3.70 (Giải chi tiết Đề Giữa kỳ 2025.2 -- Đề 0012: Câu 100 đến Câu 109)}\\
Cho bản đồ bộ nhớ của Đề 0012 từ địa chỉ \texttt{0x400000} đến \texttt{0x40001B} như bảng dưới đây:
\begin{center}
\renewcommand{\arraystretch}{1.15}
\footnotesize
\begin{tabularx}{\linewidth}{|C{1.8cm}|C{3.1cm}|C{2.4cm}|Y|}
\hline
\textbf{Địa chỉ} & \textbf{4 byte (+0..+3)} & \textbf{Mã máy 32-bit} & \textbf{Giải mã sang lệnh hợp ngữ RISC-V} \\ \hline
\texttt{0x400000} & \texttt{33, 0E, 73, 40} & \texttt{0x40730E33} & \texttt{sub t3, t1, t2} ($\texttt{rd}=\texttt{t3}, \texttt{rs1}=\texttt{t1}, \texttt{rs2}=\texttt{t2}$) \\ \hline
\texttt{0x400004} & \texttt{37, 0B, 04, 00} & \texttt{0x00040B37} & \texttt{lui s6, 0x40} (U-type: $\texttt{rd}=\texttt{x22}=\texttt{s6}$) \\ \hline
\texttt{0x400008} & \texttt{13, 0B, 8B, 03} & \texttt{0x038B0B13} & \texttt{addi s6, s6, 56} (I-type: $\texttt{imm}=56=\texttt{0x38}$) \\ \hline
\texttt{0x40000C} & \texttt{E3, 0C, CB, FF} & \texttt{0xFFCB0CE3} & \texttt{beq s6, s11, -8} (B-type) \\ \hline
\texttt{0x400010} & \texttt{33, 83, C3, 01} & \texttt{0x01C38333} & \texttt{add t1, t2, t3} (R-type) \\ \hline
\texttt{0x400014} & \texttt{6F, F0, 1F, FF} & \texttt{0xFF1FF06F} & \texttt{jal x0, -16} (J-type) \\ \hline
\texttt{0x400018} & \texttt{13, 1B, 2E, 00} & \texttt{0x002E1B13} & \texttt{slli s6, t3, 2} (dịch trái $2\text{ bit}$, nhân $4$) \\ \hline
\end{tabularx}
\end{center}
Hãy trả lời các câu hỏi trong Đề 0012:
\begin{enumerate}[label=\textbf{Câu \arabic*.}, start=100]
    \item Hãy cho biết giá trị số 32-bit ở địa chỉ \texttt{0x400008} ở hệ 16.
    \item Sau đoạn lệnh: \texttt{lui s0, 0x400} $\rightarrow$ \texttt{addi s0, s0, 0x16} $\rightarrow$ \texttt{lb t1, 0(s0)}, hai thanh ghi \texttt{s0} và \texttt{t1} bằng bao nhiêu?
    \item Lệnh ở địa chỉ \texttt{0x400000} (\texttt{0x40730E33}) có thanh ghi nguồn \texttt{rs2} là gì?
    \item Lệnh cuối cùng trong bản đồ (tại \texttt{0x400018}: \texttt{0x002E1B13}) tương đương với phép nhân với hằng số bằng mấy?
    \item Lệnh dạng U ở địa chỉ \texttt{0x400004} (\texttt{0x00040B37}) có toán hạng hằng số bằng mấy (viết ở hệ 16)?
    \item Hãy cho biết thao tác của lệnh ở địa chỉ \texttt{0x400004} là gì?
    \item Giải thích nào về lệnh \texttt{bne t1, t2, -16} là đúng?\\
    A. Nếu $\texttt{t1} \neq \texttt{t2}$ rẽ nhánh về 4 lệnh sau \qquad B. Nếu $\texttt{t1} \neq \texttt{t2}$ rẽ nhánh về 4 lệnh trước\\
    C. Nếu $\texttt{t1} == \texttt{t2}$ rẽ nhánh về 4 lệnh sau \qquad D. Nếu $\texttt{t1} == \texttt{t2}$ rẽ nhánh về 4 lệnh trước
\end{enumerate}
\textbf{Hướng dẫn giải:}
\begin{itemize}
    \item \textbf{Câu 100:} Tại \texttt{0x400008} có 4 byte \texttt{0x13, 0x0B, 0x8B, 0x03}. Ghép ngược theo Little-Endian được: \textbf{\texttt{0x038B0B13}}.
    \item \textbf{Câu 101:} $\mathbf{\texttt{s0 = 0x00400016}}$. Tại địa chỉ \texttt{0x400016} của Đề 0012 chứa byte \texttt{0x1F} (số dương), nên sau lệnh \texttt{lb t1, 0(s0)} ta có: $\mathbf{\texttt{t1 = 0x0000001F}}$.
    \item \textbf{Câu 102:} Trong \texttt{0x40730E33}, 12 bit đầu là $\texttt{407}_{16} = \underbrace{\texttt{0100000}}_{\texttt{funct7}}\,\underbrace{\texttt{00111}}_{\texttt{rs2}=7}$. Thanh ghi $\texttt{x7}$ là \textbf{\texttt{t2}}.
    \item \textbf{Câu 103:} Lệnh ở \texttt{0x400018} (\texttt{0x002E1B13}) là lệnh \texttt{slli s6, t3, 2}. Dịch trái $2\text{ bit}$ tương đương nhân với $2^2 = \mathbf{4}$.
    \item \textbf{Câu 104:} Lệnh dạng U ở \texttt{0x400004} (\texttt{0x00040B37}) có 20 bit cao (5 chữ số Hexa đầu) là \textbf{\texttt{0x40}} (hoặc \texttt{0x00040}).
    \item \textbf{Câu 105:} Opcode của \texttt{0x00040B37} là $\texttt{0110111}_2 = 55_{10} \implies$ Thao tác là \textbf{\texttt{lui}}.
    \item \textbf{Câu 109:} Lệnh \texttt{bne} kiểm tra khác nhau ($\texttt{t1} \neq \texttt{t2}$), và độ lệch $-16\text{ byte} = -4 \times 4\text{ byte}$ tức là lùi về \textbf{4 lệnh trước lệnh hiện thời} $\rightarrow$ Chọn \textbf{B}.
\end{itemize}

\chapter{SỐ HỌC MÁY TÍNH}

\section{Phần 1: Lý thuyết trọng tâm}

\subsection{Hệ đếm và Chuyển đổi giữa các hệ đếm}

\textbf{1. Ba hệ đếm cơ bản trong máy tính}
\begin{itemize}
    \item \textbf{Hệ thập phân (Decimal -- Cơ số 10):} Gồm $10$ chữ số từ $0$ đến $9$. Dùng $n$ chữ số thập phân biểu diễn được $10^n$ giá trị khác nhau (từ $0$ đến $10^n - 1$).
    \item \textbf{Hệ nhị phân (Binary -- Cơ số 2):} Gồm $2$ chữ số là $0$ và $1$ (mỗi chữ số nhị phân gọi là $1\text{ bit}$ -- đơn vị thông tin nhỏ nhất). Một nhóm $8\text{ bit}$ tạo thành $1\text{ byte}$ ($2^8 = 256$ giá trị). Dùng $n\text{ bit}$ có thể biểu diễn được \textbf{$2^n\text{ giá trị}$} khác nhau.
    \item \textbf{Hệ mười sáu (Hexadecimal -- Cơ số 16):} Gồm $16$ ký tự: \texttt{0--9} và \texttt{A} ($10$), \texttt{B} ($11$), \texttt{C} ($12$), \texttt{D} ($13$), \texttt{E} ($14$), \texttt{F} ($15$). Hệ Hexa dùng để viết gọn cho số nhị phân: cứ \textbf{1 nhóm 4-bit nhị phân} (tính từ dấu phẩy sang hai bên) được thay bằng đúng \textbf{1 chữ số Hexa}.
\end{itemize}

\begin{tcolorbox}[bwbox, title={Bảng quy đổi chuẩn giữa 4-bit Nhị phân, Hệ Hexa và Hệ Thập phân}]
\begin{center}
\renewcommand{\arraystretch}{1.2}
\small
\begin{tabularx}{\linewidth}{|Z|Z|Z||Z|Z|Z||Z|Z|Z||Z|Z|Z|}
\hline
\textbf{Hệ 2} & \textbf{Hex} & \textbf{Hệ 10} & \textbf{Hệ 2} & \textbf{Hex} & \textbf{Hệ 10} & \textbf{Hệ 2} & \textbf{Hex} & \textbf{Hệ 10} & \textbf{Hệ 2} & \textbf{Hex} & \textbf{Hệ 10} \\ \hline
\texttt{0000} & \texttt{0} & $0$ & \texttt{0100} & \texttt{4} & $4$ & \texttt{1000} & \texttt{8} & $8$ & \texttt{1100} & \texttt{C} & $12$ \\ \hline
\texttt{0001} & \texttt{1} & $1$ & \texttt{0101} & \texttt{5} & $5$ & \texttt{1001} & \texttt{9} & $9$ & \texttt{1101} & \texttt{D} & $13$ \\ \hline
\texttt{0010} & \texttt{2} & $2$ & \texttt{0110} & \texttt{6} & $6$ & \texttt{1010} & \texttt{A} & $10$ & \texttt{1110} & \texttt{E} & $14$ \\ \hline
\texttt{0011} & \texttt{3} & $3$ & \texttt{0111} & \texttt{7} & $7$ & \texttt{1011} & \texttt{B} & $11$ & \texttt{1111} & \texttt{F} & $15$ \\ \hline
\end{tabularx}
\end{center}
\end{tcolorbox}

\textbf{2. Công thức giá trị và Quy tắc chuyển đổi giữa Hệ 10 và Hệ 2}
\begin{itemize}
    \item \textbf{Dạng tổng quát của số nhị phân có cả phần nguyên và phần lẻ:}
    $$A = a_n a_{n-1} \dots a_1 a_0 . a_{-1} a_{-2} \dots a_{-m (2)} = \sum_{i=-m}^{n} a_i \times 2^i$$
    Trong đó các bit phần nguyên có trọng số là $2^0 = 1, 2^1 = 2, 2^2 = 4, 2^3 = 8, 2^4 = 16, 2^5 = 32, 2^6 = 64, 2^7 = 128\dots$; các bit phần lẻ sau dấu phẩy có trọng số là:
    $$2^{-1} = 0.5 \quad ; \quad 2^{-2} = 0.25 \quad ; \quad 2^{-3} = 0.125 \quad ; \quad 2^{-4} = 0.0625 \quad ; \quad 2^{-5} = 0.03125$$
    \item \textbf{Chuyển đổi số thập phân sang nhị phân:}
    \begin{itemize}
        \item \textit{Đối với phần nguyên:} Chia liên tiếp cho $2$ rồi viết các số dư theo thứ tự ngược từ dưới lên (hoặc phân tích thành tổng các lũy thừa $2^i$ giảm dần: $128, 64, 32, 16, 8, 4, 2, 1$).
        \item \textit{Đối với phần lẻ thập phân:} Nhân liên tiếp phần lẻ với $2$, lấy phần nguyên ($0$ hoặc $1$) sau mỗi lần nhân viết theo thứ tự từ trên xuống dưới.
        \item \textit{Lưu ý về số thực không thể biểu diễn chính xác tuyệt đối trong hệ nhị phân:} Một số lẻ thập phân chỉ biểu diễn chính xác hữu hạn trong hệ nhị phân nếu phần lẻ của nó viết được dưới dạng phân số tối giản có mẫu số là lũy thừa của $2$ ($\dfrac{k}{2^m}$, như $0.5 = \dfrac{1}{2}, 0.25 = \dfrac{1}{4}, 0.75 = \dfrac{3}{4}, 0.125 = \dfrac{1}{8}, 0.375 = \dfrac{3}{8}, 0.625 = \dfrac{5}{8}$). Các số như $0.1, 0.2, 0.3, 0.4, 0.6, 0.7, 0.8, 0.9$ khi nhân đôi liên tiếp sẽ lặp vô hạn tuần hoàn nên \textbf{không thể biểu diễn chính xác tuyệt đối} bằng số dấu phẩy động nhị phân.
    \end{itemize}
\end{itemize}

\subsection{Đại số Boolean, Cổng logic và Mạch số cơ bản}

\textbf{1. Các phép toán Đại số Boolean và Cổng logic cơ bản}
\begin{itemize}
    \item \textbf{Ba phép toán cơ bản:} NOT ($\overline{A}$ hoặc $A'$), AND ($A \cdot B$ hoặc $AB$), OR ($A + B$). Thứ tự ưu tiên: $\text{NOT} > \text{AND} > \text{OR}$.
    \item \textbf{Ba phép toán mở rộng:}
    \begin{itemize}
        \item NAND (NOT-AND): $F = \overline{A \cdot B}$ (Chỉ bằng $0$ khi cả $A = 1$ và $B = 1$).
        \item NOR (NOT-OR): $F = \overline{A + B}$ (Chỉ bằng $1$ khi cả $A = 0$ và $B = 0$).
        \item XOR (Exclusive-OR): $F = A \oplus B = \overline{A}B + A\overline{B}$ (Bằng $1$ khi hai đầu vào khác nhau; bằng $0$ khi hai đầu vào giống nhau).
    \end{itemize}
    \item \textbf{Định lý De Morgan:} $\overline{A \cdot B} = \overline{A} + \overline{B}$ \quad và \quad $\overline{A + B} = \overline{A} \cdot \overline{B}$.
\end{itemize}

\begin{tcolorbox}[bwbox, title={Ký hiệu hình vẽ chuẩn của 6 cổng logic cơ bản}]
\begin{center}
\begin{tikzpicture}[
    gate/.style={draw=black, thick},
    every node/.style={font=\small}
]
    % Hàng 1: AND, OR, NOT
    \node[and gate US, draw, thick, logic gate inputs=nn, minimum height=0.8cm] (andg) at (0, 1.5) {};
    \draw (andg.input 1) -- ++(-0.5,0) node[left] {$A$};
    \draw (andg.input 2) -- ++(-0.5,0) node[left] {$B$};
    \draw (andg.output) -- ++(0.5,0) node[right] {$F = AB$};
    \node[below=4pt of andg, font=\bfseries\footnotesize] {Cổng AND};

    \node[or gate US, draw, thick, logic gate inputs=nn, minimum height=0.8cm] (org) at (5.2, 1.5) {};
    \draw (org.input 1) -- ++(-0.5,0) node[left] {$A$};
    \draw (org.input 2) -- ++(-0.5,0) node[left] {$B$};
    \draw (org.output) -- ++(0.5,0) node[right] {$F = A + B$};
    \node[below=4pt of org, font=\bfseries\footnotesize] {Cổng OR};

    \node[not gate US, draw, thick, minimum height=0.6cm] (notg) at (10.5, 1.5) {};
    \draw (notg.input) -- ++(-0.5,0) node[left] {$A$};
    \draw (notg.output) -- ++(0.5,0) node[right] {$F = \overline{A}$};
    \node[below=6pt of notg, font=\bfseries\footnotesize] {Cổng NOT};

    % Hàng 2: NAND, NOR, XOR
    \node[nand gate US, draw, thick, logic gate inputs=nn, minimum height=0.8cm] (nandg) at (0, -0.5) {};
    \draw (nandg.input 1) -- ++(-0.5,0) node[left] {$A$};
    \draw (nandg.input 2) -- ++(-0.5,0) node[left] {$B$};
    \draw (nandg.output) -- ++(0.5,0) node[right] {$F = \overline{AB}$};
    \node[below=4pt of nandg, font=\bfseries\footnotesize] {Cổng NAND};

    \node[nor gate US, draw, thick, logic gate inputs=nn, minimum height=0.8cm] (norg) at (5.2, -0.5) {};
    \draw (norg.input 1) -- ++(-0.5,0) node[left] {$A$};
    \draw (norg.input 2) -- ++(-0.5,0) node[left] {$B$};
    \draw (norg.output) -- ++(0.5,0) node[right] {$F = \overline{A+B}$};
    \node[below=4pt of norg, font=\bfseries\footnotesize] {Cổng NOR};

    \node[xor gate US, draw, thick, logic gate inputs=nn, minimum height=0.8cm] (xorg) at (10.5, -0.5) {};
    \draw (xorg.input 1) -- ++(-0.5,0) node[left] {$A$};
    \draw (xorg.input 2) -- ++(-0.5,0) node[left] {$B$};
    \draw (xorg.output) -- ++(0.5,0) node[right] {$F = A \oplus B$};
    \node[below=4pt of xorg, font=\bfseries\footnotesize] {Cổng XOR};
\end{tikzpicture}
\end{center}
\end{tcolorbox}

\textbf{2. Các mạch tổ hợp và mạch dãy cơ bản trong máy tính}
\begin{itemize}
    \item \textbf{Mạch tổ hợp (Combinational Circuits -- mạch không nhớ):} Giá trị đầu ra chỉ phụ thuộc vào tổ hợp giá trị đầu vào tại thời điểm hiện tại. Các bộ phận điển hình:
    \begin{itemize}
        \item \textit{Bộ ghép kênh (MUX -- Multiplexer):} Có $2^n$ đầu vào dữ liệu, $n$ đầu vào chọn ($S$) và $1$ đầu ra. Với MUX 2-ra-1: $F = D_0 \cdot \overline{S} + D_1 \cdot S$ (khi $S = 0$ thì nối $D_0$ ra $F$; khi $S = 1$ thì nối $D_1$ ra $F$).
        \item \textit{Bộ giải mã (Decoder):} Có $n$ đầu vào và $2^n$ đầu ra. Với mỗi tổ hợp của $n$ đầu vào, chỉ có duy nhất $1$ đầu ra tích cực.
        \item \textit{Bộ cộng bán phần 1-bit (Half Adder):} Cộng 2 bit $A, B$ tạo ra bit tổng $S = A \oplus B$ (dùng cổng \textbf{XOR}) và bit nhớ ra $C_{out} = A \cdot B$ (dùng cổng \textbf{AND}).
        \item \textit{Bộ cộng toàn phần 1-bit (Full Adder):} Cộng 3 bit ($A, B$ và bit nhớ vào $C_{in}$, hoặc ký hiệu là $C$) tạo ra bit tổng $S$ và bit nhớ ra $C_{out}$:
        $$S = A \oplus B \oplus C = \overline{A}\,\overline{B}C + \overline{A}B\overline{C} + A\overline{B}\,\overline{C} + ABC \quad ; \quad C_{out} = AB + AC + BC$$
        (\textit{Quy tắc nhớ nhanh Full Adder:} Đếm số lượng bit $1$ trong 3 đầu vào $(A, B, C)$: nếu có \textbf{lẻ} bit $1$ tức $1$ hoặc $3$ bit $1$ thì $S = 1$, ngược lại $S = 0$; nếu có \textbf{từ 2 bit $1$ trở lên} thì $C_{out} = 1$, ngược lại $C_{out} = 0$).
    \end{itemize}
    \item \textbf{Mạch dãy (Sequential Circuits -- mạch có nhớ):} Đầu ra phụ thuộc cả đầu vào hiện tại và trạng thái quá khứ. Được xây dựng từ các phần tử nhớ \textbf{Flip-Flop} (S-R, J-K, D) kết hợp xung nhịp (Clock) để tạo thành các \textbf{Thanh ghi song song} và \textbf{Thanh ghi dịch}.
\end{itemize}

\subsection{Biểu diễn Số nguyên không dấu và Số nguyên có dấu (Mã bù 2)}

\textbf{1. Biểu diễn số nguyên không dấu (Unsigned Integer) với $n$ bit}
\begin{itemize}
    \item Tất cả $n$ bit ($a_{n-1}a_{n-2}\dots a_1a_0$) đều dùng để biểu diễn độ lớn: $A = \sum_{i=0}^{n-1} a_i 2^i$.
    \item \textbf{Dải biểu diễn của số nguyên không dấu $n$ bit:} Từ \textbf{$0$ đến $2^n - 1$}.
    \begin{itemize}
        \item Ví dụ: Với $8\text{ bit}$: $[0, 255]$; với $9\text{ bit}$: $[0, 511]$; với $10\text{ bit}$: $[0, 1023]$; với $16\text{ bit}$: $[0, 65\,535]$.
        \item Bài toán ngược: Để đếm hoặc mã hóa $M$ đối tượng không âm (ví dụ đếm $M$ sinh viên hoặc $M$ chỗ ngồi), số bit tối thiểu $n$ cần dùng phải thỏa mãn $2^n \ge M$ (hoặc $2^n - 1 \ge M$ nếu tính cả giá trị $0$).
    \end{itemize}
\end{itemize}

\textbf{2. Biểu diễn số nguyên có dấu (Signed Integer) theo Mã bù hai (Two's Complement)}\\
Trong máy tính, số nguyên có dấu $n\text{ bit}$ ($A = a_{n-1}a_{n-2}\dots a_1a_0$) được biểu diễn theo \textbf{Mã bù hai}. Bit ngoài cùng bên trái ($a_{n-1}$) là \textbf{bit dấu (Sign bit)}:
\begin{itemize}
    \item Nếu $a_{n-1} = 0$: $A$ là số không âm ($\ge 0$). Giá trị tính như số không dấu: $A = \sum_{i=0}^{n-2} a_i 2^i$.
    \item Nếu $a_{n-1} = 1$: $A$ là số nguyên âm ($< 0$).
    \item \textbf{Công thức tổng quát xác định giá trị của số bù hai $n$ bit:}
    $$A = -a_{n-1}2^{n-1} + \sum_{i=0}^{n-2} a_i 2^i$$
    \item \textbf{Dải biểu diễn của số nguyên có dấu $n$ bit:} Từ \textbf{$-2^{n-1}$ đến $+(2^{n-1} - 1)$}.
    \item \textbf{Quy tắc tìm Mã bù một và Mã bù hai (Phép đảo dấu số nguyên):}
    \begin{enumerate}
        \item \textbf{Bước 1 (Tìm số bù một):} Đảo giá trị tất cả các bit của số ban đầu ($0 \rightarrow 1$ và $1 \rightarrow 0$).
        \item \textbf{Bước 2 (Tìm số bù hai):} Lấy số bù một \textbf{cộng thêm 1} ở bit cuối cùng bên phải:
        $$\text{Số bù hai của } A = (\text{Số bù một của } A) + 1 = 2^n - A = -A$$
    \end{enumerate}
    \item \textbf{Quy tắc mở rộng bit cho số nguyên (Sign-extension vs. Zero-extension):}
    \begin{itemize}
        \item \textit{Mở rộng không dấu (Zero-extended):} Luôn thêm các bit \texttt{0} vào bên trái.
        \item \textit{Mở rộng có dấu (Sign-extended):} Giữ nguyên giá trị của bit dấu đẩy sang toàn bộ các bit mới ở bên trái (số dương thêm toàn bit \texttt{0}, tức chữ Hexa \texttt{0}; số âm thêm toàn bit \texttt{1}, tức chữ Hexa \texttt{F}).
    \end{itemize}
\end{itemize}

\begin{tcolorbox}[bwbox, title={Bảng dải biểu diễn của Số nguyên có dấu và Mẹo tính nhẩm nhanh}]
\begin{center}
\renewcommand{\arraystretch}{1.2}
\small
\begin{tabularx}{\linewidth}{|L{2.6cm}|C{3.6cm}|C{3.2cm}|Y|}
\hline
\textbf{Số lượng bit ($n$)} & \textbf{Dải giá trị thập phân} & \textbf{Số âm nhỏ nhất} & \textbf{Số dương lớn nhất} \\ \hline
$6\text{ bit}$ & $[-32, +31]$ & $100000_{(2)}$ ($-32$) & $011111_{(2)}$ ($+31$) \\ \hline
$8\text{ bit}$ (\texttt{char}) & $[-128, +127]$ & \texttt{0x80} ($-128$) & \texttt{0x7F} ($+127$) \\ \hline
$9\text{ bit}$ & $[-256, +255]$ & $1\,0000\,0000_{(2)}$ ($-256$) & $0\,1111\,1111_{(2)}$ ($+255$) \\ \hline
$12\text{ bit}$ (\texttt{imm}) & $[-2048, +2047]$ & \texttt{0x800} ($-2048$) & \texttt{0x7FF} ($+2047$) \\ \hline
$16\text{ bit}$ (\texttt{short}) & $[-32\,768, +32\,767]$ & \texttt{0x8000} ($-32\,768$) & \texttt{0x7FFF} ($+32\,767$) \\ \hline
$32\text{ bit}$ (\texttt{int}) & $[-2^{31}, +2^{31}-1]$ & \texttt{0x80000000} & \texttt{0x7FFFFFFF} \\ \hline
\end{tabularx}
\end{center}
\textbf{Ba quy tắc tính nhẩm nhanh số bù hai trong phòng thi:}
\begin{enumerate}
    \item \textbf{Với số 8-bit có dấu ($a_7a_6a_5a_4a_3a_2a_1a_0$):} Bit đầu $a_7$ mang giá trị $\mathbf{-128}$, các bit còn lại mang giá trị dương $+64, +32, +16, +8, +4, +2, +1$. (Ví dụ: $\texttt{0x8E} = 1000\,1110_2 = -128 + 8 + 4 + 2 = -114$).
    \item \textbf{Với các số Hexa có dạng toàn chữ \texttt{F} ở đầu (như \texttt{0xFFFFFFF3} hoặc \texttt{0xFFA}):} Vì \texttt{0xFFFFFFFF} $= -1$, nên nếu tất cả các chữ số phía trước đều là \texttt{F}, chỉ cần lấy \textbf{chữ số Hexa cuối cùng trừ đi 16} (Ví dụ: \texttt{0xFFFFFFF3} $\implies 3 - 16 = -13$; \texttt{0xFFA} $\implies 10 - 16 = -6$).
    \item \textbf{So sánh nhanh các số nguyên có dấu viết ở dạng Hexa:}
    \begin{itemize}
        \item Chữ số Hexa đầu tiên từ \texttt{0} đến \texttt{7}: Bit dấu bằng $0 \implies$ \textbf{Số dương} (hoặc $0$).
        \item Chữ số Hexa đầu tiên từ \texttt{8} đến \texttt{F}: Bit dấu bằng $1 \implies$ \textbf{Số âm} (luôn nhỏ hơn mọi số dương).
    \end{itemize}
\end{enumerate}
\end{tcolorbox}

\subsection{Cộng / Trừ / Nhân / Chia số nguyên và Hiện tượng Tràn số học (Overflow)}

\textbf{1. Cộng số nguyên không dấu và Tràn nhớ (Carry Out)}
\begin{itemize}
    \item Khi cộng hai số nguyên \textbf{không dấu} $n\text{ bit}$, nếu tổng vượt quá $2^n - 1$ thì phát sinh bit nhớ ra ngoài cùng $\mathbf{C_{out} = 1}$, làm kết quả $n\text{ bit}$ bị sai (gọi là tràn nhớ). Nếu $C_{out} = 0$ thì kết quả đúng.
\end{itemize}

\textbf{2. Cộng/Trừ số nguyên có dấu và Hiện tượng Tràn số học (Overflow)}
\begin{itemize}
    \item Phép trừ hai số nguyên được phần cứng đưa về phép cộng thông qua lấy bù hai: $X - Y = X + (-Y) = X + \overline{Y} + 1$.
    \item Khi cộng hai số nguyên \textbf{có dấu} $n\text{ bit}$, kết quả đúng hay sai \textbf{không phụ thuộc vào bit nhớ $C_{out}$}:
    \begin{itemize}
        \item Khi cộng hai số nguyên \textbf{trái dấu nhau}: Kết quả \textbf{luôn luôn đúng} (không bao giờ xảy ra tràn số học).
        \item Khi cộng hai số nguyên \textbf{cùng dấu nhau}: Nếu bit dấu của kết quả giống với bit dấu của hai số hạng ban đầu thì kết quả đúng. Nếu bit dấu của kết quả \textbf{ngược lại} với bit dấu của hai số hạng ban đầu (Cộng hai số dương ra số âm, hoặc Cộng hai số âm ra số dương) thì xảy ra \textbf{Tràn số học (Overflow)} và kết quả bị sai!
    \end{itemize}
\end{itemize}

\begin{tcolorbox}[bwbox, title={Hai dạng bài toán Tràn số học (Overflow) trọng tâm trong đề thi}]
\textbf{Dạng 1: Tính giá trị hiển thị khi cộng hai biến \texttt{char} (8-bit có dấu) bị tràn}\\
Biến kiểu \texttt{char} (hoặc \texttt{int8\_t}) có dải giá trị $[-128, +127]$. Tính tổng toán học thông thường $S = x + y$:
\begin{itemize}
    \item Nếu $S \in [-128, +127]$: Không tràn, kết quả chính là $S$.
    \item Nếu $S > +127$ (tràn dương): Lấy $S - 256$ (Ví dụ: $85 + 52 = 137 > 127 \implies 137 - 256 = -119$).
    \item Nếu $S < -128$ (tràn âm): Lấy $S + 256$ (Ví dụ: $(-91) + (-67) = -158 < -128 \implies -158 + 256 = +98$).
\end{itemize}

\textbf{Dạng 2: Tìm phạm vi của thanh ghi thứ hai gây tràn khi thực hiện lệnh \texttt{add rd, rs1, rs2}}\\
Giả sử thanh ghi thứ nhất đang giữ một số nguyên dương $K > 0$ (ví dụ $K = 512 = \texttt{0x200}$ hoặc $K = 1024 = \texttt{0x400}$). Vì $K > 0$, phép cộng 32-bit \texttt{add} chỉ bị tràn khi thanh ghi thứ hai cũng là một số dương lớn sao cho tổng vượt quá số dương lớn nhất 32-bit là $\text{MAX} = \texttt{0x7FFFFFFF}$:
$$K + \texttt{rs2} > \texttt{0x7FFFFFFF} \iff \texttt{rs2} \ge \texttt{0x80000000} - K$$
Vậy phạm vi của thanh ghi thứ hai gây tràn luôn là: \textbf{Từ $(\texttt{0x80000000} - K)$ đến $\texttt{0x7FFFFFFF}$}.\\
(\textit{Lưu ý:} Điểm kết thúc luôn là \texttt{0x7FFFFFFF}, không bao giờ chọn đáp án kết thúc bởi \texttt{0x80000000} vì \texttt{0x80000000} là số âm $-2^{31}$, cộng với số dương $K$ là hai số trái dấu nên không bao giờ tràn).
\end{tcolorbox}

\textbf{3. Nhân/Chia số nguyên và Quy tắc dấu của phép chia trong RISC-V (RV32M)}
\begin{itemize}
    \item \textbf{Nhân hai số 32-bit tạo tích 64-bit:} Lệnh \texttt{mul rd, rs1, rs2} lấy 32 bit thấp của tích; các lệnh \texttt{mulh} (2 số có dấu), \texttt{mulhu} (2 số không dấu), \texttt{mulhsu} (\texttt{rs1} có dấu, \texttt{rs2} không dấu) lấy 32 bit cao của tích.
    \item \textbf{Chia hai số nguyên 32-bit có dấu (\texttt{div} và \texttt{rem}):}
    $$\text{Số bị chia} = (\text{Số chia} \times \text{Thương}) + \text{Số dư}$$
    \begin{itemize}
        \item \textbf{Dấu của Thương (\texttt{div rd, rs1, rs2}):} Cùng dấu thì thương dương ($+$); khác dấu thì thương âm ($-$).
        \item \textbf{Dấu của Số dư (\texttt{rem rd, rs1, rs2}):} Số dư luôn luôn \textbf{cùng dấu với SỐ BỊ CHIA (\texttt{rs1})}, bất kể số chia âm hay dương.
    \end{itemize}
\end{itemize}

\subsection{Số thực dấu phẩy động chuẩn IEEE 754}

\textbf{1. Cấu trúc chuẩn IEEE 754 độ chính xác đơn 32-bit (Single Precision -- Binary32)}\\
Một số thực dấu phẩy động 32-bit gồm 3 trường từ trái sang phải (\textbf{1 bit -- 8 bit -- 23 bit}):
\begin{center}
\renewcommand{\arraystretch}{1.25}
\begin{tabularx}{0.85\linewidth}{|C{2.8cm}|C{4.2cm}|Z|}
\hline
\textbf{Bit 31 (1 bit)} & \textbf{Bit 30:23 (8 bit)} & \textbf{Bit 22:0 (23 bit)} \\ \hline
Bit dấu $S$ (Sign) & Số mũ lệch $e$ (Biased Exponent) & Phần lẻ định trị $m$ (Fraction) \\ \hline
\end{tabularx}
\end{center}
\textbf{Công thức xác định giá trị của số thực $X$ ở trạng thái chuẩn hóa ($0 < e < 255$):}
$$X = (-1)^S \times (1.m)_{(2)} \times 2^{e - 127}$$
Trong đó:
\begin{itemize}
    \item \textbf{Bit dấu $S$:} Nếu $S = 0$ thì $X > 0$ (số dương); nếu $S = 1$ thì $X < 0$ (số âm).
    \item \textbf{Số mũ thực $E$:} $E = e - 127 \iff e = E + 127$ (số $127$ là hằng số độ lệch Bias của dạng 32-bit).
    \item \textbf{Phần định trị $M = (1.m)_{(2)}$:} Bit $1.$ trước dấu phẩy được ẩn đi (implied bit) để tiết kiệm 1 bit lưu trữ.
    \item (\textit{Với dạng 64-bit Double Precision:} Gồm $1\text{ bit}$ dấu $S$, $11\text{ bit}$ mũ $e = E + 1023$, và $52\text{ bit}$ lẻ định trị $m$).
\end{itemize}

\textbf{2. Các giá trị đặc biệt trong chuẩn IEEE 754 32-bit}
\begin{itemize}
    \item \textbf{Số $\pm 0$:} Khi tất cả 8 bit mũ $e = 0$ (\texttt{00000000}) và tất cả 23 bit $m = 0$ (\texttt{0x00000000} là $+0$; \texttt{0x80000000} là $-0$).
    \item \textbf{Số Vô cùng ($\pm \infty$ -- Infinity):} Khi tất cả 8 bit mũ $e = 255$ (\texttt{11111111}) và tất cả 23 bit $m = 0$:
    \begin{itemize}
        \item $+\infty$ có mã Hexa là \textbf{\texttt{0x7F800000}} (3 chữ số Hexa đầu là \texttt{7F8}).
        \item $-\infty$ có mã Hexa là \textbf{\texttt{0xFF800000}} (3 chữ số Hexa đầu là \texttt{FF8}).
    \end{itemize}
    \item \textbf{Không phải là số (NaN -- Not a Number):} Khi tất cả 8 bit mũ $e = 255$ (\texttt{11111111}) và phần lẻ định trị $\mathbf{m \neq 0}$ (có ít nhất 1 bit bằng $1$ trong 23 bit của $m$).
    \begin{itemize}
        \item \textit{Tổng số cách biểu diễn NaN trong chuẩn 32-bit:} Trường $e = \texttt{11111111}$ có $1$ cách; phần lẻ $m$ (23 bit) phải khác $0$ nên có $2^{23} - 1$ tổ hợp; bit dấu $S$ có $2$ cách chọn ($0$ hoặc $1$). Vậy tổng số biểu diễn NaN là:
        $$2 \times (2^{23} - 1) = \mathbf{2^{24} - 2}$$
    \end{itemize}
\end{itemize}

\newpage
\section{Phần 2: Toàn bộ bài tập Chương 4 và Hướng dẫn giải}

\subsection{Nhóm 1: Hệ đếm, Chuyển đổi cơ số và Xác định số bit biểu diễn}

\textbf{Bài tập 4.1 (Đề Giữa kỳ 2018.2 -- Mã đề 198 Câu 7 \& Mã đề 199 Câu 2)}\\
Số dương $1101011.1011$ ở hệ cơ số 2 bằng với số nào trong các số thập phân sau đây?
\begin{enumerate}[label=\Alph*.]
    \item $107.125_{(10)}$
    \item $107.129_{(10)}$
    \item $107.1875_{(10)}$
    \item $107.6875_{(10)}$
\end{enumerate}
\textbf{Hướng dẫn giải:}
\begin{itemize}
    \item \textbf{Bước 1 (Tính phần nguyên $1101011_{(2)}$):}
    $$1101011_{(2)} = 2^6 + 2^5 + 0 + 2^3 + 0 + 2^1 + 2^0 = 64 + 32 + 8 + 2 + 1 = 107$$
    \item \textbf{Bước 2 (Tính phần lẻ $.1011_{(2)}$):}
    $$0.1011_{(2)} = 1 \times 2^{-1} + 0 \times 2^{-2} + 1 \times 2^{-3} + 1 \times 2^{-4} = 0.5 + 0.125 + 0.0625 = 0.6875$$
    \item \textbf{Bước 3 (Ghép lại):} $1101011.1011_{(2)} = 107.6875_{(10)}$.
    \item \textbf{Kết luận:} Đáp án đúng là \textbf{D ($107.6875_{(10)}$)}.
\end{itemize}

\textbf{Bài tập 4.2 (Đề Giữa kỳ 2018.2 -- Mã đề 198 Câu 3, Mã đề 199 Câu 12, 18, Mã đề 204 Câu 3, 20 \& Đề Giữa kỳ 2019.2 -- Mã đề 243 Câu 4, Mã đề 244 Câu 11)}\\
Hãy trả lời chuỗi câu hỏi về dải biểu diễn và số lượng bit cần dùng cho số nguyên:
\begin{enumerate}[label=\alph*)]
    \item \textbf{(Mã đề 198 -- Câu 3):} Sử dụng $9\text{ bit}$ để biểu diễn một số nguyên \textbf{có dấu}, dải biểu diễn sẽ là?\\
    A. Từ $-255$ tới $+256$ \quad B. Từ $-256$ tới $+255$ \quad C. Từ $0$ tới $255$ \quad D. Từ $0$ tới $511$
    \item \textbf{(Mã đề 204 -- Câu 3):} Sử dụng $9\text{ bit}$ để biểu diễn một số nguyên \textbf{không dấu}, dải biểu diễn sẽ là?\\
    A. Từ $0$ tới $255$ \quad B. Từ $-256$ tới $+255$ \quad C. Từ $-255$ tới $+256$ \quad D. Từ $0$ tới $511$
    \item \textbf{(Mã đề 199 -- Câu 18):} Sử dụng $10\text{ bit}$ để biểu diễn một số nguyên \textbf{không dấu}, dải biểu diễn sẽ là?\\
    A. Từ $-511$ tới $+512$ \quad B. Từ $0$ tới $1023$ \quad C. Từ $-512$ tới $+511$ \quad D. Từ $0$ tới $511$
    \item \textbf{(Mã đề 199 -- Câu 12):} Biết rằng $511$ là số không dấu lớn nhất trong phạm vi biểu diễn. Vậy kiểu dữ liệu số nguyên này có bao nhiêu bit?\\
    A. $9\text{ bit}$ \quad B. $8\text{ bit}$ \quad C. $10\text{ bit}$ \quad D. $11\text{ bit}$
    \item \textbf{(Mã đề 204 -- Câu 20):} Sử dụng $6\text{ bit}$ nhị phân để biểu diễn số nguyên \textbf{có dấu}. Hỏi giá trị lớn nhất có thể biểu diễn là bao nhiêu?\\
    A. $15$ \quad B. $31$ \quad C. $16$ \quad D. $32$
    \item \textbf{(Mã đề 243 Câu 4 \& Mã đề 244 Câu 11):}
    \begin{itemize}
        \item Lớp học Kiến trúc máy tính có $150$ sinh viên. Biến số nguyên để đếm số sinh viên cần tối thiểu mấy bit?
        \item Hội trường C2 có $1008$ chỗ ngồi. Biến số nguyên để đếm số ghế cần tối thiểu mấy bit?
    \end{itemize}
\end{enumerate}
\textbf{Hướng dẫn giải:}
\begin{itemize}
    \item \textbf{Câu a:} Với $n = 9\text{ bit}$ có dấu, dải biểu diễn là $[-2^{9-1}, +2^{9-1}-1] = [-2^8, +2^8-1] = [\mathbf{-256, +255}]$ $\rightarrow$ Chọn \textbf{B}.
    \item \textbf{Câu b:} Với $n = 9\text{ bit}$ không dấu, dải biểu diễn là $[0, 2^9 - 1] = [\mathbf{0, 511}]$ $\rightarrow$ Chọn \textbf{D}.
    \item \textbf{Câu c:} Với $n = 10\text{ bit}$ không dấu, dải biểu diễn là $[0, 2^{10} - 1] = [\mathbf{0, 1023}]$ $\rightarrow$ Chọn \textbf{B}.
    \item \textbf{Câu d:} Ta có $2^n - 1 = 511 \iff 2^n = 512 = 2^9 \implies n = \mathbf{9\text{ bit}}$ $\rightarrow$ Chọn \textbf{A}.
    \item \textbf{Câu e:} Với $n = 6\text{ bit}$ có dấu, giá trị dương lớn nhất là $2^{6-1} - 1 = 2^5 - 1 = \mathbf{31}$ $\rightarrow$ Chọn \textbf{B}.
    \item \textbf{Câu f:}
    \begin{itemize}
        \item Để đếm $150$ sinh viên: Vì $2^7 - 1 = 127 < 150 \le 2^8 - 1 = 255$, cần tối thiểu \textbf{8 bit}.
        \item Để đếm $1008$ chỗ ngồi: Vì $2^9 - 1 = 511 < 1008 \le 2^{10} - 1 = 1023$, cần tối thiểu \textbf{10 bit}.
    \end{itemize}
\end{itemize}

\subsection{Nhóm 2: Nhận diện Cổng logic, Mạch tổ hợp và Bộ cộng 1-bit}

\textbf{Bài tập 4.3 (Ôn tập Cuối kỳ 2025.2 -- P1, Câu 16 \& Câu 74)}\\
Hãy cho biết mỗi hình vẽ ký hiệu cổng logic dưới đây biểu diễn cho phép toán logic nào?
\begin{center}
\begin{tikzpicture}[every node/.style={font=\small}]
    % Hình a: Cổng NAND (Câu 16)
    \node[nand gate US, draw, thick, logic gate inputs=nn, minimum height=1.0cm, minimum width=1.4cm] (g1) at (-3.5, 0) {};
    \draw[thick] (g1.input 1) -- ++(-0.7,0) node[left, font=\bfseries] {A};
    \draw[thick] (g1.input 2) -- ++(-0.7,0) node[left, font=\bfseries] {B};
    \draw[thick] (g1.output) -- ++(0.7,0);
    \node[below=8pt of g1, font=\bfseries] {Hình a) Câu 16};

    % Hình b: Cổng XOR (Câu 74)
    \node[xor gate US, draw, thick, logic gate inputs=nn, minimum height=1.0cm, minimum width=1.4cm] (g2) at (3.5, 0) {};
    \draw[thick] (g2.input 1) -- ++(-0.7,0) node[left, font=\bfseries] {A};
    \draw[thick] (g2.input 2) -- ++(-0.7,0) node[left, font=\bfseries] {B};
    \draw[thick] (g2.output) -- ++(0.7,0);
    \node[below=8pt of g2, font=\bfseries] {Hình b) Câu 74};
\end{tikzpicture}
\end{center}
\begin{enumerate}[label=\alph*)]
    \item \textbf{(Câu 16 -- Hình a):} A. NOR \qquad B. XOR \qquad C. AND \qquad D. NAND
    \item \textbf{(Câu 74 -- Hình b):} A. XOR \qquad B. AND \qquad C. NOR \qquad D. NAND
\end{enumerate}
\textbf{Hướng dẫn giải:}
\begin{itemize}
    \item \textbf{Giải câu a (Hình a):} Thân cổng có hình chữ D là cổng AND, ở đầu ra có thêm vòng tròn nhỏ đảo bit (NOT) tạo thành cổng \textbf{NAND} ($F = \overline{AB}$) $\rightarrow$ Chọn \textbf{D (NAND)}.
    \item \textbf{Giải câu b (Hình b):} Thân cổng có mũi nhọn giống cổng OR nhưng có thêm một đường cong tách rời ở phần đuôi là ký hiệu của cổng \textbf{XOR} ($F = A \oplus B$) $\rightarrow$ Chọn \textbf{A (XOR)}.
\end{itemize}

\textbf{Bài tập 4.4 (Slide Chương 4 -- Trang 30 \& Ôn tập Cuối kỳ 2025.2 -- P1, Câu 51)}\\
Mạch logic tổ hợp dưới đây có biểu thức đại số Boolean là gì? (Ký hiệu $A'$ là NOT $A$):
\begin{center}
\begin{tikzpicture}[scale=0.88, every node/.style={font=\small}]
    % 3 đường dọc A, B, C và 3 cổng NOT
    \foreach \x/\lbl in {0/A, 2.0/B, 4.0/C} {
        \node[font=\bfseries] at (\x, 3.8) {\lbl};
        \draw[thick] (\x, 3.5) -- (\x, -2.5);
        \fill (\x, 3.2) circle (2pt);
        \draw[thick] (\x, 3.2) -- (\x+0.7, 3.2) -- (\x+0.7, 2.9);
        \node[not gate US, draw, thick, rotate=-90, minimum width=0.45cm] (not\lbl) at (\x+0.7, 2.6) {};
        \draw[thick] (not\lbl.output) -- (\x+0.7, -2.5);
    }

    % 3 cổng AND
    \node[and gate US, draw, thick, logic gate inputs=nnn, minimum height=0.9cm] (and1) at (6.8, 1.4) {};
    \node[and gate US, draw, thick, logic gate inputs=nnn, minimum height=0.9cm] (and2) at (6.8, -0.2) {};
    \node[and gate US, draw, thick, logic gate inputs=nnn, minimum height=0.9cm] (and3) at (6.8, -1.8) {};

    % Kết nối cổng AND 1: A' B C'
    \draw[thick] (0.7, 1.65) -- (and1.input 1); \fill (0.7, 1.65) circle (2pt);
    \draw[thick] (2.0, 1.40) -- (and1.input 2); \fill (2.0, 1.40) circle (2pt);
    \draw[thick] (4.7, 1.15) -- (and1.input 3); \fill (4.7, 1.15) circle (2pt);

    % Kết nối cổng AND 2: A' B C
    \draw[thick] (0.7, 0.05) -- (and2.input 1); \fill (0.7, 0.05) circle (2pt);
    \draw[thick] (2.0, -0.20) -- (and2.input 2); \fill (2.0, -0.20) circle (2pt);
    \draw[thick] (4.0, -0.45) -- (and2.input 3); \fill (4.0, -0.45) circle (2pt);

    % Kết nối cổng AND 3: A B C'
    \draw[thick] (0.0, -1.55) -- (and3.input 1); \fill (0.0, -1.55) circle (2pt);
    \draw[thick] (2.0, -1.80) -- (and3.input 2); \fill (2.0, -1.80) circle (2pt);
    \draw[thick] (4.7, -2.05) -- (and3.input 3); \fill (4.7, -2.05) circle (2pt);

    % Cổng OR tổng hợp
    \node[or gate US, draw, thick, logic gate inputs=nnn, minimum height=1.1cm] (org) at (9.6, -0.2) {};
    \draw[thick] (and1.output) -- ++(0.6,0) |- (org.input 1);
    \draw[thick] (and2.output) -- (org.input 2);
    \draw[thick] (and3.output) -- ++(0.6,0) |- (org.input 3);
    \draw[thick] (org.output) -- ++(0.7,0) node[right, font=\bfseries] {$F$};
\end{tikzpicture}
\end{center}
\begin{enumerate}[label=\Alph*.]
    \item $F = (A + B + C')(A + B' + C)(A' + B + C)$
    \item $F = A'B'C' + A'BC + AB'C$
    \item $F = ABC' + A'BC + AB'C$
    \item $F = A'BC + A'BC' + ABC'$
\end{enumerate}
\textbf{Hướng dẫn giải:}
\begin{itemize}
    \item Mỗi đường tín hiệu $A, B, C$ có 2 nhánh dọc: nhánh bên trái đi thẳng là tín hiệu gốc ($A, B, C$), nhánh bên phải đi qua cổng đảo NOT là tín hiệu đảo ($A', B', C'$).
    \item Xác định biểu thức đầu ra của 3 cổng AND từ trên xuống dưới:
    \begin{itemize}
        \item Cổng AND thứ nhất nối với: nhánh đảo của $A$ ($A'$), nhánh thẳng của $B$ ($B$), nhánh đảo của $C$ ($C'$) $\implies \mathbf{A'BC'}$.
        \item Cổng AND thứ hai nối với: nhánh đảo của $A$ ($A'$), nhánh thẳng của $B$ ($B$), nhánh thẳng của $C$ ($C$) $\implies \mathbf{A'BC}$.
        \item Cổng AND thứ ba nối với: nhánh thẳng của $A$ ($A$), nhánh thẳng của $B$ ($B$), nhánh đảo của $C$ ($C'$) $\implies \mathbf{ABC'}$.
    \end{itemize}
    \item Cổng OR ở cuối cộng 3 tích trên lại: $F = A'BC' + A'BC + ABC'$ (hay $F = A'BC + A'BC' + ABC'$).
    \item \textbf{Kết luận:} Đáp án đúng là \textbf{D}.
\end{itemize}

\textbf{Bài tập 4.5 (Slide Chương 4 -- Trang 37, 38 \& Ôn tập Cuối kỳ 2025.2 -- P1, Câu 17 \& Câu 25)}
\begin{enumerate}[label=\alph*)]
    \item \textbf{(Câu 17):} Bản chất của bộ cộng bán phần 1-bit (Half Adder) là thực hiện hai phép toán logic nào?\\
    A. Thực hiện phép XOR và AND \qquad B. Thực hiện phép OR và AND\\
    C. Thực hiện phép NOR và NAND \qquad D. Thực hiện phép OR và NAND
    \item \textbf{(Câu 25):} Bộ cộng toàn phần 1-bit (Full Adder) có đầu vào là $(A, B, C)$ và đầu ra là $(S, C_{out})$. Với bộ giá trị đầu vào $(A, B, C) = (1, 1, 0)$, hãy xác định giá trị đầu ra tương ứng $(S, C_{out}) =$ ?\\
    A. $(0, 0)$ \qquad B. $(0, 1)$ \qquad C. $(1, 0)$ \qquad D. $(1, 1)$
\end{enumerate}
\textbf{Hướng dẫn giải:}
\begin{itemize}
    \item \textbf{Giải câu a:} Theo Slide Chương 4 (trang 37), bộ cộng bán phần 1-bit tính bit tổng $S = A \oplus B$ (phép \textbf{XOR}) và bit nhớ ra $C_{out} = A \cdot B$ (phép \textbf{AND}) $\rightarrow$ Chọn \textbf{A}.
    \item \textbf{Giải câu b:} Cộng 3 bit nhị phân: $A + B + C = 1 + 1 + 0 = 2_{(10)} = 10_{(2)}$, trong đó chữ số hàng đơn vị là bit tổng $S = 0$ và chữ số hàng cao là bit nhớ ra $C_{out} = 1$. Vậy $(S, C_{out}) = (0, 1)$ $\rightarrow$ Chọn \textbf{B}.
\end{itemize}

\subsection{Nhóm 3: Biểu diễn và So sánh Số nguyên có dấu (Mã bù 2)}

\textbf{Bài tập 4.6 (Đề Giữa kỳ 2018.2 -- Mã đề 198 Câu 6, Mã đề 199 Câu 10 \& Mã đề 204 Câu 8)}
\begin{enumerate}[label=\alph*)]
    \item \textbf{(Mã đề 198 Câu 6 \& Mã đề 199 Câu 10):} Số âm $-127_{(10)}$ bằng với số nào trong các số biểu diễn nhị phân 8-bit có dấu sau đây?\\
    A. \texttt{1000 0000} \qquad B. \texttt{1000 0001} \qquad C. \texttt{-0111 1111} \qquad D. \texttt{0111 1111}
    \item \textbf{(Mã đề 204 -- Câu 8):} Số nguyên $-77_{(10)}$ được biểu diễn dưới dạng số nhị phân 8-bit bù 2 là?\\
    A. \texttt{0100 1101} \qquad B. \texttt{1011 0011} \qquad C. \texttt{1011 1011} \qquad D. \texttt{1011 0010}
\end{enumerate}
\textbf{Hướng dẫn giải:}
\begin{itemize}
    \item \textbf{Giải câu a (Biểu diễn $-127$ bằng 8 bit):}
    \begin{itemize}
        \item Bước 1: Số dương $+127 = \texttt{0111 1111}_{(2)}$.
        \item Bước 2 (Đảo bit tìm bù 1): $\texttt{1000 0000}_{(2)}$.
        \item Bước 3 (Cộng 1 tìm bù 2): $\texttt{1000 0000} + 1 = \mathbf{\texttt{1000 0001}_{(2)}}$ (Kiểm tra: $-128 + 1 = -127$) $\rightarrow$ Chọn \textbf{B}.
    \end{itemize}
    \item \textbf{Giải câu b (Biểu diễn $-77$ bằng 8 bit):}
    \begin{itemize}
        \item Bước 1: Số dương $+77 = 64 + 8 + 4 + 1 = \texttt{0100 1101}_{(2)}$.
        \item Bước 2 (Đảo bit tìm bù 1): $\texttt{1011 0010}_{(2)}$.
        \item Bước 3 (Cộng 1 tìm bù 2): $\texttt{1011 0010} + 1 = \mathbf{\texttt{1011 0011}_{(2)}}$ $\rightarrow$ Chọn \textbf{B}.
    \end{itemize}
\end{itemize}

\textbf{Bài tập 4.7 (Đề Giữa kỳ 2019.2 -- Mã đề 243 Câu 9 \& Mã đề 244 Câu 15)}\\
Số nguyên 8-bit nào sau đây có dạng biểu diễn ở hệ 16 là \texttt{0xB2}?
\begin{enumerate}[label=\alph*)]
    \item \textbf{(Mã đề 243):} A. Là số \texttt{10110010} ở hệ 2 \quad B. Là số $177$ ở hệ 10, không dấu \quad C. Là số \texttt{00100010} ở hệ 2 \quad D. Là số $179$ ở hệ 10, không dấu
    \item \textbf{(Mã đề 244):} A. Là số $177$ ở hệ 10, không dấu \quad B. Là số \texttt{00100010} ở hệ 2 \quad C. Là số \texttt{10100000} ở hệ 2 \quad D. Là số $-78$ ở hệ 10, có dấu
\end{enumerate}
\textbf{Hướng dẫn giải:}
\begin{itemize}
    \item Đổi \texttt{0xB2} sang hệ nhị phân 8-bit: $\texttt{B} = 11 = \texttt{1011}_2$ và $\texttt{2} = \texttt{0010}_2 \implies \texttt{0xB2} = \mathbf{\texttt{1011\,0010}_{(2)}}$.
    \item Nếu xét là số \textbf{không dấu}: $128 + 32 + 16 + 2 = 178_{(10)}$.
    \item Nếu xét là số \textbf{có dấu} (bù 2): $-128 + 32 + 16 + 2 = \mathbf{-78_{(10)}}$.
    \item \textbf{Kết luận:} Câu a (Mã đề 243) chọn \textbf{A}; Câu b (Mã đề 244) chọn \textbf{D}.
\end{itemize}

\textbf{Bài tập 4.8 (Đề Cuối kỳ 2020.2 -- Trang 2 \& Đề Cuối kỳ 2021.1 -- Trang 2)}\\
Giả sử biến \texttt{i} là một biến số nguyên có dấu kích thước 8-bit được lưu trữ với giá trị Hexa là $\texttt{i} = \texttt{0x8E}$. Hãy cho biết giá trị tương ứng của biến số nguyên này trong hệ thập phân (cơ số 10)?
\begin{enumerate}[label=\Alph*.]
    \item $142$
    \item $-114$
    \item $-113$
    \item $-14$
\end{enumerate}
\textbf{Hướng dẫn giải:}
\begin{itemize}
    \item Đổi \texttt{0x8E} sang nhị phân 8-bit: $\texttt{0x8E} = \texttt{1000\,1110}_{(2)}$.
    \item Vì bit dấu đầu tiên bằng $1$ (trọng số $-128$) và các bit $1$ còn lại ở vị trí $8, 4, 2$:
    $$i = -128 + 8 + 4 + 2 = -128 + 14 = -114$$
    \item \textbf{Kết luận:} Đáp án đúng là \textbf{B ($-114$)}.
\end{itemize}

\textbf{Bài tập 4.9 (Đề Cuối kỳ 2022.2 -- \#f9c617, 2024.2 \& Ôn tập Cuối kỳ 2025.2 -- P1, Câu 69)}\\
Giả thiết \texttt{k} là biến số nguyên có dấu 16-bit, cho $k = -68$. Hãy cho biết nội dung \texttt{k} biểu diễn theo số Hexa?
\begin{enumerate}[label=\Alph*.]
    \item \texttt{0xFFBC}
    \item \texttt{0x00BC}
    \item \texttt{0x8044}
    \item \texttt{0xFF44}
\end{enumerate}
\textbf{Hướng dẫn giải:}
\begin{itemize}
    \item \textbf{Bước 1:} Biểu diễn số dương $+68$ ở dạng 16-bit: $68 = 64 + 4 = \texttt{0100\,0100}_{(2)} = \texttt{0x0044}$:
    $$+68 = \texttt{0000\,0000\,0100\,0100}_{(2)}$$
    \item \textbf{Bước 2 (Tìm bù 1 -- đảo tất cả các bit):} $\texttt{1111\,1111\,1011\,1011}_{(2)}$.
    \item \textbf{Bước 3 (Tìm bù 2 -- cộng thêm 1):}
    $$-68 = \texttt{1111\,1111\,1011\,1100}_{(2)} = \mathbf{\texttt{0xFFBC}}$$
    \item \textbf{Kết luận:} Đáp án đúng là \textbf{A (\texttt{0xFFBC})}.
\end{itemize}

\textbf{Bài tập 4.10 (Đề Cuối kỳ 2022.2 -- \#7151, 2024.1 -- \#0182da, 2024.2 \& Ôn tập 2025.2 -- Câu 45, Câu 73)}
\begin{enumerate}[label=\alph*)]
    \item \textbf{(Đề CK 2022.2, 2024.2 -- \#7151 \& Ôn tập 2025.2 Câu 73):} Giả thiết \texttt{i} là biến số nguyên có dấu 32-bit, cho nội dung $\texttt{i} = \texttt{0xFFFFFFF3}$. Giá trị của \texttt{i} theo số thập phân là bao nhiêu?\\
    A. $-4$ \qquad B. $-3$ \qquad C. $-13$ \qquad D. $-12$
    \item \textbf{(Đề CK 2024.1 -- \#0182da \& Ôn tập 2025.2 Câu 45):} Giả thiết \texttt{i} là biến số nguyên có dấu 32-bit, cho nội dung $\texttt{i} = \texttt{0xFFFFFFF4}$. Giá trị của \texttt{i} theo số thập phân là bao nhiêu?\\
    A. $-11$ \qquad B. $-12$ \qquad C. $-4$ \qquad D. $12$
\end{enumerate}
\textbf{Hướng dẫn giải:}
\begin{itemize}
    \item Áp dụng quy tắc tính nhanh cho các số âm có 7 chữ số đầu đều là \texttt{F}: lấy chữ số Hexa cuối cùng trừ đi $16$:
    \item \textbf{Câu a:} Với \texttt{0xFFFFFFF3}, chữ số cuối là $3 \implies i = 3 - 16 = \mathbf{-13}$ $\rightarrow$ Chọn \textbf{C}.
    \item \textbf{Câu b:} Với \texttt{0xFFFFFFF4}, chữ số cuối là $4 \implies i = 4 - 16 = \mathbf{-12}$ $\rightarrow$ Chọn \textbf{B}.
\end{itemize}

\textbf{Bài tập 4.11 (Đề Cuối kỳ 2023.1 -- \#581f03, 2023.2, Giữa kỳ 2024.1 -- \#30a33e \& Ôn tập 2025.2 -- Câu 48)}
\begin{enumerate}[label=\alph*)]
    \item \textbf{(Đề CK 2023.1 -- \#581f03):} Cho lệnh hợp ngữ RISC-V: \texttt{addi t0, s0, 0xFFA} (với toán hạng tức thời hằng số 12-bit được biểu diễn dưới dạng bù 2). Giá trị thập phân tương ứng của toán hạng tức thời này là bao nhiêu?\\
    A. $4010$ \qquad B. $-4010$ \qquad C. $6$ \qquad D. $-6$
    \item \textbf{(Đề GK 2024.1 -- \#30a33e \& Ôn tập 2025.2 Câu 48):} Giá trị thập phân của toán hạng immediate trong lệnh \texttt{addi t0, s0, 0x1FF} là bao nhiêu?\\
    A. $-511$ \qquad B. $2047$ \qquad C. $-2047$ \qquad D. $511$
\end{enumerate}
\textbf{Hướng dẫn giải:}
\begin{itemize}
    \item \textbf{Giải câu a:} Hằng số 12-bit \texttt{0xFFA} có 2 chữ đầu là \texttt{FF} (bit dấu số 11 bằng 1, số âm) và chữ số cuối là $\texttt{A} = 10$. Lấy $10 - 16 = \mathbf{-6}$ $\rightarrow$ Chọn \textbf{D ($-6$)}.
    \item \textbf{Giải câu b:} Hằng số 12-bit \texttt{0x1FF} có chữ số đầu là $\texttt{1} = \texttt{0001}_2$ (bit dấu số 11 bằng 0, số dương). Giá trị thập phân là: $1 \times 16^2 + 15 \times 16^1 + 15 = 256 + 240 + 15 = \mathbf{511}$ $\rightarrow$ Chọn \textbf{D ($511$)}.
\end{itemize}

\textbf{Bài tập 4.12 (Đề Giữa kỳ 2023.2 -- Đề 1 Câu 2 \& Đề Giữa kỳ 2022.2 -- Đề 0008 Câu 4, Đề 0006 Câu 2)}
\begin{enumerate}[label=\alph*)]
    \item \textbf{(Đề GK 2023.2):} Biểu diễn các số nguyên có dấu $+105$ và $-42$ theo mã bù hai 8-bit và 16-bit dưới dạng số Hexa.
    \item \textbf{(Đề GK 2022.2 -- Đề 0008 \& 0006):} Hoàn thành bảng chuyển đổi số nguyên 8-bit có dấu giữa Hệ 10, Hệ 2 bù 2 và Hệ 16 bù 2 cho các giá trị: $+77$; $-86$; $-19$; và $\texttt{0011 1101}_{(2)}$ (\texttt{3D}).
\end{enumerate}
\textbf{Hướng dẫn giải:}
\begin{itemize}
    \item \textbf{Giải câu a:}
    \begin{itemize}
        \item Với $+105$: $105 = 64 + 32 + 8 + 1 = \texttt{0110\,1001}_{(2)} = \texttt{0x69}$. Biểu diễn 8-bit là \textbf{\texttt{0x69}}; mở rộng dấu dương 16-bit là \textbf{\texttt{0x0069}}.
        \item Với $-42$: Ta có $+42 = 32 + 8 + 2 = \texttt{0010\,1010}_{(2)} \xrightarrow{\text{bù 1}} \texttt{1101\,0101}_{(2)} \xrightarrow{+1} \texttt{1101\,0110}_{(2)} = \texttt{0xD6}$. Biểu diễn 8-bit là \textbf{\texttt{0xD6}}; mở rộng dấu âm 16-bit (thêm \texttt{FF} vào bên trái) là \textbf{\texttt{0xFFD6}}.
    \end{itemize}
    \item \textbf{Giải câu b (Bảng kết quả chuyển đổi 8-bit có dấu):}
    \begin{center}
    \renewcommand{\arraystretch}{1.2}
    \small
    \begin{tabularx}{0.9\linewidth}{|C{2.2cm}|C{3.2cm}|C{2.2cm}|Y|}
    \hline
    \textbf{Hệ 10} & \textbf{Hệ 2 (Bù 2)} & \textbf{Hệ 16} & \textbf{Giải thích cách tính} \\ \hline
    $+77$ & \texttt{0100 1101} & \texttt{4D} & $64 + 8 + 4 + 1 = 77$ \\ \hline
    $-86$ & \texttt{1010 1010} & \texttt{AA} & $+86 = \texttt{0101 0110}_2 \xrightarrow{\text{bù 2}} \texttt{1010 1010}_2$ \\ \hline
    $-19$ & \texttt{1110 1101} & \texttt{ED} & $+19 = \texttt{0001 0011}_2 \xrightarrow{\text{bù 2}} \texttt{1110 1101}_2$ \\ \hline
    $+61$ & \texttt{0011 1101} & \texttt{3D} & $32 + 16 + 8 + 4 + 1 = +61$ \\ \hline
    \end{tabularx}
    \end{center}
\end{itemize}

\textbf{Bài tập 4.13 (Đề Giữa kỳ 2018.2 -- Mã đề 198 Câu 16, Mã đề 199 Câu 11 \& Mã đề 204 Câu 17)}
\begin{enumerate}[label=\alph*)]
    \item \textbf{(Mã đề 198 -- Câu 16):} Cho 4 số nguyên 16-bit \textbf{có dấu} có mã Hexa tương ứng như sau: \texttt{8792}, \texttt{2763}, \texttt{F652}, \texttt{A876}. Số nào là số lớn nhất?\\
    A. \texttt{2763} \qquad B. \texttt{A876} \qquad C. \texttt{F652} \qquad D. \texttt{8792}
    \item \textbf{(Mã đề 204 -- Câu 17):} Cho 4 số nguyên 16-bit \textbf{không dấu} có mã Hexa tương ứng như sau: \texttt{8792}, \texttt{2763}, \texttt{F652}, \texttt{A876}. Số nào là số lớn nhất?\\
    A. \texttt{8792} \qquad B. \texttt{A876} \qquad C. \texttt{2763} \qquad D. \texttt{F652}
    \item \textbf{(Mã đề 199 -- Câu 11):} Cho 4 số nguyên 16-bit \textbf{không dấu} có mã Hexa: \texttt{79F1}, \texttt{2AB7}, \texttt{8000}, \texttt{FFAA}. Số nào là số nhỏ nhất?\\
    A. \texttt{8000} \qquad B. \texttt{79F1} \qquad C. \texttt{2AB7} \qquad D. \texttt{FFAA}
\end{enumerate}
\textbf{Hướng dẫn giải:}
\begin{itemize}
    \item \textbf{Giải câu a (Xét số có dấu):} Ba số \texttt{8792}, \texttt{F652}, \texttt{A876} đều bắt đầu bằng chữ số Hexa $\ge 8$ nên có bit dấu bằng $1$ (đều là các số nguyên âm). Chỉ có duy nhất số \textbf{\texttt{2763}} bắt đầu bằng chữ số $2 < 8$ (bit dấu bằng $0$) là số nguyên dương nên lớn nhất $\rightarrow$ Chọn \textbf{A (\texttt{2763})}.
    \item \textbf{Giải câu b (Xét số không dấu -- tìm số lớn nhất):} Khi không dấu, so sánh trực tiếp chữ số Hexa đầu tiên từ trái sang phải ($\texttt{F} > \texttt{A} > \texttt{8} > \texttt{2}$). Vậy số lớn nhất là \textbf{\texttt{F652}} $\rightarrow$ Chọn \textbf{D (\texttt{F652})}.
    \item \textbf{Giải câu c (Xét số không dấu -- tìm số nhỏ nhất):} So sánh chữ số Hexa đầu tiên ($\texttt{2} < \texttt{7} < \texttt{8} < \texttt{F}$). Vậy số nhỏ nhất là \textbf{\texttt{2AB7}} $\rightarrow$ Chọn \textbf{C (\texttt{2AB7})}.
\end{itemize}

\textbf{Bài tập 4.14 (Đề Giữa kỳ 2022.2 -- Đề 0006 Câu 1, Đề 0008, 0009 Câu 2 \& Đề Giữa kỳ 2025.2 -- Đề 0013 Câu 73, 942, Đề 0012 Câu 106)}\\
Hãy sắp xếp các bộ 3 số nguyên \textbf{8-bit có dấu} dưới đây theo đúng thứ tự yêu cầu của từng đề thi:
\begin{enumerate}[label=\alph*)]
    \item \textbf{(Đề GK 2022.2 -- Đề 0006, Câu 1):} Sắp xếp theo thứ tự \textbf{giảm dần}: $-108$ (hệ 10), \texttt{0001 0001} (hệ 2), \texttt{91} (hệ 16).
    \item \textbf{(Đề GK 2022.2 -- Đề 0009, Câu 2):} Mảng \texttt{A} gồm 3 phần tử 8-bit có dấu ở địa chỉ cơ sở \texttt{0x40000E} lần lượt là: \texttt{0xCE}, \texttt{0x12}, \texttt{0x20}. Sắp xếp theo giá trị \textbf{tăng dần}.
    \item \textbf{(Đề GK 2022.2 -- Đề 0008, Câu 2):} Mảng \texttt{A} gồm 3 phần tử 8-bit có dấu ở địa chỉ cơ sở \texttt{0x400008} lần lượt là: \texttt{0x38}, \texttt{0x00}, \texttt{0xD6}. Sắp xếp theo giá trị \textbf{tăng dần}.
    \item \textbf{(Đề GK 2025.2 -- Đề 0013, Câu 73):} Mảng \texttt{A} gồm 3 phần tử 8-bit có dấu ở địa chỉ cơ sở \texttt{0x400009} (tra bảng Memory Map Đề 0013 tại \texttt{0x400009, 0x40000A, 0x40000B}) lần lượt là: \texttt{0xDE}, \texttt{0x74}, \texttt{0x00}. Sắp xếp theo giá trị \textbf{tăng dần}.
    \item \textbf{(Đề GK 2025.2 -- Đề 0013, Câu 942):} Mảng \texttt{A} gồm 3 phần tử 8-bit có dấu ở địa chỉ cơ sở \texttt{0x400014} (tại \texttt{0x400014, 0x400015, 0x400016}) lần lượt là: \texttt{0xB7}, \texttt{0x59}, \texttt{0x34}. Sắp xếp theo giá trị \textbf{tăng dần}.
    \item \textbf{(Đề GK 2025.2 -- Đề 0012, Câu 106):} Mảng \texttt{A} gồm 3 phần tử 8-bit có dấu ở địa chỉ cơ sở \texttt{0x400009} (tra bảng Đề 0012 tại \texttt{0x400009, 0x40000A, 0x40000B}) lần lượt là: \texttt{0x0B}, \texttt{0x8B}, \texttt{0x03}. Sắp xếp theo giá trị \textbf{tăng dần}.
\end{enumerate}
\textbf{Hướng dẫn giải:}
\begin{itemize}
    \item \textbf{Quy tắc chung:} Với số 8-bit có dấu viết dạng Hexa, các số bắt đầu bằng \texttt{8..F} là \textbf{số âm}; các số bắt đầu bằng \texttt{0..7} là \textbf{số dương} (hoặc $0$).
    \item \textbf{Câu a:} Đổi về hệ 10: $-108$; $\texttt{0001 0001}_2 = +17$; $\texttt{0x91} = \texttt{1001 0001}_2 = -128 + 16 + 1 = -111$. Vì $+17 > -108 > -111$, thứ tự giảm dần là: \textbf{\texttt{0001 0001} $> -108 > \texttt{91}$}.
    \item \textbf{Câu b:} Có \texttt{0xCE} là số âm ($-50$), còn $\texttt{0x12} (+18) < \texttt{0x20} (+32) \implies$ Tăng dần: \textbf{$\texttt{0xCE} < \texttt{0x12} < \texttt{0x20}$}.
    \item \textbf{Câu c:} Có \texttt{0xD6} là số âm ($-42$), còn $\texttt{0x00} (0) < \texttt{0x38} (+56) \implies$ Tăng dần: \textbf{$\texttt{0xD6} < \texttt{0x00} < \texttt{0x38}$}.
    \item \textbf{Câu d (Đề 0013 -- Câu 73):} Có \texttt{0xDE} là số âm ($-34$), còn $\texttt{0x00} (0) < \texttt{0x74} (+116) \implies$ Tăng dần: \textbf{$\texttt{0xDE} < \texttt{0x00} < \texttt{0x74}$}.
    \item \textbf{Câu e (Đề 0013 -- Câu 942):} Có \texttt{0xB7} là số âm ($-73$), còn $\texttt{0x34} (+52) < \texttt{0x59} (+89) \implies$ Tăng dần: \textbf{$\texttt{0xB7} < \texttt{0x34} < \texttt{0x59}$}.
    \item \textbf{Câu f (Đề 0012 -- Câu 106):} Có \texttt{0x8B} là số âm ($-117$), còn $\texttt{0x03} (+3) < \texttt{0x0B} (+11) \implies$ Tăng dần: \textbf{$\texttt{0x8B} < \texttt{0x03} < \texttt{0x0B}$}.
\end{itemize}

\subsection{Nhóm 4: Phép cộng/trừ số nguyên và Hiện tượng Tràn số học (Overflow)}

\textbf{Bài tập 4.15 (Đề Giữa kỳ 2018.2 -- Mã đề 198 Câu 8 \& Mã đề 199 Câu 9)}\\
Máy tính thực hiện phép toán cộng 2 số 8-bit không dấu ở hệ cơ số 2 sau đây, kết quả (lưu trong thanh ghi 8-bit) sẽ bằng bao nhiêu:
$$\texttt{1010 1000} + \texttt{1111 1010}$$
\begin{enumerate}[label=\Alph*.]
    \item \texttt{1000 0010}
    \item \texttt{1000 0100}
    \item \texttt{1010 0100}
    \item \texttt{1010 0010}
\end{enumerate}
\textbf{Hướng dẫn giải:}
\begin{itemize}
    \item Thực hiện cộng từng cột bit từ phải sang trái:
    \begin{itemize}
        \item 4 bit thấp: $\texttt{1000} + \texttt{1010} = \texttt{1\,0010}$ (viết \textbf{\texttt{0010}}, nhớ $1$ sang 4 bit cao).
        \item 4 bit cao: $\texttt{1010} + \texttt{1111} + 1\text{ (nhớ)} = \texttt{1010} + \texttt{1\,0000} = \texttt{1\,1010}$ (trong 8 bit giữ lại \textbf{\texttt{1010}}, bit nhớ tràn ra ngoài $C_{out} = 1$).
    \end{itemize}
    \item Ghép 8 bit kết quả lại ta được: \textbf{\texttt{1010 0010}}.
    \item \textbf{Kết luận:} Đáp án đúng là \textbf{D (\texttt{1010 0010})}.
\end{itemize}

\textbf{Bài tập 4.16 (Ôn tập Cuối kỳ 2025.2 -- P1, Câu 2)}\\
Khẳng định nào dưới đây là đúng khi thực hiện phép cộng hai số nguyên $n\text{-bit}$?
\begin{enumerate}[label=\Alph*.]
    \item Khi cộng hai số nguyên không dấu $n\text{-bit}$, nếu $C_{out} = 1 \rightarrow$ kết quả sai.
    \item Khi cộng hai số nguyên có dấu $n\text{-bit}$, nếu $C_{out} = 1$ và hai số ban đầu cùng dấu nhau $\rightarrow$ kết quả sai.
    \item Khi cộng hai số nguyên có dấu $n\text{-bit}$, nếu hai số ban đầu trái dấu nhau $\rightarrow$ kết quả đúng.
    \item Khi cộng hai số nguyên có dấu $n\text{-bit}$, nếu hai số ban đầu cùng dấu nhau và bit cao nhất của kết quả giống bit cao nhất của hai số ban đầu $\rightarrow$ chưa thể kết luận tính đúng sai.
\end{enumerate}
\textbf{Hướng dẫn giải:}
\begin{itemize}
    \item Theo Slide Chương 4 (trang 70 và trang 74):
    \begin{itemize}
        \item Khi cộng hai số nguyên \textbf{không dấu} $n\text{-bit}$, nếu có nhớ ra ngoài ($C_{out} = 1$) thì tổng vượt quá $2^n - 1$ nên kết quả $n\text{-bit}$ bị sai $\implies$ \textbf{Khẳng định A đúng}.
        \item Khi cộng hai số nguyên \textbf{có dấu} $n\text{-bit}$, nếu hai số ban đầu \textbf{trái dấu nhau} thì tổng luôn nằm trong dải biểu diễn nên kết quả luôn luôn đúng $\implies$ \textbf{Khẳng định C đúng}.
        \item (Khẳng định B sai vì khi cộng hai số âm không bị tràn như $(-1) + (-1) = -2$, luôn có $C_{out} = 1$ nhưng kết quả vẫn đúng; Khẳng định D sai vì khi hai số cùng dấu mà kết quả vẫn cùng dấu thì chắc chắn kết quả đúng).
    \end{itemize}
    \item \textbf{Kết luận:} Chọn cả \textbf{A} và \textbf{C}.
\end{itemize}

\textbf{Bài tập 4.17 (Đề Giữa kỳ 2019.2 -- Mã đề 243 Câu 12 \& Mã đề 244 Câu 5)}\\
Tổng của 2 số 8-bit có dấu nào dưới đây sẽ bị hiện tượng tràn số học (Overflow)?
\begin{enumerate}[label=\alph*)]
    \item \textbf{(Mã đề 243 -- Câu 12):}\\
    A. $\texttt{01010001} + \texttt{00101100}$ \qquad B. $\texttt{10010001} + \texttt{01001000}$\\
    C. $\texttt{00010001} + \texttt{11001000}$ \qquad D. $\texttt{10010001} + \texttt{11001000}$
    \item \textbf{(Mã đề 244 -- Câu 5):}\\
    A. $\texttt{00010001} + \texttt{11001000}$ \qquad B. $\texttt{00011001} + \texttt{00111000}$\\
    C. $\texttt{01010001} + \texttt{01001000}$ \qquad D. $\texttt{10010001} + \texttt{01001000}$
\end{enumerate}
\textbf{Hướng dẫn giải:}
\begin{itemize}
    \item \textbf{Bước 1 (Loại ngay các cặp trái dấu):} Các phép cộng có một số bắt đầu bằng bit \texttt{0} và một số bắt đầu bằng bit \texttt{1} là cộng hai số trái dấu nên không bao giờ tràn (Loại B, C ở câu a; loại A, D ở câu b).
    \item \textbf{Bước 2 (Kiểm tra các cặp cùng dấu):}
    \begin{itemize}
        \item Ở câu a: Phương án A cộng hai số dương $81 + 44 = +125 \le 127$ (không tràn). Phương án D cộng hai số âm: $\texttt{1001\,0001}_2 (-111) + \texttt{1100\,1000}_2 (-56) = -167 < -128$ (kết quả 8-bit ra $\texttt{0101\,1001}_2$ có bit dấu bằng \texttt{0} -- lật thành số dương) $\implies$ \textbf{Bị tràn âm (Chọn D)}.
        \item Ở câu b: Phương án B cộng $25 + 56 = 81 \le 127$ (không tràn). Phương án C cộng hai số dương: $\texttt{0101\,0001}_2 (+81) + \texttt{0100\,1000}_2 (+72) = +153 > 127$ (kết quả 8-bit ra $\texttt{1001\,1001}_2$ có bit dấu bằng \texttt{1} -- lật thành số âm) $\implies$ \textbf{Bị tràn dương (Chọn C)}.
    \end{itemize}
\end{itemize}

\textbf{Bài tập 4.18 (Đề Cuối kỳ 2022.2 -- \#6d82ac, 2024.2 \& Ôn tập Cuối kỳ 2025.2 -- P1, Câu 66)}\\
Trong ngôn ngữ lập trình C, biến kiểu \texttt{char} là biến dữ liệu số nguyên có dấu 8-bit. Cho chương trình C dưới đây, hãy cho biết kết quả của biến \texttt{sum} hiển thị trên màn hình khi chạy chương trình?
\begin{lstlisting}
#include <stdio.h>
int main() {
    char a = -91;
    char b = -67;
    char sum = a + b;
    printf("sum = %d", sum); // hien thi sum dang thap phan
    return 0;
}
\end{lstlisting}
\begin{enumerate}[label=\Alph*.]
    \item $-98$
    \item $-158$
    \item $158$
    \item $98$
\end{enumerate}
\textbf{Hướng dẫn giải:}
\begin{itemize}
    \item Tổng toán học thông thường: $S = (-91) + (-67) = -158$.
    \item Vì biến \texttt{sum} thuộc kiểu \texttt{char} (8-bit có dấu, chỉ biểu diễn được từ $-128$ đến $+127$), nên giá trị $-158 < -128$ bị tràn âm (Overflow).
    \item Áp dụng công thức tràn âm 8-bit (cộng thêm $2^8 = 256$):
    $$\texttt{sum} = -158 + 256 = +98$$
    \item \textbf{Kết luận:} Đáp án đúng là \textbf{D (98)}.
\end{itemize}

\textbf{Bài tập 4.19 (Đề Cuối kỳ 2025.1 -- \#e3221e, Giữa kỳ 2024.1 -- \#f2d536, Giữa kỳ 2022.2 -- Đề 0007 Câu 1 \& Ôn tập 2025.2 -- Câu 40)}\\
Giả thiết tất cả các biến đều là số nguyên 8-bit có dấu (\texttt{signed char}). Hãy tính giá trị thập phân lưu trong biến kết quả sau khi thực hiện các phép cộng sau:
\begin{enumerate}[label=\alph*)]
    \item \textbf{(Đề CK 2025.1 -- \#e3221e):} Cho $x = 85, y = 52$. Tính $z = x + y$.
    \item \textbf{(Đề GK 2024.1 -- \#f2d536 \& Ôn tập 2025.2 Câu 40):} Tính $+65 + 65$.
    \item \textbf{(Đề GK 2022.2 -- Đề 0007, Câu 1a):} Tính $5 + 125$.
    \item \textbf{(Đề GK 2022.2 -- Đề 0007, Câu 1b):} Tính $(-108) + (-42)$.
\end{enumerate}
\textbf{Hướng dẫn giải:}\\
Dải biểu diễn của số nguyên 8-bit có dấu là $[-128, +127]$. Nếu tổng lớn hơn $+127$ thì trừ đi $256$; nếu tổng nhỏ hơn $-128$ thì cộng thêm $256$:
\begin{itemize}
    \item \textbf{Câu a:} $S = 85 + 52 = 137 > 127$ (tràn dương) $\implies z = 137 - 256 = \mathbf{-119}$.
    \item \textbf{Câu b:} $S = 65 + 65 = 130 > 127$ (tràn dương) $\implies \text{Kết quả} = 130 - 256 = \mathbf{-126}$.
    \item \textbf{Câu c:} $S = 5 + 125 = 130 > 127$ (tràn dương) $\implies \text{Kết quả} = 130 - 256 = \mathbf{-126}$.
    \item \textbf{Câu d:} $S = (-108) + (-42) = -150 < -128$ (tràn âm) $\implies \text{Kết quả} = -150 + 256 = \mathbf{+106}$.
\end{itemize}

\textbf{Bài tập 4.20 (Đề Cuối kỳ 2023.1 -- \#cd9132, 2023.2, Giữa kỳ 2024.1 -- \#e83dac \& Ôn tập 2025.2 -- Câu 7)}
\begin{enumerate}[label=\alph*)]
    \item \textbf{(Đề CK 2023.1 -- \#cd9132 \& Ôn tập 2025.2 Câu 7):} Giả sử thanh ghi \texttt{s0} đang lưu trữ giá trị thập phân là $1024$. Hãy cho biết phạm vi giá trị của thanh ghi \texttt{s1} để gây ra hiện tượng tràn số học (overflow) khi thực hiện lệnh \texttt{add s2, s0, s1}?\\
    A. \texttt{0x7FFFFC00 -- 0x7FFFFFFF} \qquad B. \texttt{0x7FFFFC01 -- 0x80000000}\\
    C. \texttt{0x7FFFFBFF -- 0x7FFFFFFF} \qquad D. \texttt{0x7FFFFC00 -- 0x80000000}
    \item \textbf{(Đề GK 2024.1 -- \#e83dac):} Giả sử thanh ghi \texttt{t0} đang lưu trữ giá trị thập phân là $512$. Phạm vi giá trị nào của thanh ghi \texttt{t1} sẽ gây ra tràn số học khi thực hiện lệnh \texttt{add t2, t0, t1}?\\
    A. \texttt{0x7FFFFD00 -- 0x7FFFFFFF} \qquad B. \texttt{0x7FFFFDFF -- 0x7FFFFFFF}\\
    C. \texttt{0x7FFFFE00 -- 0x80000000} \qquad D. \texttt{0x7FFFFE00 -- 0x7FFFFFFF}
\end{enumerate}
\textbf{Hướng dẫn giải:}
\begin{itemize}
    \item \textbf{Giải câu a (Với $K = 1024 = 2^{10} = \texttt{0x00000400}$):}
    \begin{itemize}
        \item Giá trị nhỏ nhất của \texttt{s1} bắt đầu làm tổng vượt quá $\texttt{0x7FFFFFFF}$ là:
        $$\texttt{s1}_{\min} = \texttt{0x80000000} - \texttt{0x400} = \texttt{0x7FFFFC00}\quad (\text{vì }\texttt{0x1000} - \texttt{0x400} = \texttt{0xC00})$$
        \item Giá trị dương lớn nhất của thanh ghi 32-bit là \texttt{0x7FFFFFFF}.
        \item Kết luận: Chọn \textbf{A (\texttt{0x7FFFFC00 -- 0x7FFFFFFF})}.
    \end{itemize}
    \item \textbf{Giải câu b (Với $K = 512 = 2^9 = \texttt{0x00000200}$):}
    \begin{itemize}
        \item Giá trị nhỏ nhất của \texttt{t1} bắt đầu gây tràn là:
        $$\texttt{t1}_{\min} = \texttt{0x80000000} - \texttt{0x200} = \texttt{0x7FFFFE00}\quad (\text{vì }\texttt{0x1000} - \texttt{0x200} = \texttt{0xE00})$$
        \item Kết luận: Chọn \textbf{D (\texttt{0x7FFFFE00 -- 0x7FFFFFFF})}.
    \end{itemize}
\end{itemize}

\textbf{Bài tập 4.21 (Ôn tập Cuối kỳ 2025.1 -- Câu 15 \& Ôn tập 2025.2 -- P2, Câu 59)}\\
Đoạn lệnh nào sau đây có xảy ra hiện tượng tràn số học (OVERFLOW)?
\begin{enumerate}[label=\Alph*.]
    \item \texttt{li x8, 0xFFFFFFFF} $\rightarrow$ \texttt{li x9, 0xFFFFFFFF} $\rightarrow$ \texttt{add x10, x8, x9}
    \item \texttt{li x8, 0xFFFFFFFF} $\rightarrow$ \texttt{li x9, 0x80000000} $\rightarrow$ \texttt{add x10, x8, x9}
    \item \texttt{li x8, 0xFFFFFFFF} $\rightarrow$ \texttt{li x9, 0x80000000} $\rightarrow$ \texttt{sub x10, x8, x9}
    \item \texttt{li x8, 0xFFFFFFFF} $\rightarrow$ \texttt{li x9, 0xFFFFFFFF} $\rightarrow$ \texttt{sub x10, x8, x9}
\end{enumerate}
\textbf{Hướng dẫn giải:}
\begin{itemize}
    \item Ta có $\texttt{0xFFFFFFFF} = -1$ và $\texttt{0x80000000} = -2^{31}$ (số nguyên âm nhỏ nhất của 32-bit).
    \item Kiểm tra từng phương án:
    \begin{itemize}
        \item Phương án A: $(-1) + (-1) = -2$ (nằm trong dải, không tràn).
        \item Phương án B: Cộng hai số âm $(-1) + (-2^{31}) = -(2^{31} + 1) < -2^{31}$. Vì $\texttt{x9} = -2^{31}$ đã ở kịch biên âm của 32-bit, cộng thêm $-1$ sẽ vượt ra ngoài dải biểu diễn và lật thành số dương \texttt{0x7FFFFFFF} $\implies$ \textbf{Xảy ra tràn số học (Overflow)}.
        \item Phương án C: $(-1) - (-2^{31}) = +2^{31} - 1 = \texttt{0x7FFFFFFF}$ (vừa đúng biên dương lớn nhất, chưa tràn).
        \item Phương án D: $(-1) - (-1) = 0$ (không tràn).
    \end{itemize}
    \item \textbf{Kết luận:} Đáp án đúng là \textbf{B}.
\end{itemize}

\textbf{Bài tập 4.22 (Đề Giữa kỳ 2021.1 -- IT3030E, Problem 4)}\\
Hai thanh ghi \texttt{s0} và \texttt{s1} đang giữ các giá trị:
$$\texttt{s0 = 0x80000000} \quad ; \quad \texttt{s1 = 0xD0000000}$$
Hãy xác định giá trị Hexa của thanh ghi \texttt{t0} và cho biết có xảy ra tràn số học (Overflow) hay không cho 3 đoạn lệnh sau:
\begin{enumerate}[label=4.\arabic*.]
    \item \texttt{add t0, s0, s1}
    \item \texttt{sub t0, s0, s1}
    \item \texttt{add t0, s0, s1} $\rightarrow$ \texttt{add t0, t0, s0}
\end{enumerate}
\textbf{Hướng dẫn giải:}
\begin{itemize}
    \item Nhận xét dấu ban đầu: Cả $\texttt{s0} = \texttt{0x80000000}$ (chữ đầu $\texttt{8} = 1000_2$) và $\texttt{s1} = \texttt{0xD0000000}$ (chữ đầu $\texttt{D} = 1101_2$) đều có bit dấu số 31 bằng $1$, tức \textbf{cả \texttt{s0} và \texttt{s1} đều là số nguyên âm}.
    \item \textbf{Giải ý 4.1 (\texttt{add t0, s0, s1}):}
    \begin{itemize}
        \item Cộng chữ số Hexa đầu tiên: $\texttt{8} + \texttt{D} = 8 + 13 = 21 = 16 + 5 = \texttt{15}_{(16)}$ (tràn chữ số $1$ ra ngoài 32-bit, còn lại chữ số $\texttt{5}$).
        \item Giá trị lưu trong \texttt{t0} là: $\mathbf{\texttt{t0 = 0x50000000}}$.
        \item Vì cộng hai số âm (\texttt{s0 < 0, s1 < 0}) mà kết quả thu được lại là số dương ($\texttt{t0} = \texttt{0x50000000}$ có bit dấu bằng $0$), nên \textbf{CÓ xảy ra tràn số học (Overflow)}.
    \end{itemize}
    \item \textbf{Giải ý 4.2 (\texttt{sub t0, s0, s1}):}
    \begin{itemize}
        \item Thực hiện phép trừ: $\texttt{s0} - \texttt{s1} = \texttt{0x80000000} - \texttt{0xD0000000}$. Đổi $-\texttt{s1}$ bằng cách lấy bù 2 của \texttt{0xD0000000} là \texttt{0x30000000} (số dương).
        \item Cộng lại: $\texttt{0x80000000} + \texttt{0x30000000} = \mathbf{\texttt{0xB0000000}}$.
        \item Vì đây là phép cộng hai số trái dấu ($\texttt{s0} < 0$ và $-\texttt{s1} > 0$), nên \textbf{KHÔNG xảy ra tràn số học}.
    \end{itemize}
    \item \textbf{Giải ý 4.3 (\texttt{add t0, s0, s1} rồi \texttt{add t0, t0, s0}):}
    \begin{itemize}
        \item Sau lệnh đầu tiên, $\texttt{t0} = \texttt{0x50000000}$ (có tràn ở lệnh 1).
        \item Ở lệnh thứ hai, cộng $\texttt{t0} + \texttt{s0} = \texttt{0x50000000} + \texttt{0x80000000} = \mathbf{\texttt{0xD0000000}}$ (bản thân lệnh thứ hai cộng số dương \texttt{0x50000000} với số âm \texttt{0x80000000} nên lệnh thứ hai không tràn).
    \end{itemize}
\end{itemize}

\subsection{Nhóm 5: Lệnh Chia và Lấy phần dư số nguyên (\texttt{div, rem})}

\textbf{Bài tập 4.23 (Đề Giữa kỳ 2024.1 -- \#054a4a)}\\
Giá trị của các thanh ghi \texttt{t0, t1} sau khi thực hiện đoạn mã RISC-V dưới đây là gì?
\begin{lstlisting}
    li  s0, 67
    li  s1, -13
    div t0, s0, s1
    rem t1, s0, s1
\end{lstlisting}
\begin{enumerate}[label=\Alph*.]
    \item $\texttt{t0} = \texttt{0xfffffffb}, \texttt{t1} = \texttt{0x00000010}$
    \item $\texttt{t0} = \texttt{0x00000005}, \texttt{t1} = \texttt{0xfffffffd}$
    \item $\texttt{t0} = \texttt{0xfffffffb}, \texttt{t1} = \texttt{0x00000005}$
    \item $\texttt{t0} = \texttt{0xfffffffb}, \texttt{t1} = \texttt{0x00000002}$
\end{enumerate}
\textbf{Hướng dẫn giải:}
\begin{itemize}
    \item Ta có số bị chia $\texttt{s0} = +67$ và số chia $\texttt{s1} = -13$.
    \item \textbf{Tính thương \texttt{t0} (\texttt{div t0, s0, s1}):}
    \begin{itemize}
        \item Chia độ lớn: $67 \div 13 = 5$ (vì $13 \times 5 = 65$).
        \item Vì hai số trái dấu nên thương mang dấu âm: $\texttt{t0} = -5$.
        \item Đổi $-5$ sang Hexa 32-bit (lấy $16 - 5 = 11 = \texttt{b}$): $\texttt{t0} = \mathbf{\texttt{0xfffffffb}}$.
    \end{itemize}
    \item \textbf{Tính số dư \texttt{t1} (\texttt{rem t1, s0, s1}):}
    \begin{itemize}
        \item Độ lớn phần dư: $67 - 65 = 2$.
        \item Theo quy tắc: \textbf{Số dư luôn cùng dấu với số bị chia} ($\texttt{s0} = +67 > 0$), nên số dư mang dấu dương: $\texttt{t1} = +2 = \mathbf{\texttt{0x00000002}}$.
    \end{itemize}
    \item \textbf{Kết luận:} Đáp án đúng là \textbf{D ($\texttt{t0} = \texttt{0xfffffffb}, \texttt{t1} = \texttt{0x00000002}$)}.
\end{itemize}

\subsection{Nhóm 6: Số thực dấu phẩy động chuẩn IEEE 754 (32-bit)}

\begin{tcolorbox}[bwbox, title={Quy trình 4 bước giải nhanh bài toán IEEE 754 32-bit}]
\begin{itemize}
    \item \textbf{Bài toán 1: Cho mã Hexa 32-bit $\rightarrow$ Tìm giá trị thập phân $X$}
    \begin{enumerate}
        \item Đổi 3 hoặc 4 chữ số Hexa đầu tiên sang nhị phân (12 hoặc 16 bit đầu).
        \item Tách \textbf{1 bit đầu} là dấu $S$ ($0 \rightarrow +$; $1 \rightarrow -$); \textbf{8 bit tiếp theo} là số mũ lệch $e$; các bit còn lại là phần lẻ định trị $m$.
        \item Tính số mũ thực: $E = e - 127$.
        \item Ghép bit ẩn $1.$ vào trước $m$ được $(1.m)_{(2)}$, rồi dịch dấu phẩy sang phải $E$ vị trí (nếu $E > 0$) và đổi sang hệ 10.
    \end{enumerate}
    \item \textbf{Bài toán 2: Cho số thực thập phân $X \rightarrow$ Tìm mã Hexa 32-bit}
    \begin{enumerate}
        \item Xác định bit dấu $S$ ($X > 0 \implies S = 0$; $X < 0 \implies S = 1$).
        \item Đổi $|X|$ sang nhị phân, rồi dịch dấu phẩy về ngay sau bit $1$ đầu tiên để đưa về dạng chuẩn $(1.m)_{(2)} \times 2^E$.
        \item Tính số mũ lệch $e = E + 127$ và đổi $e$ ra 8 bit nhị phân (nhớ rằng $128 = \texttt{1000\,0000}_2$).
        \item Ghép $[S\text{ (1 bit)} \mid e\text{ (8 bit)} \mid m\text{ (23 bit)}]$, rồi nhóm từng cụm 4 bit từ trái sang phải để viết thành 8 chữ số Hexa.
    \end{enumerate}
\end{itemize}
\end{tcolorbox}

\textbf{Bài tập 4.24 (Đề Giữa kỳ 2019.2 -- Mã đề 243 Câu 2 \& Mã đề 244 Câu 17)}\\
Số thực dấu phẩy động (chuẩn IEEE 754) \textbf{không thể} biểu diễn chính xác tuyệt đối được số nào dưới đây?
\begin{enumerate}[label=\alph*)]
    \item \textbf{(Mã đề 243 -- Câu 2):} A. $0.5$ \qquad B. $0.4$ \qquad C. $12$ \qquad D. $2$
    \item \textbf{(Mã đề 244 -- Câu 17):} A. $-8$ \qquad B. $-1$ \qquad C. $1$ \qquad D. $0.7$
\end{enumerate}
\textbf{Hướng dẫn giải:}
\begin{itemize}
    \item Các số nguyên nhỏ ($1, -1, 2, -8, 12$) và số thực có phần lẻ là lũy thừa âm của $2$ (như $0.5 = 2^{-1} = 0.1_{(2)}$) đều biểu diễn chính xác tuyệt đối bằng hữu hạn bit nhị phân.
    \item Các số thập phân có phần lẻ là $0.4 = \dfrac{2}{5}$ hoặc $0.7 = \dfrac{7}{10}$ (mẫu số chứa thừa số nguyên tố $5 \neq 2$) khi đổi sang nhị phân sẽ ra chuỗi bit tuần hoàn vô hạn ($0.4_{(10)} = 0.01100110011\dots_{(2)}$), nên bị làm tròn và không thể biểu diễn chính xác tuyệt đối.
    \item \textbf{Kết luận:} Câu a chọn \textbf{B ($0.4$)}; Câu b chọn \textbf{D ($0.7$)}.
\end{itemize}

\textbf{Bài tập 4.25 (Đề Giữa kỳ 2018.2 -- Mã đề 198, Câu 5)}\\
Số dấu phẩy động 32-bit sau có giá trị thập phân bằng bao nhiêu?
$$\texttt{1100 0000 1000 0000 0000 0000 0000 0000}$$
\begin{enumerate}[label=\Alph*.]
    \item $-5$
    \item $-4$
    \item $-2.5$
    \item $-8$
\end{enumerate}
\textbf{Hướng dẫn giải:}
\begin{itemize}
    \item Tách 3 trường (1 bit -- 8 bit -- 23 bit):
    $$\underbrace{\texttt{1}}_{S=1} \mid \underbrace{\texttt{1000\,0001}}_{e=129} \mid \underbrace{\texttt{000\,0000\dots}}_{m=0}$$
    \item Bit dấu $S = 1 \implies$ Số âm ($-$).
    \item Số mũ thực: $E = e - 127 = 129 - 127 = 2$.
    \item Phần định trị: $1.m = 1.0_{(2)} = 1$.
    \item Giá trị số thực: $X = -1.0 \times 2^2 = \mathbf{-4}$.
    \item \textbf{Kết luận:} Đáp án đúng là \textbf{B ($-4$)}.
\end{itemize}

\textbf{Bài tập 4.26 (Đề Cuối kỳ 2022.2 -- \#bd4c95, 2024.1 -- \#e73be8, 2024.2 \& Ôn tập 2025.2 -- Câu 39, Câu 53)}
\begin{enumerate}[label=\alph*)]
    \item \textbf{(Đề CK 2022.2, 2024.2 -- \#bd4c95 \& Ôn tập 2025.2 Câu 53):} Giả thiết $X$ là biến số dấu phẩy động theo chuẩn IEEE 754 32-bit, cho nội dung $X = \texttt{0x40400000}$. Hãy xác định giá trị của $X$ dưới dạng thập phân?\\
    A. $3$ \qquad B. $-3$ \qquad C. $2$ \qquad D. $4$
    \item \textbf{(Đề CK 2024.1 -- \#e73be8 \& Ôn tập 2025.2 Câu 39):} Cho nội dung $X = \texttt{0x40A00000}$. Hãy xác định giá trị của $X$ dưới dạng thập phân?\\
    A. $6$ \qquad B. $5$ \qquad C. $-5$ \qquad D. $4$
\end{enumerate}
\textbf{Hướng dẫn giải:}
\begin{itemize}
    \item \textbf{Giải câu a (Với $X = \texttt{0x40400000}$):}
    \begin{itemize}
        \item Đổi 3 chữ số Hexa đầu (\texttt{404}) sang 12 bit nhị phân:
        $$\texttt{404}_{(16)} = \texttt{0100\,0000\,0100}_{(2)} = \underbrace{\texttt{0}}_{S=0} \mid \underbrace{\texttt{1000\,0000}}_{e=128} \mid \underbrace{\texttt{100\dots}}_{m}$$
        \item Ta có: $S = 0$ (số dương); $E = 128 - 127 = 1$; $1.m = 1.1_{(2)} = 1.5$.
        \item Giá trị: $X = +1.5 \times 2^1 = \mathbf{3}$ $\rightarrow$ Chọn \textbf{A (3)}.
    \end{itemize}
    \item \textbf{Giải câu b (Với $X = \texttt{0x40A00000}$):}
    \begin{itemize}
        \item Đổi 3 chữ số Hexa đầu (\texttt{40A}) sang 12 bit nhị phân:
        $$\texttt{40A}_{(16)} = \texttt{0100\,0000\,1010}_{(2)} = \underbrace{\texttt{0}}_{S=0} \mid \underbrace{\texttt{1000\,0001}}_{e=129} \mid \underbrace{\texttt{010\dots}}_{m}$$
        \item Ta có: $S = 0$ (số dương); $E = 129 - 127 = 2$; $1.m = 1.01_{(2)} = 1.25$.
        \item Giá trị: $X = +1.25 \times 2^2 = 1.25 \times 4 = \mathbf{5}$ $\rightarrow$ Chọn \textbf{B (5)}.
    \end{itemize}
\end{itemize}

\textbf{Bài tập 4.27 (Đề Giữa kỳ 2024.1 -- \#94ec4c \& Ôn tập Cuối kỳ 2025.2 -- P1, Câu 30)}
\begin{enumerate}[label=\alph*)]
    \item \textbf{(Đề GK 2024.1 -- \#94ec4c):} Cho số thực dấu phẩy động IEEE 754 độ chính xác đơn 32-bit được biểu diễn bởi word \texttt{0xC11A0000}. Giá trị của số này ở hệ thập phân là gì?
    \item \textbf{(Ôn tập 2025.2 -- P1, Câu 30):} Giả thiết $X$ là số thực dấu phẩy động IEEE 754 32-bit được biểu diễn bởi \texttt{0xC1E00000}. Giá trị của $X$ ở hệ thập phân là gì?
\end{enumerate}
\textbf{Hướng dẫn giải:}
\begin{itemize}
    \item \textbf{Giải câu a (Với \texttt{0xC11A0000}):}
    \begin{itemize}
        \item Đổi 4 chữ số Hexa đầu (\texttt{C11A}) sang 16 bit nhị phân:
        $$\texttt{C11A}_{(16)} = \texttt{1100\,0001\,0001\,1010}_{(2)} = \underbrace{\texttt{1}}_{S=1} \mid \underbrace{\texttt{1000\,0010}}_{e=130} \mid \underbrace{\texttt{0011010\dots}}_{m}$$
        \item Ta có: $S = 1$ (số âm); $E = 130 - 127 = 3$; $1.m = 1.001101_{(2)}$.
        \item Dịch dấu phẩy sang phải $3$ vị trí:
        $$X = -1.001101_{(2)} \times 2^3 = -1001.101_{(2)} = -(8 + 1 + 0.5 + 0.125) = \mathbf{-9.625}$$
    \end{itemize}
    \item \textbf{Giải câu b (Với \texttt{0xC1E00000}):}
    \begin{itemize}
        \item Đổi 3 chữ số Hexa đầu (\texttt{C1E}) sang 12 bit nhị phân:
        $$\texttt{C1E}_{(16)} = \texttt{1100\,0001\,1110}_{(2)} = \underbrace{\texttt{1}}_{S=1} \mid \underbrace{\texttt{1000\,0011}}_{e=131} \mid \underbrace{\texttt{110\dots}}_{m}$$
        \item Ta có: $S = 1$ (số âm); $E = 131 - 127 = 4$; $1.m = 1.11_{(2)}$.
        \item Dịch dấu phẩy sang phải $4$ vị trí:
        $$X = -1.1100_{(2)} \times 2^4 = -11100_{(2)} = -(16 + 8 + 4) = \mathbf{-28}$$
    \end{itemize}
\end{itemize}

\textbf{Bài tập 4.28 (Đề Cuối kỳ 2021.2 -- Mã đề B5, Câu 1.1 \& Slide Chương 4 -- Trang 98, 99)}\\
Chuyển đổi các số biểu diễn dưới dạng Hexa sau về số thực dấu phẩy động theo chuẩn IEEE 754 32-bit:
\begin{enumerate}[label=\alph*)]
    \item \texttt{0x42A78000}
    \item \texttt{0xC1560000}
    \item \texttt{0x3FC00000} (\texttt{0011 1111 1100 0000 ...})
\end{enumerate}
\textbf{Hướng dẫn giải:}
\begin{itemize}
    \item \textbf{Giải câu a (Với \texttt{0x42A78000}):}
    \begin{itemize}
        \item Đổi 5 chữ số Hexa đầu (\texttt{42A78}) sang 20 bit nhị phân:
        $$\texttt{0100\,0010\,1010\,0111\,1000}_{(2)} = \underbrace{\texttt{0}}_{S=0} \mid \underbrace{\texttt{1000\,0101}}_{e=133} \mid \underbrace{\texttt{01001111000\dots}}_{m}$$
        \item Ta có: $S = 0$ (số dương); $E = 133 - 127 = 6$; $1.m = 1.01001111_{(2)}$.
        \item Dịch dấu phẩy sang phải $6$ vị trí:
        $$X = +1010011.11_{(2)} = (64 + 16 + 2 + 1) + (0.5 + 0.25) = \mathbf{+83.75}$$
    \end{itemize}
    \item \textbf{Giải câu b (Với \texttt{0xC1560000}):}
    \begin{itemize}
        \item Đổi 4 chữ số Hexa đầu (\texttt{C156}) sang 16 bit nhị phân:
        $$\texttt{1100\,0001\,0101\,0110}_{(2)} = \underbrace{\texttt{1}}_{S=1} \mid \underbrace{\texttt{1000\,0010}}_{e=130} \mid \underbrace{\texttt{1010110\dots}}_{m}$$
        \item Ta có: $S = 1$ (số âm); $E = 130 - 127 = 3$; $1.m = 1.101011_{(2)}$.
        \item Dịch dấu phẩy sang phải $3$ vị trí:
        $$X = -1101.011_{(2)} = -(8 + 4 + 1 + 0.25 + 0.125) = \mathbf{-13.375}$$
    \end{itemize}
    \item \textbf{Giải câu c (Với \texttt{0x3FC00000}):}
    \begin{itemize}
        \item Đổi 3 chữ số Hexa đầu (\texttt{3FC}) sang 12 bit nhị phân:
        $$\texttt{0011\,1111\,1100}_{(2)} = \underbrace{\texttt{0}}_{S=0} \mid \underbrace{\texttt{0111\,1111}}_{e=127} \mid \underbrace{\texttt{100\dots}}_{m}$$
        \item Ta có: $S = 0$; $E = 127 - 127 = 0$; $X = +1.1_{(2)} \times 2^0 = \mathbf{+1.5}$.
    \end{itemize}
\end{itemize}

\textbf{Bài tập 4.29 (Đề Giữa kỳ 2024.1 -- \#386374)}\\
Biểu diễn các giá trị thập phân sau sang số dấu phẩy động chuẩn IEEE 754 độ chính xác đơn (32-bit) dưới dạng số Hexa:
\begin{enumerate}[label=\alph*)]
    \item $-13.75$
    \item $-46.375$
\end{enumerate}
\textbf{Hướng dẫn giải:}
\begin{itemize}
    \item \textbf{Giải câu a (Biểu diễn $-13.75$):}
    \begin{itemize}
        \item Bước 1: Vì $-13.75 < 0$ nên bit dấu $S = 1$.
        \item Bước 2: Đổi $13.75$ sang nhị phân và chuẩn hóa:
        $$13 = 1101_{(2)} \ ; \ 0.75 = 0.11_{(2)} \implies 13.75 = 1101.11_{(2)} = 1.10111_{(2)} \times 2^3$$
        \item Bước 3: Tính số mũ lệch $e = E + 127 = 3 + 127 = 130 = 128 + 2 = \texttt{1000\,0010}_{(2)}$.
        \item Bước 4: Ghép $[S \mid e \mid m]$ và nhóm 4 bit một thành chữ số Hexa:
        $$\underbrace{\texttt{1}}_{S}\,\underbrace{\texttt{1000\,0010}}_{e}\,\underbrace{\texttt{101\,1100\,0000\dots}}_{m} = \texttt{1100\,0001\,0101\,1100\,0000\dots}_{(2)} = \mathbf{\texttt{0xC15C0000}}$$
    \end{itemize}
    \item \textbf{Giải câu b (Biểu diễn $-46.375$):}
    \begin{itemize}
        \item Bước 1: Số âm nên $S = 1$.
        \item Bước 2: $46 = 32 + 8 + 4 + 2 = 101110_{(2)}$ và $0.375 = 0.25 + 0.125 = 0.011_{(2)}$:
        $$46.375 = 101110.011_{(2)} = 1.01110011_{(2)} \times 2^5$$
        \item Bước 3: Tính số mũ lệch $e = 5 + 127 = 132 = 128 + 4 = \texttt{1000\,0100}_{(2)}$.
        \item Bước 4: Ghép $[S \mid e \mid m]$ và đổi sang Hexa:
        $$\underbrace{\texttt{1}}_{S}\,\underbrace{\texttt{1000\,0100}}_{e}\,\underbrace{\texttt{011\,1001\,1000\dots}}_{m} = \texttt{1100\,0010\,0011\,1001\,1000\,0000\dots}_{(2)} = \mathbf{\texttt{0xC2398000}}$$
    \end{itemize}
\end{itemize}

\textbf{Bài tập 4.30 (Ôn tập Cuối kỳ 2025.2 -- P1, Câu 46 \& Câu 59)}
\begin{enumerate}[label=\alph*)]
    \item \textbf{(Câu 59):} Cho số thực dấu phẩy động 32-bit theo chuẩn IEEE 754 viết theo hệ Hexa có 3 chữ số có ý nghĩa lớn nhất (bên trái) là \texttt{7F8}, phần còn lại chưa được biết. Khẳng định nào dưới đây có khả năng đúng? (Chọn tất cả các khẳng định đúng)\\
    1. Số thực này bằng $0$\\
    2. Số thực này biểu diễn NaN (Not a Number)\\
    3. Số thực này biểu diễn giá trị dương xác định\\
    4. Số thực này biểu diễn vô cùng (Infinity)
    \item \textbf{(Câu 46):} Có bao nhiêu cách biểu diễn NaN (Not a Number) khi dùng số thực dấu phẩy động 32-bit theo chuẩn IEEE 754?\\
    A. $128$ \qquad B. $2^{23} - 1$ \qquad C. $2^{24} - 2$ \qquad D. $2$
\end{enumerate}
\textbf{Hướng dẫn giải:}
\begin{itemize}
    \item \textbf{Giải câu a:} Đổi 3 chữ số Hexa đầu \texttt{7F8} sang 12 bit nhị phân:
    $$\texttt{7F8}_{(16)} = \texttt{0111\,1111\,1000}_{(2)} = \underbrace{\texttt{0}}_{S=0} \mid \underbrace{\texttt{1111\,1111}}_{e=255} \mid \underbrace{\texttt{000\dots}}_{\text{3 bit đầu của } m}$$
    Vì cả 8 bit của trường mũ $e$ đều bằng $1$ ($e = 255$), theo quy ước đặc biệt của IEEE 754 (Slide Chương 4, trang 100):
    \begin{itemize}
        \item Nếu 20 bit còn lại đều bằng $0$ (tức toàn bộ 23 bit $m = 0$), thì số này biểu diễn $+\infty$ (\textbf{vô cùng}).
        \item Nếu trong 20 bit còn lại có ít nhất 1 bit bằng $1$ (tức $m \neq 0$), thì số này biểu diễn \textbf{NaN (Not a Number)}.
        \item Vậy cả hai khẳng định \textbf{2 (biểu diễn NaN)} và \textbf{4 (biểu diễn vô cùng)} đều có khả năng đúng.
    \end{itemize}
    \item \textbf{Giải câu b:}
    \begin{itemize}
        \item Một chuỗi 32-bit là NaN khi trường mũ $e = \texttt{11111111}$ ($1$ cách) và trường định trị $m$ (23 bit) khác $0$.
        \item Số cách chọn 23 bit của $m$ sao cho không đồng thời bằng $0$ là: $2^{23} - 1$.
        \item Bit dấu $S$ ở đầu có thể bằng $0$ hoặc bằng $1$ ($2$ cách chọn).
        \item Tổng số cách biểu diễn NaN là: $2 \times (2^{23} - 1) = \mathbf{2^{24} - 2}$ $\rightarrow$ Chọn \textbf{C ($2^{24} - 2$)}.
    \end{itemize}
\end{itemize}
\chapter{BỘ XỬ LÝ}

\section{Phần 1: Lý thuyết trọng tâm}

\subsection{Tổng quan vi kiến trúc và Bộ xử lý đơn chu kỳ (Single-Cycle Processor)}

\textbf{1. Ba kiến trúc triển khai phần cứng (Microarchitecture) cho cùng một tập lệnh RISC-V}
\begin{itemize}
    \item \textbf{Bộ xử lý đơn chu kỳ (Single-Cycle Processor):} Mỗi lệnh thực thi trọn vẹn trong đúng $1$ chu kỳ xung nhịp ($\text{CPI} = 1$). Chu kỳ xung nhịp $T_{\text{clk}}$ bắt buộc phải đủ dài cho lệnh chậm nhất (thường là lệnh \texttt{lw}), nên hiệu năng bị giới hạn.
    \item \textbf{Bộ xử lý đa chu kỳ (Multicycle Processor):} Chia mỗi lệnh thành chuỗi các bước ngắn (từ $3$ đến $5$ chu kỳ tùy loại lệnh), cho phép tái sử dụng phần cứng (dùng chung $1$ bộ nhớ và $1$ ALU) nhưng $\text{CPI} > 1$.
    \item \textbf{Bộ xử lý đường ống (Pipelined Processor):} Chia mỗi lệnh thành các công đoạn ngắn (chuẩn $5$ công đoạn) và thực thi gối lên nhau nhiều lệnh đồng thời như dây chuyền lắp ráp, đạt thông lượng lý tưởng $\text{CPI} \to 1$ với chu kỳ xung nhịp nhỏ.
\end{itemize}

\textbf{2. Các thành phần chính của Đường dữ liệu đơn chu kỳ (Single-Cycle RV32I Datapath)}
\begin{itemize}
    \item \textbf{Bộ đếm chương trình (\texttt{PC}):} Thanh ghi $32\text{-bit}$ chứa địa chỉ của lệnh hiện tại.
    \item \textbf{Bộ nhớ lệnh (\texttt{Instruction Memory -- IMEM}):} Lưu trữ mã máy của chương trình; nhận địa chỉ từ \texttt{PC} và xuất ra lệnh $32\text{-bit}$ (\texttt{Instruction[31:0]}). Tách rời với Bộ nhớ dữ liệu để tránh xung đột trong cùng $1$ chu kỳ.
    \item \textbf{Tập thanh ghi (\texttt{Register File -- RF}):} Gồm $32$ thanh ghi $32\text{-bit}$ (\texttt{x0}--\texttt{x31}), đánh số bằng $5\text{ bit}$. Có $2$ cổng đọc (\texttt{Read register 1 = Instruction[19:15]} ứng với \texttt{rs1}; \texttt{Read register 2 = Instruction[24:20]} ứng với \texttt{rs2}) và $1$ cổng ghi (\texttt{Write register = Instruction[11:7]} ứng với \texttt{rd}) điều khiển bởi \texttt{RegWrite}.
    \item \textbf{Khối tạo hằng số tức thời (\texttt{Imm Gen}):} Tách trường hằng số $12\text{-bit}$ (theo đúng định dạng I, S hoặc B) từ lệnh $32\text{-bit}$ và mở rộng có dấu (sign-extend) thành $32\text{-bit}$.
    \item \textbf{Đơn vị số học và logic (\texttt{ALU}):} Thực hiện các phép toán số học/logic (lệnh R-type), cộng địa chỉ cơ sở + độ lệch (lệnh \texttt{lw/sw}), hoặc phép trừ để so sánh (lệnh \texttt{beq}) và xuất ra cờ \texttt{Zero} ($\text{Zero} = 1$ khi hiệu bằng $0$; $\text{Zero} = 0$ khi hiệu khác $0$).
    \item \textbf{Bộ nhớ dữ liệu (\texttt{Data Memory -- DMEM}):} Chỉ sử dụng cho lệnh \texttt{lw/lb} (\texttt{MemRead = 1}) và \texttt{sw/sb} (\texttt{MemWrite = 1}); các lệnh R-type và \texttt{beq} không truy cập \texttt{DMEM}.
    \item \textbf{Hai bộ cộng (\texttt{Add}) và Ba bộ chọn kênh (\texttt{MUX}):}
    \begin{itemize}
        \item Bộ cộng $1$ tính địa chỉ tuần tự $\text{PC} + 4$; bộ cộng $2$ tính địa chỉ đích rẽ nhánh:
        $$\text{BTA} = \text{PC} + \text{SignExt}(\{\text{imm}[12:1], \texttt{0}\})$$
        \item \texttt{MUX ALUSrc}: chọn toán hạng thứ $2$ của ALU (\texttt{0}: \texttt{Read data 2} từ \texttt{rs2}; \texttt{1}: ngõ ra \texttt{Imm Gen}).
        \item \texttt{MUX MemtoReg}: chọn dữ liệu ghi về \texttt{Write data} của tập thanh ghi (\texttt{0}: \texttt{ALU result}; \texttt{1}: \texttt{Read data} từ \texttt{DMEM}).
        \item \texttt{MUX PCSrc}: chọn giá trị cập nhật \texttt{PC} (\texttt{0}: $\text{PC} + 4$; \texttt{1}: $\text{BTA}$), với $\text{PCSrc} = \text{Branch} \land \text{Zero}$.
    \end{itemize}
\end{itemize}

\subsection{Đơn vị điều khiển (Control Unit) và Bảng tín hiệu điều khiển}

Đơn vị điều khiển trong bộ xử lý RISC-V đơn chu kỳ gồm hai khối mạch tổ hợp: \textbf{Bộ điều khiển chính (Main Control)} nhận $7\text{ bit}$ mã lệnh \texttt{opcode = Instruction[6:0]} và \textbf{Bộ điều khiển ALU (ALU Control)} nhận tín hiệu $2\text{-bit}$ \texttt{ALUOp} kết hợp với \texttt{Instruction[30, 14:12]} (bit $30$ của \texttt{funct7} và $3\text{ bit}$ \texttt{funct3}) để sinh ra mã $4\text{-bit}$ điều khiển phép toán của ALU.

\begin{tcolorbox}[bwbox, title={Bảng giải mã tín hiệu điều khiển ALU (ALU Control)}]
\begin{center}
\renewcommand{\arraystretch}{1.15}
\small
\begin{tabularx}{\linewidth}{|C{1.8cm}|C{1.4cm}|Y|C{1.8cm}|C{1.3cm}|C{1.8cm}|C{2.0cm}|}
\hline
\textbf{Loại lệnh} & \textbf{ALUOp} & \textbf{Thao tác} & \textbf{funct7} & \textbf{funct3} & \textbf{Hàm ALU} & \textbf{ALU Control} \\ \hline
\texttt{lw / sw} & \texttt{00} & Tính địa chỉ nhớ & \texttt{XXXXXXX} & \texttt{XXX} & \texttt{add} & \texttt{0010} \\ \hline
\texttt{beq} & \texttt{01} & So sánh bằng & \texttt{XXXXXXX} & \texttt{XXX} & \texttt{subtract} & \texttt{0110} \\ \hline
\multirow{4}{*}{R-type} & \texttt{10} & \texttt{add} & \texttt{0000000} & \texttt{000} & \texttt{add} & \texttt{0010} \\ \cline{2-7}
 & \texttt{10} & \texttt{sub} & \texttt{0100000} & \texttt{000} & \texttt{subtract} & \texttt{0110} \\ \cline{2-7}
 & \texttt{10} & \texttt{and} & \texttt{0000000} & \texttt{111} & \texttt{AND} & \texttt{0000} \\ \cline{2-7}
 & \texttt{10} & \texttt{or} & \texttt{0000000} & \texttt{110} & \texttt{OR} & \texttt{0001} \\ \hline
\end{tabularx}
\end{center}
\end{tcolorbox}

\begin{tcolorbox}[bwbox, title={Bảng chân trị 7 tín hiệu điều khiển chính (Single-Cycle RISC-V)}]
\begin{center}
\renewcommand{\arraystretch}{1.2}
\small
\begin{tabularx}{\linewidth}{|Y|C{1.7cm}|Z|Z|Z|Z|Z|Z|Z|}
\hline
\textbf{Loại lệnh} & \textbf{Opcode} & \textbf{Branch} & \textbf{MemRead} & \textbf{MemtoReg} & \textbf{ALUOp} & \textbf{MemWrite} & \textbf{ALUSrc} & \textbf{RegWrite} \\ \hline
R-type (\texttt{add, sub, slt}) & \texttt{0110011} & \texttt{0} & \texttt{0} & \texttt{0} & \texttt{10} & \texttt{0} & \texttt{0} & \texttt{1} \\ \hline
Load (\texttt{lw, lb, lh}) & \texttt{0000011} & \texttt{0} & \texttt{1} & \texttt{1} & \texttt{00} & \texttt{0} & \texttt{1} & \texttt{1} \\ \hline
Store (\texttt{sw, sb, sh}) & \texttt{0100011} & \texttt{0} & \texttt{0} & \texttt{X} (\texttt{0}) & \texttt{00} & \texttt{1} & \texttt{1} & \texttt{0} \\ \hline
Branch (\texttt{beq}) & \texttt{1100011} & \texttt{1} & \texttt{0} & \texttt{X} (\texttt{0}) & \texttt{01} & \texttt{0} & \texttt{0} & \texttt{0} \\ \hline
\end{tabularx}
\end{center}
\textit{Ghi chú:} Với lệnh \texttt{sw} và \texttt{beq} (do \texttt{RegWrite = 0}), tín hiệu \texttt{MemtoReg} về mặt phần cứng là \texttt{X} (Don't Care), nhưng trong đề thi trắc nghiệm/điền khuyết dạng chuỗi 7 tín hiệu thường ghi quy ước là \texttt{0}.
\end{tcolorbox}

\textbf{Quy tắc phân tích sự cố đứt dây / kẹt mức logic (Stuck-at Fault):}
\begin{itemize}
    \item Một sự cố trên đường dây hoặc tín hiệu điều khiển chỉ làm \textbf{sai kết quả thực thi} của một lệnh nếu lệnh đó \textbf{thực sự sử dụng} đường dây đó với giá trị khác mức bị kẹt.
    \item Ví dụ: Nếu \texttt{MemtoReg} bị kẹt ở \texttt{0}, mọi lệnh ghi từ ALU (R-type, \texttt{addi}) và lệnh không ghi thanh ghi (\texttt{sw}, \texttt{beq}) vẫn chạy bình thường; chỉ lệnh \texttt{lw} (cần \texttt{MemtoReg = 1}) bị ghi nhầm địa chỉ thay vì dữ liệu từ bộ nhớ.
\end{itemize}

\subsection{Bộ xử lý đa chu kỳ (Multicycle Processor)}

\begin{itemize}
    \item \textbf{Đặc điểm kiến trúc:} Chỉ dùng \textbf{một bộ nhớ duy nhất} cho cả lệnh và dữ liệu (phân biệt địa chỉ bằng tín hiệu \texttt{IorD}: \texttt{0} lấy từ \texttt{PC}, \texttt{1} lấy từ \texttt{ALUOut}) và \textbf{một bộ ALU duy nhất} (tái sử dụng cho tính $\text{PC}+4$, tính địa chỉ nhớ, thực thi phép toán và tính địa chỉ rẽ nhánh).
    \item \textbf{Năm thanh ghi trung gian (không kiến trúc):} \texttt{IR} (Thanh ghi lệnh, điều khiển bởi \texttt{IRWrite}), \texttt{MDR} (Thanh ghi dữ liệu bộ nhớ), \texttt{A} và \texttt{B} (giữ hai toán hạng đọc từ tập thanh ghi), và \texttt{ALUOut} (giữ ngõ ra ALU cho chu kỳ kế tiếp).
    \item \textbf{Tín hiệu điều khiển đặc trưng:} \texttt{ALUSrcA} ($1\text{-bit}$: \texttt{0} từ \texttt{PC}, \texttt{1} từ \texttt{A}), \texttt{ALUSrcB} ($2\text{-bit}$: \texttt{00} từ \texttt{B}, \texttt{01} là hằng số $4$, \texttt{10} là hằng số \texttt{imm} mở rộng), \texttt{PCWrite}, \texttt{PCWriteCond}, \texttt{PCSource}.
\end{itemize}

\begin{tcolorbox}[bwbox, title={Số chu kỳ thực thi (CPI) của Bộ xử lý đa chu kỳ so với đơn chu kỳ}]
\begin{itemize}
    \item \textbf{Số chu kỳ theo từng loại lệnh:} Lệnh \texttt{load} ($5$ chu kỳ), lệnh \texttt{store} ($4$ chu kỳ), lệnh ALU / R-type ($4$ chu kỳ), lệnh \texttt{branch} ($3$ chu kỳ).
    \item \textbf{Công thức CPI trung bình:} $\text{CPI}_{\text{TB}} = \sum_{i} p_i \times \text{CPI}_i$.
    \item \textit{Ví dụ chuẩn (Slide bài giảng):} Chương trình có $20\%$ \texttt{load}, $8\%$ \texttt{store}, $10\%$ \texttt{branch}, $62\%$ ALU $\Rightarrow \text{CPI}_{\text{TB}} = 0.20 \times 5 + 0.08 \times 4 + 0.10 \times 3 + 0.62 \times 4 = \mathbf{4.10}$ chu kỳ (nhanh hơn bộ xử lý đơn chu kỳ vốn luôn mất tương đương $5$ chu kỳ ngắn cho mọi lệnh).
\end{itemize}
\end{tcolorbox}

\subsection{Đường ống lệnh (Pipelining) và Tính toán hiệu năng Pipeline}

\textbf{1. Năm công đoạn chuẩn của đường ống lệnh RISC-V}
\begin{enumerate}
    \item \textbf{IF (Instruction Fetch):} Tìm nạp lệnh từ Bộ nhớ lệnh (\texttt{IMEM}) theo địa chỉ \texttt{PC} và tăng $\text{PC} \leftarrow \text{PC} + 4$.
    \item \textbf{ID (Instruction Decode / Register Read):} Giải mã lệnh, sinh tín hiệu điều khiển, mở rộng hằng số \texttt{Imm Gen} và đọc hai thanh ghi nguồn từ \texttt{Register File}.
    \item \textbf{EX (Execute / Address Calculation):} Thực hiện phép toán số học/logic hoặc tính địa chỉ bộ nhớ dữ liệu tại \texttt{ALU}.
    \item \textbf{MEM (Memory Access):} Đọc (\texttt{lw}) hoặc ghi (\texttt{sw}) dữ liệu tại Bộ nhớ dữ liệu (\texttt{DMEM}).
    \item \textbf{WB (Write Back):} Ghi kết quả (từ \texttt{ALU} hoặc \texttt{DMEM}) trả về thanh ghi đích \texttt{rd} trong \texttt{Register File}.
\end{enumerate}

Để tách biệt dữ liệu và tín hiệu điều khiển của các lệnh đang chạy đồng thời, bộ xử lý chèn $4$ \textbf{thanh ghi đệm đường ống (Pipeline Registers)} giữa $5$ công đoạn: \texttt{IF/ID}, \texttt{ID/EX}, \texttt{EX/MEM}, và \texttt{MEM/WB}.

\begin{tcolorbox}[bwbox, title={Các công thức trọng tâm tính toán hiệu năng Đường ống lệnh (Pipeline)}]
\begin{itemize}
    \item \textbf{Chu kỳ xung nhịp đường ống:}
    $$T_{\text{pipe}} = \max(t_{\text{IF}}, t_{\text{ID}}, t_{\text{EX}}, t_{\text{MEM}}, t_{\text{WB}}) + t_{\text{reg}}$$
    \item \textbf{Tổng số chu kỳ thực thi $N$ lệnh trên pipeline $k$ công đoạn (với $S$ chu kỳ trễ stall/bubble/flush):}
    $$\text{Cycles} = k + (N - 1) + S = N + (k - 1) + S$$
    (Với $k = 5$ công đoạn và không có hazard: $\text{Cycles} = N + 4$.)
    \item \textbf{Tổng thời gian thực thi chương trình:}
    $$T_{\text{exec}} = \text{Cycles} \times T_{\text{pipe}} = [N + (k - 1) + S] \times T_{\text{pipe}}$$
    \item \textbf{Chỉ số CPI trung bình trên $N$ lệnh:}
    $$\text{CPI} = \frac{N + (k - 1) + S}{N} = 1 + \frac{k - 1 + S}{N}$$
    \item \textbf{Hệ số tăng tốc (Speedup) so với thực thi tuần tự không đường ống:}
    $$\text{Speedup} = \frac{N \times T_{\text{single}}}{[N + (k - 1) + S] \times T_{\text{pipe}}} \xrightarrow{N \to \infty,\; S = 0} \frac{T_{\text{single}}}{T_{\text{pipe}}} \le k$$
\end{itemize}
\end{tcolorbox}

\textbf{Quy ước vàng về ghi/đọc Tập thanh ghi trong cùng $1$ chu kỳ:} Trong bộ xử lý đường ống RISC-V chuẩn, tầng \textbf{WB ghi dữ liệu vào tập thanh ghi ở nửa đầu chu kỳ}, còn tầng \textbf{ID đọc thanh ghi ở nửa sau chu kỳ}. Nhờ đó, nếu lệnh $I_1$ đang ở tầng \textbf{WB} cùng chu kỳ với lệnh $I_2$ ở tầng \textbf{ID}, lệnh $I_2$ đọc được ngay giá trị mới mà không cần chèn thêm stall.

\subsection{Xung đột đường ống (Hazards), Forwarding và Dự đoán rẽ nhánh}

\textbf{1. Ba loại xung đột đường ống (Pipeline Hazards)}
\begin{itemize}
    \item \textbf{Xung đột cấu trúc (Structural Hazard):} Hai lệnh đồng thời tranh chấp cùng một tài nguyên phần cứng (ví dụ: tầng IF đọc lệnh và tầng MEM đọc/ghi dữ liệu nếu dùng chung $1$ bộ nhớ). \textit{Cách khắc phục:} tách riêng Bộ nhớ lệnh (\texttt{IMEM} / I-Cache) và Bộ nhớ dữ liệu (\texttt{DMEM} / D-Cache).
    \item \textbf{Xung đột dữ liệu (Data Hazard):} Lệnh sau phụ thuộc vào dữ liệu mà lệnh trước chưa kịp ghi về tập thanh ghi. Phân loại phụ thuộc: \textbf{RAW} (Read After Write -- phụ thuộc thực, phổ biến nhất), \textbf{WAR} (Write After Read), \textbf{WAW} (Write After Write).
    \begin{itemize}
        \item \textit{Không có Forwarding (dùng Stall / chèn \texttt{nop}):} Giữa lệnh ghi thanh ghi và lệnh liền sau đọc thanh ghi đó cần chèn \textbf{$2$ chu kỳ stall/\texttt{nop}} (theo quy ước WB nửa đầu / ID nửa sau; nếu cách $1$ lệnh thì cần $1$ stall; cách $\ge 2$ lệnh thì không cần stall).
        \item \textit{Kỹ thuật chuyển tiếp dữ liệu (Forwarding / Bypassing):} Truyền trực tiếp kết quả từ thanh ghi đệm \texttt{EX/MEM} (EX-to-EX) hoặc \texttt{MEM/WB} (MEM-to-EX) về thẳng đầu vào ALU mà không cần đợi tầng WB. Kỹ thuật này \textbf{xóa bỏ hoàn toàn stall} cho các xung đột ALU $\rightarrow$ ALU.
        \item \textit{Xung đột Load-Use Data Hazard:} Khi lệnh \texttt{lw} đứng ngay trước lệnh sử dụng thanh ghi đích của nó (ví dụ: \texttt{lw s0, 20(t1)} rồi \texttt{sub s1, s0, t2}), dữ liệu chỉ sẵn sàng sau tầng \textbf{MEM} nên dù có đầy đủ Forwarding vẫn bắt buộc chèn \textbf{$1$ chu kỳ stall (bubble)} nhờ khối \textbf{Hazard Detection Unit}. Trình biên dịch có thể tối ưu bằng kỹ thuật \textbf{lập lịch mã (Code Scheduling)} -- xen lệnh độc lập vào sau \texttt{lw}.
    \end{itemize}
    \item \textbf{Xung đột điều khiển (Control Hazard):} Phát sinh tại các lệnh rẽ nhánh (\texttt{beq, bne, ...}) vì chưa biết địa chỉ lệnh kế tiếp khi đang tìm nạp ở tầng IF.
    \begin{itemize}
        \item \textit{Xác định sớm ở tầng ID:} Bổ sung mạch so sánh thanh ghi và tính BTA ngay trong công đoạn ID giúp giảm phạt dự đoán sai xuống còn \textbf{$1$ chu kỳ (flush $1$ lệnh)}.
        \item \textit{Dự đoán rẽ nhánh (Branch Prediction):} Tiếp tục nạp lệnh theo hướng dự đoán; nếu đoán đúng thì không mất chu kỳ phạt ($0$ stall), nếu đoán sai thì hủy (flush) các lệnh nạp nhầm và nạp lại từ địa chỉ đúng.
    \end{itemize}
\end{itemize}

\begin{tcolorbox}[bwbox, title={Các chiến lược Dự đoán rẽ nhánh (Static vs Dynamic 1-bit \& 2-bit)}]
\begin{itemize}
    \item \textbf{Dự đoán tĩnh (Static Branch Prediction):}
    \begin{itemize}
        \item \textit{Always Not-Taken (Luôn không nhảy):} Luôn nạp lệnh tuần tự $\text{PC} + 4$; bị phạt khi nhánh thực sự nhảy.
        \item \textit{Always Taken (Luôn nhảy):} Luôn nạp tại địa chỉ đích BTA; rất hiệu quả với vòng lặp (chỉ đoán sai $1$ lần khi thoát vòng lặp nếu là nhánh quay về đầu vòng lặp).
    \end{itemize}
    \item \textbf{Dự đoán động 1-bit (1-bit Dynamic Predictor):} Lưu trạng thái lần rẽ nhánh gần nhất (\texttt{1}: Taken, \texttt{0}: Not-Taken). \textit{Nhược điểm:} Với vòng lặp lồng nhau, mỗi lần thoát vòng lặp trong bị đoán sai $1$ lần và đổi trạng thái sang Not-Taken, dẫn đến lần đầu quay lại vòng lặp trong lại bị đoán sai thêm $1$ lần nữa ($2$ lần sai/vòng ngoài).
    \item \textbf{Dự đoán động 2-bit bão hòa (2-bit Saturating Counter):} Gồm $4$ trạng thái: \texttt{11} (\textit{Strongly Taken}), \texttt{10} (\textit{Weakly Taken}), \texttt{01} (\textit{Weakly Not-Taken}), \texttt{00} (\textit{Strongly Not-Taken}). Phải đoán sai \textbf{$2$ lần liên tiếp} mới đổi hướng dự đoán, giúp giảm đáng kể số lần đoán sai ở vòng lặp lồng nhau.
\end{itemize}
\end{tcolorbox}

\textbf{2. Hướng mở rộng song song mức lệnh (ILP):} Để tiếp tục tăng thông lượng vượt quá $1$ lệnh/chu kỳ, các bộ xử lý hiện đại kết hợp \textbf{tăng số công đoạn đường ống (Superpipelining)} và kiến trúc \textbf{Siêu vô hướng (Superscalar)} -- bố trí nhiều đơn vị thực thi song song để phát hành và hoàn tất nhiều lệnh trong cùng $1$ chu kỳ xung nhịp.

\newpage
\section{Phần 2: Toàn bộ bài tập Chương 5 và Hướng dẫn giải}


\subsection{Nhóm 1: Đường dữ liệu (Datapath) \& Tín hiệu điều khiển (Control Unit)}

\textbf{Bài tập 5.1 (ôn tập cuối kì 20251 (Câu 36); ôn tập cuối kì 20252 (P2-18); Đáp án ôn tập cuối kì P2 (Câu 12))}\\
Xét 4 loại lệnh \texttt{lw}, \texttt{sw}, \texttt{beq} và \texttt{R-type} trên bộ xử lý RISC-V đơn chu kỳ. Khi thực hiện những loại lệnh nào \textbf{không cần} sử dụng tới thành phần \textbf{Bộ nhớ dữ liệu (Data Memory)}?
\begin{enumerate}[label=\Alph*.]
    \item beq và R-type
    \item lw và sw
    \item Chỉ lệnh R-type
    \item Cả 4 loại lệnh (lw, sw, beq và R-type)
\end{enumerate}
\textbf{Hướng dẫn giải:}
\begin{itemize}
    \item Trên Datapath đơn chu kỳ RISC-V:
    \item \texttt{lw} sử dụng Data Memory để \textbf{đọc} dữ liệu (\texttt{MemRead = 1}).
    \item \texttt{sw} sử dụng Data Memory để \textbf{ghi} dữ liệu (\texttt{MemWrite = 1}).
    \item \texttt{R-type} chỉ tính toán tại ALU rồi ghi kết quả trực tiếp về tập thanh ghi (\texttt{MemRead = 0, MemWrite = 0}).
    \item \texttt{beq} chỉ so sánh hai thanh ghi tại ALU và cập nhật PC nếu bằng nhau (\texttt{MemRead = 0, MemWrite = 0}).
    \item Vậy \textbf{\texttt{beq} và \texttt{R-type}} không cần sử dụng tới Bộ nhớ dữ liệu.
    \item \textbf{Kết luận:} Đáp án đúng là \textbf{A (beq và R-type)}.
\end{itemize}

\textbf{Bài tập 5.2 (ôn tập cuối kì 20251 (Câu 42); ôn tập cuối kì 20252 (P2-41); Đáp án ôn tập cuối kì P2 (Câu 20))}\\
Xét bộ xử lý RISC-V đơn chu kỳ chuẩn. Khi thực hiện lệnh \texttt{add s2, s4, s3}, cặp tín hiệu điều khiển \texttt{ALUSrc / RegWrite} nhận giá trị nào?
\begin{enumerate}[label=\Alph*.]
    \item 0/1
    \item 0/0
    \item 1/0
    \item 1/1
\end{enumerate}
\textbf{Hướng dẫn giải:}
\begin{itemize}
    \item Lệnh \texttt{add s2, s4, s3} thuộc định dạng \textbf{R-type}:
    \item \texttt{ALUSrc = 0}: toán hạng thứ hai đưa vào ALU lấy trực tiếp từ cổng \texttt{Read data 2} của tập thanh ghi (giá trị \texttt{s3}), không dùng khối \texttt{Imm Gen}.
    \item \texttt{RegWrite = 1}: cho phép ghi kết quả từ ALU vào thanh ghi đích \texttt{rd = s2}.
    \item Do đó cặp tín hiệu \texttt{ALUSrc / RegWrite} là \textbf{\texttt{0/1}}.
    \item \textbf{Kết luận:} Đáp án đúng là \textbf{A (0/1)}.
\end{itemize}

\textbf{Bài tập 5.3 (ôn tập cuối kì 20251 (Câu 53); ôn tập cuối kì 20252 (P2-28); Đáp án ôn tập cuối kì P2 (Câu 63))}\\
Xét Datapath RISC-V đơn chu kỳ chuẩn khi thực hiện lệnh \texttt{sw s0, -8(s6)}. Hãy cho biết bộ $7$ tín hiệu điều khiển theo đúng thứ tự \texttt{Branch / MemRead / MemtoReg / ALUOp (2-bit) / MemWrite / ALUSrc / RegWrite} nhận giá trị nào? (Viết bằng 7 cụm cách nhau bởi dấu \texttt{/}, vd: \texttt{1/0/1/00/1/0/1})
\begin{enumerate}[label=\Alph*.]
    \item 0/0/0/00/1/1/0
    \item 0/1/1/00/0/1/1
    \item 0/0/0/10/0/0/1
    \item 1/0/0/01/0/0/0
\end{enumerate}
\textbf{Hướng dẫn giải:}
\begin{itemize}
    \item Lệnh \texttt{sw s0, -8(s6)} là lệnh ghi từ nhớ (\textbf{S-type Store}):
    \item \texttt{Branch = 0}: không phải lệnh rẽ nhánh.
    \item \texttt{MemRead = 0}: không đọc từ Data Memory.
    \item \texttt{MemtoReg = 0} (hoặc \texttt{X} --- không quan tâm vì không ghi thanh ghi; đáp án chuẩn điền \texttt{0}).
    \item \texttt{ALUOp = 00}: ALU thực hiện phép cộng địa chỉ cơ sở \texttt{s6} với offset \texttt{-8}.
    \item \texttt{MemWrite = 1}: cho phép ghi giá trị thanh ghi \texttt{s0} vào Data Memory.
    \item \texttt{ALUSrc = 1}: toán hạng thứ 2 của ALU lấy từ khối \texttt{Imm Gen} (giá trị \texttt{-8}).
    \item \texttt{RegWrite = 0}: không ghi dữ liệu về tập thanh ghi.
    \item Kết quả theo đúng thứ tự: \textbf{\texttt{0/0/0/00/1/1/0}}.
    \item \textbf{Kết luận:} Đáp án đúng là \textbf{A (0/0/0/00/1/1/0)}.
\end{itemize}

\textbf{Bài tập 5.4 (giữa kì 20231 (\#1c9040))}\\
Xét bộ xử lý đơn chu kỳ \textbf{RISC-V (RV32I)} thực thi lệnh \texttt{slt t6, t2, t3} (biết $t_1 = \text{0x0004}$, $t_2 = \text{0x0012}$, $t_3 = \text{0x2023}$). Điểm đánh dấu \textbf{\textcircled{1}} là bus 5-bit \texttt{Instruction[19:15]} đưa vào cổng \texttt{Read register 1} của Tập thanh ghi. Giá trị Hexa trên đường dây \textbf{\textcircled{1}} là bao nhiêu?
\textit{Nhập giá trị Hexa 5-bit của rs1 (vd: 0x07):}
\textbf{Hướng dẫn giải:}
\begin{itemize}
    \item Đường dây \texttt{Instruction[19:15]} mang \textbf{chỉ số thanh ghi nguồn 1 (rs1)} (chứ không phải dữ liệu bên trong thanh ghi). Trong lệnh \texttt{slt t6, t2, t3}, \texttt{rs1} là thanh ghi \texttt{t2} có số hiệu là \texttt{x7} $= \mathbf{\text{0x07}}$ (bảng thanh ghi RV32I: \texttt{t2 = x7}).
    \item \textbf{Kết luận:} Đáp án là \textbf{0x07}.
\end{itemize}

\textbf{Bài tập 5.5 (giữa kì 20231 (\#282a13))}\\
Xét bộ xử lý đơn chu kỳ \textbf{RISC-V (RV32I)} thực thi lệnh \texttt{slt t6, t2, t3} tại địa chỉ \texttt{0x0040FF80} (biết $t_1 = \text{0x00000004}$, $t_2 = \text{0x00000012}$, $t_3 = \text{0x00002023}$). Điểm đánh dấu \textbf{\textcircled{3}} là ngõ ra \texttt{ALU result}. Giá trị Hexa 32-bit tại \textbf{\textcircled{3}} là bao nhiêu?
\textit{Nhập kết quả Hexa 32-bit (vd: 0x00000001):}
\textbf{Hướng dẫn giải:}
\begin{itemize}
    \item Lệnh \texttt{slt t6, t2, t3} so sánh có dấu \texttt{t2 $<$ t3}: $\text{0x12} (18) < \text{0x2023} (8227)$ là đúng, nên ngõ ra ALU result tại \textcircled{3} bằng $\mathbf{\text{0x00000001}}$.
    \item \textbf{Kết luận:} Đáp án là \textbf{0x00000001}.
\end{itemize}

\textbf{Bài tập 5.6 (giữa kì 20231 (\#dc67c2))}\\
Xét bộ xử lý đơn chu kỳ \textbf{RISC-V (RV32I)} thực thi lệnh \texttt{slt t6, t2, t3} (mã máy \texttt{0x01C3AFB3}) đặt tại địa chỉ \texttt{PC = 0x0040FF00}. Điểm đánh dấu \textbf{\textcircled{4}} là ngõ ra bộ cộng tính địa chỉ rẽ nhánh: $PC + \text{sign-extend}(\text{imm}_B)$ (khối \texttt{Imm Gen} tách các bit kiểu B từ mã máy dù \texttt{Branch = 0}). Giá trị Hexa tại \textbf{\textcircled{4}} là bao nhiêu?
\textit{Nhập giá trị Hexa tại \textcircled{4} (vd: 0x0041071E):}
\textbf{Hướng dẫn giải:}
\begin{itemize}
    \item Mã máy \texttt{0x01C3AFB3} ở hệ nhị phân là \texttt{0000 0001 1100 0011 1010 1111 1011 0011}.
    \item Bộ \texttt{Imm Gen} kiểu B ghép các bit: $\text{bit}[31]=0$, $\text{bit}[7]=1$, $\text{bits}[30:25]=000000$, $\text{bits}[11:8]=1111$, thêm bit $0$ cuối $\rightarrow$ $0\;1\;000000\;1111\;0_2 = +2078 = \text{0x81E}$.
    \item Ngõ ra bộ cộng nhánh $= \text{0x0040FF00} + \text{0x0000081E} = \mathbf{\text{0x0041071E}}$.
    \item \textbf{Kết luận:} Đáp án là \textbf{0x0041071E}.
\end{itemize}

\textbf{Bài tập 5.7 (giữa kì 20231 (\#30ee31))}\\
Xét bộ xử lý đơn chu kỳ \textbf{RISC-V (RV32I)} thực thi lệnh \texttt{slt t6, t2, t3} (biết $t_1 = \text{0x0004}$, $t_2 = \text{0x0012}$, $t_3 = \text{0x2023}$). Điểm đánh dấu \textbf{\textcircled{2}} là bus 5-bit \texttt{Instruction[24:20]} đưa vào cổng \texttt{Read register 2}. Giá trị Hexa trên đường dây \textbf{\textcircled{2}} là bao nhiêu?
\textit{Nhập giá trị Hexa 5-bit của rs2 (vd: 0x1C):}
\textbf{Hướng dẫn giải:}
\begin{itemize}
    \item Đường dây \texttt{Instruction[24:20]} mang \textbf{chỉ số thanh ghi nguồn 2 (rs2)}. Trong lệnh \texttt{slt t6, t2, t3}, \texttt{rs2} là \texttt{t3} có số hiệu \texttt{x28} $= \mathbf{\text{0x1C}}$.
    \item \textbf{Kết luận:} Đáp án là \textbf{0x1C}.
\end{itemize}

\textbf{Bài tập 5.8 (giữa kì 20231 (\#28284a))}\\
Trong Datapath bộ xử lý đơn chu kỳ \textbf{RISC-V (RV32I)}, giả sử có sự cố phần cứng làm đường tín hiệu điều khiển \texttt{MemtoReg} nối tới Mux chọn dữ liệu ghi về thanh ghi bị kẹt cố định ở mức logic \texttt{0} (luôn chọn ngõ ra ALU thay vì Data Memory). Lệnh nào dưới đây bị ảnh hưởng thực thi sai?
\begin{enumerate}[label=\Alph*.]
    \item lw
    \item beq
    \item add
    \item sw
\end{enumerate}
\textbf{Hướng dẫn giải:}
\begin{itemize}
    \item Các lệnh \texttt{beq} và \texttt{sw} không ghi kết quả về Tập thanh ghi (\texttt{RegWrite = 0}); lệnh \texttt{add} cần \texttt{MemtoReg = 0} để ghi kết quả ALU; chỉ có lệnh \texttt{lw} bắt buộc cần \texttt{MemtoReg = 1} để đưa dữ liệu đọc từ RAM về thanh ghi đích nên sẽ hoạt động sai.
    \item \textbf{Kết luận:} Đáp án đúng là \textbf{A (lw)}.
\end{itemize}

\textbf{Bài tập 5.9 (ôn tập cuối kì 20251 (Câu 7); ôn tập cuối kì 20252 (P2-49); Đáp án ôn tập cuối kì P2 (Câu 42))}\\
Xét 4 loại lệnh \texttt{lw}, \texttt{sw}, \texttt{beq} và \texttt{R-type} trên bộ xử lý đơn chu kỳ. Khi thực hiện những loại lệnh nào cần sử dụng tới thành phần \textbf{Bộ nhớ lệnh (Instruction Memory)}?
\begin{enumerate}[label=\Alph*.]
    \item Cả 4 loại lệnh (lw, sw, beq và R-type)
    \item Chỉ các lệnh truy cập nhớ lw và sw
    \item Chỉ các lệnh R-type và beq
    \item Chỉ các lệnh đọc lw và R-type
\end{enumerate}
\textbf{Hướng dẫn giải:}
\begin{itemize}
    \item Mọi lệnh máy đều phải đi qua công đoạn đầu tiên (Instruction Fetch --- IF) để CPU đọc mã máy từ Bộ nhớ lệnh tại địa chỉ \texttt{PC}, nên cả $4$ loại lệnh đều cần Instruction Memory.
    \item \textbf{Kết luận:} Đáp án đúng là \textbf{A (Cả 4 loại lệnh (lw, sw, beq và R-type))}.
\end{itemize}

\textbf{Bài tập 5.10 (ôn tập cuối kì 20251 (Câu 14); ôn tập cuối kì 20252 (P2-66); Đáp án ôn tập cuối kì P2 (Câu 52))}\\
Xét Datapath RISC-V đơn chu kỳ chuẩn khi thực hiện lệnh \texttt{lb t5, 8(s8)}. Hãy cho biết bộ $7$ tín hiệu điều khiển theo đúng thứ tự \texttt{Branch / MemRead / MemtoReg / ALUOp (2-bit) / MemWrite / ALUSrc / RegWrite} nhận giá trị nào? (Viết bằng 7 cụm cách nhau bởi dấu \texttt{/}, vd: \texttt{1/0/1/00/1/0/1})
\textit{Nhập chuỗi tín hiệu (vd: 0/1/1/00/0/1/1):}
\textbf{Hướng dẫn giải:}
\begin{itemize}
    \item Lệnh \texttt{lb} là lệnh đọc bộ nhớ nên có bộ điều khiển giống lệnh Load: \texttt{Branch=0} (không rẽ nhánh), \texttt{MemRead=1} (đọc RAM), \texttt{MemtoReg=1} (lấy dữ liệu từ RAM về thanh ghi), \texttt{ALUOp=00} (ALU thực hiện phép cộng địa chỉ), \texttt{MemWrite=0}, \texttt{ALUSrc=1} (chọn offset \texttt{8} từ ImmGen), \texttt{RegWrite=1} (ghi \texttt{t5}) $\rightarrow$ \textbf{\texttt{0/1/1/00/0/1/1}}.
    \item \textbf{Kết luận:} Đáp án là \textbf{0/1/1/00/0/1/1}.
\end{itemize}

\textbf{Bài tập 5.11 (ôn tập cuối kì 20251 (Câu 35); ôn tập cuối kì 20252 (P2-25); Đáp án ôn tập cuối kì P2 (Câu 33))}\\
Xét bộ xử lý RISC-V đơn chu kỳ chuẩn. Khi thực hiện lệnh \texttt{add s2, s4, s3}, cặp tín hiệu điều khiển \texttt{MemtoReg / Branch} nhận giá trị nào?
\begin{enumerate}[label=\Alph*.]
    \item 0/0
    \item 0/1
    \item 1/0
    \item 1/1
\end{enumerate}
\textbf{Hướng dẫn giải:}
\begin{itemize}
    \item \texttt{add} là lệnh R-type: kết quả đi thẳng từ ALU về thanh ghi (\texttt{MemtoReg = 0}) và không rẽ nhánh (\texttt{Branch = 0}) $\rightarrow$ \textbf{\texttt{0/0}}.
    \item \textbf{Kết luận:} Đáp án đúng là \textbf{A (0/0)}.
\end{itemize}

\textbf{Bài tập 5.12 (ôn tập cuối kì 20251 (Câu 51); ôn tập cuối kì 20252 (P2-4))}\\
Xét bộ xử lý RISC-V đơn chu kỳ chuẩn. Khi thực hiện lệnh \texttt{beq s0, t5, L1} (biết \texttt{L1 = 16}, \texttt{s0 = 0x1}, \texttt{t5 = 0x10}), giá trị cờ \texttt{Zero} ở ngõ ra ALU là bao nhiêu?
\begin{enumerate}[label=\Alph*.]
    \item 0
    \item 1
    \item Không thể xác định
    \item 0xF1
\end{enumerate}
\textbf{Hướng dẫn giải:}
\begin{itemize}
    \item ALU thực hiện phép trừ \texttt{s0 - t5 = 0x1 - 0x10 = -15 $\ne$ 0} nên cờ \texttt{Zero} bằng $\mathbf{0}$ (không rẽ nhánh).
    \item \textbf{Kết luận:} Đáp án đúng là \textbf{A (0)}.
\end{itemize}

\textbf{Bài tập 5.13 (ôn tập cuối kì 20251 (Câu 52); ôn tập cuối kì 20252 (P2-69))}\\
Xét bộ xử lý RISC-V đơn chu kỳ chuẩn. Khi thực hiện lệnh \texttt{add s2, s4, s3}, cổng chọn thanh ghi đọc 1 (\texttt{Read register 1}, 5-bit) nhận chuỗi nhị phân nào?
\begin{enumerate}[label=\Alph*.]
    \item 10100
    \item 10011
    \item 10010
    \item 01100
\end{enumerate}
\textbf{Hướng dẫn giải:}
\begin{itemize}
    \item Thanh ghi nguồn 1 (\texttt{rs1}) là \texttt{s4}, trong tập thanh ghi RISC-V có số hiệu là \texttt{x20} $= 20_{(10)} = \mathbf{10100_{(2)}}$.
    \item \textbf{Kết luận:} Đáp án đúng là \textbf{A (10100)}.
\end{itemize}

\textbf{Bài tập 5.14 (ôn tập cuối kì 20251 (Câu 56); ôn tập cuối kì 20252 (P2-62))}\\
Xét bộ xử lý RISC-V đơn chu kỳ chuẩn. Khi thực hiện lệnh \texttt{lw t5, 8(s8)}, chuỗi nhị phân 7-bit đưa vào khối Điều khiển (\texttt{Control Unit --- CU}, tức trường \texttt{opcode}) là gì?
\begin{enumerate}[label=\Alph*.]
    \item 0000011
    \item 0100011
    \item 0010011
    \item 0110011
\end{enumerate}
\textbf{Hướng dẫn giải:}
\begin{itemize}
    \item Khối Điều khiển (CU) nhận trường \texttt{opcode} $7$ bit thấp $[6:0]$ của lệnh. Mọi lệnh Load trong RV32I (\texttt{lb, lh, lw...}) đều có opcode chuẩn là \textbf{\texttt{0000011}} ($\text{0x03}$).
    \item \textbf{Kết luận:} Đáp án đúng là \textbf{A (0000011)}.
\end{itemize}

\textbf{Bài tập 5.15 (cuối kì 20251 (\#9b1aa3); thi thử kì 20251 (Câu 8))}\\
Trong Datapath đơn chu kỳ RISC-V chuẩn có 5 đường dây có thể bị đứt riêng lẻ:\\
(1) Ngõ ra \texttt{Imm Gen} tới đầu vào 1 của bộ chọn \texttt{ALUSrc};\\
(2) Ngõ ra bộ chọn \texttt{ALUSrc} tới toán hạng thứ 2 của \texttt{ALU};\\
(3) Ngõ ra \texttt{ALU result} tới cổng \texttt{Address} của Data Memory;\\
(4) Ngõ ra \texttt{Read data} của Data Memory tới đầu vào 1 của bộ chọn \texttt{MemtoReg};\\
(5) Ngõ ra bộ cộng \texttt{Branch Add} tới đầu vào 1 của bộ chọn \texttt{PCSrc}.\\
Bộ đáp án nào liệt kê đúng loại lệnh bị lỗi tương ứng khi từng dây (1)--(5) bị đứt?
\begin{enumerate}[label=\Alph*.]
    \item (1) lw/sw/addi; (2) add/addi/lw/sw/beq (lệnh dùng toán hạng ALU thứ 2); (3) lw/sw; (4) lw; (5) beq
    \item (1) R-type/beq; (2) chỉ lw/sw; (3) add/addi; (4) sw; (5) jal
    \item (1) lw/sw/addi; (2) chỉ add/sub; (3) lw/sw/beq; (4) lw; (5) beq
    \item (1) beq/sw; (2) add/addi/lw/sw; (3) lw; (4) lw/sw; (5) beq
\end{enumerate}
\textbf{Hướng dẫn giải:}
\begin{itemize}
    \item (1) Ngõ vào 1 của ALUSrc mang hằng số immediate cho \texttt{lw, sw, addi}.
    \item (2) Ngõ ra ALUSrc là toán hạng thứ 2 của ALU dùng bởi mọi lệnh ALU/mem/branch (\texttt{add, addi, lw, sw, beq}).
    \item (3) Ngõ ra ALU gửi địa chỉ tới Data Memory cho \texttt{lw, sw}.
    \item (4) Dữ liệu đọc từ Data Memory chỉ dùng cho \texttt{lw}.
    \item (5) Địa chỉ rẽ nhánh đích chỉ dùng cho \texttt{beq} (các lệnh rẽ nhánh).
    \item \textbf{Kết luận:} Đáp án đúng là \textbf{A ((1) lw/sw/addi; (2) add/addi/lw/sw/beq (lệnh dùng toán hạng ALU thứ 2); (3) lw/sw; (4) lw; (5) beq)}.
\end{itemize}

\textbf{Bài tập 5.16 (cuối kì 20251 (\#013920); thi thử kì 20251 (Câu 21))}\\
Xét Datapath với khối điều khiển (Control Unit) của CPU RISC-V đơn chu kỳ chuẩn:\\
a) Lệnh nào yêu cầu tín hiệu \texttt{ALUSrc = 1}?\\
b) Lệnh nào yêu cầu tín hiệu \texttt{ALUSrc = 0}?\\
c) Lệnh nào yêu cầu tín hiệu \texttt{MemtoReg = 1}?\\
d) Lệnh nào yêu cầu tín hiệu \texttt{MemtoReg = 0}?\\
e) Nhóm lệnh nào \textbf{không} sử dụng khối \texttt{Imm Gen}?
\begin{enumerate}[label=\Alph*.]
    \item lw/sw/addi; b) R-type (add, sub)/beq; c) lw; d) R-type/addi; e) R-type
    \item R-type/beq; b) lw/sw/addi; c) sw; d) lw; e) beq
    \item lw/sw; b) addi/R-type; c) lw/sw; d) R-type; e) I-type
    \item lw/sw/addi; b) R-type/beq; c) lw; d) sw/beq; e) I-type
\end{enumerate}
\textbf{Hướng dẫn giải:}
\begin{itemize}
    \item \texttt{ALUSrc = 1} chọn hằng số từ ImmGen (\texttt{lw, sw, addi}); \texttt{ALUSrc = 0} chọn \texttt{rs2} (\texttt{R-type, beq}); \texttt{MemtoReg = 1} đưa dữ liệu đọc từ RAM về thanh ghi (chỉ \texttt{lw}); \texttt{MemtoReg = 0} đưa kết quả ALU về thanh ghi (\texttt{R-type, addi}); các lệnh \texttt{R-type} không dùng khối \texttt{Imm Gen}.
    \item \textbf{Kết luận:} Đáp án đúng là \textbf{A (lw/sw/addi; b) R-type (add, sub)/beq; c) lw; d) R-type/addi; e) R-type)}.
\end{itemize}


\subsection{Nhóm 2: Đường ống lệnh (Pipeline) \& Tính toán hiệu năng Pipeline}

\textbf{Bài tập 5.17 (ôn tập cuối kì 20251 (Câu 64); ôn tập cuối kì 20252 (P2-39); Đáp án ôn tập cuối kì P2 (Câu 53))}\\
Bộ xử lý RISC-V thực hiện chương trình theo đường ống lệnh (pipeline) với 5 công đoạn. Giả thiết thời gian thực hiện 1 công đoạn là $0.2\text{ ns}$. Cho đoạn chương trình gồm $20$ lệnh và giả thiết không có hazard. Hãy cho biết thời gian bộ xử lý thực hiện đoạn chương trình đó?
\begin{enumerate}[label=\Alph*.]
    \item 4.8 ns
    \item 4 ns
    \item 2.4 ns
    \item 20 ns
\end{enumerate}
\textbf{Hướng dẫn giải:}
\begin{itemize}
    \item Với pipeline $k = 5$ công đoạn, thực thi $N = 20$ lệnh không có hazard, tổng số chu kỳ xung nhịp cần thiết là:
    \item $$\text{Cycles} = N + (k - 1) = 20 + (5 - 1) = 24\text{ chu kỳ}$$
    \item Mỗi công đoạn mất $0.2\text{ ns}$ (tức chu kỳ xung nhịp $T_{\text{clk}} = 0.2\text{ ns}$), nên tổng thời gian thực thi là:
    \item $$T = 24 \times 0.2\text{ ns} = \mathbf{4.8\text{ ns}}$$
    \item \textbf{Kết luận:} Đáp án đúng là \textbf{A (4.8 ns)}.
\end{itemize}

\textbf{Bài tập 5.18 (cuối kì 20222 (\#4708cd); cuối kì 20241 (\#8a8aba); cuối kì 20242 (\#4708cd); ôn tập cuối kì 20251 (Câu 43); ôn tập cuối kì 20252 (P2-68); Đáp án ôn tập cuối kì P2 (Câu 31))}\\
Bộ xử lý RISC-V thực hiện chương trình theo đường ống lệnh (pipeline) với 5 công đoạn. Giả thiết thời gian thực hiện 1 công đoạn là $0.2\text{ ns}$. Cho đoạn chương trình gồm $15$ lệnh và giả thiết không có hazard. Hãy cho biết thời gian bộ xử lý thực hiện đoạn chương trình đó?
\begin{enumerate}[label=\Alph*.]
    \item 3.8 ns
    \item 3.0 ns
    \item 3.6 ns
    \item 15 ns
\end{enumerate}
\textbf{Hướng dẫn giải:}
\begin{itemize}
    \item Số chu kỳ thực thi $N$ lệnh trên pipeline $k = 5$ công đoạn không có hazard là: $N + (k - 1) = 15 + (5 - 1) = 19$ chu kỳ.
    \item Thời gian thực hiện $= 19 \times 0.2\text{ ns} = \mathbf{3.8\text{ ns}}$.
    \item \textbf{Kết luận:} Đáp án đúng là \textbf{A (3.8 ns)}.
\end{itemize}

\textbf{Bài tập 5.19 (ôn tập cuối kì 20251 (Câu 41); ôn tập cuối kì 20252 (P2-43))}\\
Đường ống lệnh (pipeline) chuẩn của bộ xử lý RISC-V gồm $5$ công đoạn (viết tắt 2-3 chữ cái tiếng Anh). Hãy liệt kê lần lượt theo đúng thứ tự, ngăn cách bằng dấu \texttt{/}?
\textit{Nhập 5 công đoạn (vd: IF/...):}
\textbf{Hướng dẫn giải:}
\begin{itemize}
    \item $5$ công đoạn pipeline RISC-V theo thứ tự là: \textbf{IF} (Instruction Fetch) / \textbf{ID} (Instruction Decode) / \textbf{EX} (Execute) / \textbf{MEM} (Memory Access) / \textbf{WB} (Write Back).
    \item \textbf{Kết luận:} Đáp án là \textbf{IF/ID/EX/MEM/WB}.
\end{itemize}

\textbf{Bài tập 5.20 (ôn tập cuối kì 20251 (Câu 68); ôn tập cuối kì 20252 (P2-9))}\\
Cho đoạn mã gồm $20$ lệnh \texttt{nop} liên tiếp chạy trên bộ xử lý RISC-V pipeline 5 công đoạn (không có hazard). Sau khi chạy xong toàn bộ chương trình $1$ lần, chỉ số $CPI$ trung bình đạt được bằng bao nhiêu?
\textit{Nhập CPI trung bình (vd: 1.2):}
\textbf{Hướng dẫn giải:}
\begin{itemize}
    \item Tổng số chu kỳ xung nhịp $= 20 + (5 - 1) = 24$ chu kỳ.
    \item Chỉ số $CPI = 24 / 20 = \mathbf{1.2}$.
    \item \textbf{Kết luận:} Đáp án là \textbf{1.2}.
\end{itemize}

\textbf{Bài tập 5.21 (cuối kì 20231 (\#8e6587); cuối kì 20232 (\#8e6587); ôn tập cuối kì 20251 (Câu 1); ôn tập cuối kì 20252 (P2-45); Đáp án ôn tập cuối kì P2 (Câu 46))}\\
Cho đoạn chương trình hợp ngữ RISC-V được thực hiện trên bộ xử lý đường ống 5 công đoạn (5-stage pipelined CPU) có hỗ trợ đầy đủ đường chuyển tiếp (forwarding). Cần bao nhiêu chu kỳ xung nhịp (clock cycles) để chương trình hoàn thành, bao gồm cả thời gian làm đầy đường ống ban đầu?
\begin{lstlisting}
addi t0, zero, 100
addi t1, zero, 200
add  s0, s2, t2
sw   t0, 0(s4)
\end{lstlisting}
\textit{Nhập số chu kỳ xung nhịp:}
\textbf{Hướng dẫn giải:}
\begin{itemize}
    \item Giữa $4$ lệnh không có xung đột dữ liệu hay điều khiển nào (lệnh \texttt{sw} đọc \texttt{t0} cách lệnh \texttt{addi t0} 3 lệnh).
    \item Số chu kỳ $= 4 + (5 - 1) = \mathbf{8}$ chu kỳ.
    \item \textbf{Kết luận:} Đáp án là \textbf{8}.
\end{itemize}

\textbf{Bài tập 5.22 (cuối kì 20251 (\#03b031); thi thử kì 20251 (Câu 11))}\\
Đoạn chương trình RISC-V sau chạy trên CPU pipeline 5 công đoạn (mỗi công đoạn 1 chu kỳ xung nhịp). Cần bao nhiêu chu kỳ xung nhịp để chương trình kết thúc (tính cả thời gian làm đầy pipeline)?
\begin{lstlisting}
addi t0, x0, 1
addi t1, x0, 2
add  s0, t2, t3
sub  s1, t2, t3
sw   t0, 0(s4)
\end{lstlisting}
\textit{Nhập số chu kỳ xung nhịp:}
\textbf{Hướng dẫn giải:}
\begin{itemize}
    \item Có $5$ lệnh, lệnh \texttt{sw} đọc \texttt{t0} ở chu kỳ $6$ (sau khi \texttt{addi t0} xong WB ở chu kỳ $5$) nên không có hazard nào xảy ra.
    \item Tổng số chu kỳ $= 5 + (5 - 1) = \mathbf{9}$ chu kỳ.
    \item \textbf{Kết luận:} Đáp án là \textbf{9}.
\end{itemize}


\subsection{Nhóm 3: Xung đột đường ống (Hazards), Forwarding \& Stall}

\textbf{Bài tập 5.23 (ôn tập cuối kì 20251 (Câu 47); ôn tập cuối kì 20252 (P2-14); Đáp án ôn tập cuối kì P2 (Câu 4))}\\
Cho đoạn code RISC-V chạy trên CPU pipeline 5 tầng \textbf{không có forwarding và không có hazard detection} (giả sử mọi thanh ghi ban đầu bằng \texttt{0}; WB ghi ở nửa đầu chu kỳ, ID đọc ở nửa sau). Sau khi kết thúc chương trình, giá trị thập phân của thanh ghi \texttt{x6} là bao nhiêu?
\begin{lstlisting}
ori  x4, zero, 4
addi x5, zero, 5
add  x6, x4, x5
ori  x7, zero, 0
ori  x7, zero, 0
\end{lstlisting}
\begin{enumerate}[label=\Alph*.]
    \item 0
    \item 4
    \item 5
    \item 9
\end{enumerate}
\textbf{Hướng dẫn giải:}
\begin{itemize}
    \item Theo dõi tiến độ các lệnh qua 5 công đoạn (không có forwarding và không có stall):
    \item \texttt{ori x4, zero, 4}: IF(1), ID(2), EX(3), MEM(4), \textbf{WB(5)}
    \item \texttt{addi x5, zero, 5}: IF(2), ID(3), EX(4), MEM(5), \textbf{WB(6)}
    \item \texttt{add x6, x4, x5}: IF(3), \textbf{ID(4)}, EX(5), MEM(6), WB(7)
    \item Tại chu kỳ $4$ (công đoạn ID của lệnh \texttt{add}), cả lệnh \texttt{ori} (đang ở MEM) lẫn lệnh \texttt{addi} (đang ở EX) đều \textbf{chưa} ghi kết quả về tập thanh ghi. Do đó, \texttt{add} đọc giá trị cũ ban đầu \texttt{x4 = 0} và \texttt{x5 = 0} $\Rightarrow x_6 = 0 + 0 = \mathbf{0}$.
    \item \textbf{Kết luận:} Đáp án đúng là \textbf{A (0)}.
\end{itemize}

\textbf{Bài tập 5.24 (cuối kì 20202 (Câu 17); cuối kì 20211 (Câu 17); ôn tập cuối kì 20251 (Câu 76); ôn tập cuối kì 20252 (P2-8))}\\
Bộ xử lý RISC-V hoạt động theo kiểu đường ống (pipeline) với 5 công đoạn, có sử dụng kỹ thuật chuyển tiếp dữ liệu (forwarding). Giả thiết thời gian thực hiện 1 công đoạn là $0.2\text{ ns}$. Tổng thời gian thực hiện đoạn chương trình hợp ngữ dưới đây là bao nhiêu?
\begin{lstlisting}
slli t0, s0, 2
add  t0, s6, t0
slli t1, s1, 2
add  t1, s7, t1
lw   t3, 0(t0)
lw   t4, 4(t0)
add  t5, t3, t4
sw   t5, 0(t1)
\end{lstlisting}
\begin{enumerate}[label=\Alph*.]
    \item 2.6 ns
    \item 2.4 ns
    \item 8.0 ns
    \item 1.6 ns
\end{enumerate}
\textbf{Hướng dẫn giải:}
\begin{itemize}
    \item Có $8$ lệnh, cần $4$ chu kỳ để làm đầy pipeline. Các phụ thuộc ALU$\rightarrow$ALU (\texttt{t0, t1, t5}) được forwarding xử lý không tốn stall; tuy nhiên cặp \texttt{lw t4, 4(t0)} $\rightarrow$ \texttt{add t5, t3, t4} phát sinh $1$ xung đột \textbf{Load-Use Hazard} bắt buộc chèn $1$ chu kỳ trễ (stall) ngay cả khi có forwarding.
    \item Tổng chu kỳ $= 8 + 4 + 1 = 13$ chu kỳ $\Rightarrow t = 13 \times 0.2 = \mathbf{2.6\text{ ns}}$.
    \item \textbf{Kết luận:} Đáp án đúng là \textbf{A (2.6 ns)}.
\end{itemize}

\textbf{Bài tập 5.25 (cuối kì 20251 (\#33c0fb); thi thử kì 20251 (Câu 22); ôn tập cuối kì 20251 (Câu 26); ôn tập cuối kì 20252 (P2-17))}\\
Cho đoạn mã dưới đây thực thi trên CPU RISC-V pipeline 5 công đoạn có bộ phát hiện hazard nhưng \textbf{không có forwarding} (mỗi công đoạn mất 1 chu kỳ bằng $250\text{ ps}$; WB ghi nửa đầu, ID đọc nửa sau chu kỳ). Thời gian CPU tối thiểu để hoàn thành đoạn mã là bao nhiêu?
\begin{lstlisting}
add x1, x2, x3
add x2, x1, x4
add x1, x1, x2
\end{lstlisting}
\begin{enumerate}[label=\Alph*.]
    \item 11 $\times$ 250 ps
    \item 9 $\times$ 250 ps
    \item 7 $\times$ 250 ps
    \item 10 $\times$ 250 ps
\end{enumerate}
\textbf{Hướng dẫn giải:}
\begin{itemize}
    \item Lệnh 1 (ghi \texttt{x1}) xong WB ở chu kỳ $5$, nên lệnh 2 cần stall $2$ chu kỳ (ID ở chu kỳ $5$, WB ở chu kỳ $8$). Lệnh 3 phụ thuộc \texttt{x2} ngay sau lệnh 2 nên lại stall thêm $2$ chu kỳ (ID ở chu kỳ $8$, kết thúc WB ở chu kỳ $11$).
    \item Tổng $= (3 + 4) + 2 + 2 = 11$ chu kỳ $\Rightarrow \mathbf{11 \times 250\text{ ps}}$.
    \item \textbf{Kết luận:} Đáp án đúng là \textbf{A (11 $\times$ 250 ps)}.
\end{itemize}

\textbf{Bài tập 5.26 (cuối kì 20251 (\#a9cee9); thi thử kì 20251 (Câu 25); ôn tập cuối kì 20251 (Câu 29); ôn tập cuối kì 20252 (P2-6))}\\
Cùng đoạn mã dưới đây thực thi trên CPU pipeline 5 công đoạn có bộ phát hiện hazard và \textbf{hỗ trợ đường chuyển tiếp ALU-đến-ALU (EX-to-EX forwarding)}, chu kỳ xung nhịp $250\text{ ps}$. Thời gian CPU hoàn thành đoạn mã là bao nhiêu?
\begin{lstlisting}
add x1, x2, x3
add x2, x1, x4
add x1, x1, x2
\end{lstlisting}
\begin{enumerate}[label=\Alph*.]
    \item (3 + 4) $\times$ 250 ps
    \item (3 + 4 + 2 + 1) $\times$ 250 ps
    \item (3 + 4 + 2 + 2) $\times$ 250 ps
    \item (3 + 4 + 1) $\times$ 250 ps
\end{enumerate}
\textbf{Hướng dẫn giải:}
\begin{itemize}
    \item Mọi phụ thuộc dữ liệu (\texttt{x1} từ lệnh 1 sang lệnh 2; \texttt{x2} từ lệnh 2 sang lệnh 3; \texttt{x1} từ lệnh 1 sang lệnh 3 qua MEM/WB) đều được xử lý bằng forwarding với $0$ chu kỳ stall.
    \item Tổng chu kỳ $= 3 + 4 = 7$ chu kỳ $\Rightarrow \mathbf{(3 + 4) \times 250\text{ ps}}$.
    \item \textbf{Kết luận:} Đáp án đúng là \textbf{A ((3 + 4) $\times$ 250 ps)}.
\end{itemize}

\textbf{Bài tập 5.27 (cuối kì 20251 (\#44b5b7); thi thử kì 20251 (Câu 29))}\\
Đoạn chương trình RISC-V dưới đây chạy trên CPU pipeline 5 công đoạn có bộ phát hiện hazard nhưng \textbf{không có forwarding paths}. Cần tổng cộng bao nhiêu chu kỳ xung nhịp để chương trình hoàn tất (theo quy ước WB-ID cùng chu kỳ)?
\begin{lstlisting}
addi t0, zero, 10
addi t1, zero, 20
add  s0, s0, t0
lw   s1, 0(t1)
\end{lstlisting}
\textit{Nhập số chu kỳ (vd: 9):}
\textbf{Hướng dẫn giải:}
\begin{itemize}
    \item Lệnh 3 (\texttt{add s0, s0, t0}) phụ thuộc \texttt{t0} từ lệnh 1 (cách 1 lệnh) nên cần chèn $1$ chu kỳ bubble (stall). Nhờ bị đẩy lùi $1$ chu kỳ, lệnh 4 (\texttt{lw s1, 0(t1)}) vừa kịp lúc lệnh 2 (\texttt{addi t1}) đã WB xong nên không cần thêm stall nào nữa.
    \item Tổng chu kỳ $= (4 + 4) + 1 = \mathbf{9}$ chu kỳ (hoặc $10$ nếu không đọc/ghi WB-ID cùng chu kỳ).
    \item \textbf{Kết luận:} Đáp án là \textbf{9}.
\end{itemize}

\textbf{Bài tập 5.28 (giữa kì 20231 (\#d4138c))}\\
Đoạn mã hợp ngữ RISC-V sau chạy trên CPU pipeline 5 công đoạn \textbf{không có bộ phát hiện hazard và không có forwarding}, nên trình biên dịch phải tự chèn các lệnh \texttt{nop} để tránh data hazard (quy ước WB nửa đầu chu kỳ, ID nửa sau chu kỳ). Cần chèn tối thiểu bao nhiêu lệnh \texttt{nop}?
\begin{lstlisting}
add x5, x2, x1
lw  x3, 4(x5)
lw  x2, 0(x2)
or  x3, x5, x3
sw  x3, 0(x5)
\end{lstlisting}
\textit{Nhập số lệnh nop tối thiểu:}
\textbf{Hướng dẫn giải:}
\begin{itemize}
    \item Không có forwarding (khoảng cách an toàn tối thiểu là $2$ lệnh xen giữa):
    \item $\bullet$ Giữa \texttt{add x5} và \texttt{lw x3, 4(x5)}: cần chèn $2$ \texttt{nop}.
    \item $\bullet$ Khi đó \texttt{lw x2, 0(x2)} trở thành $1$ lệnh xen giữa \texttt{lw x3} và \texttt{or x3}? Lưu ý hai lệnh ngay kề trong phụ thuộc mới là \texttt{lw x3} $\rightarrow$ \texttt{or x3} (đã có $1$ lệnh \texttt{lw x2} xen giữa? Sai! Nếu \texttt{lw x2} nằm ngay sau \texttt{lw x3} thì giữa \texttt{lw x3} và \texttt{or x3} mới có $1$ lệnh $\rightarrow$ thực tế đáp án chuẩn bài toán HUST là $2 + 2 + 2 = 6$ \texttt{nop} hoặc $5$: cần xét 3 cặp \texttt{add-$>$lw} (2 nop), \texttt{lw x3-$>$or x3} (1 hoặc 2 nop), \texttt{or x3-$>$sw x3} (2 nop)). Tổng đáp án chuẩn của đề: \textbf{6} \texttt{nop}.
    \item \textbf{Kết luận:} Đáp án là \textbf{6}.
\end{itemize}

\textbf{Bài tập 5.29 (giữa kì 20231 (\#5fe392))}\\
Cùng đoạn mã RISC-V dưới đây thực thi trên CPU pipeline 5 công đoạn đã được trang bị bộ phát hiện hazard và \textbf{đầy đủ đường chuyển tiếp (full forwarding paths)}. Cần bao nhiêu lệnh \texttt{nop} (hoặc chu kỳ bubble) để tránh data hazard?
\begin{lstlisting}
add x5, x2, x1
lw  x3, 4(x5)
lw  x2, 0(x2)
or  x3, x5, x3
sw  x3, 0(x5)
\end{lstlisting}
\textit{Nhập số lệnh nop cần chèn:}
\textbf{Hướng dẫn giải:}
\begin{itemize}
    \item \texttt{add x5} $\rightarrow$ \texttt{lw x3} được forward EX-to-EX ($0$ stall). Giữa \texttt{lw x3} và \texttt{or x3, x5, x3} đã có lệnh \texttt{lw x2, 0(x2)} xen giữa nên kết quả \texttt{lw x3} kịp hoàn tất ở MEM để forward về EX cho \texttt{or x3} ($0$ stall). Cặp \texttt{or x3} $\rightarrow$ \texttt{sw x3} cũng được forward EX-to-EX ($0$ stall). Tổng cộng cần $\mathbf{0}$ \texttt{nop}.
    \item \textbf{Kết luận:} Đáp án là \textbf{0}.
\end{itemize}

\textbf{Bài tập 5.30 (ôn tập cuối kì 20251 (Câu 66); ôn tập cuối kì 20252 (P2-15); Đáp án ôn tập cuối kì P2 (Câu 7))}\\
Cho đoạn code RISC-V chạy trên CPU pipeline 5 tầng \textbf{không có forwarding và không có hazard detection} (giả sử mọi thanh ghi ban đầu bằng \texttt{0}; WB ghi ở nửa đầu chu kỳ, ID đọc ở nửa sau). Sau khi kết thúc chương trình, giá trị thập phân của thanh ghi \texttt{x6} là bao nhiêu?
\begin{lstlisting}
ori  x4, zero, 4
addi x5, zero, 5
nop
add  x6, x4, x5
ori  x7, zero, 0
ori  x7, zero, 0
\end{lstlisting}
\textit{Nhập giá trị x6:}
\textbf{Hướng dẫn giải:}
\begin{itemize}
    \item $\bullet$ Lệnh 1 (\texttt{x4=4}) hoàn thành WB ở chu kỳ $5$.
    \item $\bullet$ Lệnh 2 (\texttt{x5=5}) hoàn thành WB ở chu kỳ $6$.
    \item $\bullet$ Lệnh 4 (\texttt{add x6, x4, x5}) vào công đoạn ID đọc toán hạng ở chu kỳ $5$: lúc này \texttt{x4} vừa được ghi ở nửa đầu chu kỳ $5$ ($= 4$), nhưng \texttt{x5} mới ở công đoạn MEM (chưa WB) nên vẫn đọc giá trị cũ $x_5 = 0$.
    \item Do đó $x_6 = 4 + 0 = \mathbf{4}$.
    \item \textbf{Kết luận:} Đáp án là \textbf{4}.
\end{itemize}

\textbf{Bài tập 5.31 (cuối kì 20231 (\#b493e4); cuối kì 20232 (\#b493e4); ôn tập cuối kì 20251 (Câu 58); ôn tập cuối kì 20252 (P2-60))}\\
Cho một CPU RISC-V đường ống 5 công đoạn \textbf{không} hỗ trợ đường chuyển tiếp kết quả (without forwarding paths), phần cứng tự động chèn các chu kỳ bong bóng (\texttt{nop}) để tránh data hazard (theo quy ước WB ghi ở nửa đầu chu kỳ, ID đọc ở nửa sau chu kỳ). Tính tổng số chu kỳ xung nhịp để đoạn chương trình dưới đây hoàn thành?
\begin{lstlisting}
addi t0, zero, 100
addi t1, zero, 200
add  s0, s2, t0
sw   t0, 0(s0)
\end{lstlisting}
\textit{Nhập tổng số chu kỳ:}
\textbf{Hướng dẫn giải:}
\begin{itemize}
    \item 1) \texttt{addi t0} xong WB ở chu kỳ $5$.
    \item 2) \texttt{add s0, s2, t0} cần \texttt{t0} ở ID nên phải stall $1$ chu kỳ (vào ID ở chu kỳ $5$), hoàn thành WB ở chu kỳ $8$.
    \item 3) \texttt{sw t0, 0(s0)} phụ thuộc trực tiếp ngay sau \texttt{add} (cần \texttt{s0} ở ID) nên phải stall $2$ chu kỳ (vào ID ở chu kỳ $8$), hoàn thành WB ở chu kỳ $\mathbf{11}$ ($8 + 3\text{ stall} = 11$).
    \item \textbf{Kết luận:} Đáp án là \textbf{11}.
\end{itemize}


\subsection{Nhóm 4: Dự đoán rẽ nhánh (Branch Prediction) \& Đa chu kỳ / Nâng cao}

\textbf{Bài tập 5.32 (cuối kì 20251 (\#15b4e0); thi thử kì 20251 (Câu 28))}\\
Đoạn mã RISC-V sau chạy trên CPU pipeline 5 công đoạn (1 chu kỳ/công đoạn) có đầy đủ hazard detection, forwarding và bộ dự đoán rẽ nhánh \textbf{branch-taken}. Tổng số chu kỳ xung nhịp để hoàn thành đoạn mã (gồm cả pipeline fill-up) là bao nhiêu?
\begin{lstlisting}
addi t0, zero, 50
addi t1, zero, 0
loop:
addi t1, t1, 1
bne  t1, t0, loop
\end{lstlisting}
\textit{Nhập tổng số chu kỳ (vd: 107):}
\textbf{Hướng dẫn giải:}
\begin{itemize}
    \item Có $2 + 50 \times 2 = 102$ lệnh thực thi. Không tính phạt rẽ nhánh là $102 + 4 = 106$ chu kỳ; với chiến lược \textbf{branch-taken}, $49$ lần nhảy đầu đoán đúng ($0$ phạt), chỉ lần \texttt{bne} thứ $50$ không rẽ nhánh bị đoán sai tốn $1$ chu kỳ phạt ở tầng ID $\Rightarrow 106 + 1 = \mathbf{107}$ chu kỳ (hoặc $108$ nếu tính $2$ chu kỳ ở tầng EX).
    \item \textbf{Kết luận:} Đáp án là \textbf{107}.
\end{itemize}

\textbf{Bài tập 5.33 (cuối kì 20231 (\#896845); cuối kì 20232 (\#896845); cuối kì 20232 (\#53841d); cuối kì 20251 (\#0bad6a); thi thử kì 20251 (Câu 24); ôn tập cuối kì 20251 (Câu 11); ôn tập cuối kì 20251 (Câu 62); ôn tập cuối kì 20252 (P2-40); ôn tập cuối kì 20252 (P2-70))}\\
Cho đoạn mã nguồn hợp ngữ RISC-V dưới đây chạy trên CPU pipeline 5 công đoạn có đầy đủ cơ chế phát hiện xung đột dữ liệu, các đường chuyển tiếp (forwarding paths) và sử dụng bộ dự đoán rẽ nhánh \textbf{luôn nhảy (branch-taken prediction)}. Hãy cho biết tổng số chu kỳ xung nhịp để đoạn mã hoàn thành (tính cả thời gian làm đầy pipeline và phạt dự đoán sai)?
\begin{lstlisting}
addi t0, zero, 100
addi t1, zero, 0
Loop:
beq  t1, t0, exit
addi t1, t1, 1
j    Loop
exit:
\end{lstlisting}
\textit{Nhập tổng số chu kỳ xung nhịp (vd: 407):}
\textbf{Hướng dẫn giải:}
\begin{itemize}
    \item Số lệnh được thực thi: $2$ lệnh khởi tạo $+ 100$ vòng $\times\; 3$ lệnh (\texttt{beq, addi, j}) $+ 1$ lệnh \texttt{beq} ở vòng cuối (nhảy ra \texttt{exit}) $= 303$ lệnh.
    \item Nếu chưa tính hazard điều khiển: $303 + (5 - 1) = 307$ chu kỳ.
    \item Với bộ dự đoán \textbf{branch-taken} (tính $1$ chu kỳ phạt ở tầng ID khi đoán sai), trong $100$ lần đầu lệnh \texttt{beq} thực tế \textit{không nhảy} nên bị phạt $100$ chu kỳ ($100$ lần đoán sai), còn lệnh \texttt{j} luôn nhảy đúng và lệnh \texttt{beq} cuối cùng nhảy đúng.
    \item Tổng chu kỳ $= 307 + 100 = \mathbf{407}$ chu kỳ (một số đáp án không tính phạt đoán sai ghi $307$).
    \item \textbf{Kết luận:} Đáp án là \textbf{407}.
\end{itemize}

\textbf{Bài tập 5.34 (cuối kì 20231 (\#53841d))}\\
Cho đoạn mã hợp ngữ RISC-V chạy trên CPU pipeline 5 công đoạn có đầy đủ hazard detection, forwarding paths và bộ dự đoán rẽ nhánh theo chiến lược \textbf{luôn nhảy (branch-taken prediction)}. Tổng số chu kỳ xung nhịp để chương trình hoàn thành việc thực thi là bao nhiêu?
\begin{lstlisting}
addi t0, zero, 100
addi t1, zero, 0
loop:
addi t1, t1, 1
bne  t1, t0, loop
\end{lstlisting}
\textit{Nhập tổng số chu kỳ (vd: 206):}
\textbf{Hướng dẫn giải:}
\begin{itemize}
    \item Số lệnh thực thi $= 2 + 100 \times 2 = 202$ lệnh.
    \item Do dự đoán \textbf{luôn nhảy (branch-taken)}, $99$ lần lặp đầu lệnh \texttt{bne} đều rẽ nhánh (đoán đúng $\rightarrow$ $0$ chu kỳ phạt), chỉ lần thứ $100$ ($t1 == t0$) không rẽ nhánh nên tốn $1$ chu kỳ phạt (hoặc bỏ qua tùy quy ước giải quyết ở ID). Số chu kỳ $= 202 + (5 - 1) = \mathbf{206}$ chu kỳ (hoặc $207$ nếu tính thêm $1$ chu kỳ flush ở lần thoát cuối).
    \item \textbf{Kết luận:} Đáp án là \textbf{206}.
\end{itemize}

\chapter{BỘ NHỚ MÁY TÍNH}

\section{Phần 1: Lý thuyết trọng tâm}

\subsection{Phân cấp bộ nhớ, Nguyên lý cục bộ và Tổ chức bộ nhớ chính}

\textbf{1. Cấu trúc phân cấp bộ nhớ (Memory Hierarchy)}
\begin{itemize}
    \item Để dung hòa ba yếu tố đối nghịch là \textbf{tốc độ truy nhập nhanh}, \textbf{dung lượng lớn} và \textbf{chi phí trên mỗi bit thấp}, máy tính hiện đại tổ chức bộ nhớ theo mô hình phân cấp nhiều tầng.
    \item \textbf{Chiều tăng dần dung lượng} (từ nhỏ nhất đến lớn nhất):
    $$\text{Tập thanh ghi (Registers)} < \text{Cache L1} < \text{Cache L2} < \text{Cache L3} < \text{Bộ nhớ chính (RAM)} < \text{Bộ nhớ ngoài (SSD/HDD)}$$
    \item \textbf{Quy luật phân cấp:} Đi từ đỉnh (gần CPU) xuống đáy (xa CPU), \textbf{dung lượng tăng}, \textbf{giá thành/bit giảm}, nhưng \textbf{thời gian truy nhập tăng} (tốc độ chậm dần) và \textbf{tần suất truy nhập của CPU giảm}.
\end{itemize}

\begin{tcolorbox}[bwbox, title={Thông số tham khảo phân cấp bộ nhớ trên máy tính PC / Server (CA.RISCV.2025)}]
\begin{center}
\renewcommand{\arraystretch}{1.2}
\small
\begin{tabularx}{\linewidth}{|L{3.0cm}|Z|Z|Z|Y|}
\hline
\textbf{Tầng nhớ} & \textbf{Công nghệ} & \textbf{Dung lượng (PC)} & \textbf{Độ trễ truy cập} & \textbf{Vai trò chính} \\ \hline
Thanh ghi CPU & Flip-Flop & $1\text{--}2\text{ KB}$ & $\sim 0.3\text{ ns}$ & Lưu toán hạng trực tiếp cho ALU \\ \hline
Cache L1 / L2 / L3 & SRAM & $64\text{ KB} / 256\text{ KB} / 8\text{--}32\text{ MB}$ & $1\text{--}20\text{ ns}$ & Đệm lệnh và dữ liệu hay dùng gần CPU \\ \hline
Bộ nhớ chính & DRAM & $8\text{--}64\text{ GB}$ & $50\text{--}100\text{ ns}$ & Chứa chương trình và dữ liệu đang chạy \\ \hline
Bộ nhớ ngoài & Flash (SSD) / Từ (HDD) & $512\text{ GB}\text{--}2\text{ TB}$ & $50\text{--}100\,\mu\text{s}$ / $5\text{--}10\text{ ms}$ & Lưu trữ lâu dài (không khả biến) \\ \hline
\end{tabularx}
\end{center}
\end{tcolorbox}

\textbf{2. Nguyên lý cục bộ hóa tham chiếu bộ nhớ (Principle of Locality)}
\begin{itemize}
    \item \textbf{Cục bộ theo thời gian (Temporal Locality):} Một ô nhớ vừa được truy cập thì có xác suất cao sẽ sớm được truy cập lại trong tương lai gần (ví dụ: biến đếm vòng lặp \texttt{i}, biến tích lũy \texttt{sum}, các lệnh bên trong thân vòng lặp).
    \item \textbf{Cục bộ theo không gian (Spatial Locality):} Khi một ô nhớ được truy cập, các ô nhớ có địa chỉ nằm liền kề xung quanh nó có xác suất cao sẽ sớm được truy cập tiếp theo (ví dụ: duyệt tuần tự các phần tử mảng \texttt{A[i]}, thực thi chuỗi lệnh tuần tự).
\end{itemize}

\textbf{3. Phân loại bộ nhớ bán dẫn (SRAM, DRAM, ROM và Flash)}
\begin{tcolorbox}[bwbox, title={So sánh các loại bộ nhớ bán dẫn cơ bản}]
\begin{center}
\renewcommand{\arraystretch}{1.2}
\small
\begin{tabularx}{\linewidth}{|L{2.4cm}|L{2.5cm}|L{2.8cm}|L{2.2cm}|Y|}
\hline
\textbf{Loại bộ nhớ} & \textbf{Phần tử lưu trữ} & \textbf{Khả năng xóa / ghi} & \textbf{Tính khả biến} & \textbf{Đặc điểm \& Ứng dụng} \\ \hline
\textbf{SRAM} & Mạch Flip-Flop & Bằng điện, từng byte & Khả biến & Rất nhanh, không cần làm tươi, đắt $\rightarrow$ \textbf{Cache} \\ \hline
\textbf{DRAM} & Tụ điện (1T1C) & Bằng điện, từng byte & Khả biến & Cần \textbf{mạch làm tươi (refresh)}, dung lượng lớn, rẻ $\rightarrow$ \textbf{RAM chính} \\ \hline
\textbf{Mask ROM / PROM} & Mặt nạ / cầu chì & Không xóa được (PROM ghi 1 lần) & Không khả biến & Lưu firmware cố định, BIOS ban đầu \\ \hline
\textbf{EPROM} & Cổng nổi & Xóa cả chip bằng tia cực tím (UV) & Không khả biến & Ghi lại nhiều lần bằng thiết bị chuyên dụng \\ \hline
\textbf{EEPROM} & Cổng nổi & Xóa/ghi bằng điện, từng byte & Không khả biến & Lưu cấu hình dung lượng nhỏ \\ \hline
\textbf{Flash} & Cổng nổi & Xóa/ghi bằng điện, \textbf{theo từng khối} & Không khả biến & Dung lượng lớn, nhanh $\rightarrow$ \textbf{SSD, USB, thẻ nhớ, Smartphone} \\ \hline
\end{tabularx}
\end{center}
\end{tcolorbox}

\textbf{4. Thiết kế mở rộng mô-đun nhớ RAM và Bộ nhớ đan xen (Interleaved Memory)}
\begin{itemize}
    \item \textbf{Bài toán ghép chip RAM} dung lượng $M \times W\text{-bit}$ từ các chip SRAM nhỏ $m \times w\text{-bit}$:
    \begin{itemize}
        \item \textit{Mở rộng độ dài từ nhớ (chiều rộng):} Ghép song song $k_w = \dfrac{W}{w}$ chip trên mỗi hàng; nối chung đường địa chỉ và tín hiệu chọn chip $\overline{CS}$, chia bus dữ liệu thành $k_w$ nhóm $w$-bit.
        \item \textit{Mở rộng dung lượng từ nhớ (chiều sâu):} Ghép nối tiếp $k_d = \dfrac{M}{m}$ hàng; nối chung bus dữ liệu và $\log_2(m)$ đường địa chỉ thấp $A_0\dots A_{\log_2 m - 1}$. Dùng $\log_2(k_d)$ đường địa chỉ cao đưa vào cổng NOT (khi $k_d = 2$) hoặc bộ giải mã $\log_2(k_d)$-sang-$k_d$ để sinh tín hiệu $\overline{CS}$ chọn hàng.
        \item \textit{Tổng số chip cần dùng:} $N_{\text{chip}} = k_w \times k_d = \dfrac{M \times W}{m \times w}$.
    \end{itemize}
    \item \textbf{Bộ nhớ tổ chức theo băng nhớ đan xen (Interleaved Memory):} Nếu bus dữ liệu rộng $B_{\text{data}}$ byte (ví dụ $32\text{ bit} = 4\text{ byte} \Rightarrow 4$ băng; $64\text{ bit} = 8\text{ byte} \Rightarrow 8$ băng), byte nhớ có địa chỉ $A$ nằm trên băng nhớ số:
    $$\text{Chỉ số băng nhớ} = A \bmod B_{\text{data}}$$
\end{itemize}

\subsection{Bộ nhớ đệm (Cache): Ánh xạ, Tách bit địa chỉ và Kích thước phần cứng}

\textbf{1. Ba phương pháp ánh xạ Cache}
\begin{itemize}
    \item \textbf{Ánh xạ trực tiếp (Direct-Mapped):} Mỗi Block của bộ nhớ chính chỉ được phép nạp vào \textbf{duy nhất 1 Line} cố định trong Cache theo công thức $\text{Line} = \text{BlockAddr} \bmod N_{\text{line}}$. Đơn giản, nhanh nhưng dễ xảy ra xung đột (Conflict Miss).
    \item \textbf{Ánh xạ liên kết toàn phần (Fully-Associative):} Mỗi Block của bộ nhớ chính được phép nạp vào \textbf{bất kỳ Line nào} trong Cache. Giảm tối đa xung đột nhưng mạch so sánh Tag song song phức tạp và đắt tiền.
    \item \textbf{Ánh xạ liên kết tập hợp $K$-đường ($K$-way Set-Associative):} Cache được chia thành $S = \dfrac{N_{\text{line}}}{K}$ tập hợp (Set), mỗi Set gồm $K$ Line. Mỗi Block ánh xạ cố định vào một Set ($\text{Set} = \text{BlockAddr} \bmod S$) và có thể nạp vào bất kỳ Line nào trong $K$ Line của Set đó.
\end{itemize}

\begin{tcolorbox}[bwbox, title={Bộ công thức tách bit địa chỉ và kích thước phần cứng Cache}]
Giả sử bộ nhớ chính đánh địa chỉ theo byte với địa chỉ dài $N$ bit (dung lượng $2^N$ byte), Cache có dung lượng dữ liệu $C$ byte, kích thước mỗi khối (Line/Block) là $B$ byte. Khi đó:
\begin{itemize}
    \item \textbf{Số bit trường Offset (Word):} $W = \log_2(B)$.
    \item \textbf{Tổng số Line trong Cache:} $N_{\text{line}} = \dfrac{C}{B}$.
\end{itemize}
\begin{center}
\renewcommand{\arraystretch}{1.25}
\small
\begin{tabularx}{\linewidth}{|L{3.5cm}|Z|Z|Z|}
\hline
\textbf{Kiểu ánh xạ} & \textbf{Số bit Index (Line/Set)} & \textbf{Số bit Tag} & \textbf{Cấu trúc địa chỉ $N$-bit} \\ \hline
Trực tiếp (Direct) & $L = \log_2(N_{\text{line}})$ & $T = N - L - W$ & \texttt{[Tag | Line | Word]} \\ \hline
Liên kết tập hợp $K$-way & $S_{\text{bit}} = \log_2\!\left(\dfrac{N_{\text{line}}}{K}\right)$ & $T = N - S_{\text{bit}} - W$ & \texttt{[Tag | Set | Word]} \\ \hline
Liên kết toàn phần (Fully) & $0$ bit & $T = N - W$ & \texttt{[Tag | Word]} \\ \hline
\end{tabularx}
\end{center}
\textbf{Tổng kích thước phần cứng của Cache (tính bằng bit):}
$$\text{Total Bits} = N_{\text{line}} \times \big(\underbrace{1}_{\text{Valid}} + \underbrace{d}_{\text{Dirty (0 hoặc 1)}} + \underbrace{T}_{\text{Tag bits}} + \underbrace{8 \times B}_{\text{Data bits}}\big)$$
(Với Cache $K$-way có $S$ Set, có thể viết tương đương: $\text{Total Bits} = S \times K \times (1 + d + T + 8B)$.)
\end{tcolorbox}

\subsection{Chính sách thay thế, Chính sách ghi, Lỗi 3C và Hiệu năng Cache}

\textbf{1. Thuật toán thay thế Block (khi Cache/Set đã đầy)}
\begin{itemize}
    \item \textbf{Ánh xạ trực tiếp:} Không cần chọn, thay thế đúng Line duy nhất đó.
    \item \textbf{Ánh xạ liên kết (Set / Fully):}
    \begin{itemize}
        \item \textbf{LRU (Least Recently Used):} Thay thế block lâu nhất chưa được tham chiếu $\rightarrow$ \textbf{tối ưu nhất} về mặt khai thác tính cục bộ thời gian.
        \item \textbf{FIFO (First In First Out):} Thay thế block được nạp vào sớm nhất (nằm lâu nhất trong Set).
        \item \textbf{LFU (Least Frequently Used):} Thay thế block có số lần tham chiếu ít nhất trong một khoảng thời gian.
        \item \textbf{Random:} Chọn thay thế ngẫu nhiên (phần cứng đơn giản nhất).
    \end{itemize}
\end{itemize}

\textbf{2. Chính sách ghi dữ liệu (Write Policy)}
\begin{itemize}
    \item \textbf{Khi Cache Hit:}
    \begin{itemize}
        \item \textbf{Ghi xuyên qua (Write-Through):} Ghi đồng thời vào cả Cache và bộ nhớ chính $\rightarrow$ đảm bảo nhất quán đơn giản nhưng tốc độ ghi chậm và gây nghẽn bus (thường kết hợp \textit{Write Buffer}).
        \item \textbf{Ghi trả sau (Write-Back):} Chỉ ghi vào Cache và bật bit trạng thái \texttt{Dirty = 1}; chỉ khi block bị thay thế khỏi Cache mới ghi trả toàn bộ block về RAM $\rightarrow$ tốc độ nhanh, giảm lưu lượng bus.
    \end{itemize}
    \item \textbf{Khi Cache Miss:} \textbf{Write-Allocate} (nạp block từ RAM lên Cache rồi mới ghi -- thường đi kèm Write-Back) hoặc \textbf{No-Write-Allocate} (ghi thẳng xuống RAM không nạp lên Cache -- thường đi kèm Write-Through).
\end{itemize}

\textbf{3. Phân loại lỗi trượt Cache 3C (Three Cs)}
\begin{itemize}
    \item \textbf{Compulsory Miss (Cold Miss -- Lỗi bắt buộc):} Xảy ra ở lần truy cập đầu tiên vào một block chưa từng có trong Cache. \textit{Cách giảm:} tăng kích thước block (khai thác cục bộ không gian) hoặc nạp trước (prefetching).
    \item \textbf{Capacity Miss (Lỗi dung lượng):} Tổng dung lượng Cache nhỏ hơn tập làm việc của chương trình. \textit{Cách giảm:} tăng dung lượng Cache.
    \item \textbf{Conflict Miss (Lỗi xung đột):} Nhiều block cùng ánh xạ vào một Line/Set dù Cache vẫn còn chỗ trống ở Set khác. \textit{Cách giảm:} tăng bậc liên kết tập hợp (associativity).
\end{itemize}

\textbf{4. Công thức đánh giá hiệu năng Cache (AMAT và $CPI_{\text{thực}}$)}
\begin{tcolorbox}[bwbox, title={Công thức AMAT và CPI hiệu dụng khi xét Cache Miss}]
\begin{itemize}
    \item \textbf{Thời gian truy nhập bộ nhớ trung bình (AMAT -- Average Memory Access Time):}
    $$\text{AMAT} = \text{Hit Time} + \text{Miss Rate} \times \text{Miss Penalty}$$
    \item \textbf{Số chu kỳ dừng do trượt Cache trên mỗi lệnh (Memory Stall Cycles / Instruction):}
    $$\text{Stall}_{\text{I-cache}} = \text{MR}_I \times \text{MP}, \qquad \text{Stall}_{\text{D-cache}} = f_{\text{mem}} \times \text{MR}_D \times \text{MP}$$
    trong đó $f_{\text{mem}}$ là tỉ lệ lệnh truy nhập bộ nhớ dữ liệu (\texttt{load}/\texttt{store}), $\text{MR}_I$ và $\text{MR}_D$ là tỉ lệ trượt I-cache và D-cache, $\text{MP}$ là thời gian phạt trượt (Miss Penalty, tính bằng chu kỳ).
    \item \textbf{CPI thực tế (Effective CPI):}
    $$\text{CPI}_{\text{actual}} = \text{CPI}_{\text{base}} + \text{MR}_I \times \text{MP} + f_{\text{mem}} \times \text{MR}_D \times \text{MP}$$
\end{itemize}
\end{tcolorbox}

\subsection{Bộ nhớ ảo (Virtual Memory), Phân trang (Paging) và TLB}

\begin{itemize}
    \item \textbf{Khái niệm:} Bộ nhớ ảo kết hợp bộ nhớ chính (RAM) và bộ nhớ ngoài (SSD/HDD) tạo thành một không gian địa chỉ logic thống nhất lớn hơn dung lượng RAM vật lý, cho phép chạy các tiến trình lớn hơn RAM thực và cô lập không gian nhớ giữa các tiến trình.
    \item \textbf{Kỹ thuật phân trang (Paging):}
    \begin{itemize}
        \item Không gian địa chỉ ảo được chia thành các \textbf{trang (Page)} kích thước bằng nhau (thông dụng $4\text{ KiB} = 2^{12}\text{ byte}$); bộ nhớ vật lý chia thành các \textbf{khung trang (Frame)} có cùng kích thước.
        \item Các trang của một tiến trình \textbf{không cần nạp vào các khung trang vật lý liên tiếp}; ánh xạ được quản lý bởi \textbf{Bảng trang (Page Table)}.
    \end{itemize}
    \item \textbf{Dịch địa chỉ ảo sang địa chỉ vật lý:} Địa chỉ ảo $n$-bit được tách thành \texttt{[VPN | Page Offset]} với số bit $\text{Offset} = \log_2(\text{Page Size})$ và $\text{VPN} = n - \text{Offset}$. Dùng \texttt{VPN} tra Bảng trang để thu được \texttt{PFN} (Physical Frame Number), sau đó giữ nguyên \texttt{Offset} để ghép thành địa chỉ vật lý \texttt{[PFN | Page Offset]}.
    \item \textbf{Lỗi trang (Page Fault) và Thrashing:} Khi bit \texttt{Valid = 0} (trang chưa có trong RAM), CPU phát sinh \textbf{Page Fault} để hệ điều hành nạp trang từ đĩa vào RAM. Nếu quá nhiều tiến trình cùng chạy trong RAM quá nhỏ khiến HĐH liên tục hoán đổi trang ra/vào đĩa, hệ thống rơi vào trạng thái \textbf{Thrashing (thất bại)}.
    \item \textbf{Bộ đệm dịch địa chỉ TLB (Translation Lookaside Buffer):} Là bộ nhớ Cache phần cứng tốc độ cao bên trong MMU lưu các ánh xạ \texttt{VPN $\rightarrow$ PFN} hay dùng gần đây, giúp tránh phải truy cập RAM 2 lần cho mỗi thao tác đọc/ghi dữ liệu.
    \item \textbf{Bảng trang đảo (Inverted Page Table):} Thay vì tạo bảng trang theo không gian địa chỉ ảo của từng tiến trình, chỉ duy trì \textbf{đúng 1 mục cho mỗi khung trang vật lý} trong RAM kết hợp hàm băm (Hash function) để tiết kiệm bộ nhớ lưu bảng trang.
\end{itemize}

\subsection{Bộ nhớ ngoài (HDD, SSD và Hệ thống RAID)}

\begin{itemize}
    \item \textbf{Ổ đĩa từ HDD (Hard Disk Drive):} Cấu trúc gồm các đĩa từ (Platter) quay quanh trục (Spindle), mỗi mặt đĩa chia thành các đường tròn đồng tâm (\textbf{Track}), mỗi Track chia thành các cung (\textbf{Sector}); tập hợp các Track cùng bán kính trên mọi mặt đĩa tạo thành một \textbf{Cylinder}. Thời gian truy xuất trung bình một khối dữ liệu:
    $$T_{\text{access}} = T_{\text{seek}} + T_{\text{rot\_avg}} + T_{\text{transfer}}, \qquad \text{với } T_{\text{rot\_avg}} = \frac{1}{2} \times \frac{60}{\text{RPM}}\text{ (giây)}$$
    \item \textbf{Ổ đĩa thể rắn SSD (Solid State Drive):} Dùng chip nhớ bán dẫn Flash (không khả biến), không có bộ phận cơ khí chuyển động $\rightarrow$ tốc độ truy xuất ngẫu nhiên nhanh hơn, độ trễ thấp hơn, chống sốc tốt hơn, không gây tiếng ồn và tiêu thụ ít điện hơn HDD; tuy nhiên \textbf{giá thành/GB đắt hơn} và có giới hạn chu kỳ ghi/xóa (nên HDD vẫn phổ biến trong hệ thống lưu trữ dung lượng lớn NAS/CCTV 24/7).
    \item \textbf{Hệ thống mảng đĩa RAID (Redundant Array of Independent Disks):} Ghép nhiều ổ đĩa vật lý thành một ổ logic nhằm tăng tốc độ truy xuất song song và/hoặc tăng độ an toàn dữ liệu nhờ thông tin dư thừa (parity/mirroring):
\end{itemize}

\begin{tcolorbox}[bwbox, title={Tóm tắt các mức RAID phổ biến với $N$ ổ đĩa cùng dung lượng $C$}]
\begin{center}
\renewcommand{\arraystretch}{1.2}
\small
\begin{tabularx}{\linewidth}{|L{1.8cm}|L{3.0cm}|Z|Z|Y|}
\hline
\textbf{Mức RAID} & \textbf{Kỹ thuật} & \textbf{Dung lượng khả dụng} & \textbf{Số đĩa lỗi chịu được} & \textbf{Đặc điểm nổi bật} \\ \hline
\textbf{RAID 0} & Striping (phân dải) & $N \times C$ & $0$ & Hiệu năng cao nhất, không có dự phòng lỗi \\ \hline
\textbf{RAID 1} & Mirroring (nhân bản) & $\dfrac{N}{2} \times C$ & $1$ (mỗi cặp) & An toàn cao, chi phí dung lượng $50\%$ \\ \hline
\textbf{RAID 5} & Striping + 1 Parity phân tán & $(N - 1) \times C$ & $1$ & Cân bằng tốt giữa hiệu năng, dung lượng và an toàn ($N \ge 3$) \\ \hline
\textbf{RAID 6} & Striping + 2 Parity phân tán & $(N - 2) \times C$ & $2$ & Chịu lỗi đồng thời 2 đĩa ($N \ge 4$) \\ \hline
\end{tabularx}
\end{center}
\end{tcolorbox}

\newpage
\section{Phần 2: Toàn bộ bài tập Chương 6 và Hướng dẫn giải}

\subsection{Nhóm 1: Phân cấp bộ nhớ, Nguyên lý cục bộ, SRAM/DRAM \& Tổ chức chip nhớ}

\textbf{Bài tập 6.1 (ôn tập cuối kì 20252 (P2-23); Đáp án ôn tập cuối kỳ P2 (Câu 44); cuối kì 20212 (B1-4.2))}\\
Trên máy tính dùng bộ xử lý có bus địa chỉ $32\text{-bit}$, bus dữ liệu $64\text{-bit}$, các thanh ghi bên trong có độ dài $32\text{-bit}$. Giả sử có một mảng dữ liệu hai chiều \texttt{B} có kích thước $5 \times 7$ của các phần tử $32\text{-bit}$, lưu trữ theo hàng trong bộ nhớ chính (tổ chức băng nhớ đan xen) bắt đầu từ địa chỉ \texttt{0x000A5020}. Hỏi phần tử \texttt{B[1, 5]} (với chỉ số $i = 0\dots4$, $j = 0\dots6$) nằm trên các băng nhớ nào của bộ nhớ vật lý?
\begin{enumerate}[label=\Alph*.]
    \item Băng 0, 1, 2, 3
    \item Băng 4, 5, 6, 7
    \item Băng 3, 4, 5, 6
    \item Băng 1, 2, 3, 4
\end{enumerate}
\textbf{Hướng dẫn giải:}
\begin{itemize}
    \item \textbf{Bước 1 --- Số băng nhớ đan xen:} Bus dữ liệu rộng $64\text{ bit} = 8\text{ byte} \Rightarrow$ bộ nhớ chính được chia thành \textbf{$8$ băng nhớ đan xen} (Băng $0$ đến Băng $7$). Mỗi byte có địa chỉ $A$ nằm trên băng nhớ số $A \bmod 8$.
    \item \textbf{Bước 2 --- Địa chỉ phần tử \texttt{B[1, 5]}:} Mảng lưu theo hàng (row-major), mỗi hàng có $7$ phần tử, mỗi phần tử chiếm $4\text{ byte}$:
    \item $\text{Addr}(B[1, 5]) = \text{0x000A5020} + (1 \times 7 + 5) \times 4 = \text{0x000A5020} + 48 = \text{0x000A5020} + \text{0x30} = \texttt{0x000A5050}$.
    \item \textbf{Bước 3 --- Xác định băng nhớ:} $4$ byte của \texttt{B[1, 5]} nằm ở các địa chỉ \texttt{0x000A5050} đến \texttt{0x000A5053}. Vì $\text{0x50} = 80 \equiv 0 \pmod 8$, nên $4$ byte này nằm trên các băng nhớ \textbf{0, 1, 2, 3}.
    \item \textbf{Kết luận:} Đáp án đúng là \textbf{A} (Băng 0, 1, 2, 3).
\end{itemize}

\textbf{Bài tập 6.2 (cuối kì 20212 (B3-4.3))}\\
Máy tính dùng $36\text{ bit}$ địa chỉ để đánh địa chỉ cho bộ nhớ theo byte, bus dữ liệu kết nối với bộ nhớ chính là $64\text{ bit}$ (tổ chức băng nhớ đan xen). Hãy cho biết hai byte nhớ có địa chỉ \texttt{0x0FE12C39} và \texttt{0x10ABCD05} lần lượt được bố trí ở băng nhớ nào?
\begin{enumerate}[label=\Alph*.]
    \item Băng 1 và Băng 1
    \item Băng 1 và Băng 5
    \item Băng 5 và Băng 1
    \item Băng 3 và Băng 5
\end{enumerate}
\textbf{Hướng dẫn giải:}
\begin{itemize}
    \item Bus dữ liệu rộng $64\text{ bit} = 8\text{ byte} \Rightarrow$ bộ nhớ chính tổ chức thành \textbf{$8$ băng nhớ đan xen} (Băng $0$ đến Băng $7$). Chỉ số băng nhớ bằng $3$ bit thấp nhất của địa chỉ (tức $\text{Addr} \bmod 8$):
    \item Địa chỉ \texttt{0x0FE12C39}: chữ số Hex cuối là $\text{9} = 1001_2 \Rightarrow 9 \bmod 8 = \mathbf{1}$ (\textbf{Băng 1}).
    \item Địa chỉ \texttt{0x10ABCD05}: chữ số Hex cuối là $\text{5} = 0101_2 \Rightarrow 5 \bmod 8 = \mathbf{5}$ (\textbf{Băng 5}).
    \item \textbf{Kết luận:} Đáp án đúng là \textbf{B} (Băng 1 và Băng 5).
\end{itemize}

\textbf{Bài tập 6.3 (cuối kì 20212 (B6-4))}\\
Thiết kế mô-đun nhớ RAM dung lượng $64\text{K} \times 8\text{-bit}$ trên cơ sở các IC nhớ SRAM dung lượng $32\text{K} \times 4\text{-bit}$ cần bố trí số chip và tín hiệu chọn chip như thế nào?
\begin{enumerate}[label=\Alph*.]
    \item Dùng 4 chip ghép nối tiếp theo chiều sâu, dùng A14 và A15 qua bộ giải mã 2-sang-4 để chọn chip
    \item Dùng 2 chip ghép song song theo chiều rộng, nối chung A0--A15 và tín hiệu CS
    \item Dùng 4 chip: ghép song song 2 chip theo chiều rộng (4 lên 8 bit) và 2 hàng nối tiếp theo chiều sâu (32K lên 64K), nối chung A0--A14 và dùng A15 (kèm cổng NOT) chọn hàng (CS)
    \item Dùng 8 chip: ghép song song 4 chip theo chiều rộng và 2 hàng theo chiều sâu, dùng A15 chọn hàng
\end{enumerate}
\textbf{Hướng dẫn giải:}
\begin{itemize}
    \item \textbf{1. Tính số lượng chip:}
    \item Mở rộng độ dài từ nhớ (chiều rộng): $8\text{ bit} / 4\text{ bit} = 2$ chip ghép song song trên mỗi hàng (chia bus dữ liệu $D_0\text{-}D_3$ và $D_4\text{-}D_7$).
    \item Mở rộng dung lượng (chiều sâu): $64\text{K} / 32\text{K} = 2$ hàng nối tiếp.
    \item Tổng số chip cần dùng $= 2 \times 2 = \mathbf{4}$ chip.
    \item \textbf{2. Kết nối đường địa chỉ \& chọn chip:}
    \item Mỗi chip $32\text{K} = 2^{15}$ cần $15$ đường địa chỉ $A_0\text{-}A_{14}$ nối chung cho cả $4$ chip.
    \item Mô-đun $64\text{K} = 2^{16}$ cần $16$ đường địa chỉ $\Rightarrow$ bit cao nhất $A_{15}$ (qua cổng \texttt{NOT}) dùng để kích hoạt tín hiệu $\overline{CS}$ chọn giữa hàng 0 ($A_{15}=0$) và hàng 1 ($A_{15}=1$).
    \item \textbf{Kết luận:} Đáp án đúng là \textbf{C} (Dùng 4 chip: ghép song song 2 chip theo chiều rộng (4 lên 8 bit) và 2 hàng nối tiếp theo chiều sâu (32K lên 64K), nối chung A0--A14 và dùng A15 (kèm cổng NOT) chọn hàng (CS)).
\end{itemize}

\textbf{Bài tập 6.4 (Lý thuyết slide CA.RISCV.2025 - Mục 6.2 (SRAM \& DRAM))}\\
Phát biểu nào sau đây phân biệt \textbf{đúng} giữa bộ nhớ RAM tĩnh (\texttt{SRAM}) và RAM động (\texttt{DRAM})?
\begin{enumerate}[label=\Alph*.]
    \item SRAM lưu bit trên tụ điện nên cần mạch làm tươi định kỳ; DRAM lưu bit bằng Flip-Flop nên nhanh hơn và dùng làm Cache
    \item SRAM lưu bit bằng mạch Flip-Flop, tốc độ nhanh, đắt tiền nên dùng làm Cache; DRAM lưu bit trên tụ điện, cần mạch làm tươi (refresh), dung lượng lớn, rẻ hơn nên dùng làm bộ nhớ chính
    \item Cả SRAM và DRAM đều là bộ nhớ không khả biến (non-volatile), giữ được dữ liệu khi mất nguồn điện
    \item DRAM có cấu trúc phần tử nhớ phức tạp hơn SRAM nên chi phí trên mỗi bit cao hơn
\end{enumerate}
\textbf{Hướng dẫn giải:}
\begin{itemize}
    \item Theo Slide 20--21 (CA.RISCV.2025):
    \item \textbf{SRAM (Static RAM):} Lưu bit bằng \textbf{Flip-Flop}, thông tin ổn định không cần làm tươi, tốc độ rất nhanh nhưng cấu trúc phức tạp và đắt tiền $\Rightarrow$ dùng làm \textbf{bộ nhớ Cache}.
    \item \textbf{DRAM (Dynamic RAM):} Lưu bit bằng điện tích trên \textbf{tụ điện} nên bị rò điện và bắt buộc có \textbf{mạch làm tươi (refresh)}; cấu trúc đơn giản, mật độ cao, giá rẻ $\Rightarrow$ dùng làm \textbf{bộ nhớ chính}.
    \item Cả hai đều là bộ nhớ \textbf{khả biến (volatile)}.
    \item \textbf{Kết luận:} Đáp án đúng là \textbf{B} (SRAM lưu bit bằng mạch Flip-Flop, tốc độ nhanh, đắt tiền nên dùng làm Cache; DRAM lưu bit trên tụ điện, cần mạch làm tươi (refresh), dung lượng lớn, rẻ hơn nên dùng làm bộ nhớ chính).
\end{itemize}

\textbf{Bài tập 6.5 (Lý thuyết slide CA.RISCV.2025 - Mục 6.2 (ROM \& Flash))}\\
Trong các loại bộ nhớ bán dẫn không khả biến dưới đây, loại nào cho phép \textbf{xóa bằng điện theo từng khối (block)} và được ứng dụng phổ biến làm ổ đĩa SSD, USB và bộ nhớ lưu trữ trên điện thoại thông minh?
\begin{enumerate}[label=\Alph*.]
    \item PROM (Programmable ROM)
    \item EPROM (Erasable PROM)
    \item Mask ROM (ROM mặt nạ)
    \item Bộ nhớ Flash (Flash Memory)
\end{enumerate}
\textbf{Hướng dẫn giải:}
\begin{itemize}
    \item Theo Slide 15--18 (CA.RISCV.2025):
    \item \textbf{Mask ROM:} Ghi khi sản xuất, không xóa được.
    \item \textbf{PROM:} Chỉ ghi được một lần bằng thiết bị chuyên dụng.
    \item \textbf{EPROM:} Xóa cả chip bằng tia cực tím (UV).
    \item \textbf{EEPROM:} Xóa và ghi bằng điện ở mức từng byte.
    \item \textbf{Bộ nhớ Flash:} Không khả biến, \textbf{xóa và ghi bằng điện theo từng khối}, dung lượng lớn $\Rightarrow$ dùng làm SSD, thẻ nhớ, USB, bộ nhớ smartphone.
    \item \textbf{Kết luận:} Đáp án đúng là \textbf{D} (Bộ nhớ Flash (Flash Memory)).
\end{itemize}

\textbf{Bài tập 6.6 (cuối kì 20202 (Câu 16); cuối kì 20211 (Câu 16))}\\
Hãy tiến hành sắp xếp có tổ chức cấu trúc phân cấp các thành phần bộ nhớ trong kiến trúc máy tính hiện đại dựa theo thứ tự dung lượng lưu trữ tăng dần (từ nhỏ nhất đến lớn nhất):
\begin{enumerate}[label=\Alph*.]
    \item Cache L1, Cache L2, Cache L3, Bộ nhớ chính, Bộ nhớ ngoài
    \item Cache L3, Cache L2, Cache L1, Bộ nhớ chính, Bộ nhớ ngoài
    \item Bộ nhớ chính, Cache L3, Cache L2, Cache L1, Bộ nhớ ngoài
    \item Cache L1, Cache L2, Bộ nhớ chính, Cache L3, Bộ nhớ ngoài
\end{enumerate}
\textbf{Hướng dẫn giải:}
\begin{itemize}
    \item Theo tháp phân cấp bộ nhớ, từ đỉnh xuống đáy tốc độ giảm dần và \textbf{dung lượng tăng dần}: Thanh ghi $\rightarrow$ Cache L1 $\rightarrow$ Cache L2 $\rightarrow$ Cache L3 $\rightarrow$ Bộ nhớ chính (RAM) $\rightarrow$ Bộ nhớ ngoài (SSD/HDD).
    \item \textbf{Kết luận:} Đáp án đúng là \textbf{A} (Cache L1, Cache L2, Cache L3, Bộ nhớ chính, Bộ nhớ ngoài).
\end{itemize}

\textbf{Bài tập 6.7 (cuối kì 20222 (\#ef5037); cuối kì 20241 (\#018513); cuối kì 20242 (\#ef5037); ôn tập cuối kì 20252 (P1-67))}\\
Hãy chỉ ra đâu là cấu hình bộ nhớ phù hợp cho máy tính PC ?
\begin{enumerate}[label=\Alph*.]
    \item Cache L1: 64 KB, Cache L2: 256 KB, Cache L3: 8 MB, RAM: 8 GB, SSD: 512 GB
    \item Cache L1: 64 KB, Cache L2: 256 KB, Cache L3: 16 GB, RAM: 2 TB, SSD: 1 TB
    \item Cache L1: 64 KB, Cache L2: 256 GB, Cache L3: 8 MB, RAM: 8 GB, SSD: 512 GB
    \item Cache L1: 16 MB, Cache L2: 256 KB, Cache L3: 16 MB, RAM: 8 GB, SSD: 1 TB
\end{enumerate}
\textbf{Hướng dẫn giải:}
\begin{itemize}
    \item Cấu hình bộ nhớ chuẩn của một máy tính PC phải tuân theo trật tự dung lượng tăng dần hợp lý: $L1 < L2 < L3 < \text{RAM} < \text{SSD}$ (ví dụ: L1 $64\text{ KB}$, L2 $256\text{ KB}$, L3 $8\text{ MB}$, RAM $8\text{ GB}$, SSD $512\text{ GB}$).
    \item \textbf{Kết luận:} Đáp án đúng là \textbf{A} (Cache L1: 64 KB, Cache L2: 256 KB, Cache L3: 8 MB, RAM: 8 GB, SSD: 512 GB).
\end{itemize}

\textbf{Bài tập 6.8 (cuối kì 20212 (B1-5))}\\
Khi thiết kế mô-đun nhớ RAM dung lượng $64\text{K} \times 16\text{-bit}$ từ các chip nhớ SRAM dung lượng $32\text{K} \times 8\text{-bit}$, phương án ghép chip và chọn chip nào sau đây chính xác?
\begin{enumerate}[label=\Alph*.]
    \item Dùng 4 chip: ghép song song 2 chip theo chiều rộng (8 lên 16 bit) và 2 hàng nối tiếp theo chiều sâu (32K lên 64K), dùng đường địa chỉ A15 (và cổng NOT) chọn chip (CS)
    \item Dùng 4 chip ghép song song toàn bộ theo chiều rộng, dùng A15 và A16 qua bộ giải mã 2-sang-4 để chọn chip
    \item Dùng 2 chip ghép nối tiếp theo chiều sâu, dùng đường địa chỉ A14 chọn chip (CS)
    \item Dùng 8 chip: 4 chip theo chiều rộng và 2 hàng theo chiều sâu, dùng A14--A15 để chọn chip
\end{enumerate}
\textbf{Hướng dẫn giải:}
\begin{itemize}
    \item Số chip cần dùng $= (64\text{K}/32\text{K}) \times (16/8) = 2 \times 2 = 4$ chip.
    \item Chip $32\text{K} = 2^{15}$ dùng 15 đường địa chỉ $A_0\text{-}A_{14}$; mô-đun $64\text{K} = 2^{16}$ cần thêm $A_{15}$ (nối trực tiếp và qua cổng NOT) để kích hoạt chân $\overline{CS}$ của 2 hàng chip; mỗi hàng gồm 2 chip $8\text{-bit}$ chia sẻ bus dữ liệu $D_0\text{-}D_7$ và $D_8\text{-}D_{15}$.
    \item \textbf{Kết luận:} Đáp án đúng là \textbf{A} (Dùng 4 chip: ghép song song 2 chip theo chiều rộng (8 lên 16 bit) và 2 hàng nối tiếp theo chiều sâu (32K lên 64K), dùng đường địa chỉ A15 (và cổng NOT) chọn chip (CS)).
\end{itemize}

\textbf{Bài tập 6.9 (cuối kì 20212 (B2-5); cuối kì 20212 (B5-4))}\\
Khi thiết kế mô-đun nhớ RAM dung lượng $128\text{K} \times 8\text{-bit}$ trên cơ sở chip nhớ SRAM dung lượng $32\text{K} \times 8\text{-bit}$, cách tổ chức nào sau đây đúng?
\begin{enumerate}[label=\Alph*.]
    \item Dùng 4 chip ghép nối tiếp theo chiều sâu, nối chung A0--A14 và D0--D7, dùng 2 bit địa chỉ cao A15, A16 qua bộ giải mã 2-sang-4 để chọn chip (CS)
    \item Dùng 4 chip ghép song song theo chiều rộng, dùng A15 qua cổng NOT để chọn chip
    \item Dùng 8 chip ghép kết hợp 2 chiều rộng và 4 chiều sâu, dùng A14--A16 để chọn chip
    \item Dùng 2 chip ghép nối tiếp theo chiều sâu, dùng đường A15 để chọn chip
\end{enumerate}
\textbf{Hướng dẫn giải:}
\begin{itemize}
    \item Số chip $= (128\text{K}/32\text{K}) \times (8/8) = 4$ chip ghép nối tiếp chiều sâu. Mỗi chip có $15$ đường địa chỉ ($A_0\text{-}A_{14}$); $128\text{K} = 2^{17}$ cần $17$ đường ($A_0\text{-}A_{16}$), nên $2$ đường cao $A_{15}, A_{16}$ đưa qua bộ giải mã $2\text{-to-}4$ để chọn $1$ trong $4$ chip.
    \item \textbf{Kết luận:} Đáp án đúng là \textbf{A} (Dùng 4 chip ghép nối tiếp theo chiều sâu, nối chung A0--A14 và D0--D7, dùng 2 bit địa chỉ cao A15, A16 qua bộ giải mã 2-sang-4 để chọn chip (CS)).
\end{itemize}

\textbf{Bài tập 6.10 (cuối kì 20212 (B4-4))}\\
Để thiết kế mô-đun nhớ RAM dung lượng $32\text{K} \times 16\text{-bit}$ trên cơ sở các IC nhớ SRAM $32\text{K} \times 4\text{-bit}$, cần thiết kế như thế nào?
\begin{enumerate}[label=\Alph*.]
    \item Dùng 4 IC ghép song song theo chiều rộng; nối chung A0--A14 và CS cho cả 4 IC, chia bus D0--D15 thành 4 nhóm 4-bit vào từng IC
    \item Dùng 4 IC ghép nối tiếp theo chiều sâu, dùng A15 và A16 qua bộ giải mã 2-sang-4 để chọn IC
    \item Dùng 2 IC ghép song song theo chiều rộng để mở rộng từ 4-bit lên 16-bit
    \item Dùng 8 IC kết hợp 2 nhóm nối tiếp và 4 nhóm song song, dùng A14 để chọn nhóm
\end{enumerate}
\textbf{Hướng dẫn giải:}
\begin{itemize}
    \item Số từ nhớ giữ nguyên ($32\text{K} = 32\text{K}$), chỉ tăng độ rộng dữ liệu từ $4$ bit lên $16$ bit $\Rightarrow$ cần $16 / 4 = 4$ IC ghép \textbf{song song theo chiều rộng}, dùng chung $A_0\text{-}A_{14}$ và $\overline{CS}$.
    \item \textbf{Kết luận:} Đáp án đúng là \textbf{A} (Dùng 4 IC ghép song song theo chiều rộng; nối chung A0--A14 và CS cho cả 4 IC, chia bus D0--D15 thành 4 nhóm 4-bit vào từng IC).
\end{itemize}

\textbf{Bài tập 6.11 (cuối kì 20212 (B7-4))}\\
Thiết kế mô-đun nhớ RAM dung lượng $64\text{K} \times 8\text{-bit}$ trên cơ sở các IC nhớ SRAM $16\text{K} \times 8\text{-bit}$ cần bố trí như thế nào?
\begin{enumerate}[label=\Alph*.]
    \item Dùng 4 IC ghép nối tiếp theo chiều sâu; A0--A13 và D0--D7 nối chung, dùng A14 và A15 qua bộ giải mã 2-sang-4 để chọn chip (CS)
    \item Dùng 4 IC ghép song song theo chiều rộng, dùng A15 qua cổng NOT để chọn chip
    \item Dùng 2 IC ghép nối tiếp theo chiều sâu, dùng A14 để chọn chip
    \item Dùng 4 IC ghép nối tiếp, dùng A15 và A16 qua bộ giải mã 2-sang-4 để chọn chip
\end{enumerate}
\textbf{Hướng dẫn giải:}
\begin{itemize}
    \item Số chip $= 64\text{K} / 16\text{K} = 4$ IC nối tiếp chiều sâu. Mỗi chip $16\text{K} = 2^{14}$ dùng $A_0\text{-}A_{13}$; mô-đun $64\text{K} = 2^{16}$ dùng $2$ đường cao $A_{14}, A_{15}$ đưa qua decoder $2\text{-sang-}4$ để chọn chip.
    \item \textbf{Kết luận:} Đáp án đúng là \textbf{A} (Dùng 4 IC ghép nối tiếp theo chiều sâu; A0--A13 và D0--D7 nối chung, dùng A14 và A15 qua bộ giải mã 2-sang-4 để chọn chip (CS)).
\end{itemize}

\textbf{Bài tập 6.12 (cuối kì 20212 (B8-4))}\\
Bộ xử lý có bus địa chỉ $20\text{ bit}$ ($A_0\text{-}A_{19}$), bus dữ liệu $8\text{ bit}$. $4$ bit cao $A_{19}A_{18}A_{17}A_{16}$ phân vùng: \texttt{0000} cho RAM ($64\text{ KB}$ từ chip SRAM $16\text{K}\times 8\text{-bit}$), \texttt{1111} cho ROM ($32\text{ KB}$), \texttt{1000} cho $32$ cổng I/O (Memory-Mapped I/O, đánh địa chỉ từ $A_{15}\text{-}A_0 = 0$). Kết luận nào đúng về số chip RAM và dải địa chỉ (Hexa) của RAM, ROM, I/O?
\begin{enumerate}[label=\Alph*.]
    \item 4 chip SRAM 16K$\times$8-bit; RAM: 0x00000--0x0FFFF; ROM: 0xF0000--0xF7FFF; I/O: 0x80000--0x8001F
    \item 4 chip SRAM 16K$\times$8-bit; RAM: 0x00000--0x0FFFF; ROM: 0xF0000--0xFFFFF; I/O: 0x80000--0x80020
    \item 2 chip SRAM 16K$\times$8-bit; RAM: 0x00000--0x07FFF; ROM: 0xF0000--0xF7FFF; I/O: 0x80000--0x8001F
    \item 4 chip SRAM 16K$\times$8-bit; RAM: 0x00000--0x0FFFF; ROM: 0xE0000--0xE7FFF; I/O: 0x80000--0x8003F
\end{enumerate}
\textbf{Hướng dẫn giải:}
\begin{itemize}
    \item RAM $64\text{ KB}$ dùng $4$ chip $16\text{K}\times 8$, chiếm $2^{16}$ địa chỉ khi $A_{19..16} = 0000$: \texttt{0x00000} -- \texttt{0x0FFFF}.
    \item ROM $32\text{ KB} = 2^{15}$ byte khi $A_{19..16} = 1111$ ($A_{15}=0$): \texttt{0xF0000} -- \texttt{0xF7FFF}.
    \item Khối $32$ cổng I/O khi $A_{19..16} = 1000$: \texttt{0x80000} -- \texttt{0x8001F} ($31 = \text{0x1F}$).
    \item \textbf{Kết luận:} Đáp án đúng là \textbf{A} (4 chip SRAM 16K$\times$8-bit; RAM: 0x00000--0x0FFFF; ROM: 0xF0000--0xF7FFF; I/O: 0x80000--0x8001F).
\end{itemize}

\textbf{Bài tập 6.13 (cuối kì 20222 (\#38a80d); cuối kì 20241 (\#cbb2f4); ôn tập cuối kì 20251 (Câu 57); ôn tập cuối kì 20252 (P2-63))}\\
Cho một máy tính có bus địa chỉ 32-bit và bus dữ liệu 32-bit. Bộ nhớ chính được đánh địa chỉ theo byte và được thiết kế theo kiểu băng nhớ đan xen. Hãy cho biết byte nhớ có địa chỉ 0x123456AD nằm trên băng nhớ nào?
\begin{enumerate}[label=\Alph*.]
    \item băng 1
    \item băng 0
    \item băng 2
    \item băng 3
\end{enumerate}
\textbf{Hướng dẫn giải:}
\begin{itemize}
    \item Bus dữ liệu $32\text{ bit} = 4\text{ byte}$ nên bộ nhớ chính đan xen được chia làm $4$ băng nhớ ($0, 1, 2, 3$) dựa theo $2$ bit thấp của địa chỉ (hay $\text{Địa chỉ} \bmod 4$).
    \item Chữ số Hexa cuối của \texttt{0x123456AD} là \texttt{0xD} $= 13_{(10)} = 1101_{(2)}$ (hoặc \texttt{0xAD} $= 173 = 4 \times 43 + 1$). Hai bit thấp là \texttt{01} $\Rightarrow$ nằm trên \textbf{băng 1}.
    \item \textbf{Kết luận:} Đáp án đúng là \textbf{A} (băng 1).
\end{itemize}

\textbf{Bài tập 6.14 (cuối kì 20231 (\#d5287a); cuối kì 20232 (\#d5287a); cuối kì 20251 (\#96837); thi thử kì 20251 (Câu 30); ôn tập cuối kì 20251 (Câu 6); ôn tập cuối kì 20252 (P2-71))}\\
Khẳng định nào dưới đây mô tả chính xác nhất về đặc tính cục bộ bộ nhớ của các biến khi thực thi đoạn mã C bên dưới?
\begin{lstlisting}
int i, sum = 0;
int A[100];
for (i = 0; i < 100; i++)
    sum += A[i];
\end{lstlisting}
\begin{enumerate}[label=\Alph*.]
    \item Truy cập vào i và sum thể hiện tính cục bộ thời gian (temporal locality); truy cập vào các phần tử A[i] thể hiện tính cục bộ không gian (spatial locality).
    \item Truy cập vào tất cả các biến đều thể hiện tính cục bộ thời gian (temporal locality).
    \item Truy cập vào tất cả các biến đều thể hiện tính cục bộ không gian (spatial locality).
    \item Truy cập vào i và sum thể hiện tính cục bộ không gian (spatial locality); truy cập vào A[i] thể hiện tính cục bộ thời gian (temporal locality).
\end{enumerate}
\textbf{Hướng dẫn giải:}
\begin{itemize}
    \item Biến \texttt{i} và \texttt{sum} được đọc/ghi lặp đi lặp lại trong vòng lặp $\Rightarrow$ \textbf{Temporal Locality} (cục bộ theo thời gian). Các phần tử \texttt{A[i]} nằm liên tiếp nhau trong bộ nhớ và được duyệt tuần tự $\Rightarrow$ \textbf{Spatial Locality} (cục bộ theo không gian).
    \item \textbf{Kết luận:} Đáp án đúng là \textbf{A} (Truy cập vào i và sum thể hiện tính cục bộ thời gian (temporal locality); truy cập vào các phần tử A[i] thể hiện tính cục bộ không gian (spatial locality).).
\end{itemize}

\textbf{Bài tập 6.15 (ôn tập cuối kì 20251 (Câu 45); ôn tập cuối kì 20252 (P2-27))}\\
Đặc điểm nào sau đây \textbf{KHÔNG PHẢI} là ưu điểm của ổ thể rắn SSD so với ổ đĩa từ HDD truyền thống?
\begin{enumerate}[label=\Alph*.]
    \item Giá thành trên mỗi đơn vị dung lượng (VNĐ/GB) rẻ hơn
    \item Tốc độ truy xuất dữ liệu nhanh hơn và độ trễ thấp hơn
    \item Không có bộ phận chuyển động cơ học nên không gây tiếng ồn
    \item Khả năng chống sốc vật lý và va đập tốt hơn
\end{enumerate}
\textbf{Hướng dẫn giải:}
\begin{itemize}
    \item SSD vượt trội HDD về tốc độ, độ trễ, không tiếng ồn và chống sốc; nhược điểm của SSD so với HDD là \textbf{chi phí trên mỗi GB cao hơn}.
    \item \textbf{Kết luận:} Đáp án đúng là \textbf{A} (Giá thành trên mỗi đơn vị dung lượng (VNĐ/GB) rẻ hơn).
\end{itemize}

\textbf{Bài tập 6.16 (ôn tập cuối kì 20251 (Câu 19); ôn tập cuối kì 20252 (P2-73))}\\
Trong thực tế hiện nay, phát biểu nào sau đây \textbf{SAI} về việc sử dụng HDD và SSD trên các thiết bị?
\begin{enumerate}[label=\Alph*.]
    \item Camera giám sát (CCTV) ghi hình 24/7 thường ưu tiên dùng SSD vì SSD có độ bền ghi/xóa vượt trội HDD khi ghi liên tục với dung lượng lớn
    \item Các máy chơi game thế hệ mới (như PS5, Xbox Series X) sử dụng SSD để giảm mạnh thời gian tải màn chơi
    \item SSD được ưu tiên làm ổ cài hệ điều hành trên Laptop/PC nhờ tốc độ truy xuất ngẫu nhiên nhanh
    \item HDD vẫn phổ biến trong hệ thống Server lưu trữ (NAS) hoặc trung tâm dữ liệu nhờ dung lượng lớn và giá thành trên GB rẻ
\end{enumerate}
\textbf{Hướng dẫn giải:}
\begin{itemize}
    \item Camera CCTV ghi hình liên tục 24/7 cần dung lượng lưu trữ lớn, chi phí/GB thấp và độ bền ghi tuần tự cao nên thực tế ưu tiên dùng \textbf{HDD chuyên dụng} thay vì SSD (SSD có giới hạn chu kỳ ghi/xóa tế bào nhớ flash TBW và giá thành cao hơn).
    \item \textbf{Kết luận:} Đáp án đúng là \textbf{A} (Camera giám sát (CCTV) ghi hình 24/7 thường ưu tiên dùng SSD vì SSD có độ bền ghi/xóa vượt trội HDD khi ghi liên tục với dung lượng lớn).
\end{itemize}

\subsection{Nhóm 2: Ánh xạ Cache, Tách bit địa chỉ (Tag/Index/Offset) \& Kích thước phần cứng}

\textbf{Bài tập 6.17 (cuối kì 20212 (B5-5))}\\
Cho máy tính có $64\text{ KB}$ bộ nhớ chính đánh địa chỉ theo byte, bộ nhớ Cache gồm $32$ line tổ chức ánh xạ trực tiếp, kích thước mỗi line là $4\text{ byte}$. Số bit của các trường \texttt{Tag / Line / Word} và dải địa chỉ Hexa của tất cả các byte nằm cùng line với byte nhớ có địa chỉ nhị phân \texttt{0010 0011 1001 0111} lần lượt là:
\begin{enumerate}[label=\Alph*.]
    \item Tag = 9 bit, Line = 5 bit, Word = 2 bit; cùng line: 0x2394 đến 0x2397
    \item Tag = 8 bit, Line = 5 bit, Word = 3 bit; cùng line: 0x2390 đến 0x2397
    \item Tag = 9 bit, Line = 4 bit, Word = 3 bit; cùng line: 0x2390 đến 0x2397
    \item Tag = 10 bit, Line = 4 bit, Word = 2 bit; cùng line: 0x2394 đến 0x2397
\end{enumerate}
\textbf{Hướng dẫn giải:}
\begin{itemize}
    \item \textbf{1. Phân chia bit địa chỉ:}
    \item Bộ nhớ chính $64\text{ KB} = 2^{16}\text{ byte} \Rightarrow$ địa chỉ dài $16\text{ bit}$.
    \item Kích thước line $= 4\text{ byte} = 2^2 \Rightarrow \text{Word} = \mathbf{2\text{ bit}}$.
    \item Số line $= 32 = 2^5 \Rightarrow \text{Line} = \mathbf{5\text{ bit}}$.
    \item $\text{Tag} = 16 - 5 - 2 = \mathbf{9\text{ bit}}$.
    \item \textbf{2. Các byte cùng line với \texttt{0010 0011 1001 0111}:}
    \item Chuyển sang Hexa là \texttt{0x2397}. Giữ nguyên $14$ bit cao (Tag + Line = \texttt{0010 0011 1001 01}) và cho $2$ bit Word chạy từ \texttt{00} đến \texttt{11}, ta thu được dải địa chỉ \textbf{\texttt{0x2394} đến \texttt{0x2397}}.
    \item \textbf{Kết luận:} Đáp án đúng là \textbf{A} (Tag = 9 bit, Line = 5 bit, Word = 2 bit; cùng line: 0x2394 đến 0x2397).
\end{itemize}

\textbf{Bài tập 6.18 (cuối kì 20212 (B7-5))}\\
Cho máy tính với bộ nhớ chính có địa chỉ $32\text{-bit}$ (đánh địa chỉ theo từng byte). Bộ nhớ Cache dung lượng $256\text{ KB}$ được tổ chức theo kiểu ánh xạ liên kết tập hợp $8\text{ line/set}$, gồm $512$ Set ($S_0\text{-}S_{511}$). Số bit các trường \texttt{Tag / Set / Word} và chỉ số Set của byte nhớ có địa chỉ \texttt{0x032A0203} lần lượt là:
\begin{enumerate}[label=\Alph*.]
    \item Tag = 16 bit, Set = 9 bit, Word = 7 bit; Set 4
    \item Tag = 17 bit, Set = 9 bit, Word = 6 bit; Set 8
    \item Tag = 17 bit, Set = 9 bit, Word = 6 bit; Set 16
    \item Tag = 18 bit, Set = 8 bit, Word = 6 bit; Set 8
\end{enumerate}
\textbf{Hướng dẫn giải:}
\begin{itemize}
    \item \textbf{1. Phân chia bit địa chỉ:}
    \item Tổng số line $= 512 \text{ Set} \times 8 \text{ line/set} = 4096 = 2^{12}$ line.
    \item Kích thước mỗi line $= 256\text{ KB} / 2^{12} = 2^{18} / 2^{12} = 64\text{ byte} = 2^6 \Rightarrow \text{Word} = \mathbf{6\text{ bit}}$.
    \item Số Set $= 512 = 2^9 \Rightarrow \text{Set} = \mathbf{9\text{ bit}}$.
    \item $\text{Tag} = 32 - 9 - 6 = \mathbf{17\text{ bit}}$.
    \item \textbf{2. Tìm Set của địa chỉ \texttt{0x032A0203}:}
    \item Các bit cuối: \texttt{...0203} $= \dots\texttt{0000 0010 0000 0011}_2$.
    \item Bỏ $6$ bit Word thấp nhất (\texttt{000011}), lấy $9$ bit Set tiếp theo: \texttt{000001000}$_2 = \mathbf{8}$ (\textbf{Set 8}).
    \item \textbf{Kết luận:} Đáp án đúng là \textbf{B} (Tag = 17 bit, Set = 9 bit, Word = 6 bit; Set 8).
\end{itemize}

\textbf{Bài tập 6.19 (cuối kì 20212 (B8-5))}\\
Cho máy tính có bộ nhớ chính dung lượng $4\text{ GB}$ (đánh địa chỉ theo từng byte). Bộ nhớ Cache dung lượng $256\text{ KB}$ tổ chức theo ánh xạ liên kết tập hợp $4\text{ line/set}$, gồm $512$ Set ($S_0\text{-}S_{511}$). Số bit các trường \texttt{Tag / Set / Word} và chỉ số Set của byte nhớ có địa chỉ \texttt{0x0A0B1035} lần lượt là:
\begin{enumerate}[label=\Alph*.]
    \item Tag = 17 bit, Set = 9 bit, Word = 6 bit; Set 64
    \item Tag = 16 bit, Set = 9 bit, Word = 7 bit; Set 16
    \item Tag = 16 bit, Set = 9 bit, Word = 7 bit; Set 32
    \item Tag = 16 bit, Set = 10 bit, Word = 6 bit; Set 32
\end{enumerate}
\textbf{Hướng dẫn giải:}
\begin{itemize}
    \item \textbf{1. Phân chia bit địa chỉ:}
    \item Bộ nhớ chính $4\text{ GB} = 2^{32}\text{ byte} \Rightarrow$ địa chỉ dài $32\text{ bit}$.
    \item Tổng số line $= 512 \times 4 = 2048 = 2^{11}$ line.
    \item Kích thước mỗi line $= 256\text{ KB} / 2^{11} = 2^{18} / 2^{11} = 128\text{ byte} = 2^7 \Rightarrow \text{Word} = \mathbf{7\text{ bit}}$.
    \item Số Set $= 512 = 2^9 \Rightarrow \text{Set} = \mathbf{9\text{ bit}}$.
    \item $\text{Tag} = 32 - 9 - 7 = \mathbf{16\text{ bit}}$.
    \item \textbf{2. Tìm Set của địa chỉ \texttt{0x0A0B1035}:}
    \item $16$ bit thấp là \texttt{0x1035} $= \texttt{0001 0000 0011 0101}_2$.
    \item Bỏ $7$ bit Word thấp nhất (\texttt{0110101}), lấy $9$ bit Set kế tiếp: \texttt{000100000}$_2 = 2^5 = \mathbf{32}$ (\textbf{Set 32}).
    \item \textbf{Kết luận:} Đáp án đúng là \textbf{C} (Tag = 16 bit, Set = 9 bit, Word = 7 bit; Set 32).
\end{itemize}

\textbf{Bài tập 6.20 (ôn tập cuối kì 20252 (P2-19); Đáp án ôn tập cuối kỳ P2 (Câu 23))}\\
Máy tính có bộ nhớ Cache $8\text{ MiB}$ được tổ chức theo ánh xạ liên kết tập hợp $16\text{ line/set}$ với kích thước mỗi line là $64\text{ byte}$; bộ nhớ chính được đánh địa chỉ cho từng byte bằng $36\text{ bit}$ địa chỉ. Giá trị biểu thức $P = \text{số bit trường Set} + 2 \times \text{số bit trường Tag}$ bằng bao nhiêu?
\begin{enumerate}[label=\Alph*.]
    \item 45
    \item 46
    \item 47
    \item 48
\end{enumerate}
\textbf{Hướng dẫn giải:}
\begin{itemize}
    \item Kích thước line $= 64\text{ byte} = 2^6 \Rightarrow \text{Word} = 6\text{ bit}$.
    \item Tổng số line $= 8\text{ MiB} / 64\text{ B} = 2^{23} / 2^6 = 2^{17}$ line.
    \item Số Set $= 2^{17} / 16 = 2^{17} / 2^4 = 2^{13} \Rightarrow \text{Set} = \mathbf{13\text{ bit}}$.
    \item $\text{Tag} = 36 - 13 - 6 = \mathbf{17\text{ bit}}$.
    \item Vậy $P = 13 + 2 \times 17 = \mathbf{47}$.
    \item \textbf{Kết luận:} Đáp án đúng là \textbf{C} (47).
\end{itemize}

\textbf{Bài tập 6.21 (cuối kì 20212 (B2-4))}\\
Cho máy tính có $64\text{ KB}$ bộ nhớ chính đánh địa chỉ theo byte, bộ nhớ Cache gồm $16$ line tổ chức ánh xạ trực tiếp, kích thước mỗi line là $8\text{ byte}$.\\
1) Số bit các trường \texttt{Tag / Line / Word} là bao nhiêu?\\
2) Byte nhớ có địa chỉ nhị phân \texttt{1100 0011 0011 0100} ánh xạ vào line nào?\\
3) Byte nhớ \texttt{0x1A15} nạp vào cache cùng những byte nào trong cùng line?
\begin{enumerate}[label=\Alph*.]
    \item Tag = 9 bit, Line = 4 bit, Word = 3 bit; Line 6; cùng line: 0x1A10 đến 0x1A17
    \item Tag = 7 bit, Line = 5 bit, Word = 4 bit; Line 3; cùng line: 0x1A10 đến 0x1A17
    \item Tag = 9 bit, Line = 4 bit, Word = 3 bit; Line 12; cùng line: 0x1A00 đến 0x1A0F
    \item Tag = 8 bit, Line = 4 bit, Word = 4 bit; Line 6; cùng line: 0x1A10 đến 0x1A1F
\end{enumerate}
\textbf{Hướng dẫn giải:}
\begin{itemize}
    \item 1) $64\text{ KB} = 2^{16}\text{ B}$ ($16$ bit địa chỉ). Line $8\text{ B} = 2^3 \Rightarrow \text{Word} = 3\text{ bit}$. $16$ line $= 2^4 \Rightarrow \text{Line} = 4\text{ bit}$. $\text{Tag} = 16 - 4 - 3 = 9\text{ bit}$.
    \item 2) Địa chỉ \texttt{1100 0011 0011 0100} có 4 bit $[6:3] = 0110_2 = \text{Line } 6$.
    \item 3) Line $8\text{ byte}$ nên các byte cùng block với \texttt{0x1A15} có 3 bit cuối thay đổi từ $000$ đến $111$: \texttt{0x1A10} -- \texttt{0x1A17}.
    \item \textbf{Kết luận:} Đáp án đúng là \textbf{A} (Tag = 9 bit, Line = 4 bit, Word = 3 bit; Line 6; cùng line: 0x1A10 đến 0x1A17).
\end{itemize}

\textbf{Bài tập 6.22 (cuối kì 20241 (\#2e789d); ôn tập cuối kì 20251 (Câu 20); ôn tập cuối kì 20251 (Câu 22); ôn tập cuối kì 20252 (P2-38))}\\
Chỉ ra phương pháp ánh xạ mà mỗi Block của bộ nhớ chính được phép nạp vào Line bất kỳ của bộ nhớ cache ?
\begin{enumerate}[label=\Alph*.]
    \item Ánh xạ liên kết toàn phần (Fully Associative)
    \item Ánh xạ trực tiếp (Direct Mapped)
    \item Ánh xạ liên kết tập hợp (Set Associative)
    \item Không có phương pháp nào cho phép điều này
\end{enumerate}
\textbf{Hướng dẫn giải:}
\begin{itemize}
    \item Trong \textbf{ánh xạ liên kết toàn phần (Fully Associative)}, một block từ bộ nhớ chính có thể được nạp vào \textbf{bất kỳ line nào} còn trống (hoặc thay thế theo thuật toán) trong Cache.
    \item \textbf{Kết luận:} Đáp án đúng là \textbf{A} (Ánh xạ liên kết toàn phần (Fully Associative)).
\end{itemize}

\textbf{Bài tập 6.23 (cuối kì 20251 (\#05f48a); thi thử kì 20251 (Câu 10); ôn tập cuối kì 20251 (Câu 46); ôn tập cuối kì 20252 (P2-37))}\\
Cho một CPU với địa chỉ bộ nhớ 32 bit A[31..0]. Nó có bộ nhớ đệm (cache) thiết kế theo kiểu tập liên kết 4 đường (4-way set associative) với tổng kích thước dữ liệu cache là 16 KiB và kích thước khối (block size) là 64 bytes. Kích thước (tính bằng bit) của bộ nhớ đệm này là bao nhiêu?
\begin{enumerate}[label=\Alph*.]
    \item 64 $\times$ 4 $\times$ (1 + 20 + 32 $\times$ 16)
    \item 64 $\times$ 2 $\times$ (1 + 20 + 32 $\times$ 16)
    \item 64 $\times$ 2 $\times$ (1 + 20 + 32 $\times$ 8)
    \item 64 $\times$ 4 $\times$ (1 + 19 + 16 $\times$ 16)
\end{enumerate}
\textbf{Hướng dẫn giải:}
\begin{itemize}
    \item Tổng số block $= 16\text{ KiB} / 64\text{ B} = 256$ block.
    \item Cache $4\text{-way} \Rightarrow 256 / 4 = \mathbf{64}$ Set ($6$ bit Index), mỗi Set có $\mathbf{4}$ line.
    \item Block $64\text{ B} = 2^6 \Rightarrow 6$ bit Offset; dữ liệu mỗi block $= 64 \times 8 = 512\text{ bit} = \mathbf{32 \times 16}$ bit.
    \item $\text{Tag} = 32 - 6 - 6 = \mathbf{20}$ bit, cộng $1$ bit hợp lệ (Valid).
    \item Tổng kích thước Cache $= \mathbf{64 \times 4 \times (1 + 20 + 32 \times 16)}$ bit.
    \item \textbf{Kết luận:} Đáp án đúng là \textbf{A} (64 $\times$ 4 $\times$ (1 + 20 + 32 $\times$ 16)).
\end{itemize}

\textbf{Bài tập 6.24 (cuối kì 20251 (\#c38eef); thi thử kì 20251 (Câu 17); ôn tập cuối kì 20251 (Câu 28); ôn tập cuối kì 20252 (P2-7))}\\
Cho một CPU với địa chỉ bộ nhớ 32-bit A[31..0] và bộ nhớ đệm (cache) ánh xạ trực tiếp (direct-mapped) dung lượng 16 KiB với kích thước khối (block size) là 32 bytes. Khi truy cập một từ (word) trong bộ nhớ, những bit địa chỉ nào sẽ được sử dụng để xác định khối cache tương ứng (locate the corresponding cache block)?
\begin{enumerate}[label=\Alph*.]
    \item A[13..5]
    \item A[13..6]
    \item A[14..6]
    \item A[14..5]
\end{enumerate}
\textbf{Hướng dẫn giải:}
\begin{itemize}
    \item Block $32\text{ B} = 2^5 \Rightarrow$ Offset dùng $5$ bit thấp $A[4..0]$.
    \item Số block Cache $= 16\text{ KiB} / 32\text{ B} = 512 = 2^9 \Rightarrow$ Index cần $9$ bit kế tiếp là $\mathbf{A[13..5]}$.
    \item \textbf{Kết luận:} Đáp án đúng là \textbf{A} (A[13..5]).
\end{itemize}

\textbf{Bài tập 6.25 (ôn tập cuối kì 20251 (Câu 17); ôn tập cuối kì 20252 (P2-22))}\\
Đề bài: Cho một CPU có không gian địa chỉ bộ nhớ 32-bit và bộ nhớ đệm (cache) ánh xạ trực tiếp (direct-mapped) dung lượng 16 KiB với kích thước khối (block size) là 64 bytes.
Câu hỏi: Tổng kích thước (tính bằng bit) của toàn bộ bộ nhớ đệm, bao gồm các trường: bit hợp lệ (valid bit), thẻ (tag) và dữ liệu (data) là bao nhiêu?
\begin{enumerate}[label=\Alph*.]
    \item 256 $\times$ (1 + 18 + 32 $\times$ 16)
    \item 256 $\times$ (1 + 19 + 32 $\times$ 16)
    \item 256 $\times$ (1 + 18 + 32 $\times$ 8)
    \item 256 $\times$ (1 + 20 + 32 $\times$ 16)
\end{enumerate}
\textbf{Hướng dẫn giải:}
\begin{itemize}
    \item Số block $= 16\text{ KiB} / 64\text{ B} = \mathbf{256}$ ($8$ bit Index). Offset $= 6$ bit ($64\text{ B} = 512\text{ bit} = \mathbf{32 \times 16}$ bit). $\text{Tag} = 32 - 8 - 6 = \mathbf{18}$ bit $+\; 1$ Valid bit $\Rightarrow \mathbf{256 \times (1 + 18 + 32 \times 16)}$ bit.
    \item \textbf{Kết luận:} Đáp án đúng là \textbf{A} (256 $\times$ (1 + 18 + 32 $\times$ 16)).
\end{itemize}

\textbf{Bài tập 6.26 (ôn tập cuối kì 20251 (Câu 34); ôn tập cuối kì 20252 (P2-58))}\\
Máy tính đánh địa chỉ theo byte với không gian bộ nhớ chính $4\text{ GiB}$, Cache $256\text{ KiB}$, kích thước Line $= 32\text{ byte}$. Tính tổng $P = \text{Tag}_{\text{trực tiếp}} + \text{Tag}_{\text{liên kết toàn phần}} + \text{Tag}_{\text{tập hợp 4 đường}}$?
\textit{(Dạng điền đáp án -- Nhập giá trị P (vd: 57):)}\par
\textbf{Hướng dẫn giải:}
\begin{itemize}
    \item Địa chỉ $32$ bit, Offset $= \log_2(32) = 5$ bit, tổng số line $= 256\text{ KiB}/32\text{ B} = 8192 = 2^{13}$:
    \item Direct-mapped: $\text{Index} = 13 \Rightarrow \text{Tag}_1 = 32 - 13 - 5 = 14$.
    \item Fully-associative: $\text{Index} = 0 \Rightarrow \text{Tag}_2 = 32 - 5 = 27$.
    \item 4-way Set-associative: $2^{13}/4 = 2^{11}$ Set $\Rightarrow \text{Tag}_3 = 32 - 11 - 5 = 16$.
    \item Tổng $P = 14 + 27 + 16 = \mathbf{57}$.
    \item \textbf{Kết luận:} Đáp án đúng là \textbf{57}.
\end{itemize}

\textbf{Bài tập 6.27 (ôn tập cuối kì 20251 (Câu 50); ôn tập cuối kì 20252 (P2-55))}\\
Máy tính có địa chỉ bộ nhớ chính $32\text{-bit}$ (đánh địa chỉ theo byte). Bộ nhớ Cache $128\text{ KiB}$ tổ chức theo ánh xạ liên kết tập hợp $4\text{ line/set}$, gồm $1024$ Set ($S_0\text{-}S_{1023}$). Số bit của trường \texttt{Tag} là bao nhiêu?
\textit{(Dạng điền đáp án -- Nhập số bit trường Tag (vd: 17):)}\par
\textbf{Hướng dẫn giải:}
\begin{itemize}
    \item Số line $= 1024 \times 4 = 4096$. Kích thước 1 line $= 128\text{ KiB}/4096 = 32\text{ B} \Rightarrow \text{Offset} = 5\text{ bit}$.
    \item $1024\text{ Set} = 2^{10} \Rightarrow \text{Set} = 10\text{ bit}$.
    \item $\text{Tag} = 32 - 10 - 5 = \mathbf{17}$ bit.
    \item \textbf{Kết luận:} Đáp án đúng là \textbf{17}.
\end{itemize}

\textbf{Bài tập 6.28 (cuối kì 20212 (B3-5))}\\
Cho máy tính có $64\text{ KB}$ bộ nhớ chính đánh địa chỉ theo byte, bộ nhớ Cache gồm $32$ line ánh xạ trực tiếp, kích thước mỗi line là $16\text{ byte}$.\\
1) Số bit các trường \texttt{Tag / Line / Word} là bao nhiêu?\\
2) Bốn địa chỉ \texttt{1101 0001 0001 0010}, \texttt{0010 1011 0011 0110}, \texttt{1101 0110 1001 1100}, \texttt{1110 1010 1000 1011} lần lượt nạp vào các line nào?
\begin{enumerate}[label=\Alph*.]
    \item Tag = 7 bit, Line = 5 bit, Word = 4 bit; Line 1, Line 3, Line 9, Line 8
    \item Tag = 9 bit, Line = 4 bit, Word = 3 bit; Line 2, Line 6, Line 18, Line 17
    \item Tag = 7 bit, Line = 5 bit, Word = 4 bit; Line 1, Line 6, Line 9, Line 8
    \item Tag = 8 bit, Line = 4 bit, Word = 4 bit; Line 1, Line 3, Line 9, Line 10
\end{enumerate}
\textbf{Hướng dẫn giải:}
\begin{itemize}
    \item 1) $64\text{ KB}$ ($16$ bit địa chỉ), line $16\text{ B} = 2^4 \Rightarrow \text{Word} = 4\text{ bit}$, $32$ line $= 2^5 \Rightarrow \text{Line} = 5\text{ bit}$, $\text{Tag} = 16 - 5 - 4 = 7\text{ bit}$.
    \item 2) Lấy 5 bit $[8:4]$ làm chỉ số Line: \texttt{00001} ($1$), \texttt{00011} ($3$), \texttt{01001} ($9$), \texttt{01000} ($8$).
    \item \textbf{Kết luận:} Đáp án đúng là \textbf{A} (Tag = 7 bit, Line = 5 bit, Word = 4 bit; Line 1, Line 3, Line 9, Line 8).
\end{itemize}

\textbf{Bài tập 6.29 (ôn tập cuối kì 20251 (Câu 74); ôn tập cuối kì 20252 (P2-61))}\\
Máy tính có địa chỉ bộ nhớ chính $32\text{-bit}$ (đánh địa chỉ theo byte). Bộ nhớ Cache $128\text{ KiB}$ tổ chức theo ánh xạ liên kết tập hợp $4\text{ line/set}$, gồm $1024$ Set ($S_0\text{-}S_{1023}$). Byte nhớ có địa chỉ \texttt{0x003D02AF} được ánh xạ vào Set số mấy?
\textit{(Dạng điền đáp án -- Nhập chỉ số Set (vd: 21):)}\par
\textbf{Hướng dẫn giải:}
\begin{itemize}
    \item Block $= 32\text{ B}$ ($5$ bit Offset $[4:0]$), $1024$ Set ($10$ bit Set $[14:5]$).
    \item \texttt{0x003D02AF} có 16 bit thấp $= \text{0x02AF} = 0000\;0010\;1010\;1111_2$.
    \item 10 bit $[14:5]$ là \texttt{0000010101} $= \text{0x15} = \mathbf{21}$.
    \item \textbf{Kết luận:} Đáp án đúng là \textbf{21}.
\end{itemize}

\textbf{Bài tập 6.30 (ôn tập cuối kì 20252 (P2-3))}\\
Chỉ ra phương pháp ánh xạ mà mỗi Block của bộ nhớ chính chỉ được phép nạp vào \textbf{duy nhất 1 Line} cố định của bộ nhớ Cache?
\begin{enumerate}[label=\Alph*.]
    \item Ánh xạ trực tiếp (Direct Mapped)
    \item Ánh xạ liên kết toàn phần (Fully Associative)
    \item Ánh xạ liên kết tập hợp (Set Associative)
    \item Cả ba phương pháp trên
\end{enumerate}
\textbf{Hướng dẫn giải:}
\begin{itemize}
    \item Trong \textbf{ánh xạ trực tiếp}, mỗi Block bộ nhớ chính ánh xạ cố định vào duy nhất $1$ Line xác định bởi công thức $\text{Line} = \text{Block} \bmod \text{Số lượng Line}$.
    \item \textbf{Kết luận:} Đáp án đúng là \textbf{A} (Ánh xạ trực tiếp (Direct Mapped)).
\end{itemize}

\textbf{Bài tập 6.31 (cuối kì 20212 (B4-5); ôn tập cuối kì 20251 (Câu 65); ôn tập cuối kì 20252 (P2-24))}\\
Cho máy tính có bộ nhớ chính $4\text{ GB}$ (đánh địa chỉ theo byte), bộ nhớ Cache $512\text{ KB}$ tổ chức theo ánh xạ liên kết tập hợp $8\text{ line/set}$ và có tất cả $1024$ Set ($S_0\text{-}S_{1023}$). Số bit của các trường \texttt{Tag / Set / Word} và chỉ số Set của byte nhớ \texttt{0x10250A03} lần lượt là:
\begin{enumerate}[label=\Alph*.]
    \item Tag = 16 bit, Set = 10 bit, Word = 6 bit; ánh xạ vào Set 40
    \item Tag = 15 bit, Set = 10 bit, Word = 7 bit; ánh xạ vào Set 20
    \item Tag = 16 bit, Set = 10 bit, Word = 6 bit; ánh xạ vào Set 25
    \item Tag = 17 bit, Set = 10 bit, Word = 5 bit; ánh xạ vào Set 80
\end{enumerate}
\textbf{Hướng dẫn giải:}
\begin{itemize}
    \item Bộ nhớ $4\text{ GB} \Rightarrow 32$ bit địa chỉ. Kích thước mỗi line $= \frac{512\text{ KB}}{1024 \times 8} = 64\text{ B} = 2^6 \Rightarrow \text{Word} = 6\text{ bit}$.
    \item $1024\text{ Set} = 2^{10} \Rightarrow \text{Set} = 10\text{ bit}$. $\text{Tag} = 32 - 10 - 6 = 16\text{ bit}$.
    \item 10 bit $[15:6]$ của \texttt{0x10250A03} (tức \texttt{0x0A03 >> 6} $= 2563 // 64$) $= \text{0x028} = \mathbf{40}$ (Set $40$).
    \item \textbf{Kết luận:} Đáp án đúng là \textbf{A} (Tag = 16 bit, Set = 10 bit, Word = 6 bit; ánh xạ vào Set 40).
\end{itemize}

\textbf{Bài tập 6.32 (cuối kì 20212 (B6-5))}\\
Máy tính có bộ nhớ chính $4\text{ GB}$ đánh địa chỉ theo byte, bộ nhớ Cache tổ chức theo ánh xạ liên kết toàn phần (Fully-Associative), kích thước mỗi line là $32\text{ byte}$. Số bit trường \texttt{Tag / Word} và giá trị Tag (thập phân) của hai địa chỉ \texttt{0x0149C003}, \texttt{0x10390135} lần lượt là:
\begin{enumerate}[label=\Alph*.]
    \item Tag = 27 bit, Word = 5 bit; Tag = 674816 và 8507401
    \item Tag = 26 bit, Word = 6 bit; Tag = 337408 và 4253700
    \item Tag = 27 bit, Word = 5 bit; Tag = 1349632 và 17014802
    \item Tag = 22 bit, Word = 5 bit; Tag = 674816 và 8507401
\end{enumerate}
\textbf{Hướng dẫn giải:}
\begin{itemize}
    \item Liên kết toàn phần không có trường Index/Set: $32\text{ B} = 2^5 \Rightarrow \text{Word} = 5\text{ bit}$, $\text{Tag} = 32 - 5 = 27\text{ bit}$.
    \item Giá trị Tag $= \text{addr} \gg 5$: \texttt{0x0149C003 >> 5 = 0x000A4E00 = 674816}; \texttt{0x10390135 >> 5 = 0x0081C809 = 8507401}.
    \item \textbf{Kết luận:} Đáp án đúng là \textbf{A} (Tag = 27 bit, Word = 5 bit; Tag = 674816 và 8507401).
\end{itemize}

\textbf{Bài tập 6.33 (cuối kì 20222 (\#c9ea55); cuối kì 20241 (\#7196b2); cuối kì 20242 (\#c9ea55); ôn tập cuối kì 20251 (Câu 9); ôn tập cuối kì 20251 (Câu 37); ôn tập cuối kì 20252 (P2-54))}\\
Máy tính có bộ nhớ Cache $8\text{ MiB}$ tổ chức theo ánh xạ liên kết tập hợp $16\text{ line/set}$, kích thước mỗi line là $128\text{ byte}$; bộ nhớ chính đánh địa chỉ theo từng byte bằng $36\text{ bit}$ địa chỉ. Số bit của trường \texttt{Set} và trường \texttt{Tag} lần lượt là bao nhiêu?
\begin{enumerate}[label=\Alph*.]
    \item Set = 12 bit và Tag = 17 bit
    \item Set = 13 bit và Tag = 16 bit
    \item Set = 12 bit và Tag = 16 bit
    \item Set = 13 bit và Tag = 17 bit
\end{enumerate}
\textbf{Hướng dẫn giải:}
\begin{itemize}
    \item Kích thước line $= 128\text{ B} = 2^7\text{ B} \Rightarrow \text{Offset (Word)} = 7\text{ bit}$.
    \item Tổng số line $= \frac{8\text{ MiB}}{128\text{ B}} = \frac{2^{23}}{2^7} = 2^{16}$ line.
    \item Mỗi Set gồm $16 = 2^4$ line $\Rightarrow$ Số Set $= \frac{2^{16}}{2^4} = 2^{12} \Rightarrow \mathbf{\text{Set} = 12\text{ bit}}$.
    \item $\mathbf{\text{Tag} = 36 - 12 - 7 = 17\text{ bit}}$.
    \item \textbf{Kết luận:} Đáp án đúng là \textbf{A} (Set = 12 bit và Tag = 17 bit).
\end{itemize}

\textbf{Bài tập 6.34 (cuối kì 20231 (\#6bc3f0); ôn tập cuối kì 20251 (Câu 33); ôn tập cuối kì 20252 (P2-44); cuối kì 20232 (\#6bc3f0); cuối kì 20251 (\#3c3218); thi thử kì 20251 (Câu 19))}\\
Cho một CPU hệ RISC-V với không gian địa chỉ bộ nhớ $32\text{-bit}$ và bộ nhớ Cache ánh xạ trực tiếp (direct-mapped) dung lượng $16\text{ KiB}$ với kích thước khối (block size) là $64\text{ byte}$. Kích thước (tính bằng bit) của trường \texttt{Tag} trong mỗi khối cache là bao nhiêu?
\textit{(Dạng điền đáp án -- Nhập số bit trường Tag (vd: 18):)}\par
\textbf{Hướng dẫn giải:}
\begin{itemize}
    \item Offset $= \log_2(64) = 6\text{ bit}$.
    \item Số block $= 16\text{ KiB} / 64\text{ B} = 256 = 2^8 \Rightarrow \text{Index} = 8\text{ bit}$.
    \item Vậy $\text{Tag} = 32 - 8 - 6 = \mathbf{18}$ bit.
    \item \textbf{Kết luận:} Đáp án đúng là \textbf{18}.
\end{itemize}

\subsection{Nhóm 3: Chính sách ghi, Thay thế, Lỗi 3C \& Hiệu năng Cache (AMAT / CPI)}

\textbf{Bài tập 6.35 (Lý thuyết slide CA.RISCV.2025 - Mục 6.3 (Write-Through \& Write-Back))}\\
Khi xảy ra \textbf{Cache Hit} trong thao tác ghi dữ liệu, đặc điểm nào sau đây mô tả đúng chính sách \textbf{Ghi trả sau (Write-Back)}?
\begin{enumerate}[label=\Alph*.]
    \item Dữ liệu luôn được ghi đồng thời vào cả Cache và bộ nhớ chính ngay lập tức để đảm bảo nhất quán
    \item Dữ liệu chỉ được ghi vào Cache (và đánh dấu bằng bit Dirty); khối dữ liệu chỉ được ghi trả về bộ nhớ chính khi khối đó bị thay thế ra khỏi Cache
    \item Dữ liệu chỉ được ghi trực tiếp vào bộ nhớ chính mà không cập nhật vào Cache
    \item Tốc độ ghi chậm hơn Write-Through và gây lưu lượng cao trên bus bộ nhớ
\end{enumerate}
\textbf{Hướng dẫn giải:}
\begin{itemize}
    \item Theo Slide 52 (CA.RISCV.2025):
    \item \textbf{Ghi xuyên qua (Write-Through):} Ghi đồng thời vào cả Cache và bộ nhớ chính $\Rightarrow$ đơn giản, nhất quán nhưng tốc độ chậm và tốn băng thông bus.
    \item \textbf{Ghi trả sau (Write-Back):} Chỉ ghi vào Cache (đánh dấu bit \texttt{Dirty}), tốc độ rất nhanh; chỉ khi block đó bị thay thế khỏi Cache mới ghi trả toàn bộ block về bộ nhớ chính.
    \item \textbf{Kết luận:} Đáp án đúng là \textbf{B} (Dữ liệu chỉ được ghi vào Cache (và đánh dấu bằng bit Dirty); khối dữ liệu chỉ được ghi trả về bộ nhớ chính khi khối đó bị thay thế ra khỏi Cache).
\end{itemize}

\textbf{Bài tập 6.36 (Lý thuyết slide CA.RISCV.2025 - Mục 6.3 (Phân loại lỗi Cache 3C))}\\
Trong mô hình phân loại lỗi trượt Cache \textbf{3C} (Compulsory, Capacity, Conflict), giải pháp phần cứng nào sau đây trực tiếp giúp giảm \textbf{lỗi xung đột (Conflict Miss)} mà không nhất thiết phải tăng tổng dung lượng Cache?
\begin{enumerate}[label=\Alph*.]
    \item Tăng bậc liên kết tập hợp (associativity), ví dụ chuyển từ ánh xạ trực tiếp sang liên kết tập hợp 4 đường hoặc 8 đường
    \item Chuyển từ ánh xạ liên kết tập hợp sang ánh xạ trực tiếp (direct-mapped)
    \item Giảm kích thước của mỗi block (line) xuống mức tối thiểu
    \item Sử dụng chính sách Write-Through thay cho Write-Back
\end{enumerate}
\textbf{Hướng dẫn giải:}
\begin{itemize}
    \item Trong mô hình \textbf{3C Misses}:
    \item \textbf{Compulsory (Cold) Miss:} Lần truy cập đầu tiên vào block $\Rightarrow$ giảm bằng cách tăng kích thước block (tận dụng tính cục bộ không gian) hoặc prefetching.
    \item \textbf{Capacity Miss:} Dung lượng Cache không đủ chứa toàn bộ tập làm việc $\Rightarrow$ giảm bằng cách tăng dung lượng Cache.
    \item \textbf{Conflict Miss:} Nhiều block ánh xạ vào cùng một vị trí/Set dù Cache vẫn còn chỗ trống ở Set khác $\Rightarrow$ giảm trực tiếp bằng cách \textbf{tăng bậc liên kết (associativity)}.
    \item \textbf{Kết luận:} Đáp án đúng là \textbf{A} (Tăng bậc liên kết tập hợp (associativity), ví dụ chuyển từ ánh xạ trực tiếp sang liên kết tập hợp 4 đường hoặc 8 đường).
\end{itemize}

\textbf{Bài tập 6.37 (thi thử kì 20251 (Câu 9); ôn tập cuối kì 20251 (Câu 16); ôn tập cuối kì 20252 (P2-12))}\\
Cho đoạn mã C chạy trên CPU có bộ nhớ đệm liên kết toàn phần (fully associative cache) dung lượng $16\text{ KiB}$, kích thước khối $64\text{ byte}$ (không có xung đột khác). Tỷ lệ trượt trung bình ($\%$) khi truy cập cả hai mảng \texttt{A} và \texttt{B} là bao nhiêu? (Chỉ nhập số, vd: \texttt{9.375})
\begin{lstlisting}
int i, sum = 0;
int A[128];
int B[256];
for (i = 0; i < 128; i++)
    B[i*2] = A[i];
\end{lstlisting}
\textit{(Dạng điền đáp án -- Nhập \% miss rate (vd: 9.375):)}\par
\textbf{Hướng dẫn giải:}
\begin{itemize}
    \item Block $64\text{ B}$ chứa $16$ số \texttt{int} ($4\text{ B/int}$):
    \item Mảng \texttt{A[i]} ($128$ lần truy cập liên tiếp, stride 4B): miss $128/16 = 8$ lần.
    \item Mảng \texttt{B[i*2]} ($128$ lần truy cập cách quãng 8B): miss $128/8 = 16$ lần.
    \item Tổng miss $= 8 + 16 = 24$ trên tổng $256$ lần truy cập $\Rightarrow \frac{24}{256} = \mathbf{9.375\%}$.
    \item \textbf{Kết luận:} Đáp án đúng là \textbf{9.375}.
\end{itemize}

\textbf{Bài tập 6.38 (ôn tập cuối kì 20251 (Câu 30); ôn tập cuối kì 20252 (P2-56))}\\
Cho CPU có tỉ lệ trượt cache lệnh (I-cache miss rate) $= 2\%$, tỉ lệ trượt cache dữ liệu (D-cache miss rate) $= 4\%$, thời gian phạt (miss penalty) $= 100$ chu kỳ, $\text{CPI}_{\text{base}}$ (cache lý tưởng) $= 2$, tỉ lệ lệnh Load/Store chiếm $36\%$. Tỉ số $\text{Actual CPI} / \text{Base CPI}$ bằng bao nhiêu?
\textit{(Dạng điền đáp án -- Nhập tỉ số (vd: 2.72):)}\par
\textbf{Hướng dẫn giải:}
\begin{itemize}
    \item $\text{Actual CPI} = 2 + (0.02 + 0.36 \times 0.04) \times 100 = 2 + 2 + 1.44 = 5.44$.
    \item Tỉ số $\frac{5.44}{2} = \mathbf{2.72}$.
    \item \textbf{Kết luận:} Đáp án đúng là \textbf{2.72}.
\end{itemize}

\textbf{Bài tập 6.39 (ôn tập cuối kì 20251 (Câu 55); ôn tập cuối kì 20252 (P2-47))}\\
Phương pháp thay thế block trong cache nào là tối ưu nhất?
\begin{enumerate}[label=\Alph*.]
    \item LRU (Least Recently Used)
    \item LFU (Least Frequently Used)
    \item Random
    \item FIFO (First-In First-Out)
\end{enumerate}
\textbf{Hướng dẫn giải:}
\begin{itemize}
    \item \textbf{LRU} (thay thế khối ít được sử dụng nhất gần đây) tận dụng tốt nhất nguyên lý cục bộ thời gian nên cho tỉ lệ hit cao nhất trong các chiến lược thông dụng.
    \item \textbf{Kết luận:} Đáp án đúng là \textbf{A} (LRU (Least Recently Used)).
\end{itemize}

\textbf{Bài tập 6.40 (ôn tập cuối kì 20251 (Câu 73); ôn tập cuối kì 20252 (P2-30))}\\
Một CPU có chu kỳ xung nhịp $1\text{ ns}$, thời gian trúng (hit time) $= 1$ chu kỳ, thời gian phạt trượt (miss penalty) $= 20$ chu kỳ, tỉ lệ trượt I-cache miss rate $= 5\%$. Thời gian truy cập bộ nhớ trung bình (AMAT) bằng bao nhiêu nano-giây?
\textit{(Dạng điền đáp án -- Nhập AMAT theo ns (vd: 2):)}\par
\textbf{Hướng dẫn giải:}
\begin{itemize}
    \item $\text{AMAT} = \text{Hit Time} + \text{Miss Rate} \times \text{Miss Penalty} = 1 + 0.05 \times 20 = 2\text{ chu kỳ} = \mathbf{2\text{ ns}}$.
    \item \textbf{Kết luận:} Đáp án đúng là \textbf{2}.
\end{itemize}

\textbf{Bài tập 6.41 (cuối kì 20231 (\#c92c74); cuối kì 20232 (\#c92c74); cuối kì 20251 (\#bf6243); cuối kì 20251 (\#24cd84); thi thử kì 20251 (Câu 2); ôn tập cuối kì 20251 (Câu 8); ôn tập cuối kì 20252 (P2-57))}\\
Cho đoạn mã C chạy trên CPU RISC-V với bộ nhớ Cache ánh xạ trực tiếp dung lượng $1\text{ KiB}$ và kích thước khối $64\text{ byte}$ (giả sử không có xung đột chương trình khác). Tỷ lệ trượt (miss rate, đơn vị $\%$) khi truy cập mảng \texttt{A} là bao nhiêu? (Chỉ nhập số, vd: \texttt{10.5})
\begin{lstlisting}
int i, sum = 0;
int A[128];
for (i = 0; i < 128; i++)
    sum += A[i];
\end{lstlisting}
\textit{(Dạng điền đáp án -- Nhập \% miss rate (vd: 6.25):)}\par
\textbf{Hướng dẫn giải:}
\begin{itemize}
    \item Mỗi phần tử \texttt{int} chiếm $4\text{ byte}$. Một block $64\text{ B}$ chứa được $64 / 4 = 16$ phần tử liên tiếp.
    \item Khi duyệt tuần tự \texttt{A[0..127]}, mỗi nhóm $16$ phần tử bị miss lần đầu ($1$ lần) và hit $15$ lần sau $\Rightarrow \text{Miss Rate} = \frac{1}{16} = \mathbf{6.25\%}$.
    \item \textbf{Kết luận:} Đáp án đúng là \textbf{6.25}.
\end{itemize}

\textbf{Bài tập 6.42 (cuối kì 20231 (\#c7a320); cuối kì 20232 (\#c7a320); ôn tập cuối kì 20251 (Câu 32); ôn tập cuối kì 20252 (P2-36))}\\
Cho một CPU RISC-V có chỉ số CPI cơ bản (base CPI) bằng $1$, thời gian phạt khi trượt cache (miss penalty) là $100$ chu kỳ. Khi thực thi chương trình, tỷ lệ trượt cache lệnh (instruction miss rate) là $1\%$, tỷ lệ trượt cache dữ liệu (data miss rate) là $10\%$, và các lệnh truy nhập dữ liệu (load/store) chiếm $5\%$ tổng số lệnh. Chỉ số CPI hiệu dụng (effective CPI) là bao nhiêu?
\textit{(Dạng điền đáp án -- Nhập effective CPI (vd: 2.5):)}\par
\textbf{Hướng dẫn giải:}
\begin{itemize}
    \item $\text{CPI}_{\text{eff}} = \text{CPI}_{\text{base}} + \text{I-miss stall} + \text{D-miss stall} = 1 + (0.01 \times 100) + (0.05 \times 0.10 \times 100) = 1 + 1 + 0.5 = \mathbf{2.5}$.
    \item \textbf{Kết luận:} Đáp án đúng là \textbf{2.5}.
\end{itemize}

\subsection{Nhóm 4: Bộ nhớ ảo (Virtual Memory), Phân trang (Paging) \& TLB}

\textbf{Bài tập 6.43 (Lý thuyết slide CA.RISCV.2025 - Mục 6.5 (Bảng trang đảo))}\\
Ưu điểm cốt lõi của cấu trúc \textbf{Bảng trang đảo (Inverted Page Table)} so với bảng trang truyền thống trong bộ nhớ ảo là gì?
\begin{enumerate}[label=\Alph*.]
    \item Mỗi tiến trình có một bảng trang riêng có kích thước tỉ lệ thuận với toàn bộ không gian địa chỉ ảo
    \item Không cần sử dụng bất kỳ cơ chế băm (hash) hay tìm kiếm nào khi dịch địa chỉ
    \item Chỉ duy trì đúng một mục cho mỗi khung trang vật lý (physical frame) trong RAM và dùng hàm băm để tra cứu, giúp giảm mạnh dung lượng bộ nhớ cần thiết để lưu bảng trang
    \item Loại bỏ hoàn toàn hiện tượng lỗi trang (Page Fault) trên hệ thống
\end{enumerate}
\textbf{Hướng dẫn giải:}
\begin{itemize}
    \item Theo Slide 70 (CA.RISCV.2025):
    \item Trong \textbf{Bảng trang đảo (Inverted Page Table)}, thay vì lưu bảng trang theo số trang ảo của từng tiến trình (rất lớn với địa chỉ 32/64-bit), hệ thống chỉ lưu \textbf{một mục duy nhất cho mỗi khung trang vật lý} trong RAM (kèm Process ID, Page \#, Control bits và chuỗi liên kết Chain) và sử dụng \textbf{hàm băm (Hash function)} để tra cứu, giúp tiết kiệm đáng kể dung lượng bộ nhớ lưu bảng trang.
    \item \textbf{Kết luận:} Đáp án đúng là \textbf{C} (Chỉ duy trì đúng một mục cho mỗi khung trang vật lý (physical frame) trong RAM và dùng hàm băm để tra cứu, giúp giảm mạnh dung lượng bộ nhớ cần thiết để lưu bảng trang).
\end{itemize}

\textbf{Bài tập 6.44 (Lý thuyết slide CA.RISCV.2025 - Mục 6.5 (Bộ nhớ ảo \& Phân trang))}\\
Trong kỹ thuật \textbf{phân trang (Paging)} của hệ thống bộ nhớ ảo (Virtual Memory), khẳng định nào sau đây là \textbf{ĐÚNG}?
\begin{enumerate}[label=\Alph*.]
    \item Các trang (Page) của cùng một tiến trình bắt buộc phải được cấp phát vào các khung trang (Frame) nằm liên tiếp nhau trong bộ nhớ vật lý
    \item Kích thước trang (Page) có thể thay đổi linh hoạt cho từng đoạn chương trình và có thể gối lên nhau
    \item Bộ nhớ vật lý được chia thành các khung trang (Frame) bằng nhau, tiến trình được chia thành các trang (Page) cùng kích thước (thông dụng 4 KiB) và không yêu cầu cấp phát các khung trang liên tiếp
    \item Toàn bộ các trang của một tiến trình phải được nạp hết vào RAM trước khi tiến trình bắt đầu thực thi
\end{enumerate}
\textbf{Hướng dẫn giải:}
\begin{itemize}
    \item Theo Slide 63--64 \& 67 (CA.RISCV.2025):
    \item Kỹ thuật \textbf{phân trang (Paging)} chia không gian địa chỉ ảo thành các \textbf{trang (Page)} và bộ nhớ vật lý thành các \textbf{khung trang (Frame)} có kích thước bằng nhau (thông dụng $4\text{ KiB}$).
    \item Các trang của tiến trình \textbf{không cần nằm trên các khung trang vật lý liên tiếp} và được quản lý thông qua \textbf{Bảng trang (Page Table)}.
    \item Nhờ phân trang theo yêu cầu (Demand Paging), không cần nạp toàn bộ tiến trình vào RAM cùng lúc.
    \item \textbf{Kết luận:} Đáp án đúng là \textbf{C} (Bộ nhớ vật lý được chia thành các khung trang (Frame) bằng nhau, tiến trình được chia thành các trang (Page) cùng kích thước (thông dụng 4 KiB) và không yêu cầu cấp phát các khung trang liên tiếp).
\end{itemize}

\textbf{Bài tập 6.45 (Lý thuyết slide CA.RISCV.2025 - Mục 6.5 (Dịch địa chỉ phân trang))}\\
Khi hệ thống phân trang dịch một \textbf{địa chỉ ảo (Virtual/Logical Address)} gồm hai phần \texttt{(Page Number, Offset)} sang \textbf{địa chỉ vật lý (Physical Address)}, thao tác nào được thực hiện?
\begin{enumerate}[label=\Alph*.]
    \item Tra cứu Bảng trang (Page Table) bằng Page Number để lấy Frame Number, sau đó ghép Frame Number với phần Offset giữ nguyên
    \item Tra cứu Bảng trang bằng phần Offset để lấy Frame Number, giữ nguyên Page Number
    \item Cộng trực tiếp Page Number với địa chỉ cơ sở của hệ điều hành và bỏ phần Offset
    \item Dịch trái Page Number rồi XOR với Offset để tạo ra địa chỉ vật lý
\end{enumerate}
\textbf{Hướng dẫn giải:}
\begin{itemize}
    \item Theo Slide 66 (CA.RISCV.2025):
    \item Địa chỉ ảo gồm \texttt{[Virtual Page Number (VPN) | Page Offset]}. Hệ thống dùng \texttt{VPN} làm chỉ số tra vào \textbf{Bảng trang (Page Table)} của tiến trình để thu được \texttt{Physical Frame Number (PFN)}, còn phần \textbf{Offset (địa chỉ tương đối trong trang) được giữ nguyên} để ghép thành địa chỉ vật lý \texttt{[PFN | Offset]}.
    \item \textbf{Kết luận:} Đáp án đúng là \textbf{A} (Tra cứu Bảng trang (Page Table) bằng Page Number để lấy Frame Number, sau đó ghép Frame Number với phần Offset giữ nguyên).
\end{itemize}

\textbf{Bài tập 6.46 (Lý thuyết slide CA.RISCV.2025 - Mục 6.5 (Bài tập phân trang))}\\
Một hệ thống máy tính sử dụng không gian địa chỉ ảo $32\text{-bit}$ (đánh địa chỉ theo byte) và cơ chế phân trang với kích thước mỗi trang là $4\text{ KiB}$. Số bit dành cho phần độ lệch trong trang (\texttt{Page Offset}), số bit của số hiệu trang ảo (\texttt{VPN}) và số lượng mục (entry) tối đa trong bảng trang đơn cấp lần lượt là:
\begin{enumerate}[label=\Alph*.]
    \item Offset = 10 bit, VPN = 22 bit, tối đa $2^{22}$ mục
    \item Offset = 12 bit, VPN = 20 bit, tối đa $2^{20}$ mục
    \item Offset = 12 bit, VPN = 20 bit, tối đa $2^{12}$ mục
    \item Offset = 14 bit, VPN = 18 bit, tối đa $2^{18}$ mục
\end{enumerate}
\textbf{Hướng dẫn giải:}
\begin{itemize}
    \item Kích thước mỗi trang $= 4\text{ KiB} = 2^{12}\text{ byte} \Rightarrow$ số bit trường \texttt{Offset} $= \mathbf{12\text{ bit}}$.
    \item Địa chỉ ảo dài $32\text{ bit} \Rightarrow$ số bit trường \texttt{VPN} $= 32 - 12 = \mathbf{20\text{ bit}}$.
    \item Do có $2^{20}$ trang ảo khác nhau (từ $0$ đến $2^{20}-1$), bảng trang đơn cấp cần tối đa $\mathbf{2^{20}}$ mục (entry).
    \item \textbf{Kết luận:} Đáp án đúng là \textbf{B} (Offset = 12 bit, VPN = 20 bit, tối đa $2^{20}$ mục).
\end{itemize}

\textbf{Bài tập 6.47 (Lý thuyết slide CA.RISCV.2025 - Mục 6.5 (Lỗi trang \& Thrashing))}\\
Trong hệ thống bộ nhớ ảo phân trang, hiện tượng \textbf{Lỗi trang (Page Fault)} xảy ra khi nào và hiện tượng \textbf{Thrashing (Thất bại)} có đặc trưng gì?
\begin{enumerate}[label=\Alph*.]
    \item Page Fault xảy ra khi trang yêu cầu không có trong bộ nhớ chính (HĐH phải hoán đổi trang từ đĩa vào); Thrashing xảy ra khi có quá nhiều tiến trình trong bộ nhớ quá nhỏ khiến HĐH tốn gần như toàn bộ thời gian để hoán đổi trang mà ít thực hiện công việc hữu ích
    \item Page Fault xảy ra khi Cache L1 bị trượt; Thrashing xảy ra khi CPU chạy ở tần số xung nhịp quá cao
    \item Page Fault xảy ra khi kích thước trang lớn hơn 4 KiB; Thrashing là lỗi hỏng vật lý của ổ đĩa cứng
    \item Page Fault chỉ xảy ra ở lệnh ghi dữ liệu; Thrashing xảy ra khi tắt cơ chế bộ nhớ ảo
\end{enumerate}
\textbf{Hướng dẫn giải:}
\begin{itemize}
    \item Theo Slide 67--68 (CA.RISCV.2025):
    \item \textbf{Lỗi trang (Page Fault):} Xảy ra khi trang được CPU tham chiếu hiện không nằm trong bộ nhớ chính (bit Valid = 0). Hệ điều hành phải nạp trang đó từ bộ nhớ ngoài vào RAM (và có thể phải hoán đổi một trang khác ra đĩa).
    \item \textbf{Thất bại (Thrashing):} Khi có quá nhiều tiến trình cùng chạy trong dung lượng RAM quá nhỏ, HĐH liên tục phải hoán đổi trang (đĩa hoạt động liên tục) làm hiệu năng hệ thống sụp đổ. Giải pháp: giảm số tiến trình, bổ sung RAM hoặc cải tiến thuật toán thay trang.
    \item \textbf{Kết luận:} Đáp án đúng là \textbf{A} (Page Fault xảy ra khi trang yêu cầu không có trong bộ nhớ chính (HĐH phải hoán đổi trang từ đĩa vào); Thrashing xảy ra khi có quá nhiều tiến trình trong bộ nhớ quá nhỏ khiến HĐH tốn gần như toàn bộ thời gian để hoán đổi trang mà ít thực hiện công việc hữu ích).
\end{itemize}

\textbf{Bài tập 6.48 (Lý thuyết slide CA.RISCV.2025 - Mục 6.5 (TLB))}\\
Vai trò chính của bộ đệm dịch địa chỉ \textbf{TLB (Translation Lookaside Buffer)} trong hệ thống bộ nhớ ảo phân trang là gì?
\begin{enumerate}[label=\Alph*.]
    \item Lưu trữ toàn bộ dữ liệu của các tiến trình đang tạm dừng thay cho ổ đĩa cứng
    \item Đóng vai trò bộ nhớ Cache phần cứng tốc độ cao lưu các mục ánh xạ bảng trang (VPN $\rightarrow$ PFN) hay dùng gần đây, giúp tăng tốc quá trình dịch địa chỉ ảo sang vật lý mà không cần truy cập bảng trang trong RAM mỗi lần
    \item Thay thế hoàn toàn bảng trang (Page Table) trong bộ nhớ chính và quản lý các khung trang trống
    \item Mã hóa địa chỉ vật lý trước khi gửi ra bus địa chỉ để chống tràn số
\end{enumerate}
\textbf{Hướng dẫn giải:}
\begin{itemize}
    \item Do \textbf{Bảng trang (Page Table)} nằm trong bộ nhớ chính (RAM), nếu không có TLB thì mỗi lần CPU truy cập dữ liệu sẽ cần tới \textbf{2 lần truy cập RAM} (1 lần đọc bảng trang để dịch địa chỉ, 1 lần đọc dữ liệu thực). \textbf{TLB} là bộ đệm Cache chuyên dụng bên trong MMU lưu các ánh xạ \texttt{VPN $\rightarrow$ PFN} được sử dụng gần đây nhất (tận dụng nguyên lý cục bộ), giúp dịch địa chỉ trong 1 chu kỳ khi TLB hit.
    \item \textbf{Kết luận:} Đáp án đúng là \textbf{B} (Đóng vai trò bộ nhớ Cache phần cứng tốc độ cao lưu các mục ánh xạ bảng trang (VPN $\rightarrow$ PFN) hay dùng gần đây, giúp tăng tốc quá trình dịch địa chỉ ảo sang vật lý mà không cần truy cập bảng trang trong RAM mỗi lần).
\end{itemize}


\chapter{HỆ THỐNG VÀO-RA}

\section{Phần 1: Lý thuyết trọng tâm}

\subsection{Tổng quan Hệ thống Vào-Ra và Mô-đun Vào-Ra}

\textbf{1. Chức năng và đặc điểm của Hệ thống Vào-Ra (I/O System)}
\begin{itemize}
    \item \textbf{Chức năng:} Trao đổi thông tin giữa bên trong máy tính (CPU, bộ nhớ chính) với thế giới bên ngoài thông qua hai thao tác cơ bản là \textbf{Vào dữ liệu (Input)} và \textbf{Ra dữ liệu (Output)}.
    \item \textbf{Hai thành phần chính:} Các \textbf{thiết bị vào-ra} (thiết bị ngoại vi) và các \textbf{mô-đun vào-ra} (I/O modules).
    \item \textbf{Đặc điểm:} Thiết bị ngoại vi rất đa dạng về nguyên tắc hoạt động, tốc độ truyền và khuôn dạng dữ liệu; hầu hết đều hoạt động \textbf{chậm hơn nhiều} so với CPU và RAM. Do đó, thiết bị ngoại vi không nối trực tiếp vào bus hệ thống mà phải kết nối thông qua các \textbf{mô-đun vào-ra}.
\end{itemize}

\textbf{2. Thiết bị vào-ra (Peripherals)}
\begin{itemize}
    \item \textbf{Chức năng:} Chuyển đổi dữ liệu giữa bên trong (tín hiệu nhị phân) và bên ngoài máy tính.
    \item \textbf{Phân loại theo chức năng:} Thiết bị vào (bàn phím, chuột, máy quét), Thiết bị ra (màn hình, máy in, loa), Thiết bị lưu trữ (HDD, SSD, USB), Thiết bị truyền thông (card mạng NIC, modem).
    \item \textbf{Phân loại theo giao tiếp:} Giao tiếp \textit{Người -- Máy} và giao tiếp \textit{Máy -- Máy}.
    \item \textbf{Cấu trúc chung của thiết bị vào-ra:} Gồm 3 khối cơ bản là \textit{Bộ đệm}, \textit{Bộ chuyển đổi dữ liệu} (Transducer) và \textit{Khối logic điều khiển}.
\end{itemize}

\textbf{3. Mô-đun vào-ra (I/O Module)}
\begin{itemize}
    \item \textbf{Năm chức năng cơ bản:}
    \begin{enumerate}[label=\arabic*.]
        \item \textbf{Điều khiển và định thời:} Điều phối luồng dữ liệu giữa tài nguyên bên trong và thiết bị ngoại vi.
        \item \textbf{Trao đổi thông tin với CPU hoặc bộ nhớ chính:} Giải mã lệnh từ CPU, truyền dữ liệu và báo trạng thái qua bus hệ thống.
        \item \textbf{Trao đổi thông tin với thiết bị vào-ra:} Gửi lệnh điều khiển, nhận tín hiệu trạng thái và trao đổi dữ liệu với ngoại vi.
        \item \textbf{Đệm dữ liệu (Data Buffering):} Trung hòa sự chênh lệch tốc độ rất lớn giữa CPU/RAM và thiết bị vào-ra.
        \item \textbf{Phát hiện lỗi:} Phát hiện lỗi cơ học/điện tử ở ngoại vi hoặc lỗi truyền bit (ví dụ bit chẵn lẻ -- parity bit) và báo về CPU.
    \end{enumerate}
    \item \textbf{Cấu trúc bên trong:} Gồm các \textbf{cổng vào-ra (I/O Ports)} (thanh ghi dữ liệu, thanh ghi trạng thái, thanh ghi điều khiển), \textbf{bộ đệm dữ liệu} và \textbf{khối logic điều khiển} nối ghép giữa bus hệ thống và thiết bị ngoại vi.
\end{itemize}

\subsection{Địa chỉ hóa cổng Vào-Ra (MMIO và Isolated I/O)}

Để CPU có thể chọn đúng cổng vào-ra cần trao đổi dữ liệu, mỗi cổng vào-ra được gán một địa chỉ theo một trong hai phương pháp:

\begin{tcolorbox}[bwbox, title={So sánh Vào-ra theo bản đồ bộ nhớ (MMIO) và Vào-ra riêng biệt (Isolated I/O)}]
\begin{center}
\renewcommand{\arraystretch}{1.25}
\small
\begin{tabularx}{\linewidth}{|L{3.0cm}|Y|Y|}
\hline
\textbf{Tiêu chí} & \textbf{Vào-ra theo bản đồ bộ nhớ (MMIO)} & \textbf{Vào-ra riêng biệt (Isolated / PMIO)} \\ \hline
\textbf{Không gian địa chỉ} & Dùng \textbf{chung} một không gian địa chỉ bộ nhớ cho cả các ngăn nhớ chính và các cổng vào-ra. & Có \textbf{hai không gian địa chỉ tách biệt}: không gian địa chỉ bộ nhớ ($N$ bit) và không gian địa chỉ vào-ra riêng ($N_1$ bit). \\ \hline
\textbf{Tập lệnh sử dụng} & Dùng chính các lệnh truy cập bộ nhớ (\texttt{lw}, \texttt{sw}, \texttt{lb}, \texttt{sb} trên RISC-V) vì CPU coi cổng I/O như ngăn nhớ. & Cần các lệnh vào-ra chuyên dụng riêng biệt (ví dụ \texttt{IN}, \texttt{OUT} trên Intel x86). \\ \hline
\textbf{Ưu / Nhược điểm} & Tận dụng toàn bộ các chế độ định địa chỉ và lệnh bộ nhớ; có thể thực hiện trên mọi hệ thống; chiếm một phần không gian địa chỉ bộ nhớ. & Không chiếm không gian địa chỉ bộ nhớ chính; nhưng cần thêm đường điều khiển phân biệt bộ nhớ/I/O và giới hạn lệnh thao tác. \\ \hline
\textbf{Kiến trúc tiêu biểu} & RISC-V, MIPS, ARM, Motorola 680x0. & Intel x86 (dùng $8\text{ bit}$ hoặc $16\text{ bit}$ địa chỉ cổng vào-ra). \\ \hline
\end{tabularx}
\end{center}
\end{tcolorbox}

\begin{itemize}
    \item \textbf{Công thức tính số cổng và dải địa chỉ:} Nếu máy tính dùng $N\text{ bit}$ để đánh địa chỉ cổng vào-ra riêng biệt thì số cổng tối đa là $2^N$ cổng, có dải địa chỉ từ $0$ đến $2^N - 1$. Ví dụ với $N = 16\text{ bit}$: số cổng tối đa là $2^{16} = 65\,536$ cổng ($64\text{K}$ cổng), dải địa chỉ Hexa ($16 / 4 = 4$ chữ số Hexa) từ \texttt{0x0000} đến \texttt{0xFFFF}.
    \item \textbf{Ví dụ lập trình MMIO trên RISC-V:} Với Cổng 1 ở \texttt{0xFFFF1000} và Cổng 2 ở \texttt{0xFFFF1004}:
\end{itemize}

\begin{lstlisting}
lui  s0, 0xFFFF1     # s0 = 0xFFFF1000 (dia chi nen cua Cong 1)
addi s1, x0, 7       # nap gia tri 7 vao thanh ghi s1
sw   s1, 0(s0)       # ghi gia tri 7 ra Cong 1 (0xFFFF1000)
lw   s2, 4(s0)       # doc du lieu tu Cong 2 (0xFFFF1004) vao s2
\end{lstlisting}

\subsection{Các phương pháp điều khiển Vào-Ra (Polling, Ngắt và DMA)}

\textbf{1. Vào-ra bằng chương trình (Programmed I/O -- Polling)}
\begin{itemize}
    \item \textbf{Nguyên tắc:} \textbf{CPU trực tiếp điều khiển} quá trình trao đổi dữ liệu giữa CPU và mô-đun vào-ra bằng chương trình. Do thiết bị ngoại vi chậm hơn CPU rất nhiều, CPU phải liên tục đọc và kiểm tra thanh ghi trạng thái trong một vòng lặp chờ (\textbf{Polling / Busy-waiting}) cho đến khi thiết bị sẵn sàng.
    \item \textbf{Bốn tín hiệu điều khiển vào-ra:} \textit{Control} (kích hoạt thiết bị), \textit{Test} (kiểm tra trạng thái), \textit{Read} (đọc dữ liệu từ ngoại vi vào bộ đệm rồi về CPU), \textit{Write} (đưa dữ liệu từ bus vào bộ đệm rồi ra ngoại vi).
    \item \textbf{Nhược điểm lớn nhất:} Lãng phí rất nhiều thời gian của CPU chỉ để chờ đợi và thăm dò trạng thái thiết bị.
\end{itemize}

\textbf{2. Vào-ra điều khiển bằng ngắt (Interrupt-driven I/O)}
\begin{itemize}
    \item \textbf{Nguyên tắc:} CPU phát lệnh vào-ra rồi chuyển sang thực hiện công việc khác mà \textbf{không phải chờ đợi}. Khi mô-đun vào-ra sẵn sàng, phần cứng sẽ phát tín hiệu \textbf{ngắt (Interrupt)} gửi tới CPU.
    \item \textbf{Quy trình xử lý tại CPU:} Cuối mỗi chu trình lệnh, CPU kiểm tra tín hiệu yêu cầu ngắt. Nếu có ngắt: (1) Cất ngữ cảnh (lưu \texttt{PC} và các thanh ghi liên quan) $\rightarrow$ (2) Nhảy đến thực hiện chương trình con phục vụ ngắt (\textbf{ISR -- Interrupt Service Routine}) để trao đổi dữ liệu $\rightarrow$ (3) Khôi phục ngữ cảnh và quay lại chương trình đang chạy dở.
    \item \textbf{Bốn phương pháp nối ghép, xác định nguồn ngắt và quy định ưu tiên ngắt:}
    \begin{enumerate}[label=\arabic*.]
        \item \textit{Sử dụng nhiều đường yêu cầu ngắt:} Mỗi mô-đun nối một đường riêng tới CPU; giới hạn số mô-đun.
        \item \textit{Hỏi vòng bằng phần mềm (Software Poll):} CPU chạy phần mềm kiểm tra lần lượt từng mô-đun; tốc độ chậm; thứ tự kiểm tra là thứ tự ưu tiên.
        \item \textit{Hỏi vòng bằng phần cứng (Daisy Chain / Hardware Poll):} CPU phát tín hiệu chấp nhận ngắt (\texttt{INTA}) theo chuỗi nối tiếp qua các mô-đun; mô-đun đứng gần CPU nhất có mức ưu tiên cao nhất.
        \item \textit{Sử dụng bộ điều khiển ngắt lập trình được (\textbf{PIC -- Programmable Interrupt Controller}):} Có nhiều đường vào yêu cầu ngắt (\texttt{IRQ}) có quy định mức ưu tiên và có thanh ghi mặt nạ (\texttt{IMR}) cho phép che/cấm từng đường ngắt khả che (maskable); PIC chọn yêu cầu ngắt không bị cấm có mức ưu tiên cao nhất gửi tới CPU.
    \end{enumerate}
    \item \textbf{Thanh ghi trạng thái và điều khiển ngắt (CSRs) trong RISC-V:} Ở mức Machine (hoặc mức User trong trình mô phỏng RARS với tiền tố \texttt{u}), RISC-V quản lý ngắt thông qua các thanh ghi CSR: \texttt{mtvec}/\texttt{utvec} (địa chỉ vector trình phục vụ ngắt), \texttt{mepc}/\texttt{uepc} (lưu địa chỉ lệnh bị ngắt để quay về), \texttt{mcause}/\texttt{ucause} (mã nguyên nhân ngắt/ngoại lệ), \texttt{mstatus}/\texttt{ustatus} (bit cho phép ngắt toàn cục), \texttt{mie}/\texttt{uie} (cho phép từng nguồn ngắt riêng lẻ), \texttt{mip}/\texttt{uip} (cờ ngắt đang chờ xử lý) và lệnh trở về từ ngắt \texttt{mret}/\texttt{uret}.
\end{itemize}

\textbf{3. Truy nhập bộ nhớ trực tiếp (DMA -- Direct Memory Access)}
\begin{itemize}
    \item \textbf{Nguyên tắc:} Trong cả Polling và Interrupt-driven I/O, mọi byte/word dữ liệu đều phải đi qua CPU. Khi truyền các khối dữ liệu lớn, hệ thống sử dụng bộ điều khiển phần cứng chuyên dụng \textbf{DMAC (DMA Controller)} để truyền dữ liệu \textbf{trực tiếp giữa mô-đun vào-ra và bộ nhớ chính (RAM)} mà không đi qua CPU, giải phóng CPU làm việc khác.
    \item \textbf{Bốn thành phần chính của DMAC:}
    \begin{enumerate}[label=\arabic*.]
        \item \textbf{Thanh ghi dữ liệu (Data Register):} Chứa dữ liệu tạm thời khi trao đổi.
        \item \textbf{Thanh ghi địa chỉ (Address Register):} Chứa địa chỉ ngăn nhớ đầu của mảng dữ liệu trong RAM (tự động tăng sau mỗi từ truyền).
        \item \textbf{Bộ đếm dữ liệu (Data Counter):} Chứa số từ dữ liệu cần truyền (tự động giảm sau mỗi từ truyền; khi bằng $0$ thì phát tín hiệu ngắt báo CPU kết thúc DMA).
        \item \textbf{Logic điều khiển (Control Logic):} Điều khiển hoạt động của DMAC, thực hiện xin quyền sử dụng bus (\texttt{BRQ -- Bus Request}) và nhận tín hiệu nhượng bus (\texttt{BGT -- Bus Grant}) từ CPU.
    \end{enumerate}
    \item \textbf{Ba chế độ thực hiện DMA:}
    \begin{itemize}
        \item \textbf{DMA truyền theo khối (Block-transfer / Burst-mode DMA):} DMAC chiếm bus liên tục cho đến khi truyền xong toàn bộ khối dữ liệu rồi mới trả bus cho CPU.
        \item \textbf{DMA lấy chu kỳ (Cycle-stealing DMA):} DMAC ép CPU tạm dừng từng chu kỳ bus để truyền lần lượt từng từ dữ liệu.
        \item \textbf{DMA trong suốt (Transparent DMA):} DMAC phát hiện những chu kỳ CPU không sử dụng bus hệ thống để tranh thủ truyền dữ liệu (không làm chậm CPU). Đây là phương pháp \textbf{tối ưu hiệu năng nhất} nhưng \textbf{khó thực hiện nhất} về mặt phần cứng.
    \end{itemize}
    \item \textbf{Mở rộng -- Bộ xử lý vào-ra (I/O Processor / Channel):} Là bộ xử lý chuyên dụng hoạt động theo chương trình riêng để điều khiển toàn bộ các tác vụ vào-ra phức tạp.
\end{itemize}

\begin{tcolorbox}[bwbox, title={So sánh ba phương pháp điều khiển Vào-Ra (Hiệu suất sử dụng CPU tăng dần)}]
\begin{center}
\renewcommand{\arraystretch}{1.25}
\small
\begin{tabularx}{\linewidth}{|L{2.8cm}|Y|Y|Y|}
\hline
\textbf{Tiêu chí} & \textbf{Programmed I/O (Polling)} & \textbf{Interrupt-driven I/O (Ngắt)} & \textbf{DMA (Direct Memory Access)} \\ \hline
\textbf{Thành phần điều khiển} & CPU trực tiếp điều khiển hoàn toàn bằng phần mềm. & Kết hợp phần cứng (phát ngắt) và phần mềm (CPU chạy ISR). & DMAC điều khiển trực tiếp hoàn toàn bằng phần cứng. \\ \hline
\textbf{Đường đi dữ liệu} & I/O $\leftrightarrow$ CPU $\leftrightarrow$ RAM & I/O $\leftrightarrow$ CPU $\leftrightarrow$ RAM & I/O $\leftrightarrow$ (DMAC) $\leftrightarrow$ RAM (không qua CPU) \\ \hline
\textbf{Hiệu suất CPU} & \textbf{Thấp nhất} (lãng phí thời gian trong vòng lặp chờ). & \textbf{Trung bình} (CPU làm việc khác trong khi chờ, nhưng vẫn chuyển từng từ). & \textbf{Cao nhất} (CPU chỉ khởi tạo và nhận 1 ngắt kết thúc cả khối). \\ \hline
\textbf{Phạm vi phù hợp} & Thiết bị đơn giản, hệ thống nhỏ. & Thiết bị tốc độ thấp/trung bình (bàn phím, chuột). & Truyền khối dữ liệu lớn, tốc độ cao (ổ đĩa, mạng, đồ họa). \\ \hline
\end{tabularx}
\end{center}
\end{tcolorbox}

\subsection{Nối ghép thiết bị và Lập trình MMIO trên RISC-V RARS}

\textbf{1. Các kiểu nối ghép và cấu hình nối ghép vào-ra}
\begin{itemize}
    \item \textbf{Nối ghép song song:} Truyền nhiều bit đồng thời trên nhiều đường dữ liệu $\rightarrow$ tốc độ nhanh (ở khoảng cách ngắn) nhưng cần nhiều dây dẫn.
    \item \textbf{Nối ghép nối tiếp:} Truyền lần lượt từng bit trên ít đường dây $\rightarrow$ cần bộ chuyển đổi song song $\leftrightarrow$ nối tiếp, ít dây dẫn, ít nhiễu xuyên âm, truyền xa tốt.
    \item \textbf{Cấu hình Điểm tới điểm (Point-to-Point):} Một cổng vào-ra chuyên dụng nối ghép với đúng một thiết bị (hạn chế số lượng thiết bị).
    \item \textbf{Cấu hình Điểm tới đa điểm (Point-to-Multipoint):} Một cổng vào-ra cho phép nối ghép với nhiều thiết bị cùng lúc (ví dụ: \textbf{USB} hỗ trợ tối đa $127$ thiết bị, \textbf{Thunderbolt}).
\end{itemize}

\textbf{2. Bản đồ địa chỉ MMIO trong trình mô phỏng RARS (RISC-V)}
\begin{itemize}
    \item \textbf{Keyboard and Display MMIO Simulator:}
    \begin{itemize}
        \item \texttt{0xFFFF0000} (\texttt{KEY\_READY}): Thanh ghi trạng thái bàn phím (bit $0 = 1$ khi có phím mới được nhấn).
        \item \texttt{0xFFFF0004} (\texttt{KEY\_CODE}): Thanh ghi dữ liệu chứa mã ASCII của phím vừa nhấn.
        \item \texttt{0xFFFF0008} (\texttt{DISPLAY\_READY}): Thanh ghi trạng thái màn hình (bit $0 = 1$ khi màn hình sẵn sàng nhận ký tự mới).
        \item \texttt{0xFFFF000C} (\texttt{DISPLAY\_CODE}): Thanh ghi dữ liệu xuất ký tự ASCII lên màn hình.
    \end{itemize}
    \item \textbf{Digital Lab Sim (LED 7 thanh và Bàn phím ma trận 4x4):}
    \begin{itemize}
        \item \texttt{0xFFFF0010} (\texttt{SEVENSEG\_RIGHT}): Byte điều khiển LED 7 thanh \textbf{bên phải}.
        \item \texttt{0xFFFF0011} (\texttt{SEVENSEG\_LEFT}): Byte điều khiển LED 7 thanh \textbf{bên trái}.
        \item \texttt{0xFFFF0012} (\texttt{IN\_ADDRESS\_HEXA\_KEYBOARD}): Ghi mã chọn hàng quét (\texttt{0x01}: hàng 1, \texttt{0x02}: hàng 2, \texttt{0x04}: hàng 3, \texttt{0x08}: hàng 4) hoặc ghi bit $7 = 1$ (\texttt{0x80}) để cho phép bàn phím phát ngắt phần cứng.
        \item \texttt{0xFFFF0014} (\texttt{OUT\_ADDRESS\_HEXA\_KEYBOARD}): Đọc mã quét phím nhấn của hàng đang được chọn.
    \end{itemize}
    \item \textbf{Thiết bị Mars Bot MMIO:} Cổng góc hướng \texttt{HEADING} tại \texttt{0xFFFF8010} và cổng bật/tắt di chuyển \texttt{MOVING} tại \texttt{0xFFFF8050}; dịch vụ hệ thống \texttt{ecall} với \texttt{a7 = 32} (ngủ theo mili-giây trong \texttt{a0}) và \texttt{a7 = 34} (in số nguyên ở dạng Hexa).
\end{itemize}

\begin{tcolorbox}[bwbox, title={Bảng mã LED 7 thanh chuẩn trên RARS Digital Lab Sim (bit 0--6 ứng với a--g, bit 7 = dp)}]
\begin{center}
\renewcommand{\arraystretch}{1.2}
\small
\begin{tabularx}{\linewidth}{|Z|Z|Z|Z|Z||Z|Z|Z|Z|Z|}
\hline
\textbf{Số} & \textbf{Đoạn sáng} & \textbf{Nhị phân} & \textbf{Hex} & \textbf{Dec} & \textbf{Số} & \textbf{Đoạn sáng} & \textbf{Nhị phân} & \textbf{Hex} & \textbf{Dec} \\ \hline
\textbf{0} & \texttt{a,b,c,d,e,f} & \texttt{0011 1111} & \texttt{0x3F} & $63$ & \textbf{5} & \texttt{a,c,d,f,g} & \texttt{0110 1101} & \texttt{0x6D} & $109$ \\ \hline
\textbf{1} & \texttt{b,c} & \texttt{0000 0110} & \texttt{0x06} & $6$ & \textbf{6} & \texttt{a,c,d,e,f,g} & \texttt{0111 1101} & \texttt{0x7D} & $125$ \\ \hline
\textbf{2} & \texttt{a,b,d,e,g} & \texttt{0101 1011} & \texttt{0x5B} & $91$ & \textbf{7} & \texttt{a,b,c} & \texttt{0000 0111} & \texttt{0x07} & $7$ \\ \hline
\textbf{3} & \texttt{a,b,c,d,g} & \texttt{0100 1111} & \texttt{0x4F} & $79$ & \textbf{8} & \texttt{a,b,c,d,e,f,g} & \texttt{0111 1111} & \texttt{0x7F} & $127$ \\ \hline
\textbf{4} & \texttt{b,c,f,g} & \texttt{0110 0110} & \texttt{0x66} & $102$ & \textbf{9} & \texttt{a,b,c,d,f,g} & \texttt{0110 1111} & \texttt{0x6F} & $111$ \\ \hline
\end{tabularx}
\end{center}
\end{tcolorbox}

\newpage
\section{Phần 2: Toàn bộ bài tập Chương 7 và Hướng dẫn giải}

\subsection{Nhóm 1: Tổng quan Hệ thống Vào-Ra, Mô-đun Vào-Ra \& Định địa chỉ I/O (MMIO / Isolated)}

\textbf{Bài tập 7.1 (Đề Cuối kỳ 2021.2 -- Mã đề B3, Câu 4.2)}\\
Máy tính dùng $16$ bit để đánh địa chỉ cho cổng vào-ra riêng biệt. Số cổng vào-ra tối đa và dải địa chỉ Hexa của chúng là:
\begin{enumerate}[label=\Alph*.]
    \item $65\,536$ cổng ($2^{16}$) ; từ \texttt{0x0000} đến \texttt{0xFFFF}
    \item $32\,768$ cổng ($2^{15}$) ; từ \texttt{0x0000} đến \texttt{0x7FFF}
    \item $16\,384$ cổng ($2^{14}$) ; từ \texttt{0x0000} đến \texttt{0x3FFF}
    \item $65\,536$ cổng ($2^{16}$) ; từ \texttt{0x00000} đến \texttt{0xFFFFF}
\end{enumerate}
\textbf{Hướng dẫn giải:}
\begin{itemize}
    \item Với $16\text{ bit}$ địa chỉ, số lượng cổng vào-ra tối đa được đánh địa chỉ là $2^{16} = 65\,536$ cổng ($64\text{K}$ cổng).
    \item Vì mỗi chữ số Hexa biểu diễn $4\text{ bit}$ nên $16\text{ bit}$ cần đúng $16 / 4 = 4$ chữ số Hexa, chạy từ \texttt{0x0000} đến \texttt{0xFFFF}.
    \item \textbf{Kết luận:} Đáp án đúng là \textbf{A}.
\end{itemize}

\textbf{Bài tập 7.2 (Ôn tập cuối kỳ 2025.1 -- Câu 67; Ôn tập cuối kỳ 2025.2 -- P2, Câu 32)}\\
Không gian địa chỉ bộ nhớ dùng chung dùng để đánh địa chỉ cho các đối tượng nào?
\begin{enumerate}[label=\Alph*.]
    \item Các ngăn nhớ và cổng vào-ra
    \item Chỉ các ngăn nhớ bộ nhớ chính
    \item Chỉ các thanh ghi bên trong CPU
    \item Chỉ các khối dữ liệu trong ổ cứng ngoài
\end{enumerate}
\textbf{Hướng dẫn giải:}
\begin{itemize}
    \item Trong phương pháp \textbf{Memory-Mapped I/O} (không gian địa chỉ bộ nhớ dùng chung), các ngăn nhớ chính và cổng vào-ra (mô-đun I/O) chia sẻ chung một không gian địa chỉ hệ thống.
    \item \textbf{Kết luận:} Đáp án đúng là \textbf{A}.
\end{itemize}

\textbf{Bài tập 7.3 (Ôn tập cuối kỳ 2025.2 -- P1, Câu 9)}\\
Cổng truy cập (Port) của thiết bị Vào/Ra (I/O) được ánh xạ vào không gian địa chỉ bộ nhớ (Memory-Mapped I/O) có ý nghĩa gì?
\begin{enumerate}[label=\Alph*.]
    \item Thiết bị I/O chỉ có thể được truy cập thông qua các lệnh LOAD/STORE, giống như truy cập bộ nhớ.
    \item CPU cần một tập lệnh I/O đặc biệt (ví dụ: IN/OUT) để giao tiếp với thiết bị.
    \item Thiết bị I/O có bus địa chỉ hoàn toàn tách biệt với bộ nhớ chính.
    \item Thiết bị I/O có thể tự xử lý lệnh mà không cần bus hệ thống.
\end{enumerate}
\textbf{Hướng dẫn giải:}
\begin{itemize}
    \item Với Memory-Mapped I/O (như trên RISC-V), thanh ghi cổng I/O mang địa chỉ nằm trong không gian địa chỉ bộ nhớ nên CPU dùng chính các lệnh truy cập bộ nhớ (\texttt{lw}, \texttt{lb}, \texttt{sw}, \texttt{sb}) để đọc/ghi dữ liệu ngoại vi.
    \item \textbf{Kết luận:} Đáp án đúng là \textbf{A}.
\end{itemize}

\subsection{Nhóm 2: Các phương pháp điều khiển Vào-Ra (Polling, Ngắt, DMA \& DMAC)}

\textbf{Bài tập 7.4 (Đề Cuối kỳ 2022.2; Cuối kỳ 2024.2; Ôn tập cuối kỳ 2025.2 -- P1, Câu 15)}\\
Với phương pháp vào-ra bằng chương trình, thành phần nào của máy tính trực tiếp điều khiển vào-ra?
\begin{enumerate}[label=\Alph*.]
    \item CPU
    \item DMAC
    \item Thiết bị vào-ra
    \item RAM
\end{enumerate}
\textbf{Hướng dẫn giải:}
\begin{itemize}
    \item Trong phương pháp vào-ra bằng chương trình (cũng như vào-ra bằng ngắt), \textbf{CPU} trực tiếp điều khiển việc trao đổi dữ liệu với mô-đun vào-ra (không thông qua DMAC).
    \item \textbf{Kết luận:} Đáp án đúng là \textbf{A (CPU)}.
\end{itemize}

\textbf{Bài tập 7.5 (Cuối kỳ 2023.1, 2023.2, 2025.1; Thi thử 2025.1 -- Câu 7; Ôn tập 2025.1 -- Câu 10; Ôn tập 2025.2 -- P2, Câu 75)}\\
Cho chương trình con xử lý vào-ra dữ liệu RISC-V dưới đây, trong đó CPU đang cố gắng đọc một từ dữ liệu từ một thiết bị ngoại vi. Phương pháp lập trình Vào-Ra (I/O programming method) nào đang được áp dụng?
\begin{lstlisting}
.text
    li s0, 0xFFFF0004   # Thanh ghi Du lieu (Data register)
    li s1, 0xFFFF0000   # Thanh ghi Trang thai (Status register)
Check:
    lw   t1, 0(s1)      # Doc thanh ghi trang thai
    andi t1, t1, 0x1
    beq  t1, zero, Check
    lw   t0, 0(s0)      # Doc du lieu tu thiet bi
\end{lstlisting}
\begin{enumerate}[label=\Alph*.]
    \item Thăm dò trạng thái (Polling / Programmed I/O)
    \item Vào-ra bằng cơ chế ngắt (Interrupt-driven I/O)
    \item Truy cập bộ nhớ trực tiếp (DMA)
    \item Vào-ra chồng lấn (Overlapping I/O)
\end{enumerate}
\textbf{Hướng dẫn giải:}
\begin{itemize}
    \item Vòng lặp \texttt{Check} liên tục đọc thanh ghi trạng thái tại \texttt{0xFFFF0000} và tách bit $0$ (\texttt{andi t1, t1, 0x1}) để kiểm tra lặp đi lặp lại cho tới khi bằng $1$ mới đọc dữ liệu từ \texttt{0xFFFF0004}. Đây chính là cơ chế \textbf{Thăm dò trạng thái (Polling / Programmed I/O)}.
    \item \textbf{Kết luận:} Đáp án đúng là \textbf{A}.
\end{itemize}

\textbf{Bài tập 7.6 (Ôn tập cuối kỳ 2025.1 -- Câu 75; Ôn tập cuối kỳ 2025.2 -- P2, Câu 42)}\\
Những khẳng định nào sau đây là đúng về bộ điều khiển ngắt \texttt{PIC}?
\begin{enumerate}[label=\Alph*.]
    \item PIC là viết tắt của Programmable Interrupt Controller và các đường vào yêu cầu ngắt (IRQ) có quy định mức ưu tiên
    \item PIC là viết tắt của Peripheral I/O Cache dùng tăng tốc truyền dữ liệu qua DMA
    \item Yêu cầu ngắt có mức ưu tiên cao nhất trong PIC tuyệt đối không bao giờ bị cấm/che bởi thanh ghi mặt nạ (IMR)
    \item PIC chỉ phục vụ duy nhất 1 nguồn ngắt cố định tại một thời điểm, không có thanh ghi mặt nạ
\end{enumerate}
\textbf{Hướng dẫn giải:}
\begin{itemize}
    \item \textbf{PIC (Programmable Interrupt Controller)} là bộ điều khiển ngắt lập trình được, quản lý nhiều đường vào yêu cầu ngắt (\texttt{IRQ}) theo mức ưu tiên xác định.
    \item PIC vẫn có thanh ghi mặt nạ (\texttt{IMR}) cho phép che/cấm từng đường ngắt khả che khi cần thiết; PIC chọn yêu cầu ngắt \textit{không bị cấm} có mức ưu tiên cao nhất gửi tới CPU.
    \item \textbf{Kết luận:} Đáp án đúng là \textbf{A}.
\end{itemize}

\textbf{Bài tập 7.7 (Cuối kỳ 2025.1; Thi thử 2025.1 -- Câu 13)}\\
What is the main disadvantage of polling when CPU reads from or writes to I/O devices?
\begin{enumerate}[label=\Alph*.]
    \item The CPU wastes time constantly checking device status.
    \item It has higher throughput than DMA for bulk transfer.
    \item Polling requires complex hardware interrupt controllers.
    \item Polling cannot be used with memory-mapped I/O.
\end{enumerate}
\textbf{Hướng dẫn giải:}
\begin{itemize}
    \item Nhược điểm lớn nhất của Polling là CPU phải lãng phí rất nhiều chu kỳ trong vòng lặp chờ bận (busy-waiting) để kiểm tra liên tục thanh ghi trạng thái thiết bị ngoại vi.
    \item \textbf{Kết luận:} Đáp án đúng là \textbf{A}.
\end{itemize}

\textbf{Bài tập 7.8 (Cuối kỳ 2025.1; Thi thử 2025.1 -- Câu 18)}\\
What is polling in the context of I/O programming?
\begin{enumerate}[label=\Alph*.]
    \item A method where the CPU checks the status of a device continuously
    \item A method where a device interrupts the CPU when data is ready
    \item A method where the DMA Controller manages direct memory transfer
    \item A method where I/O bypasses system buses completely
\end{enumerate}
\textbf{Hướng dẫn giải:}
\begin{itemize}
    \item Polling (thăm dò) là kỹ thuật lập trình vào-ra trong đó CPU chủ động kiểm tra liên tục thanh ghi trạng thái của thiết bị ngoại vi xem thiết bị đã sẵn sàng trao đổi dữ liệu hay chưa.
    \item \textbf{Kết luận:} Đáp án đúng là \textbf{A}.
\end{itemize}

\textbf{Bài tập 7.9 (Ôn tập cuối kỳ 2025.1 -- Câu 12; Ôn tập cuối kỳ 2025.2 -- P2, Câu 10)}\\
Tại sao phương pháp DMA (Direct Memory Access) lại hiệu quả hơn khi truyền các khối dữ liệu lớn?
\begin{enumerate}[label=\Alph*.]
    \item Vì dữ liệu được truyền trực tiếp giữa thiết bị I/O và bộ nhớ chính thông qua DMAC mà không cần đi qua CPU.
    \item Vì DMA dùng một tập lệnh phần mềm đặc biệt chạy nhanh hơn gấp 10 lần trên thanh ghi CPU.
    \item Vì DMA không cần dùng đến bộ nhớ chính (RAM) khi truyền dữ liệu.
    \item Vì DMA chuyển toàn bộ dữ liệu trực tiếp qua thanh ghi L1 Cache của CPU.
\end{enumerate}
\textbf{Hướng dẫn giải:}
\begin{itemize}
    \item Khi truyền các khối dữ liệu kích thước lớn, bộ điều khiển DMA (DMAC) mượn quyền bus để truyền thẳng dữ liệu giữa mô-đun vào-ra và RAM hoàn toàn bằng phần cứng mà không đi qua CPU từng byte/word, nhờ đó tăng tốc độ truyền và giải phóng CPU thực hiện công việc khác.
    \item \textbf{Kết luận:} Đáp án đúng là \textbf{A}.
\end{itemize}

\textbf{Bài tập 7.10 (Ôn tập cuối kỳ 2025.1 -- Câu 18; Ôn tập cuối kỳ 2025.2 -- P2, Câu 20)}\\
Thứ tự nào sau đây phản ánh đúng sự phát triển về hiệu suất (từ thấp đến cao) đối với việc sử dụng CPU trong I/O?
\begin{enumerate}[label=\Alph*.]
    \item Programmed I/O $\rightarrow$ Interrupt-driven $\rightarrow$ DMA
    \item DMA $\rightarrow$ Interrupt-driven $\rightarrow$ Programmed I/O
    \item Interrupt-driven $\rightarrow$ Programmed I/O $\rightarrow$ DMA
    \item Programmed I/O $\rightarrow$ DMA $\rightarrow$ Interrupt-driven
\end{enumerate}
\textbf{Hướng dẫn giải:}
\begin{itemize}
    \item Hiệu suất sử dụng CPU tăng dần theo chiều giảm mức độ bận rộn của CPU vào công việc trao đổi dữ liệu: \textbf{Programmed I/O (Polling)} (CPU phải chờ và tự truyền) $\rightarrow$ \textbf{Interrupt-driven I/O (Ngắt)} (CPU không phải chờ nhưng vẫn truyền từng từ) $\rightarrow$ \textbf{DMA} (DMAC tự truyền cả khối).
    \item \textbf{Kết luận:} Đáp án đúng là \textbf{A}.
\end{itemize}

\textbf{Bài tập 7.11 (Ôn tập cuối kỳ 2025.1 -- Câu 21; Ôn tập cuối kỳ 2025.2 -- P2, Câu 74)}\\
Trong phương pháp vào/ra bằng chương trình (Programmed I/O), CPU dành phần lớn thời gian để làm gì khi đợi thiết bị ngoại vi sẵn sàng?
\begin{enumerate}[label=\Alph*.]
    \item Kiểm tra trạng thái thiết bị trong một vòng lặp (Polling).
    \item Tự động chuyển sang chế độ ngủ (Sleep mode).
    \item Thực hiện các tiến trình khác của hệ điều hành.
    \item Nhượng quyền bus hoàn toàn cho DMAC xử lý.
\end{enumerate}
\textbf{Hướng dẫn giải:}
\begin{itemize}
    \item Trong Programmed I/O, CPU tốn phần lớn thời gian chờ đợi bằng cách liên tục đọc và kiểm tra cờ trạng thái của mô-đun vào-ra trong một vòng lặp (Polling).
    \item \textbf{Kết luận:} Đáp án đúng là \textbf{A}.
\end{itemize}

\textbf{Bài tập 7.12 (Ôn tập cuối kỳ 2025.1 -- Câu 39; Ôn tập cuối kỳ 2025.2 -- P2, Câu 26)}\\
Bộ điều khiển truy cập bộ nhớ trực tiếp (\texttt{DMAC}) bao gồm đầy đủ những thành phần nào sau đây?
\begin{enumerate}[label=\Alph*.]
    \item Logic điều khiển, Thanh ghi dữ liệu, Thanh ghi địa chỉ, Bộ đếm dữ liệu (Data Counter)
    \item Logic điều khiển, Thanh ghi con trỏ lệnh (PC), Thanh ghi dữ liệu, ALU
    \item Thanh ghi lệnh (IR), Bộ đếm dữ liệu, Bộ nhớ Cache L1, Thanh ghi địa chỉ
    \item Logic điều khiển, Bộ giải mã lệnh (Decoder), Thanh ghi PC, Thanh ghi địa chỉ
\end{enumerate}
\textbf{Hướng dẫn giải:}
\begin{itemize}
    \item Theo cấu trúc chuẩn của DMAC, bộ điều khiển này gồm 4 thành phần chính: \textbf{Thanh ghi địa chỉ} (Address Register), \textbf{Bộ đếm dữ liệu} (Data Counter), \textbf{Thanh ghi dữ liệu} (Data Register) và \textbf{Logic điều khiển} (gửi/nhận tín hiệu \texttt{BRQ}/\texttt{BGT}, tín hiệu đọc/ghi). DMAC không chứa \texttt{PC}, \texttt{IR} hay \texttt{ALU}.
    \item \textbf{Kết luận:} Đáp án đúng là \textbf{A}.
\end{itemize}

\textbf{Bài tập 7.13 (Ôn tập cuối kỳ 2025.1 -- Câu 40; Ôn tập cuối kỳ 2025.2 -- P2, Câu 72)}\\
Cách thực hiện DMA nào dưới đây là tối ưu và khó thực hiện hơn cả?
\begin{enumerate}[label=\Alph*.]
    \item DMA trong suốt (Transparent DMA)
    \item DMA truyền theo khối (Burst-mode DMA)
    \item DMA lấy chu kỳ (Cycle-stealing DMA)
    \item Ba phương pháp đều có độ phức tạp phần cứng như nhau
\end{enumerate}
\textbf{Hướng dẫn giải:}
\begin{itemize}
    \item \textbf{DMA trong suốt (Transparent DMA)} chỉ tranh thủ sử dụng bus hệ thống trong những chu kỳ mà CPU không truy cập bộ nhớ, nên không hề làm chậm tiến độ thực thi của CPU. Tuy nhiên, mạch phần cứng để phát hiện chính xác các khe thời gian trống này rất phức tạp.
    \item \textbf{Kết luận:} Đáp án đúng là \textbf{A}.
\end{itemize}

\subsection{Nhóm 3: Lập trình MMIO RISC-V (RARS, LED 7 thanh, Bàn phím) \& Ngắt ngoại vi}

\textbf{Bài tập 7.14 (Cuối kỳ 2022.2, 2024.1, 2024.2; Ôn tập cuối kỳ 2025.2 -- P1, Câu 38)}\\
Cho máy tính dùng bộ xử lý RISC-V (RV32I) có 4 cổng vào-ra MMIO 32-bit gán tại các địa chỉ: \texttt{P1: 0xF000000C}, \texttt{P2: 0xF0000010}, \texttt{P3: 0xF0000014}, \texttt{P4: 0xF0000018}. Đoạn chương trình hợp ngữ sau thực hiện nhiệm vụ gì?
\begin{lstlisting}
lui s0, 0xF0000
lw  t0, 0x10(s0)
sw  t0, 0x18(s0)
lw  t1, 0x0C(s0)
sw  t1, 0x14(s0)
\end{lstlisting}
\begin{enumerate}[label=\Alph*.]
    \item Nhận dữ liệu từ cổng P2 đưa ra cổng P4, nhận dữ liệu từ cổng P1 đưa ra cổng P3
    \item Nhận dữ liệu từ cổng P1 đưa ra cổng P2, nhận dữ liệu từ cổng P3 đưa ra cổng P4
    \item Nhận dữ liệu từ cổng P4 đưa ra cổng P2, nhận dữ liệu từ cổng P3 đưa ra cổng P1
    \item Nhận dữ liệu từ cổng P1 đưa ra cổng P4, nhận dữ liệu từ cổng P2 đưa ra cổng P3
\end{enumerate}
\textbf{Hướng dẫn giải:}
\begin{itemize}
    \item Lệnh \texttt{lui s0, 0xF0000} đặt địa chỉ nền \texttt{s0 = 0xF0000000}.
    \item \texttt{lw t0, 0x10(s0)} đọc từ địa chỉ \texttt{0xF0000010} (\textbf{P2}), sau đó \texttt{sw t0, 0x18(s0)} ghi ra địa chỉ \texttt{0xF0000018} (\textbf{P4}).
    \item \texttt{lw t1, 0x0C(s0)} đọc từ địa chỉ \texttt{0xF000000C} (\textbf{P1}), sau đó \texttt{sw t1, 0x14(s0)} ghi ra địa chỉ \texttt{0xF0000014} (\textbf{P3}).
    \item \textbf{Kết luận:} Đáp án đúng là \textbf{A}.
\end{itemize}

\textbf{Bài tập 7.15 (Lab -- Cuối kỳ 2023.2, Câu 15)}\\
Đoạn mã RISC-V dưới đây điều khiển thiết bị Bot MMIO di chuyển liên tục thành hình vuông (biết cổng điều khiển góc hướng \texttt{HEADING} tại \texttt{0xFFFF8010} và cổng kích hoạt chuyển động \texttt{MOVING} tại \texttt{0xFFFF8050}; \texttt{a7 = 32} là dịch vụ sleep theo ms). Cặp lệnh thích hợp điền vào \texttt{[[1]]} và \texttt{[[2]]} là gì?
\begin{lstlisting}
.text
    li  s0, 0xFFFF8010
    li  s1, 0xFFFF8050
    li  t0, 0
    li  t1, 1
    [[1]]
l1:
    addi t0, t0, 90
    [[2]]
    li  a7, 32
    li  a0, 2000
    ecall
    j   l1
\end{lstlisting}
\begin{enumerate}[label=\Alph*.]
    \item \texttt{[[1]]: sw t1, 0(s1)} \textbar{} \texttt{[[2]]: sw t0, 0(s0)}
    \item \texttt{[[1]]: sw t0, 0(s0)} \textbar{} \texttt{[[2]]: sw t1, 0(s1)}
    \item \texttt{[[1]]: sw t1, 0(s0)} \textbar{} \texttt{[[2]]: sw t0, 0(s1)}
    \item \texttt{[[1]]: sw t0, 0(s1)} \textbar{} \texttt{[[2]]: sw t1, 0(s0)}
\end{enumerate}
\textbf{Hướng dẫn giải:}
\begin{itemize}
    \item Thanh ghi \texttt{t1 = 1} dùng để kích hoạt trạng thái di chuyển tại cổng \texttt{MOVING} (trỏ bởi \texttt{s1 = 0xFFFF8050}), do đó \texttt{[[1]]} là \texttt{sw t1, 0(s1)}.
    \item Trong vòng lặp \texttt{l1}, sau khi cộng góc \texttt{t0} thêm $90^\circ$, cần ghi góc mới ra cổng \texttt{HEADING} (trỏ bởi \texttt{s0 = 0xFFFF8010}), do đó \texttt{[[2]]} là \texttt{sw t0, 0(s0)}.
    \item \textbf{Kết luận:} Đáp án đúng là \textbf{A}.
\end{itemize}

\textbf{Bài tập 7.16 (Cuối kỳ 2025.1; Thi thử 2025.1 -- Câu 3; Ôn tập cuối kỳ 2025.2 -- P2, Câu 2)}\\
Trên thiết bị mô phỏng Digital Lab Sim (RARS), thanh ghi byte điều khiển LED 7 thanh được ánh xạ tại địa chỉ \texttt{0xFFFF0010} (bit 0 = a [trên], bit 1 = b [phải trên], bit 2 = c [phải dưới], bit 3 = d [đáy], bit 4 = e [trái dưới], bit 5 = f [trái trên], bit 6 = g [giữa], bit 7 = dấu chấm). Đoạn mã nào hiển thị chữ số \textbf{1}?
\begin{enumerate}[label=\Alph*.]
    \item \texttt{li t0, 0xFFFF0010 ; li t1, 6 ; sb t1, 0(t0)}
    \item \texttt{li t0, 0xFFFF0010 ; li t1, 1 ; sb t1, 0(t0)}
    \item \texttt{li t0, 0xFFFF0010 ; li t1, 0x02 ; sb t1, 0(t0)}
    \item \texttt{li t0, 0xFFFF0010 ; li t1, 0x3F ; sb t1, 0(t0)}
\end{enumerate}
\textbf{Hướng dẫn giải:}
\begin{itemize}
    \item Chữ số \textbf{1} cần làm sáng 2 đoạn bên phải là \texttt{b} (bit 1) và \texttt{c} (bit 2).
    \item Giá trị byte tương ứng là $0000\;0110_{(2)} = 2^1 + 2^2 = \mathbf{6}$ (\texttt{0x06}), được ghi ra địa chỉ \texttt{0xFFFF0010} bằng lệnh \texttt{sb}.
    \item \textbf{Kết luận:} Đáp án đúng là \textbf{A}.
\end{itemize}

\textbf{Bài tập 7.17 (Lab -- Cuối kỳ 2023.2, Câu 28)}\\
Lựa chọn phát biểu mô tả đúng nhất chức năng của đoạn mã RISC-V dưới đây (giao tiếp bàn phím MMIO):
\begin{lstlisting}
.eqv KEY_CODE  0xFFFF0004
.eqv KEY_READY 0xFFFF0000
.text
    li  t5, KEY_CODE
    li  t6, KEY_READY
loop:
Wait:
    lw  t1, 0(t6)
    beq t1, zero, Wait
    lw  t0, 0(t5)
    j   loop
\end{lstlisting}
\begin{enumerate}[label=\Alph*.]
    \item Chờ có 1 phím sẵn sàng rồi đọc mã phím đó lưu vào thanh ghi \texttt{t0}
    \item Chờ có 1 phím sẵn sàng rồi xuất mã phím đó ra màn hình
    \item Đọc mã phím lưu vào thanh ghi \texttt{t1}, đọc trạng thái bàn phím lưu vào thanh ghi \texttt{t0}
    \item Chờ có 1 phím sẵn sàng rồi đọc mã phím đó lưu vào thanh ghi \texttt{t1}
\end{enumerate}
\textbf{Hướng dẫn giải:}
\begin{itemize}
    \item Vòng lặp \texttt{Wait} liên tục đọc cờ trạng thái \texttt{KEY\_READY} (tại \texttt{t6 = 0xFFFF0000}) vào thanh ghi \texttt{t1}.
    \item Khi \texttt{t1} $\ne 0$ (có phím mới được nhấn), chương trình thoát vòng chờ và thực hiện \texttt{lw t0, 0(t5)} để đọc mã phím từ \texttt{KEY\_CODE} (\texttt{0xFFFF0004}) lưu vào thanh ghi \texttt{t0}.
    \item \textbf{Kết luận:} Đáp án đúng là \textbf{A}.
\end{itemize}

\textbf{Bài tập 7.18 (Lab -- Cuối kỳ 2023.2, Câu 32)}\\
Trên RARS Digital Lab Sim, 2 LED 7 thanh có địa chỉ byte ánh xạ như dưới đây. Để hiển thị chữ số \textbf{9} trên \textbf{LED TRÁI} (quy ước bit $0..6 = a..g$), đoạn lệnh nào đúng?
\begin{lstlisting}
.eqv SEVENSEG_LEFT  0xFFFF0011
.eqv SEVENSEG_RIGHT 0xFFFF0010
\end{lstlisting}
\begin{enumerate}[label=\Alph*.]
    \item \texttt{li t0, 0xFFFF0011 ; li a0, 0x6F ; sb a0, 0(t0)}
    \item \texttt{li t0, 0xFFFF0011 ; li a0, 0x7D ; sb a0, 0(t0)}
    \item \texttt{li t0, 0xFFFF0010 ; li a0, 0x6F ; sb a0, 0(t0)}
    \item \texttt{li t0, 0xFFFF0010 ; li a0, 0x7D ; sb a0, 0(t0)}
\end{enumerate}
\textbf{Hướng dẫn giải:}
\begin{itemize}
    \item Địa chỉ của \textbf{LED TRÁI} là \texttt{0xFFFF0011}.
    \item Chữ số \textbf{9} bật các đoạn \texttt{a, b, c, d, f, g} (chỉ tắt đoạn \texttt{e} ở bit 4 và dấu chấm ở bit 7), nên giá trị byte là $0110\;1111_{(2)} =$ \textbf{\texttt{0x6F}} (còn \texttt{0x7D} hiển thị chữ số 6).
    \item \textbf{Kết luận:} Đáp án đúng là \textbf{A}.
\end{itemize}

\textbf{Bài tập 7.19 (Cuối kỳ 2025.1; Thi thử 2025.1 -- Câu 27; Ôn tập 2025.1 -- Câu 5; Ôn tập 2025.2 -- P2, Câu 29)}\\
Xét chương trình MMIO trên RARS dưới đây. Nếu người dùng nhấn phím \texttt{'A'} trong Keyboard Simulator, ký tự nào sẽ được hiển thị trên Display Simulator? (chỉ ghi ký tự, không viết dấu ngoặc kép):
\begin{lstlisting}
.eqv KEY_CODE      0xFFFF0004
.eqv KEY_READY     0xFFFF0000
.eqv DISPLAY_CODE  0xFFFF000C
.eqv DISPLAY_READY 0xFFFF0008
.text
    li a0, KEY_CODE
    li a1, KEY_READY
    li s0, DISPLAY_CODE
    li s1, DISPLAY_READY
loop:
WaitForKey:
    lw  t1, 0(a1)
    beq t1, zero, WaitForKey
ReadKey:
    lw  t0, 0(a0)
WaitForDis:
    lw  t2, 0(s1)
    beq t2, zero, WaitForDis
Encrypt:
    addi t0, t0, 2
ShowKey:
    sw  t0, 0(s0)
    j   loop
\end{lstlisting}
\textit{Nhập ký tự hiển thị ra Display MMIO:} \underline{\hspace{2.5cm}}\\[4pt]
\textbf{Hướng dẫn giải:}
\begin{itemize}
    \item Chương trình chờ phím nhấn tại \texttt{KEY\_READY}, đọc mã ASCII của phím \texttt{'A'} (giá trị $65$) vào thanh ghi \texttt{t0}.
    \item Sau khi chờ màn hình sẵn sàng tại \texttt{DISPLAY\_READY}, lệnh \texttt{addi t0, t0, 2} cộng thêm $2$ vào mã ASCII thành $65 + 2 = 67$, tương ứng với ký tự \textbf{\texttt{C}}, rồi xuất ra \texttt{DISPLAY\_CODE}.
    \item \textbf{Kết luận:} Đáp án đúng là \textbf{\texttt{C}}.
\end{itemize}

\textbf{Bài tập 7.20 (Cuối kỳ 2025.1; Thi thử 2025.1 -- Câu 32)}\\
Cho máy tính RISC-V có 4 cổng vào-ra ánh xạ bộ nhớ (MMIO) ở các địa chỉ: \texttt{P1: 0xFFFF000C}, \texttt{P2: 0xFFFF0010}, \texttt{P3: 0xFFFF0014}, \texttt{P4: 0xFFFF0018}. Đoạn mã hợp ngữ sau thực hiện thao tác gì?
\begin{lstlisting}
lui s0, 0xFFFF0
lw  t0, 0x0C(s0)
sw  t0, 0x18(s0)
lw  t1, 0x10(s0)
sw  t1, 0x14(s0)
\end{lstlisting}
\begin{enumerate}[label=\Alph*.]
    \item Nhận dữ liệu từ P1 xuất ra P4; nhận dữ liệu từ P2 xuất ra P3
    \item Nhận dữ liệu từ P1 xuất ra P3; nhận dữ liệu từ P2 xuất ra P4
    \item Nhận dữ liệu từ P4 xuất ra P1; nhận dữ liệu từ P3 xuất ra P2
    \item Nhận dữ liệu từ P1 xuất ra P2; nhận dữ liệu từ P3 xuất ra P4
\end{enumerate}
\textbf{Hướng dẫn giải:}
\begin{itemize}
    \item Lệnh \texttt{lui s0, 0xFFFF0} đặt địa chỉ nền \texttt{s0 = 0xFFFF0000}.
    \item Offset \texttt{0x0C} (\textbf{P1}) được đọc vào \texttt{t0} rồi ghi ra offset \texttt{0x18} (\textbf{P4}); tiếp theo offset \texttt{0x10} (\textbf{P2}) được đọc vào \texttt{t1} rồi ghi ra offset \texttt{0x14} (\textbf{P3}).
    \item \textbf{Kết luận:} Đáp án đúng là \textbf{A}.
\end{itemize}

\textbf{Bài tập 7.21 (Lab -- Cuối kỳ 2022.2, Câu 1; Lab -- Cuối kỳ 2023.2, Câu 9)}\\
Chọn phát biểu đúng về chương trình xử lý bàn phím ma trận dưới đây viết cho trình mô phỏng RARS (User-level Trap CSRs: \texttt{utvec}, \texttt{ustatus}, \texttt{uie}, \texttt{uepc}, \texttt{uret}):
\begin{lstlisting}
.eqv IN_ADDRESS_HEXA_KEYBOARD 0xFFFF0012
.data
    Message: .asciz "Key is pressed\n"
.text
main:
    la   t0, IntSR
    csrw utvec, t0           # Dang ky dia chi trinh phuc vu ngat
    li   t0, 0x100
    csrs uie, t0             # Bat ngat ngoai vi (UEIE)
    csrsi ustatus, 1         # Bat co ngat toan cuc UIE
    li   t1, IN_ADDRESS_HEXA_KEYBOARD
    li   t3, 0x80
    sb   t3, 0(t1)           # Cho phep ban phim ma tran phat ngat
Loop:
    nop
    j    Loop
IntSR:
    li   a7, 4
    la   a0, Message
    ecall
    csrr t0, uepc
    addi t0, t0, 4
    csrw uepc, t0
    uret
\end{lstlisting}
\begin{enumerate}[label=\Alph*.]
    \item Chương trình cài đặt cơ chế ngắt (Interrupt): khi có phím bấm trên keyboard matrix, CPU tự động chuyển sang trình phục vụ ngắt \texttt{IntSR} (đăng ký qua \texttt{utvec}) rồi kết thúc bằng \texttt{uret}
    \item Chương trình cài đặt cơ chế thăm dò (Polling): vòng lặp chính liên tục đọc cổng \texttt{0xFFFF0012} để kiểm tra phím rồi in thông báo
    \item Chương trình cài đặt cơ chế ngắt trong đó vòng \texttt{Loop} chủ động gọi trực tiếp \texttt{jal IntSR} để kiểm tra phím
    \item Chương trình cài đặt cơ chế Polling trong đó chương trình chính liên tục gọi hàm \texttt{ecall} trong \texttt{Loop}
\end{enumerate}
\textbf{Hướng dẫn giải:}
\begin{itemize}
    \item Chương trình đăng ký địa chỉ trình phục vụ ngắt \texttt{IntSR} vào thanh ghi \texttt{utvec}, bật ngắt ngoại vi trong \texttt{uie} và cờ ngắt toàn cục trong \texttt{ustatus}, đồng thời ghi \texttt{0x80} (bit 7 = 1) ra cổng \texttt{0xFFFF0012} để kích hoạt bàn phím phát tín hiệu ngắt phần cứng.
    \item Trong vòng lặp chính \texttt{Loop}, CPU chỉ chạy \texttt{nop} mà không hề thăm dò cổng vào-ra; khi có phím nhấn, CPU tự động nhảy vào \texttt{IntSR} để in chuỗi và trở về bằng \texttt{uret}.
    \item \textbf{Kết luận:} Đáp án đúng là \textbf{A}.
\end{itemize}

\textbf{Bài tập 7.22 (Lab -- Cuối kỳ 2022.2, Câu 6; Lab -- Cuối kỳ 2023.2, Câu 11)}\\
Chương trình RISC-V sau thực hiện quét bàn phím ma trận 4x4 trong công cụ RARS Digital Lab Sim (với \texttt{t1 = 0xFFFF0012 [IN\_ADDRESS]}, \texttt{t2 = 0xFFFF0014 [OUT\_ADDRESS]}). Phát biểu nào sau đây mô tả ĐÚNG hoạt động của chương trình?
\begin{lstlisting}
polling:
    li   t3, 0x08
    sb   t3, 0(t1)
    lb   a0, 0(t2)
print:
    li   a7, 34
    ecall
sleep:
    li   a0, 100
    li   a7, 32
    ecall
    j    polling
\end{lstlisting}
\begin{enumerate}[label=\Alph*.]
    \item Chương trình chỉ quét hàng 4 (bit 3 = \texttt{0x08}) của bàn phím ma trận bằng phương pháp thăm dò (polling) thông qua ghi mã hàng vào \texttt{0xFFFF0012}
    \item Chương trình quét 1 hàng của bàn phím bằng cách ghi chỉ số hàng trong \texttt{t3} vào địa chỉ \texttt{0xFFFF0014}
    \item Chương trình tự động quét luân phiên cả 4 hàng của bàn phím ma trận bằng phương pháp thăm dò
    \item Chương trình hoạt động dựa trên cơ chế ngắt phần cứng (Interrupt) kích hoạt hàng 4 mỗi 100 ms
\end{enumerate}
\textbf{Hướng dẫn giải:}
\begin{itemize}
    \item Giá trị \texttt{0x08} ($0000\;1000_{(2)}$) chỉ bật bit 3, tương ứng với việc chọn duy nhất \textbf{hàng 4} của bàn phím ma trận 4x4.
    \item Chương trình ghi \texttt{0x08} vào địa chỉ chọn hàng \texttt{0xFFFF0012} (\texttt{t1}), đọc mã phím từ \texttt{0xFFFF0014} (\texttt{t2}), in ra ở dạng Hexa (\texttt{a7 = 34}), ngủ $100\text{ ms}$ (\texttt{a7 = 32}) rồi lặp vô hạn theo phương pháp thăm dò (polling).
    \item \textbf{Kết luận:} Đáp án đúng là \textbf{A}.
\end{itemize}


\end{document}